%& -job-name=Internal
\documentclass[a4paper,landscape]{memoir}
\usepackage{chlewrics}
\usepackage{multicol}
\setlength{\columnsep}{6mm}
\setsecnumdepth{chapter}
\begin{document}
\frontmatter
\chapterstyle{thatcher}

\begin{titlingpage}
\centering
\vspace*{\fill}
{\LARGE\bfseries Internal and Non-Fiction\\\Huge Songbook\par}
\bigskip
{\large A private notebook, set to music\par}
\vspace*{\fill}
{\normalsize Directed and Curated by \bfseries Carlos Eŭgenĭo Thompson Pinzón\par}
{\normalsize Produced by \bfseries Suno\par}
\bigskip
{\normalsize CC-BY\par}
{\normalsize 2026\par}
\vspace*{\fill}
\end{titlingpage}

\clearpage\begin{multicols*}{2}
\tableofcontents*
\end{multicols*}
\mainmatter
\renewcommand\chaptername{Album~}

\part{Personal Chronicle}

\chapter{Earlier Thoughts I (1988–1992)}
\albumcover{covers/internal/earlier-thoughts-i}
The earliest layer of the whole songbook: a teenager's hatred-poem in three languages (\emph{Je hais}, \emph{Odio}, \emph{Jag avskyr}), a Swedish cycle of waking up to black eyes and static (\emph{Despertar}), and the first travel notes of a life that hadn't started moving yet. None of it was written to be sung -- it was journal entries and language-class exercises that happened to survive four decades, waiting for a voice.
\clearpage\begin{multicols*}{4}
\songtitle{Je hais ma haine}

\tag*{1988}\tag{electroswing}\tag{original-lyrics}\tag{french-lyrics}

\begin{notes}
A playful electroswing piece that layers vintage swagger over modern club energy. %
The lyrics begin with “Je hais ce que je sens pour toi,”, giving a quick sense of the song’s tone.
\end{notes}

\begin{style}
Electroswing with a high-tempo dance beat at 124 BPM in F minor. %
The track features a prominent brass section with trumpets and trombones playing syncopated staccato stabs. %
A synthesized upright bass provides a walking bassline that interacts with a four-on-the-floor kick drum and a crisp, high-frequency snare. %
A bright, processed female vocal performs repetitive melodic hooks with heavy compression and a slight slapback delay. %
Glitchy record scratches and vinyl crackle effects are layered over the percussion. %
A swing-style piano plays rhythmic chords on the off-beats. %
The arrangement uses sudden filter sweeps and momentary silences for rhythmic emphasis.
\end{style}

\begin{lyrics}

\songpart{Intro}
\annotation{four-on-the-floor kick drum, vinyl crackle, staccato brass stabs}\\
Je hais ce que je sens pour toi,\\
Et je hais ce que tu me fais sentir.\\
\annotation{record scratch}\\
Je hais t'aimer, et je me hais pour ça.\\
Je hais te haïr, et je me hais pour ça.\\
\annotation{full band enters, walking synth bass, swing piano}\\
Je hais ce que je fais pour toi\\
Et je hais ce que je ne fais pas pour toi

\songpart{Chorus}
\annotation{processed female vocals, bright brass section}\\
Je hais ce que tu me fais...\\
Et je hais toute ma haine.

\songpart{Outro}
\annotation{filter sweep on brass}\\
Je hais tant ce que je hais...\\
Que je finis en haïssant tout !

\end{lyrics}


\songtitle{Dig}

\tag*{1991}\tag{original-lyrics}\tag{electroswing}\tag{swedish-lyrics}\tag{dance}

\begin{notes}
A playful electroswing piece that layers vintage swagger over modern club energy. %
The lyrics begin with “Dig vill jag skriva om.”, giving a quick sense of the song’s tone.
\end{notes}

\begin{style}
Electroswing with a high-tempo dance beat at 124 BPM in F minor. %
The track features a prominent brass section with trumpets and trombones playing syncopated staccato stabs. %
A synthesized upright bass provides a walking bassline that interacts with a four-on-the-floor kick drum and a crisp, high-frequency snare. %
A bright, processed female vocal performs repetitive melodic hooks with heavy compression and a slight slapback delay. %
Glitchy record scratches and vinyl crackle effects are layered over the percussion. %
A swing-style piano plays rhythmic chords on the off-beats. %
The arrangement uses sudden filter sweeps and momentary silences for rhythmic emphasis.
\end{style}

\begin{lyrics}

\songpart{Intro}
\annotation{four-on-the-floor kick drum, vinyl crackle, staccato brass stabs}\\
Dig vill jag skriva om.\\
Dig har jag svrivit förr.\\
\annotation{record scratch}\\
Men dig kan jag inte tänka något om.\\
\annotation{full band enters, walking synth bass, swing piano}\\
Inte när du har blivit så länge bort.\\
Inte när ögonen dina har lämnat skenet i minet mitt.

\songpart{Chorus}
\annotation{processed female vocals, bright brass section}\\
Inte heller när ingenting är ditt, sedan du inte finns.

\songpart{Outro}
\annotation{filter sweep on brass}\\
Ögonen dina, svarta ögon.  Dem vill jag att bli mina.

\end{lyrics}


\songtitle{Si pienso en ti}

\tag*{1990}\tag{original-source}\tag{house-pop}\tag{dubstep}\tag{spanish-lyrics}

\begin{notes}
A cross-genre track that balances house pulse with dubstep intensity. %
The lyrics begin with “Si pienso en una persona,”, giving a quick sense of the song’s tone.
\end{notes}

\begin{style}
House-pop with dubstep drops, swung four-on-the-floor verses, half-time chorus lift, and Latin American female vocals; verse stays airy with clipped percussion and a rubbery bass pulse, pre-chorus strips to handclaps and filtered pads, chorus opens into stacked harmonies and a heavy wobble bass hit. %
Use playful lead doubles, glossy ad-lib tails, reverse swells before drops, and a bright, punchy mix with intimate close-mic detail., female vocals, latin, soft, dubstep, playful
\end{style}

\begin{lyrics}

\songpart{Verse 1}
Si pienso en una persona,\\
pienso en tu rostro.

Si pienso en una voz,\\
busco la tuya en mi memoria.

Si pienso en una mirada,\\
la imagino de tus ojos.

Si pienso en lo que nunca fue,\\
aparece tu sombra.

\songpart{Pre-Chorus}
Tu imagen viene,\\
sin pedir permiso.

Tus ojos, tu cabello,\\
tu boca, tu piel.

Y yo me quedo ahí,\\
sin saber qué hacer.

\songpart{Chorus}
Si pienso en ti, ahí estás tú\\
Si pienso en ti, ahí estás tú\\
Lo que más deseo, eres tú\\
Lo que más deseo, eres tú

Si pienso en ti, ahí estás tú\\
Si pienso en ti, ahí estás tú\\
No sé borrarte de mí\\
No sé borrarte de mí

\songpart{Verse 2}
Si pienso en mi dolor,\\
recuerdo tus palabras.

Aquellos saludos\\
que tú no respondiste.

Si pienso en mi vida,\\
vuelvo a encontrarte.

En unas esperanzas\\
que aún no he perdido.

En unos ojos abiertos\\
mirando a todas partes,\\
exceptuándome a mí.

\songpart{Pre-Chorus}
Y una voz dulce\\
con palabras ásperas.

Un nombre que temo\\
pronunciar despacio.

Un rostro que no puedo\\
soltar de mi pecho.

\songpart{Chorus}
Si pienso en ti, ahí estás tú\\
Si pienso en ti, ahí estás tú\\
Lo que más deseo, eres tú\\
Lo que más deseo, eres tú

Si pienso en ti, ahí estás tú\\
Si pienso en ti, ahí estás tú\\
No sé borrarte de mí\\
No sé borrarte de mí

\songpart{Bridge}
Si pienso en cualquier cosa,\\
tu figura cae primero.

Visible o escondida,\\
siempre llegas tú.

Y aunque quiera negarlo,\\
me sigues encontrando.

Tu nombre en mi garganta,\\
tu recuerdo en mi piel.

\songpart{Final Chorus}
Si pienso en ti, ahí estás tú\\
Si pienso en ti, ahí estás tú\\
Lo que más deseo, eres tú\\
Lo que más deseo, eres tú

Si pienso en ti, ahí estás tú\\
Si pienso en ti, ahí estás tú\\
En todo lo que soy\\
siempre estás tú

\end{lyrics}


\songtitle{Si pienso en ti}

\tag*{1991}
\begin{notes}
Lyrics adapted by SUNO from an original 1991 poem by Carlos Thompson.
\end{notes}
\begin{style}
Latin American female vocals; Cinematic indie-pop instrumental with a slow, swaying pulse and soft syncopated percussion; verse texture stays intimate with muted piano, brushed shakers, and airy pad haze, pre-drop strips to filtered keys and heartbeat kick, chorus opens with warm string swells and a glimmering lead motif, Breathy vocal chops, reverse swells, and sunlit chime accents, Wide, polished, photoreal mix with a golden glow, natural, soft, minimal, warm
\end{style}
\begin{lyrics}
\songpart{Verse 1}
Si pienso en una persona,\\
pienso en tu rostro.

Si pienso en una voz,\\
busco la tuya en mi memoria.

Si pienso en una mirada,\\
la imagino de tus ojos.

Si pienso en lo que nunca fue,\\
aparece tu sombra.

\songpart{Pre-Chorus}
Tu imagen viene,\\
sin pedir permiso.

Tus ojos, tu cabello,\\
tu boca, tu piel.

Y yo me quedo ahí,\\
sin saber qué hacer.

\songpart{Chorus}
Si pienso en ti, ahí estás tú\\
Si pienso en ti, ahí estás tú\\
Lo que más deseo, eres tú\\
Lo que más deseo, eres tú

Si pienso en ti, ahí estás tú\\
Si pienso en ti, ahí estás tú\\
No sé borrarte de mí\\
No sé borrarte de mí

\songpart{Verse 2}
Si pienso en mi dolor,\\
recuerdo tus palabras.

Aquellos saludos\\
que tú no respondiste.

Si pienso en mi vida,\\
vuelvo a encontrarte.

En unas esperanzas\\
que aún no he perdido.

En unos ojos abiertos\\
mirando a todas partes,\\
exceptuándome a mí.

\songpart{Pre-Chorus}
Y una voz dulce\\
con palabras ásperas.

Un nombre que temo\\
pronunciar despacio.

Un rostro que no puedo\\
soltar de mi pecho.

\songpart{Chorus}
Si pienso en ti, ahí estás tú\\
Si pienso en ti, ahí estás tú\\
Lo que más deseo, eres tú\\
Lo que más deseo, eres tú

Si pienso en ti, ahí estás tú\\
Si pienso en ti, ahí estás tú\\
No sé borrarte de mí\\
No sé borrarte de mí

\songpart{Bridge}
Si pienso en cualquier cosa,\\
tu figura cae primero.

Visible o escondida,\\
siempre llegas tú.

Y aunque quiera negarlo,\\
me sigues encontrando.

Tu nombre en mi garganta,\\
tu recuerdo en mi piel.

\songpart{Final Chorus}
Si pienso en ti, ahí estás tú\\
Si pienso en ti, ahí estás tú\\
Lo que más deseo, eres tú\\
Lo que más deseo, eres tú

Si pienso en ti, ahí estás tú\\
Si pienso en ti, ahí estás tú\\
En todo lo que soy\\
siempre estás tú
\end{lyrics}


\songtitle{Travel Notes}

\tag*{1989}\tag{house}\tag{original-source}\tag*{english-lyrics}

\begin{notes}
A cross-genre track that balances house pulse with dubstep intensity. %
The lyrics begin with “Chlewey closed the school gate”, giving a quick sense of the song’s tone.
\end{notes}

\begin{style}
House with dubstep drops, playful high-pitched female vocals, Latin-American pronunciation, four-on-the-floor pulse, rubbery bass hits, verse kept sparse and bouncy, pre-chorus filters up with risers and snapped claps, chorus opens into glittering synth stabs and a wobbling sub drop, whispered ad-libs and stacked doubles on the hook, bright punchy mix with glossy top end, playful, electronic, dubstep, latin, poetry, soft, female vocals
\end{style}

\begin{lyrics}

\songpart{Verse 1}
Chlewey closed the school gate\\
No undergrad, not yet\\
He wanted one last trip\\
One last math quest

WordPerfect screen at night\\
XT clone in the room\\
Pages full of stories\\
Then gone, too soon

\songpart{Pre-Chorus}
Free time on his hands\\
No class, no train\\
Just a new notebook\\
And a half-lit brain

First page, first breath\\
Ink on the line\\
What was lost in the past\\
Started shining again

\songpart{Chorus}
Travel notes, travel notes\\
Turn the page, let it glow\\
Travel notes, travel notes\\
Where the old ghosts go (go, go)

China on his tongue\\
Stockholm in his bones\\
Travel notes, travel notes\\
He wrote it all alone (alone)

\songpart{Verse 2}
July took him east\\
For the China sky\\
Math Olympiads trip\\
With a writer’s eye

Hèla smiled at him\\
Tunis in her name\\
For one bright moment\\
He forgot the same

Adriana’s shadow\\
Slipped out of sight\\
Then the memory came back\\
By the end of the night

\songpart{Pre-Chorus}
Then north came the cold\\
Blue water, white streets\\
Two and a half years\\
With a new heartbeat

Scripts turned to letters\\
Pages learned to stay\\
From handwritten lines\\
To an electronic way

\songpart{Chorus}
Travel notes, travel notes\\
Turn the page, let it glow\\
Travel notes, travel notes\\
Where the old ghosts go (go, go)

China on his tongue\\
Stockholm in his bones\\
Travel notes, travel notes\\
He wrote it all alone (alone)

\songpart{Bridge}
Dates got worn away\\
Revision by revision\\
Some years blurred together\\
In the half-light vision

Still the words kept moving\\
Still the hand kept true\\
Lost in the timeline\\
But alive in the blue

\songpart{Chorus}
Travel notes, travel notes\\
Turn the page, let it glow\\
Travel notes, travel notes\\
Where the old ghosts go (go, go)

China on his tongue\\
Stockholm in his bones\\
Travel notes, travel notes\\
He wrote it all alone (alone)

\songpart{Outro}
First year, last page\\
Then the rest unfolds\\
In a notebook’s shadow\\
The memory holds (holds)

\end{lyrics}


\songtitle{Adriana Tal Vez}
\tag*{1990}\tag{bhangra-hip-hop}\tag{spanish-lyrics}
\begin{notes}
Lyrics by SUNO based on a 1990 text by Carlos Thompson.
\end{notes}

\begin{style}
Emotional bhangra-hip-hop with a mid-tempo dhol groove, syncopated claps, and a rolling bassline; verse rides sparse with tight rap phrasing and hand drum accents, pre-chorus opens with lifted backing chants and snare pushes, chorus hits bigger with dhol hits, harmonized lead, and call-response repeats. %
Vocal is close-mic and aching in verses, doubled on hooks, with short delay throws on the name. %
Ear candy: reversed cymbal swells into the chorus, flute fragments between lines, and a brief breakdown before the final lift. %
Bright, punchy, and heartfelt., rap, emotional, german hip hop, bhangra
\end{style}

\begin{lyrics}
\songpart{Verse 1}
Adriana, yo no sé\\
si esto tuvo un lugar.\\
Tal vez nunca vi la puerta\\
o no la supe tocar.

Yo quería darte un hola,\\
una risa al pasar.\\
Y en la cara de los días\\
tu silencio fue mi mar.

\songpart{Pre-Chorus}
Si yo hubiera pedido\\
una amistad de verdad,\\
si tú hubieras querido\\
cruzar conmigo esta ciudad...

Yo habría guardado\\
tu nombre al despertar,\\
pero me quedé tan cerca\\
que no me pude lanzar.

\songpart{Chorus}
Adriana, tal vez fue así\\
(tal vez fue así)\\
Adriana, te quise a ti\\
(te quise a ti)\\
Si no hubo un adiós,\\
fue por mí, fue por mí.\\
Adriana, tal vez fue así.

\songpart{Verse 2}
No te conozco del todo,\\
y aun así te sentí.\\
Como un mensaje borrado\\
que quedó dentro de mí.

Me pregunto si era amor\\
o te deseo nada más.\\
No quiero palabras vacías\\
si me escuchas de verdad.

Y si es culpa de mis miedos\\
que no llegues a leer,\\
yo cargo con ese ruido\\
que no supe detener.

\songpart{Pre-Chorus}
Quiero saber qué piensas,\\
quiero tu respuesta aquí.\\
Dime por qué no fui capaz\\
de despedirme de ti.

Quizás nunca te salude,\\
quizás ése fue el final,\\
pero en mi pecho tu sombra\\
sigue dando vueltas más.

\songpart{Chorus}
Adriana, tal vez fue así\\
(tal vez fue así)\\
Adriana, te quise a ti\\
(te quise a ti)\\
Si no hubo un adiós,\\
fue por mí, fue por mí.\\
Adriana, tal vez fue así.

\songpart{Bridge}
No te pido que regreses,\\
solo quiero comprender,\\
si lo nuestro fue un camino\\
o una puerta sin abrir.

Si me miras desde lejos,\\
si me piensas una vez,\\
déjame caer en calma\\
con lo que no puede ser.

\songpart{Final Chorus}
Adriana, tal vez fue así\\
(tal vez fue así)\\
Adriana, te quise a ti\\
(te quise a ti)\\
Si no hubo un adiós,\\
fue por mí, fue por mí.\\
Adriana, tal vez fue así.

Adriana, dime por qué\\
(dime por qué)\\
yo me quedo sin saber.\\
Chao, y piensa en mí.

\end{lyrics}


\songtitle{Si pienso}

\tag*{1990}\tag{trance}\tag{original-work}\tag{spanish-lyrics}

\begin{notes}
Original work by Carlos Thompson.\\
In Spanish. %
Trance meditation on memory and obsessive thought patterns.
\end{notes}

\begin{style}
trance, latin pop, soprano lead vocals, latin-american pronunciation, bassline melody lead, airy vocal phrasing, delay throws, harmonized chorus, sidechain pumping, glossy synth pads, arpeggiated trance plucks, four-on-the-floor kick, gated reverb snare, 128 BPM, bittersweet longing, neon melancholy
\end{style}

\begin{lyrics}
Si pienso en una persona; pienso en tu rostro...\\
Si pienso en una voz; busco la tuya en mi memoria...\\
Si pienso en una mirada; la imagino de tus ojos...\\
Si pienso en lo que nunca fue; tu imagen aparece, tus ojos, tu cabello, tu boca, tu piel...\\
Si pienso en qué deseo; primero estás tú, porque tú eres lo que más deseo...\\
Si pienso en mi dolor; recuerdo aquellas palabras que tú me dijiste, aquellos saludos que no respondiste...\\
Si pienso en mi vida; necesariamente tengo que pensar en ti, en las esperanzas que aún no he perdido, en las imágenes que de ti aún guardo, en unos ojos abiertos mirando a todas partes exceptuándome a mi, en una voz dulce con palabras ásperas, en un nombre que temo pronunciar, un rostro que no puedo olvidar, tantas cosas que desearía no recordar...\\
Pero no puedo...\\
Si pienso en cualquier cosa, ahí, visible o escondida en algún lado, ahí estás tú.
\end{lyrics}


\songtitle{Iros de risk}

\tag*{1989}\tag{trance}\tag{original-work}\tag{spanish-lyrics}

\begin{notes}
Original work.\\
Experimental trance composition with fragmented narrative structure.
\end{notes}

\begin{style}
trance, latin pop, soprano lead vocals, latin-american pronunciation, bassline melody lead, airy vocal phrasing, delay throws, harmonized chorus, sidechain pumping, glossy synth pads, arpeggiated trance plucks, four-on-the-floor kick, gated reverb snare, 128 BPM, bittersweet longing, neon melancholy
\end{style}

\begin{lyrics}
Iros de risk, nefa he sabido como entrar, solo vi y temí, e iros de risk.  Decidí y ahí se va.  Ce ça, aill be, je bi.\\
O.K., she's 'ere and aim 'ere.\\
Pero no será\\
no será en aim sure
'n aim' sure\\
sure\\
shit\\
shes 'ere\\
en' aim 'ere\\
ce nothink\\
nada!
\end{lyrics}


\songtitle{Je hais}

\tag*{1988}\tag{original-lyrics}\tag{hard-rock}

\subsection{Notes}

\subsection{Style}
Heavy metal with electronic instrumentals. %
deep bass voice.

\begin{lyrics}
Je hais ce que je sens pour toi,\\
et je hais ce que tu me fais sentir.\\
Je hais t'aimer, et je me hais pour ça.\\
Je hais te haïr, et je me hais pour ça.\\
Je hais ce que je fais pour toi,\\
et je hais ce que je ne fais pas pour toi.\\
Je hais ce que tu me fais...\\
Et je hais tout mon haine.\\
Je hais tant ce que je hais ...\\
que je finisen haisant tout.
\end{lyrics}


\songtitle{Je hais}

\tag*{1988}\tag{hard-rock}\tag{original-lyrics}\tag{french-lyrics}\tag{je-hais-1988}

\begin{notes}
See \emph{\ref{sec:je-hais}} on page \pageref{sec:je-hais}.
\end{notes}

\begin{style}
hard rock, male voice, powerful, industrial
\end{style}
\begin{lyrics}
Je hais ce que je sens pour toi,\\
et je hais ce que tu me fais sentir.\\
Je hais t'aimer, et je me hais pour ça.\\
Je hais te haïr, et je me hais pour ça.\\
Je hais ce que je fais pour toi,\\
et je hais ce que je ne fais pas pour toi.\\
Je hais ce que tu me fais...\\
Et je hais tout mon haine.\\
Je hais tant ce que je hais ...\\
que je finisen haisant tout.
\end{lyrics}


\songtitle{La haine que je te dois}

\tag*{1988}\tag{pop-rock}\tag{original-lyrics}\tag{french-lyrics}\tag{je-hais-1988}

\begin{notes}
See \emph{\ref{sec:je-hais}} on page \pageref{sec:je-hais}.
\end{notes}

\begin{style}
pop rock, 1970s
\end{style}
\begin{lyrics}
Je hais ce que je sens pour toi,\\
et je hais ce que tu me fais sentir.\\
Je hais t'aimer, et je me hais pour ça.\\
Je hais te haïr, et je me hais pour ça.\\
Je hais ce que je fais pour toi,\\
et je hais ce que je ne fais pas pour toi.\\
Je hais ce que tu me fais...\\
Et je hais tout mon haine.\\
Je hais tant ce que je hais ...\\
que je finisen haisant tout.
\end{lyrics}


\songtitle{Al despertar}

\tag*{1991}\tag{original-source}\tag{despertar-1991}\tag{ai-adapted}\tag{electronic-dubstep}

\begin{notes}
Based on the original \emph{Despertar} poem by Carlos Thompson, adaptep to song structure in 2026 with help of Gemini.
\end{notes}
\begin{style}
Style: electronic, dubstep built-ups, house, 90 bpm\\
synthesizers, bass, guitars\\
male baritone vocals
\end{style}
\begin{lyrics}
\songpart{Intro - spoken softly}
He creído verte en los sueños perdidos en mi memoria al despertar.\\
He creído verte en rostros ajenos, extraños, perdidos entre los demás.

\songpart{Verse 1 - half spoken, half sung}
He creído olvidarte y ver lo que vi en ti en otros lugares donde tu huella no ha podido estar.\\
Donde el aire no ha podido ser tocado por tu aliento.\\
Donde tu nombre tenga otro sentido y otro color.

\annotation{Build}\\
He creído por momentos escapar.\\
Qué error.

\annotation{Sung}\\
Estás tan en mí que en todo recuerdo está tu nombre o la imagen perdida de tus ojos.

\songpart{Chorus - emotional, melodic}
Y quiero seguir buscándote,\\
y quiero seguir olvidándote,\\
y quiero seguir amándote,\\
y quiero seguir odiándote.

Quiero seguir siendo quien fui cuando tú existías.\\
Quiero seguir pensando cómo verte sin ser visto, cómo ser visto sin volver a verte.

\songpart{Verse 2 - restrained, intimate}
Tan en mí que en todo rostro sólo veo lo que vi en ti.\\
Para dejar tu huella grabada, para buscar el toque de tu aliento.

Te he traído en mí a donde el verde se acaba por segunda vez (tercera, diré.)\\
A donde la nostalgia crece y no me permite olvidarte.

\annotation{Build}\\
No a ti, en realidad.\\
Sino a la imagen de los sueños perdidos en mi memoria al despertar.

Es difícil convencerme de que no estás aquí, por más que lo sé.

\songpart{Chorus - stronger}
Y quiero seguir buscándote,\\
y quiero seguir olvidándote,\\
y quiero seguir amándote,\\
y quiero seguir odiándote.

Quiero regresar al tiempo en que no eras recuerdo, sino una realidad.\\
Por más que esa realidad lejos estuvo de ser mía.

\songpart{Bridge - spoken over ambient instruments}
Y cada escaso segundo que pasa estoy más lejos de ese tiempo.\\
Y cada escaso segundo que pasa estás más lejos de existir en mí.

Sólo queda buscarte en las ramas muertas, en árboles vacíos deshojados por el otoño.\\
Ese otoño del que has hablado y del cual no quiero creer que conozca a tus ojos.

\annotation{Breakdown - whispered}\\
Y son otros rostros y otros muros los que guardan tu imagen.\\
Otras imágenes las que adquieren tu nombre.

Eres un recuerdo.\\
Un fantasma de algo que no existe más.

Más que donde no quiero verte.\\
Más que donde quiero verte.

\annotation{Final Chorus - exhausted, fading}\\
Y quiero seguir buscándote…\\
aunque no estés aquí.

Y quiero seguir olvidándote…\\
por más que no pueda.

Quiero seguir siendo quien fui cuando tú existías.

\songpart{Outro - whispered}
Muros amarillos como las hojas que no tuvo este otoño.

Y los sueños terminan perdiéndose en mi memoria al despertar.

Y aún no sé dónde estás.\\
No sé si existes.\\
No sé si has existido alguna vez.

No sé si tu nombre y tu historia no serán un invento de mi imaginación.\\
O si no será el otoño el que trae tu imagen.

\annotation{Fade out}
\end{lyrics}


\songtitle{Message trop tard}

\tag*{1988}\tag{musical-theater}\tag{original-source}\tag{french-lyrics}\tag{je-hais-1988}

\begin{notes}
Lyrics by Gemini on a 1988 poem by Carlos Thompson and context by the author, high interaction.
(See \emph{Je hais}).
\end{notes}
\begin{style}
French musical theater ballad. %
Tenor male vocals with a wide vibrato and dramatic delivery. %
Orchestral arrangement featuring a grand piano playing arpeggiated chords and a string section providing sustained legato pads and melodic counterpoint. %
The arrangement builds in intensity with a staccato string ostinato and orchestral percussion including timpani rolls and cymbal swells. %
Key of C minor. %
Tempo is 72 BPM. %
Time signature is 4/4. %
The vocal performance uses dynamic shifts from soft, breathy delivery to powerful, belted chest voice. %
High-register violin lines double the vocal melody during the climax.
\end{style}
\begin{lyrics}
\annotation{Instrumental Intro}\\
\annotation{Melancholic acoustic guitar strumming with a soft, lonely piano melody}

\songpart{Verse 1}
Un regard volé près des casiers, le cœur qui déraille\\
Mes mots s'emmêlent, mes mains tremblent, c'est la pagaille\\
Je voulais juste te plaire, être celui qu'on remarque\\
Mais ma maladresse dessine les contours d'une marque\\
Je t'idéalise dans le noir, j'invente notre histoire

\songpart{Pre-Chorus}
\annotation{Bass guitar enters}\annotation{Whisper}\\
Un message envoyé trop tard, une attente infinie\\
Je perds le contrôle de mes nuits, de ma vie

\songpart{Chorus}
\annotation{grand piano arpeggios, legato string pads}\\
Je hais ce que je sens pour toi,\\
Et je hais ce que tu me fais sentir.\\
Je hais t'aimer, et je me hais pour ça.\\
Je hais te haïr, et je me hais pour ça.

\songpart{Verse 2}
\annotation{Music drops back to piano and bass}\\
Puis est venu le mot de trop, la confession de trop\\
Tu as reculé d'un pas, tu as dit "c'est trop"\\
Tu as tracé une ligne, posé tes limites sacrées\\
Et mon ego s'effondre, ma fierté est brisée\\
Ton refus légitime devient mon obsession

\songpart{Chorus}
\annotation{grand piano arpeggios, legato string pads}\\
Je hais ce que je sens pour toi\\
Et je hais ce que tu me fais sentir\\
Je hais t'aimer et je me hais pour ça\\
Je hais te haïr et je me hais pour ça

\songpart{Bridge}
\annotation{staccato string ostinato enters}\\
Je hais ce que je fais pour toi\\
Et je hais ce que je ne fais pas pour toi\\
\annotation{timpani roll, cymbal swell}\\
Et je hais ce que je ne fais pas pour toi.\\
Je hais ce que tu me fais...\\
Et je hais toute ma haine.

\annotation{Guitar Solo}\\
\annotation{Dramatic, crying guitar solo building up to a chaotic climax}

\songpart{Bridge Climax}
\annotation{Vocals become strained, desperate, almost screaming}\\
Je vois mon ombre t'étouffer, mes colères te traquer\\
Ce n'est plus de l'amour, c'est du poison fabriqué !\\
Je hais tant ce que je hais...\\
Que je finis en haïssant tout !

\songpart{Outro}
\annotation{All instruments abruptly cut out except for a naked, fading piano}\\
Je lâche prise.\\
Je m'efface.\\
Mon premier amour...\\
Ma première prison.\\
\annotation{Fade out to silence}
\end{lyrics}


\songtitle{Svarta Ögon}

\subsection{Notes}
Original text by Carlos Thompson ca. %
1991.

\subsection{Style}
dubstep, deep baritone vocals

\begin{lyrics}ögona dina, svarta ögon.\\
Dom vill jà å bli mina.
\end{lyrics}


\songtitle{Svarta Ögon Skena}

\tag*{1991}\tag{electronic-dubstep}\tag{original-lyrics}\tag{despertar-1991}\tag{swedish-lyrics}

\begin{notes}
Poem written by Carlos Thompson in 1991.
\end{notes}
\begin{style}
dubstep, tenor vocals, trance, industrial
\end{style}
\begin{lyrics}
Dej vill jà skriva om.\\
Dej har jà svrivì förr.\\
Men dej kan jà inte tänka nåt om.\\
Inte när'u har blivì så länge bort.\\
Inte när ögona dina har lämnà skenè i minè mitt.\\
Inte heller när ingè é ditt, sen du inte finns.

ögonen dina, svarta ögon. Dem vill jag att bli mina.

Dig vill jag skriva om. Dig har jag skrivit förr. Men dig kan jag inte tänka något om. %
Inte när du har blivit så länge bort. Inte när ögonen dina har lämnat skenet i minet mitt. %
Inte heller när inget är ditt, sedan du inte finns.
\end{lyrics}


\songtitle{Bilden}

\tag*{1991}\tag{original-lyrics}\tag{electronic-dubstep}

\subsection{Notes}
Free form poem writen by Carlos Thompson in 1991.


\subsection{Style}
electronic house background, conversational voice

\begin{lyrics}
\songpart{Intro (instrumental)}

\annotation{Conversational}\\
\annotation{Tenor vocals}\\
Kan inte stå ut länge.  Kan inte.\\
Hej!\\
Känner'u mej?\\
Kommer'u ihåg den där killen som nån gång sa dej jà älskar'ej?\\
Den som nån gång sa dej jà vill dej va' min?\\
Har'u inte glömt'en?\\
För dé har jà inte.\\
Jà har inte kunnà glömma tjejen, jà har sagt å älska'na.\\
Jà kommer fortfarande ihåg'na.\\
Hennes bild ligger kvar innerst hos mej.\\
Bilden jà har försökt finna hos så många andra.\\
Bilden jà har trott finna hos så många andra.\\
Bilden som inte finns hos nån anann.\\
Bilden som inte é en bild...\\
bara ett namn...  ditt.\\
amnè...\\
bara minnè.\\
Minnè om en tjej, jà har sagt å älska'na.\\
Minnè om hennes bild.\\
Minnè om minnè.\\
Bilden om ett minn, jà har redan glömt.\\
Bilden om dej.

\songpart{Outro (instrumental)}

\end{lyrics}


\songtitle{Vakning}

\tag*{1991}\tag{dubstep}\tag{original-lyrics}\tag{despertar-1991}\tag{multilingual}

\begin{notes}
Written in 1991 by Carlos Thompson. Adapted to song structure and production annotations in 2026 by the author.
\end{notes}
\begin{style}
electronic, energetic, synthesizers, percussions, dubstep, multilingual, smooth voice, male vocals
\end{style}
\begin{lyrics}
\songpart{Intro}
\annotation{Instrumental}

\annotation{Spoken word}\\
\annotation{baritone vocals, whisper, fast, smooth, Swedish}\\
Och drömarna går till slut vilse in i mit minne vid vakningen.\\
Och ännu kan jag inte var du håller till, kan varken om du är till eller om ditt namn och din historia inte är en uppfinning från min fantasi och om det är hösten som bär din bild.

\annotation{Instrumental bridge, drums, cymbals}

\songpart{Verse 1}
\annotation{baritone vocals, Swedish}\\
Dina förbanade svarta ögon!\\
Skiten jag inte vill tänka till.\\
\annotation{Spanish}
Tú eres única, el problema soy yo quien quiere olvidarte buscándote en otros rostros. 
\annotation{Swedish}\\
Du är unik.  Problemet är jag, som vill glöma bort dig genom att söka dig hos andra ansikten.

\annotation{Instrumental bridge}

\songpart{Verse 2}
\annotation{baritone vocals, Spanish}\\
Y trato de volver a repetir:\\
No existes!\\
\annotation{English}
But I do not like to hate you.  You know.  I really hate you, and I still like you, but it is the felling I hate.

\annotation{Instrumental bridge}

\songpart{Outro}
\annotation{baritone vocals, Swedish, shouting}\\
Och Adriana är inte till.\\
\annotation{whisper, instruments off}\\
Punkt.
\end{lyrics}


\songtitle{Jag avskyr}
\tag*{1991}\tag{dubstep}\tag{original-lyrics}\tag{swedish-lyrics}\tag{je-hais-1988}

\begin{notes}
Translation of \emph{Odio}/\emph{Je hais} into Swedish by Carlos Thompson.

See \emph{\ref{sec:je-hais}} on page \pageref{sec:je-hais}.
\end{notes}
\begin{style}
Dubstep epic with half-time wobble drops, syncopated sub punches, and tight snare snaps; intro starts with distant bass growls and clipped breathy textures, pre-drop builds with rising filters and metallic ticks, chorus hits with huge detuned synth wobble and marching toms. %
Glitched vocal chops, reverse swells, dust-sparkle risers, cinematic low-end, bright and brutal mix., rich, complex, light, deep, dramatic, soft, natural, dubstep, female voice, intimate, almost conversational
\end{style}
\begin{lyrics}Jag avskyr det, jag känner om dej,\\
och avskyr det, du får mej att känna.\\
Jag avskyr det att älska dej, och avskyr mej därför.\\
Jag avskyr det att avsky dej, och avskyr mej därför.\\
Jag avskyr det, jag gör för dej,\\
och avskyr det, jag inte gör för dej.\\
Jag avskyr det, du gör åt mej,\\
och avskyr hela min avsky.\\
Jag avskyr det, jag avskyr, så\\
att jag slutar och avskyr allt.
\end{lyrics}

\songtitle{Vaina que crece}

\tag*{1992}\tag{art-pop}\tag{graeville}\tag{original-work}\tag{spanish-lyrics}

\begin{notes}
Original work.\\
Theme: Mental health, accountability, and the need for structured support systems.
\end{notes}

\begin{style}
Art pop electrónico con pulso medio, groove sincopado y un crecimiento gradual de tensión; verso en capas mínimas con bajo filtrado y piano seco, pre-coro abre espacio con armonías aireadas y percusión más marcada, coro entra ancho con sintes brillantes y baterías firmes. %
Voz íntima al frente con dobles suaves, coros apilados en el gancho y pequeñas colas de delay en palabras clave. %
Riser reverso antes del coro, chasquidos digitales entre frases y un estallido de textura en el final. %
Mezcla limpia, grande y luminosa.
\end{style}

\begin{lyrics}
\songpart{Verse 1}
Una pequeña vaina comienza\\
y poco a poco adquiere forma\\
se elabora\\
en silencio

Llega la distancia\\
el cambio\\
el recuerdo\\
de un obscuro viaje al Despertar

La inspiración que cede al Græville\\
y se pierde\\
en un par de experimentos estúpidos\\
hasta que duele más la espera

\songpart{Pre-Chorus}
Y vuelve a moverse\\
bajo la piel\\
vuelve a romperse\\
vuelve a nacer

La crisis llama\\
sin pedir perdón\\
y lo que era polvo\\
toma voz

\songpart{Chorus}
Una nueva historia\\
un nuevo sueño\\
una nueva historia\\
crece en el pecho

Que crece\\
que engrandece\\
que rompe el miedo\\
y vuelve a empezar

Una nueva historia\\
un nuevo sueño\\
una nueva historia\\
viene a estallar

\songpart{Verse 2}
Perdida en un abismo electrónico\\
la tarde se queda sin nombre\\
y todo gira\\
sin dirección

Entre chispas\\
y gestos torpes\\
se arma el vacío\\
se parte el sol

Y de repente\\
sin pedir permiso\\
la idea respira\\
se levanta sola

\songpart{Pre-Chorus}
Se junta el pulso\\
con el dolor\\
se enciende el borde\\
del corazón

La crisis llama\\
sin pedir perdón\\
y lo que era ruina\\
toma voz

\songpart{Chorus}
Una nueva historia\\
un nuevo sueño\\
una nueva historia\\
crece en el pecho

Que crece\\
que engrandece\\
que rompe el miedo\\
y vuelve a empezar

Una nueva historia\\
un nuevo sueño\\
una nueva historia\\
viene a estallar

\songpart{Bridge}
Y llega la distancia\\
ya no pesa igual\\
el recuerdo no manda\\
puede soltar

De un oscuro viaje al Despertar\\
nace la ruta\\
nace el umbral

\songpart{Final Chorus}
Una nueva historia\\
un nuevo sueño\\
una nueva historia\\
crece en el pecho

Que crece\\
que engrandece\\
que rompe el miedo\\
y vuelve a empezar

Una nueva historia\\
un nuevo sueño\\
una nueva historia\\
viene a estallar
\end{lyrics}


\end{multicols*}

\chapter{Earlier Thoughts II (1995–1998)}
\albumcover{covers/internal/earlier-thoughts-ii}
The university years: blank pages that stayed blank, ink-stained tempers, a mind drifting into its own clouds mid-sentence. \emph{Siento}, \emph{Rabia en la boca} and \emph{Esfera entintada} catalogue a temper still learning its own edges, while \emph{En la nube} and \emph{Circulo convexo} watch the same mind float somewhere else entirely. Three languages, one notebook, one restless kid.
\clearpage\begin{multicols*}{4}
\songtitle{Siento}

\tag*{1996}\tag{original-lyrics}\tag{glam-metal}

\begin{notes}
A distinctive arrangement that supports the song's storytelling and emotional arc. %
The lyrics begin with “Siento haberte amado de una forma tan profunda”, giving a quick sense of the song’s tone.
\end{notes}

\begin{style}
glam metal
\end{style}

\begin{lyrics}

Siento haberte amado de una forma tan profunda\\
que hoy te odio.\\
Siento dolor y siento frío\\
Siento el vacío que me explota por dentro\\
y da rabia.\\
Siento haberte querido\\
y haber fallado en decirlo\\
y siento grandemente que te hayas ido\\
sin palabras de despedida\\
sin una maldita forma de hacerme saber\\
si algo de lo que siento has sentido\\
si en algo\\
alguna vez\\
me has querido.

\end{lyrics}


\songtitle{No sé, pero vuelve}

\tag{original-source}\tag{house}\tag{dubstep}\tag{spanish-lyrics}\tag{female-vocals}

\begin{notes}
A cross-genre track that balances house pulse with dubstep intensity. %
The lyrics begin with “¿Cómo me siento contigo?”, giving a quick sense of the song’s tone.
\end{notes}

\begin{style}
House with dubstep bass hits and a playful bounce, mid-tempo four-on-the-floor verses; pre-chorus opens into filtered synth lifts and vocal stacks; chorus drops with wobbling bass, clap snaps, and chantable hooks. %
Soft peachy female vocals, lightly doubled in verses, wide harmonies on the hook, short delay throws on key phrases, glossy risers and a reverse swell into each drop, bright and polished mix with a club-ready low end., female vocals, soft, dubstep, playful
\end{style}

\begin{lyrics}

\songpart{Verse 1}
¿Cómo me siento contigo?\\
Si ni yo lo tengo claro\\
Quiero sentarme a tu lado\\
Y decirlo sin disparos

Que yo misma me enredo\\
Cuando intento ser sincera\\
Te miro, me deshago\\
Y la boca se me queda

\songpart{Pre-Chorus}
Pero te pienso, te pienso\\
Aunque me cueste admitirlo\\
Si digo poco, me pierdo\\
Si digo todo, me caigo contigo

Y aunque no sepa ordenarlo\\
Sé que no quiero soltarte\\
Hay algo tuyo en mis manos\\
Que me hace temblar al tocarte

\songpart{Chorus}
La necesito\\
Más de lo que puedo decir\\
La necesito\\
Y me da miedo descubrir\\
Si eres amiga\\
O si te tengo que soltar\\
La necesito\\
Y no sé cómo respirar

(La necesito)\\
(La necesito)

\songpart{Verse 2}
Habrá un día, tal vez\\
Que te pueda llamar amiga\\
O convencerme por fin\\
De bajarte de esta herida

Pero hoy no me sale\\
Ponerte en una respuesta\\
Hoy te pienso demasiado\\
Y me pierdo en la pregunta abierta

\songpart{Pre-Chorus}
Y si te hablo, me tiemblo\\
Y si callo, me persigues\\
Yo no sé qué estoy buscando\\
Pero tú ya lo intuyes

Déjame decirlo lento\\
Sin prometerte de más\\
Que entre lo que siento y lo que temo\\
Ya no me queda paz

\songpart{Chorus}
La necesito\\
Más de lo que puedo decir\\
La necesito\\
Y me da miedo descubrir\\
Si eres amiga\\
O si te tengo que soltar\\
La necesito\\
Y no sé cómo respirar

(La necesito)\\
(La necesito)

\songpart{Bridge}
No quiero mentirme\\
Ni pedirte un papel\\
Solo quiero verte\\
Y no romperme otra vez

Si eres puerto o tormenta\\
Si te quedas o te vas\\
Yo no tengo las respuestas\\
Pero te quiero mirar

\songpart{Final Chorus}
La necesito\\
Más de lo que puedo decir\\
La necesito\\
Y me da miedo descubrir\\
Si eres amiga\\
O si te tengo que soltar\\
La necesito\\
Y no sé cómo respirar

(La necesito)\\
(La necesito)

\end{lyrics}


\songtitle{Líneas en blanco}

\tag{original-source}\tag{house}\tag{dubstep}\tag{spanish-lyrics}\tag{female-vocals}\tag{salsa}

\begin{notes}
A cross-genre track that balances house pulse with dubstep intensity. %
The lyrics begin with “Y me encuentro otra vez”, giving a quick sense of the song’s tone.
\end{notes}

\begin{style}
House with dubstep drops and playful female vocals, four-on-the-floor pulse and syncopated bass wobble; verse stays airy with sparse kick, plucked keys, and filtered handclaps, pre-chorus opens with rising synth sweeps and stacked breaths, chorus hits wide with sub drops and chopped vocal hooks. %
Soft close-mic lead, doubled on the hook, tiny ad-libs and echo throws, with reversed swells and glittery bell flickers between phrases. %
Bright, glossy, and buoyant mix with a punchy low end., soft, playful, dubstep, female vocals, salsa
\end{style}

\begin{lyrics}

\songpart{Verse 1}
Y me encuentro otra vez\\
ante líneas en blanco\\
otra ocasión me llama\\
y yo me acerco despacio

Son palabras de relleno\\
mientras limpio mi mente\\
y ordeno ese huracán\\
que me baila de frente

De suscesos pasados\\
hay mariposas en guerra\\
revolotean en mi cabeza\\
y no les cierro la puerta

La belleza de unos ojos verdes\\
que me acompañaban\\
ojos ajenos, tan cercanos\\
y tan lejos al alba

\songpart{Pre-Chorus}
Y algo tiembla aquí dentro\\
quiere salir, quiere hablar\\
pero el miedo le pone\\
su cerrojo lunar

Una luna sonriente\\
me guiña desde Orión\\
y yo sigo juntando\\
lo que rompe mi voz

\songpart{Chorus}
Líneas en blanco\\
líneas en blanco\\
yo me suelto despacio\\
(yo me suelto despacio)

Líneas en blanco\\
líneas en blanco\\
lo que callo se vuelve canto\\
(lo que callo se vuelve canto)

\songpart{Verse 2}
Una mariposa de Orión\\
cruza lento mis cielos\\
y me deja en la piel\\
un temblor tan sincero

Es el sentimiento de abandono\\
o la salsa que no sé bailar\\
tantas cosas enredadas\\
pidiendo forma al pasar

Se agolpan y se juntan\\
como vidrio y algodón\\
quieren gritarle al destino\\
quieren nombre y dirección

Pero el filtro de mis temores\\
les baja la intención\\
y sólo quedan destellos\\
sobre mi pluma en expansión

\songpart{Pre-Chorus}
Y algo tiembla aquí dentro\\
quiere salir, quiere hablar\\
pero el miedo le pone\\
su cerrojo lunar

Una luna sonriente\\
me guiña desde Orión\\
y yo sigo juntando\\
lo que rompe mi voz

\songpart{Chorus}
Líneas en blanco\\
líneas en blanco\\
yo me suelto despacio\\
(yo me suelto despacio)

Líneas en blanco\\
líneas en blanco\\
lo que callo se vuelve canto\\
(lo que callo se vuelve canto)

\songpart{Bridge}
Déjame caer\\
sin pedir perdón\\
si me tiembla la mano\\
también tiembla el corazón

Y si hoy no digo todo\\
me acerco un poco más\\
porque hasta el polvo en silencio\\
aprende a brillar

\songpart{Final Chorus}
Líneas en blanco\\
líneas en blanco\\
yo me suelto despacio\\
(yo me suelto despacio)

Líneas en blanco\\
líneas en blanco\\
lo que callo se vuelve canto\\
(lo que callo se vuelve canto)

\end{lyrics}


\songtitle{Rabia en la boca}

\tag{original-source}\tag{house}\tag{dubstep}\tag{spanish-lyrics}\tag{female-vocals}

\begin{notes}
A cross-genre track that balances house pulse with dubstep intensity. %
The lyrics begin with “No es sabor amargo en la boca”, giving a quick sense of the song’s tone.
\end{notes}

\begin{style}
House-dubstep hybrid with bright four-on-the-floor verses and a half-time bass drop in the chorus; playful peachy female vocals, clipped rhythmic phrasing, stacked doubles on the hook, whispered ad-libs in the pre-chorus, riser sweeps into each drop, reverse cymbal tails and glossy sidechain pump, polished and punchy with a wide club-ready mix, soft, playful, female vocals, dubstep
\end{style}

\begin{lyrics}

\songpart{Verse 1}
No es sabor amargo en la boca\\
ni dolor que me parta el pecho\\
es un nudo de rabia\\
desconsuelo\\
y un “qué pasó” por dentro

Miro atrás\\
y veo lo que pudo ser\\
lo que el tiempo fue borrando\\
sin pedir permiso\\
sin volver

Y me quedo ahí\\
con la mano temblando\\
viendo nombres en el polvo\\
de todo lo que nunca fue

\songpart{Pre-Chorus}
Y le echo la culpa al otro\\
para no mirarme de frente\\
y me culpo a mí\\
por la torpeza\\
por elegir mal otra vez

Se me sube la rabia\\
se me rompe la razón\\
y cuando por fin depende de mí\\
me traiciona el corazón

\songpart{Chorus}
Da rabia\\
da rabia\\
me muerde la voz

Da rabia\\
da rabia\\
lo que no pasó

Y por orgullo tonto\\
cierro la puerta hoy\\
(da rabia)\\
pero por dentro sé\\
que dolió peor

\songpart{Verse 2}
Me acuerdo de esas horas\\
de las cosas que dejé\\
de las palabras viejas\\
que juraban volver

Todo aquello que mató el tiempo\\
sin tocarlo con la mano\\
todo aquello que guardé\\
como un vidrio en el costado

Y ahora frente a mí\\
cuando toca decidir\\
sale ese orgullo idiota\\
que me hace mentir

\songpart{Pre-Chorus}
Y le echo la culpa al otro\\
para no mirarme de frente\\
y me culpo a mí\\
por la torpeza\\
por elegir mal otra vez

Se me sube la rabia\\
se me rompe la razón\\
y cuando por fin depende de mí\\
me traiciona el corazón

\songpart{Chorus}
Da rabia\\
da rabia\\
me muerde la voz

Da rabia\\
da rabia\\
lo que no pasó

Y por orgullo tonto\\
cierro la puerta hoy\\
(da rabia)\\
pero por dentro sé\\
que dolió peor

\songpart{Bridge}
Ya no quiero esconderme\\
detrás de mi versión\\
ni ponerle al otro el peso\\
de mi propia decisión

Si el pasado se me rompe\\
que se rompa de una vez\\
yo también fui quien soltó todo\\
yo también dije “después”

\songpart{Chorus}
Da rabia\\
da rabia\\
me muerde la voz

Da rabia\\
da rabia\\
lo que no pasó

Y por orgullo tonto\\
abro la puerta hoy\\
(da rabia)\\
aunque me tiemble el pulso\\
ya me elegí yo

\end{lyrics}


\songtitle{Esfera entintada}

\tag*{1996}\tag{house}\tag{original-source}\tag{spanish-lyrics}

\begin{notes}
A cross-genre track that balances house pulse with dubstep intensity. %
The lyrics begin with “Palabras en hemorragia”, giving a quick sense of the song’s tone.
\end{notes}

\begin{style}
House with dubstep drops, playful swung groove and half-time bass hits; verse stays intimate and minimal with clipped synth plucks, pre-chorus opens into rising chords and filtered tension, chorus lands with a wide sub drop and bright vocal doubles, bridge strips to pads and heartbeat kick before the final lift. %
Soft female vocals with Latin-American pronunciation, airy stacks on hooks, delay throws on key words, reversed swells and glittery risers between lines, glossy and dark-bright mix., dubstep, female vocals, soft, playful, latin
\end{style}

\begin{lyrics}

\songpart{Verse 1}
Palabras en hemorragia\\
salen de una esfera entintada\\
me flotan, me hechizan\\
me suben por dentro, calladas

Explotan en un remolino\\
quieren tapar mi alma\\
y no salen de verdad\\
sólo la esconden, la enredan

\songpart{Pre-Chorus}
Aún no entiendo qué pasó\\
qué me pasó, por qué dolió\\
mi mundo se está deshaciendo\\
en un universo negro

Pero hay destellos\\
hay una niña triste\\
que en mis recuerdos vive\\
y todavía insiste

\songpart{Chorus}
Magia negra, magia negra\\
me atrapó, me atrapó\\
Magia negra, magia negra\\
y mi voz se rompió

Magia negra, magia negra\\
no me suelta, no\\
Magia negra, magia negra\\
salió de mi corazón

\songpart{Verse 2}
Desdibujó la sonrisa\\
de una amiga en un segundo\\
y al reír, al sonreír\\
me dejó sin suelo alguno

Una mano cálida aprieta la mía\\
y me oprime sin razón\\
cuando su vida ya no cabe\\
en el borde de su dolor

\songpart{Pre-Chorus}
Y yo no sé cargar\\
con lo que no pude ver\\
sus problemas son tan grandes\\
que me olvido de entender

Pero hay destellos\\
hay una niña triste\\
que en mis recuerdos vive\\
y todavía insiste

\songpart{Chorus}
Magia negra, magia negra\\
me atrapó, me atrapó\\
Magia negra, magia negra\\
y mi voz se rompió

Magia negra, magia negra\\
no me suelta, no\\
Magia negra, magia negra\\
salió de mi corazón

\songpart{Bridge}
No venía de ella\\
no venía de su luz\\
era mi propia sombra\\
mordiéndome la cruz

Lo sano se me pudre\\
lo bueno se me va\\
y en medio del silencio\\
me vuelvo a despertar

\songpart{Final Chorus}
Magia negra, magia negra\\
me atrapó, me atrapó\\
Magia negra, magia negra\\
y mi voz se rompió

Magia negra, magia negra\\
no me suelta, no\\
Magia negra, magia negra\\
salió de mi corazón

\end{lyrics}


\songtitle{No sé, vos sabés}

\tag{original-source}\tag{latin-house}\tag{dubstep}\tag{spanish-lyrics}\tag{female-vocals}\tag{rap}\tag{blues}

\begin{notes}
A cross-genre track that balances house pulse with dubstep intensity. %
The lyrics begin with “Hay veces en la existencia”, giving a quick sense of the song’s tone.
\end{notes}

\begin{style}
Latin house-dubstep with a rolling dembow pulse and half-time bass drops; verse rides sparse kicks, palm-muted synth stabs, and a bouncy low-end groove; pre-chorus lifts with filtered claps and rising vocal stacks; chorus hits wider with chopped lead synths and gang doubles. %
Soft close-mic female vocals, airy doubles, playful ad-libs, delay throws, and bright punchy mix with shimmering risers and reversed swells., female vocals, blues, rap, german hip hop, playful, soft, latin, dubstep, rock
\end{style}

\begin{lyrics}

\songpart{Verse 1}
Hay veces en la existencia\\
de cualquier ser, o la mía,\\
cuando pasa un mal momento\\
y otro ya se espera.\\
La vida no se derrumba\\
si uno se agarra\\
con las uñas\\
a cualquier aliciente.

A veces, por miedo,\\
por no soltar la pelea,\\
uno deja seguir\\
lo que igual duele.\\
Y en medio de todo,\\
sin buscarlo,\\
una mirada,\\
una sonrisa,\\
un saludo.

\songpart{Pre-Chorus}
Y listo,\\
la vida debe continuar.\\
Los problemas,\\
después.\\
Y esto otro va pa' vos,\\
si vos.

\songpart{Chorus}
No sé, no sé\\
si te amo.\\
No sé, no sé\\
si puedo olvidarte.\\
No sé, no sé\\
si pueda perdonarte.\\
Vos sabés, vos sabés\\
que yo te miro y tiemblo.

\songpart{Verse 2}
No sé de dónde viniste\\
ni por qué lo hiciste,\\
y lo que menos comprendo\\
es por qué me importas.\\
No sé si pueda llegar\\
el día de los dos,\\
dos extraños\\
como viejos amigos\\
en la calle.

No sé\\
si pueda dejar de hacerlo,\\
si lo que siento\\
se me va a partir.\\
No sé\\
lo que estás pensando,\\
ni sé\\
cómo me vi\\
tan estúpido\\
haciendo lo que hice.

\songpart{Pre-Chorus}
Solo espero,\\
poder verte un día\\
y sin decir\\
ni una palabra,\\
vos entiendas\\
que todo sigue,\\
que todo sigue,\\
aunque me queme.

\songpart{Chorus}
No sé, no sé\\
si te amo.\\
No sé, no sé\\
si puedo olvidarte.\\
No sé, no sé\\
si pueda perdonarte.\\
Vos sabés, vos sabés\\
que yo te miro y tiemblo.

\songpart{Bridge}
Y si te veo\\
hoy de frente,\\
que no falten\\
las ganas de hablar.\\
Un gesto corto,\\
una salida,\\
y la herida\\
se pone a bailar.

No sé,\\
no sé,\\
no sé,\\
pero vuelvo.

\songpart{Chorus}
No sé, no sé\\
si te amo.\\
No sé, no sé\\
si puedo olvidarte.\\
No sé, no sé\\
si pueda perdonarte.\\
Vos sabés, vos sabés\\
que yo te miro y tiemblo.

\end{lyrics}


\songtitle{Ninety-Six Static}

\tag*{1996}\tag{industrial-pop}\tag{original-work}

\begin{notes}
\end{notes}

\begin{style}
midtempo drum and bass, industrial pop, baritone lead vocals, spoken word bridge, English Swedish Spanish, 96 BPM, syncopated breakbeat, sub bass drops, metallic synth stabs, sampler glitch edits, tape saturation, plate reverb, detuned keyboard wash, gated snare, call and response chorus, digital nostalgia, archival dread
\end{style}

\begin{lyrics}
\songpart{Verse 1}
Ninety-six, the static on the line\\
Searching for a truth that I can define\\
Leaving old groups, moving through the haze\\
Tracing out the rhythm of these shifting days

\songpart{Pre-Chorus}
Signals in the dark, the pulse is in the wire\\
The network hums, setting hearts on fire\\
Caught in the loop, the rhythm starts to grow\\
Everything I write, the only truth I know

\songpart{Chorus}
Echoes in the code, drifting in the light\\
Caught in the rhythm of the digital night\\
Voices in the static, secrets in the screen\\
Living in the spaces of the unseen
(Everything is real, everything is true)

\songpart{Verse 2}
Tal vez escribir sea solo una forma de no estar solo\\
O una manera de decir que me muero de frío\\
Tracing every movement, the pulse is getting thin\\
Watching where the outer world meets the world within

\songpart{Pre-Chorus}
Signals in the dark, the pulse is in the wire\\
The network hums, setting hearts on fire\\
Caught in the loop, the rhythm starts to grow\\
Everything I write, the only truth I know

\songpart{Chorus}
Echoes in the code, drifting in the light\\
Caught in the rhythm of the digital night\\
Voices in the static, secrets in the screen\\
Living in the spaces of the unseen

\songpart{Bridge}
Jag är rädd att jag skriver för att inte behöva leva\\
El dolor es un hábito que no puedo dejar\\
The keyboard is a mirror, the monitor is cold\\
Spinning out the stories that never will be told

\songpart{Verse 3}
No sé si lo que escribo es verdad o es una máscara\\
Quizás la sinceridad es solo otro tipo de miedo\\
Digital echoes, hanging in the air\\
Lost in the archive, caught in the snare

\songpart{Pre-Chorus}
Signals in the dark, the pulse is in the wire\\
The network hums, setting hearts on fire\\
Caught in the loop, the rhythm starts to grow\\
Everything I write, the only truth I know

\songpart{Chorus}
Echoes in the code, drifting in the light\\
Caught in the rhythm of the digital night\\
Voices in the static, secrets in the screen\\
Living in the spaces of the unseen

\songpart{Outro}
Todo esto se perderá
(Just words on a screen)\\
Todo esto es lo que soy
(The echo of a dream)\\
Fading away...\\
Fade into the archive.
\end{lyrics}


\songtitle{En la nube}

\tag*{1996}\tag{new-age}\tag{original-source}\tag{spanish-lyrics}

\begin{notes}
Original text by Carlos Thompson (1996).
\end{notes}

\begin{style}
epic, choral, synthesizers, new age
\end{style}

\begin{lyrics}
Y terminamos perdidos en un espiral amorfo, enorme, obscuro. %
En un laberinto sin salida y sin entrada. %
Pero estamos dentro. Envueltos. %
Tratando de distinguir las sombras en la obscuridad. %
Temiendo mirar a quien está a nuestro lado y descubrir que estamos solos.\\
y escuchamos una voz.\\
Una voz amiga.\\
Y no sólo distinguimos las sombras ahora sino la luz también.\\
Una luz que no enceguece.

\end{lyrics}


\songtitle{Three Notebooks}

\tag*{1995}\tag{chamber-rock}\tag{autobiographical}

\begin{notes}
A semi-autobiographical piece rooted in 1995, charting the emergence of the Chlewey online presence and the early connections to GIDES and new digital communities. %
It blends nostalgia with the thrill of a new world opening.
\end{notes}

\begin{style}
chamber rock, punk rock, 168 BPM, close-mic baritone, whispered delivery, piano ostinato, spiccato strings, tom-heavy rock drums, distorted bass guitar, parallel compression, tape saturation, slapback delay, dry chamber room, minor key, menacing desire, rapid 16th-note drums
\end{style}

\begin{lyrics}
\songpart{Verse 1}
Late ninety-four, the crisis stayed\\
A heavy cloud on Chlewey’s days\\
Then GIDES called, and opened doors\\
To code, and labs, and something more

First page online, in three tongues wide\\
A personal home page, set with pride\\
And email came like a new-line spark\\
A way to reach beyond the dark

\songpart{Pre-Chorus}
He typed his name\\
Then sent it out\\
To every screen\\
On every route

\songpart{Chorus}
Chlewey on the line\\
Chlewey on the line\\
Three languages, one sky\\
Chlewey on the line

Chlewey on the line\\
Chlewey on the line\\
He found a world alive\\
Chlewey on the line

\songpart{Verse 2}
In June, he taught the C class there\\
For windows graphics, fresh and bare\\
Vitral formed from GIDES hands\\
A split, a spark, a shift in plans

Héctor Díaz wanted more\\
Not just code stacked on a floor\\
He wanted voices in the room\\
And nerds to bloom, and nerves to loosen too

\songpart{Pre-Chorus}
New faces close\\
New rules to learn\\
A different pulse\\
A turn, a turn

\songpart{Chorus}
Chlewey on the line\\
Chlewey on the line\\
Three languages, one sky\\
Chlewey on the line

Chlewey on the line\\
Chlewey on the line\\
He found a world alive\\
Chlewey on the line

\songpart{Bridge}
Objetivo Calidad\\
Opened up the missing door\\
A life from Amistad days\\
He thought was gone forevermore

And the notebook from ninety\\
That old diary at last\\
Ran through one date, then two, then three\\
Till a third one came so fast

\songpart{Final Chorus}
Chlewey on the line\\
Chlewey on the line\\
Three languages, one sky\\
Chlewey on the line

Chlewey on the line\\
Chlewey on the line\\
He found a world alive\\
Chlewey on the line
\end{lyrics}


\songtitle{Mirá}

\tag*{1996}\tag{original-work}\tag{musical-theater}

\begin{notes}
Original work from 1996.\\
Musical theater piece using conversational narrative framing.
\end{notes}

\begin{style}
Musical theater showtune with folk-pop influences. %
Baritone male vocals deliver a narrative performance with clear diction and theatrical phrasing. %
The arrangement features a bright acoustic guitar strumming eighth-note patterns, a melodic accordion providing counter-melodies and sustained chords, and a double bass playing on beats 1 and 3. %
A light percussion section includes a tambourine on the backbeat and subtle shaker. %
The tempo is 115 BPM in 4/4 time, likely in the key of G major. %
The harmonic structure follows a diatonic progression typical of contemporary musical theater, with a clean, dry production style that emphasizes vocal clarity and acoustic instrument separation.
\end{style}

\begin{lyrics}

\songpart{Verse 1}
\annotation{strummed acoustic guitar, double bass, accordion}
Amiga, ¿puedo llamarte amiga? mirá! tengo un problema, tú lo conoces o lo has oído en parte, estoy seguro.  Resulta que, hace tiempo ya, encontré a una persona, tú la conoces, su nombre; si no lo sabes lo descubrirás. %
\annotation{tambourine enters on backbeat} Hace algún otro tiempo descubrí que la amaba, o al menos eso creí, o, ¿sabes?, creo que ni eso creí pero lo hice.  Y también hace algún tiempo, se lo hice saber.  No sé si me creyó 

\end{lyrics}


\songtitle{Círculo Convexo}

\tag*{1997}\tag{original-lyrics}\tag{new-age-techno}

\subsection{Notes}

\subsection{Style}
New age techno with industrial notes.\\
Deep bass voice.\\
Synthesizers, percussion, bass, electric guitars , powerful bass, aggressive drums, electronic elements

\begin{lyrics}
\annotation{synthesizer riff}

\songpart{segment 1}
\annotation{Spoken}\annotation{Fast Delivery}\\
\annotation{male bass vocals}\\
Círculo convexo. %
Parábola\\
Un atardecer contigo y una noche larga, extrañándote.\\
¡Vívo! ¡Lo sé!\\
Vivo pensando, pensándote, extrañándote esperando ese momento especial para verte... %
sentirte\\
Tu cara y tu piel entre mis manos\\
Tu dulce voz susurrándome al oído\\
Tu mano resbalando por mi pecho

\annotation{percussion  riff}

\songpart{segment 2}
\annotation{melodic}\\
\annotation{male bass vocals}\\
Y volvió… lloviendo en todas partes\\
y escampando en tu mirada

\annotation{guitar riff}

\songpart{segment 3}
\annotation{melodic}\\
\annotation{male bass vocals}\\
Y es así como me absorbes tanto\\
absorción positiva\\
que mi mente sólo piensa en tu presencia\\
y el temor de tu partida

\songpart{segment 4}
\annotation{Spoken}\annotation{Fast Delivery}\\
\annotation{male bass vocals}\\
Círculo convexo.  Parábola.\\
Palabras que cruzan mi mente\\
una mente sin sentido\\
una mente en blanco que juega con mi vida y la oculta\\
que esconde lo que pienso\\
y salta juguetona entre realidades.

\annotation{synthesizer riff}

\songpart{segment 5}
\annotation{Spoken}\annotation{Fast Delivery}\\
\annotation{female alto vocals}\\
Era una gota de lluvia\\
un grandioso ser que quería abarcar al mundo\\
jugando con la luz y descomponiéndola en sus colores.

Era un mundo en sí misma\\
un destino de algo grandioso\\
abarcando el mundo para hacerlo fértil\\
para alimentar una cadena de pequeños mundos.

Era un sueño que se acercaba\\
dándose cuenta de la inmensidad de la tierra a donde llegaba.

Era mucho…  y al acercarse a su destino tan poco

Era…

\annotation{synthesizer riff}\\
\annotation{fade}
\end{lyrics}


\songtitle{Papelito Chiquito}

\subsection{Notes}

\subsection{Style}
Trance-dubstep with variable BPM, swelling and compressing between driving four-on-the-floor lifts and heavy half-time drops; verse rides airy arpeggios and tight low-end pulses, pre-chorus narrows to clipped drums and tense filter movement, chorus hits with distorted bass surges and echo-stacked soprano replies. %
Tenor lead in verses, deep baritone doubles on key lines, soprano echoes on the hook words, with delay throws, reverse swells, and hard sidechain pumping. %
Bright, glossy top end over a dark, weighty center., deep, trance, dubstep, theme

\begin{lyrics}
\songpart{Verse 1}
Pedí un papelito chiquito\\
para escribir chiquito\\
¡Y equibo tiquito!\\
Estuve feliz\\
y me gustó\\
Felizmente feliz

\songpart{Pre-Chorus}
Con esa emoción\\
de descubrir algo nuevo\\
de encontrar lo perdido\\
aunque siga perdido\\
ya no me importa

\songpart{Chorus}
Estuve feliz\\
estuve feliz\\
Felizmente feliz\\
ya no me importa\\
Estuve feliz\\
(Estuve feliz)\\
Felizmente feliz\\
(estuve feliz)

\songpart{Verse 2}
Podría salir\\
en este momento\\
y vivir en un mundo\\
enigmático y desconocido\\
podría probar\\
un nuevo rumbo\\
pero me da miedo\\
me da miedo

\songpart{Pre-Chorus}
Que al escaparme\\
me olvide de ti\\
que el aire me lleve\\
y me deje vacío\\
sin tu nombre

\songpart{Chorus}
Estuve feliz\\
estuve feliz\\
Felizmente feliz\\
ya no me importa\\
Estuve feliz\\
(Estuve feliz)\\
Felizmente feliz\\
(estuve feliz)

\songpart{Bridge}
Es imposible\\
establecer un nuevo puente\\
y un nuevo lugar\\
imposible terminar\\
imposible concluir\\
pero sigo aquí\\
palabra a palabra

\songpart{Final Chorus}
Estuve feliz\\
estuve feliz\\
Felizmente feliz\\
ya no me importa\\
Estuve feliz\\
(Estuve feliz)\\
Felizmente feliz\\
(estuve feliz)
\end{lyrics}


\songtitle{Septiembre}

\tag*{1996}\tag{dubstep}\tag{original-lyrics}\tag{spanish-lyrics}

\begin{notes}
Adaptation of a postcard written in 1996 for September's Love and Friendship month.
\end{notes}
\begin{style}
Dubstep epic with a tense half-time groove and swinging syncopated drums; verse sections ride sparse sub pulses and brittle ticks, pre-drop strips to filtered growls and a rising snare ladder, chorus slams with massive wobble bass, brass-like synth hits, and a clipped vocal hook texture. %
Breathy choirs, reverse swells, dust-hit risers, and shimmering bell sparks thread the transitions. %
Wide, cinematic, glossy mix with sharp low-end impact and photoreal detail., rich, dubstep, light, soft, natural, deep, dramatic, complex, Colombian pronunciation
\end{style}
\begin{lyrics}
\songpart{Intro}
\annotation{Instrumental}

\songpart{Bridge}
\annotation{melodic speech, breathing, whispering}\\
La existencia humana está llena de pequeños detalles\\
que no por pequeños grandiosos;\\
de detalles y momentos que hacen por la alegría y la euforia\\
o tienden por la tristeza y la depresión.

Y es así como uno está vivo,\\
se siente vivo;\\
como una roca cayendo a veces, pero roca viva\\
o como un globo que flota jugetón entre las nubes.

Y pensar que en estos detalles y momentos,\\
es siempre indispensable\\
una amiga como tú.

(como tú)

(como tú)
\end{lyrics}
\songtitle{Monoestable}

\tag{original-work}\tag{alternative-rock}\tag{spanish-lyrics}

\begin{notes}
Original work in Spanish.\\
Explores state dynamics and internal oscillation; themed around instability and emotional transformation.
\end{notes}

\begin{style}
Dark alternative pop-rock in Spanish with a tense mid-tempo pulse, jagged palm-muted guitars, punchy toms, and a low synth drone; verse stays tight and claustrophobic, pre-chorus opens with rising backing vocals and filtered percussion, chorus lands with a stark shouted hook and layered doubles, bridge strips to voice and bass before the final lift. %
Close-mic lead vocal, dry and intimate in verses, wider harmonies and short delay throws on the hook word, with vinyl crackle and reverse swells between sections; bright top end over a heavy, compressed low end.
\end{style}

\begin{lyrics}
\songpart{Verse 1}
Teoría del chaos, teoría de catástrofes\\
es el gatillo que dispara al monoestable\\
¿Cómo se puede estar bien?\\
¿Cómo se puede caer así?

Sombras que se cierran\\
sobre la mente\\
que encierran, que aprisionan\\
que revientan

\songpart{Pre-Chorus}
Y no se abre la puerta\\
y no baja la tensión\\
todo gira hacia adentro\\
todo rompe alrededor

\songpart{Chorus}
Y es ella\\
siempre ella\\
Ella y el vacío\\
Ella y el vacío

Y es ella\\
siempre ella\\
Me deja sin nombre\\
me deja sin hilo

\songpart{Verse 2}
Lleva el pulso en la nuca\\
como un vidrio bajo piel\\
me promete que regreso\\
y me pierde otra vez

En la esquina de mis días\\
se me quiebra la razón\\
y lo que callo me muerde\\
desde el fondo del montón

\songpart{Pre-Chorus}
Y no se abre la puerta\\
y no baja la tensión\\
todo gira hacia adentro\\
todo rompe alrededor

\songpart{Chorus}
Y es ella\\
siempre ella\\
Ella y el vacío\\
Ella y el vacío

Y es ella\\
siempre ella\\
Me deja sin nombre\\
me deja sin hilo

\songpart{Bridge}
Dime dónde guardo el miedo\\
si me ocupa todo el cuerpo\\
dime cómo sigue vivo\\
lo que nace de un incendio

\songpart{Chorus}
Y es ella\\
siempre ella\\
Ella y el vacío\\
Ella y el vacío

Y es ella\\
siempre ella\\
Me deja sin nombre\\
me deja sin hilo
\end{lyrics}


\songtitle{Lluvia En Tu Mirada}

\tag*{1996}\tag{pop-rock}\tag{spanish-lyrics}
\begin{notes}
Lyrics by SUNO based on two line poem by Carlos Thompson.
\end{notes}
\begin{style}
Indie pop-rock uptempo with brisk offbeat guitar strums, tight kick-snare drive, and a chiming bassline; verses ride lean and restless, pre-chorus opens with rising harmonies, chorus lands brighter with stacked vocals and a singable hook, Soft ad-lib echoes on key lines, subtle reverse swells into the drop, and a final spoken outro over stripped drums, Bright, airy, and close-mic with an emotional lift, uplifting, tone
\end{style}
\begin{lyrics}
\songpart{Verse 1}
Caminé por días grises\\
con los bolsillos vacíos\\
la tarde caía lenta\\
sobre mis pasos fríos

Tu nombre en la ventana\\
me cambió la dirección\\
y aunque el cielo estaba roto\\
yo seguí la canción

\songpart{Pre-Chorus}
Porque hay algo en tu risa\\
que me vuelve a levantar\\
cuando todo se deshace\\
tú me enseñas a mirar

\songpart{Chorus}
Y volvió… lloviendo en todas partes\\
y escampando en tu mirada\\
Y volvió… lloviendo en todas partes\\
y escampando en tu mirada
(escampando en tu mirada)

\songpart{Verse 2}
Las calles iban cansadas\\
la ropa olía a temporal\\
pero al verte de frente\\
se me olvidó el mal

Tus manos hacen abrigo\\
donde el frío quiso entrar\\
y en medio de tanto invierno\\
me das otro lugar

\songpart{Pre-Chorus}
Porque hay algo en tu risa\\
que me vuelve a levantar\\
cuando todo se deshace\\
tú me enseñas a mirar

\songpart{Chorus}
Y volvió… lloviendo en todas partes\\
y escampando en tu mirada\\
Y volvió… lloviendo en todas partes\\
y escampando en tu mirada
(escampando en tu mirada)

\songpart{Bridge}
Si el mundo pesa, tú no\\
si el día muerde, tú no\\
te acercas y todo cae\\
despacio, sin rumor

Y yo, que iba a perderme\\
ya sé por dónde andar\\
si me nombras, se me enciende\\
la forma de respirar

\songpart{Chorus}
Y volvió… lloviendo en todas partes\\
y escampando en tu mirada\\
Y volvió… lloviendo en todas partes\\
y escampando en tu mirada
(escampando en tu mirada)

\songpart{Outro}
Tu existencia me arregla el día.\\
Me deja la casa adentro.\\
Me hace mejor el aire.\\
Y cuando estás, todo amaina.
\end{lyrics}


\input{temp/vaina_en_graeville}
\end{multicols*}

\chapter{Earlier Thoughts: Undated and Later Revisits}
\albumcover{covers/internal/earlier-thoughts-undated}
Four pieces that belong to the earlier-thoughts family in spirit but not in paperwork -- undated fragments, or lines written decades later that circle straight back to a cold wall, a paper staircase, a question he was still asking himself long after he should have had the answer.
\clearpage\begin{multicols*}{4}
\songtitle{Pared fría}

\tag*{2025}\tag{trance}\tag{drift}\tag{damon}\tag{original-work}\tag{spanish-lyrics}

\begin{notes}
Original work.\\
Co-written with Grok.\\
Based on character Damon from the short story \emph{Drift} by Carlos Thompson and ChatGPT.
\end{notes}

\begin{style}
trance music, dubstep club, 140 BPM, deep bass male vocals, close-mic whisper, breathy low register, clipped synth stabs, metallic percussion, sidechain pumping, half-time breakdown, gated reverb snare, stereo delay throws, tape saturation, analog synth pads, pulsing sub-bass, bittersweet longing, restrained ache
\end{style}

\begin{lyrics}
\songpart{Verse 1}
Sorprende mi silencio\\
cuando ya no sé nombrarme\\
y las palabras se esconden\\
en los pliegues de mi mente

Mi grito pide salida\\
pero se queda en la cara\\
frío de pared, piedra dura\\
pegado a mi frente

\songpart{Pre-Chorus}
Y me pregunto por qué\\
dejo pasar la vida\\
por qué mi boca se rompe\\
antes de decir lo que siento

Por qué entrego mi herida\\
al capricho del tiempo\\
sin abrir la puerta\\
sin pedir auxilio

\songpart{Chorus}
Sólo quédate aquí\\
sólo quédate aquí\\
no me preguntes por qué\\
sólo quédate aquí

Dame un hombro, por favor\\
dame un hombro, por favor\\
que pueda llorar sin hablar\\
sólo quédate aquí

\songpart{Verse 2}
No busco una explicación\\
ni un mapa de mis sombras\\
a veces necesito un alma\\
que no me suelte la mano

Que no intente descifrarme\\
que no me pida razones\\
que vea temblar mi pecho\\
y se quede a mi lado

\songpart{Pre-Chorus}
Y me pregunto por qué\\
me trago cada caída\\
si el cuerpo ya me delata\\
si la pena me nombra

Por qué dejo que el reloj\\
me atraviese despacio\\
sin decirle a los míos\\
cuánto me falta

\songpart{Chorus}
Sólo quédate aquí\\
sólo quédate aquí\\
no me preguntes por qué\\
sólo quédate aquí

Dame un hombro, por favor\\
dame un hombro, por favor\\
que pueda llorar sin hablar\\
sólo quédate aquí

\songpart{Bridge}
Si me ves callar\\
no me apartes\\
si me ves caer\\
no me sueltes

Yo también quiero saber\\
cómo se nombra el dolor\\
pero hoy sólo necesito\\
tu presencia y tu calor

\songpart{Final Chorus}
Sólo quédate aquí\\
sólo quédate aquí\\
no me preguntes por qué\\
sólo quédate aquí

Dame un hombro, por favor\\
dame un hombro, por favor\\
que pueda llorar sin hablar\\
sólo quédate aquí
\end{lyrics}


\songtitle{Papel en la escalera}

\tag*{2026}\tag{dubstep}\tag{spanish-lyrics}
\begin{notes}

\end{notes}
\begin{style}
Dubstep hymn with halftime wobble bass, stomping kick-snare, and tense syncopated synth pulses; verses stay sparse with close-mic male vocal and clipped doubles, pre-chorus opens into stacked harmony lifts, chorus drops into a wide sub-heavy chant with vowel hooks. %
Use reversed swells, glitch ticks, and filtered risers between lines; final chorus adds choir-like gang layers and a last bass surge. %
Bright, punchy mix with a cold spacious center., hard rock, glam metal, lyrical, dubstep.
vocals: male baritone lead, contained, use pitch and intonation for emphasis, no screaming, no shouting, central colombian pronunciation: seseo.
\end{style}
\begin{lyrics}
\songpart{Verse 1}
Lo que se nombra queda.
Lo que hiere, también.
Yo escribo para ocultarme
y para verme la piel.

Digo menos de lo que siento,
por cuidar lo que soy.
Pero en cada frase rota
va mi pulso y mi voz.

\songpart{Pre-Chorus}
Hay verdades torcidas
que me saben hablar.
Y en mis manos temblorosas
algo empieza a despertar.

\songpart{Chorus}
No se va, no se va
lo que toca la verdad.
No se va, no se va
aunque cambie de señal.

No se va, no se va
lo que escribe mi lugar.
No se va, no se va
yo lo vuelvo a levantar.

\songpart{Verse 2}
Bajo gradas, cruzo miradas,
y el aire dice más.
Un saludo breve en medio
de la duda y la señal.

Yo no sé qué significa
su silencio al pasar,
pero el tiempo lo acomoda
y me lo deja mirar.

\songpart{Pre-Chorus}
Mientras habla el frente,
mi cabeza cae atrás.
Y en el margen del cuaderno
algo mío vuelve a entrar.

\songpart{Chorus}
No se va, no se va
lo que toca la verdad.
No se va, no se va
aunque cambie de señal.

No se va, no se va
lo que escribe mi lugar.
No se va, no se va
yo lo vuelvo a levantar.

\songpart{Bridge}
Hay un fraude en mis palabras,
sí, lo puedo reconocer.
Pero entre la frase armada
vive algo que sí fue.

Lo que cae de tantas sombras
todavía quiere hablar,
y aun si tiembla mi respuesta
no me puedo abandonar.

\songpart{Final Chorus}
No se va, no se va
lo que toca la verdad.
No se va, no se va
aunque cambie de señal.

No se va, no se va
lo que escribe mi lugar.
No se va, no se va
yo lo vuelvo a levantar.

No se va, no se va
lo que duele al respirar.
No se va, no se va
yo lo vuelvo a levantar.
\end{lyrics}


\songtitle{La duda - por qué}

\tag*{2026}\tag{melodic-dubstep}\tag{multilingual}\tag{original-source}

\begin{notes}
\end{notes}
\begin{style}
Emotional melodic dubstep with cinematic atmosphere and controlled dynamics. 

One consistent male lead voice across all languages (Spanish, French, German, English), baritone, soft and intimate tone, no shouting, no aggressive metal-style delivery. %
Vocals should remain melodic and connected, not spoken-word unless explicitly marked.

Secondary background voice used only for echoes, delayed responses, and canon harmonies. %
female soprano vocal, distant, airy, and complementary.

Use overlapping multilingual vocal echoes with stereo delay and canon layering. %
Maintain smooth rhythmic phrasing without unnatural pauses between short lines.

Allow intentional pauses only during instrumental fills and dubstep transitions. %
Include clear instrumental breaks with wobble bass, glitch effects, reverse swells, chopped vocal fragments, and half-time drops.

Deep sub bass, syncopated drums, cinematic low-end, bright atmospheric synths, restrained emotional intensity, dramatic but controlled mix.
\end{style}
\begin{lyrics}
\songpart{Instrumental Intro}
\annotation{distant bass drones, breathy textures, slow filter rise}

\songpart{Verse}
\annotation{Lead voice, soft sung male voice, intimate, connected phrasing}\\
En serio,\\
la duda existe,\\
la duda de quién soy,\\
la duda de qué pienso,\\
la duda de por qué lo hago,\\
la duda de quién es.

\annotation{Lead voice} La duda, \annotation{Echo voice, female soprano vocals}  (la duda)\\
\annotation{Lead voice} la duda, \annotation{Echo voice} (la duda)\\
\annotation{Lead voice} la duda...\\
sólo es la duda de quién soy, (¿quien soy?)\\
y de ahí la de qué hago,\\
y el eterno "¿por qué?" (¿por qué?)

\annotation{pause - instrumental fill}

\annotation{Dubstep Bridge}\\
\annotation{heavy wobble drop, chopped vocal echoes}\\
(la duda)\\
(la duda)

\annotation{Lead voice}\\
¿Por qué?\\
(¿Por qué?)

\annotation{Lead voice - French}\\
Pourquoi ?\\
(Pourquoi ?)

(Pourquoi ?)

\annotation{Lead voice, German}\\
Warum?
(Warum?)

Warum?\\
(Warum?)

(Warum?)

\annotation{Lead voice}\\
¿Por qué?

(la duda existe,)\\
(la duda de quién soy,)\\
(la duda de qué pienso,)\\
(la duda de por qué lo hago,)

\annotation{Lead voice, English}\\
Why?
(Why?)

(la duda)

(la duda)

(en serio)

(¿Por qué?)

\annotation{instrumental break}\\
\annotation{glitched synths, no vocals, half-time drums}

\annotation{Final}\\
¿Por qué? (la duda)
\end{lyrics}


\songtitle{Espectador de mi vida}

\subsection{Notes}

\subsection{Style}
8-bit chiptune. %
Square wave lead melody, pulse wave bass, and noise channel percussion. %
140 BPM, C Major, 4/4 time. %
The lead melody features rapid arpeggios and vibrato typical of the NES sound chip. %
The bassline uses a constant eighth-note pulse. %
Percussion consists of short noise bursts for the snare and lower-frequency noise for the kick. %
The arrangement uses a repetitive loop structure with a high-pitched counter-melody entering in the second half.

\begin{lyrics}
\songpart{Verse 1}
Sigo de pie\\
sin mirar atrás\\
con este escudo\\
de pura resignación

Me digo: “puedo”\\
y vuelvo a aguantar\\
pero por dentro\\
ya no sé quién soy

\songpart{Pre-Chorus}
Y dejo caer\\
mi voz al suelo\\
como si nada\\
pudiera salvarla hoy

\songpart{Chorus}
Soy espectador de mi vida\\
soy espectador, y me duele\\
(soy espectador)\\
La veo partir\\
de mis manos vacías\\
y sigo aquí\\
sobreviviendo

Soy espectador de mi vida\\
soy espectador, y me rompe\\
(soy espectador)\\
Quiero detener\\
lo que me derriba\\
pero me quedo\\
sobreviviendo

\songpart{Verse 2}
Veo mis sueños\\
irse sin mí\\
como luces frías\\
que no pude cuidar

Una se enciende\\
otra se va\\
y yo tan lejos\\
de poderla abrazar

Me cae el mundo\\
sobre la piel\\
me va borrando\\
la fuerza y la fe

\songpart{Pre-Chorus}
Y aunque lo sé\\
no doy un paso\\
me quedo mudo\\
mirando perder

\songpart{Chorus}
Soy espectador de mi vida\\
soy espectador, y me duele\\
(soy espectador)\\
La veo partir\\
de mis manos vacías\\
y sigo aquí\\
sobreviviendo

Soy espectador de mi vida\\
soy espectador, y me rompe\\
(soy espectador)\\
Quiero detener\\
lo que me derriba\\
pero me quedo\\
sobreviviendo

\songpart{Bridge}
Si alguna vez\\
miro de frente\\
tal vez descubra\\
que aún puedo caer

Tal vez caer\\
sea despertar\\
tal vez romperme\\
sea empezar a ser

\annotation{Final Chorus}\\
Soy espectador de mi vida\\
soy espectador, y me duele\\
(soy espectador)\\
No quiero seguir\\
con el alma rendida\\
quiero salir\\
de este silencio

Soy espectador de mi vida\\
ya no me alcanza con sobrevivir\\
(ya no más)\\
Si tengo que arder\\
para estar con vida\\
déjame arder\\
y decidir
\end{lyrics}


\end{multicols*}

\chapter{Gr\ae ville}
\albumcover{covers/internal/graeville}
\emph{Græville} was a private-language experiment -- an invented town, half diary and half fiction, that let a teenager say things he couldn't say plainly. A cup of cold chocolate, a room full of stuffed animals, a barrier that never quite explains itself: the songs keep the experiment's original distance intact.
\clearpage\begin{multicols*}{4}
\songtitle{Chocolate Frío}

\tag*{2026}\tag{latin-pop}\tag{spanish-lyrics}

\begin{notes}
\end{notes}
\begin{style}
Latin pop with a jaunty vallenato-cumbia pulse, mid-tempo swing and hand percussion; verse rides close-mic guitar and a bouncing bassline, pre-chorus strips to claps and a rising guiro, chorus opens with accordion stabs and gang vocal hooks. %
Lead vocal stays intimate and playful with stacked doubles on the hook, ad-lib echoes on punch lines, and a few delay throws on the title phrase. %
Ear candy: reversed breath swell into the chorus, a school-bell chime between lines, and a short riser before the final lift. %
Bright, colorful mix with punchy low end and glossy top.
\end{style}
\begin{lyrics}
\songpart{Verso 1}
Mamá vendrá de nuevo\\
a levantarme, piensa Ámila\\
quien medio abre los ojos\\
apenas lo suficiente\\
para ver cuando su mamá entra

Pero la mamá\\
simplemente trata\\
de levantarle las cobijas\\
que Ámila agarra con fuerza

¿Por qué los padres siempre tienen\\
que despertarlo a uno\\
cuando uno está durmiendo?

\songpart{Pre-Coro}
Que el chocolate se está infriando\\
dice mamá\\
ya voy más tardecito\\
yo estoy dormida

\songpart{Coro}
¡Ay, mami, deja que se infrie!\\
ya voy más tardecito\\
¡Ay, mami, deja que se infrie!\\
yo estoy dormida\\
¡Ay, mami, deja que se infrie!\\
pero no me dejes sin bus

\songpart{Verso 2}
Aunque siempre es chévere\\
ir en el carro de papi\\
hasta el colegio

No sé por qué\\
a otros les gustan los viernes\\
yo los odio

Era el quinto viernes en línea\\
que amanecía enguayabada\\
y fuera de eso\\
para que la recibieran con cálculo

\songpart{Pre-Coro}
¡Ay, mami, deja que se infrie!\\
ya voy más tardecito\\
fue lo primero que se le ocurrió decir\\
sin pensar

Que eran las mismas palabras\\
que utilizó el jueves\\
y el miércoles\\
y el martes

\songpart{Coro}
¡Ay, mami, deja que se infrie!\\
ya voy más tardecito\\
¡Ay, mami, deja que se infrie!\\
yo estoy dormida\\
¡Ay, mami, deja que se infrie!\\
pero no me dejes sin bus

\songpart{Verso 3}
Pero al fin\\
con un dolor de cabeza ni el hijuemadre\\
los ojos cerrados\\
y tanteando el camino hasta la cocina

Llega a tomar el café frío\\
mientras su hermanita\\
sale corriendo\\
a subirse al bus

Sale de la ducha\\
para controlar la temperatura\\
antes de volverse a meter\\
bajo el chorro

\songpart{Puente}
Con media hora de retraso\\
llega al colegio\\
personalizado uniforme\\
medias caídas

Maquillaje por encima del reglamentario\\
pero necesario\\
para ocultar las ojeras\\
y los rastros de la fiesta de anoche

\songpart{Último coro}
¡Ay, mami, deja que se infrie!\\
ya voy más tardecito\\
¡Ay, mami, deja que se infrie!\\
yo estoy dormida\\
¡Ay, mami, deja que se infrie!\\
pero al fin voy pa’l colegio
\end{lyrics}


\songtitle{Cuarto de Peluches}

\subsection{Notes}

\subsection{Style}
German hip hop with Spanish lyrics, moody boom-bap swing and half-time drums; verse rides tense bass pulses and dry snares, pre-chorus strips to muted chords and breathy pauses, chorus hits with stacked unison hooks and tight gang doubles. %
Two vocal leads trade lines with whispered asides, reverse swells into the hook, tape-stop turns before the bridge, and a wide, punchy, intimate mix., rap, german hip hop

\begin{lyrics}
\songpart{Verse 1}
Graeville 4\\
atrapados de verdad\\
el cuarto era jaula\\
y también era verdad

muy lejos de ayuda\\
muy cerca del final\\
si salíamos corriendo\\
podíamos morir igual

Amila perdida\\
o perdidos para ella\\
nadie dice su nombre\\
sin temblar por dentro de ellas

no pienses en ella\\
no pienso en ella\\
pero tu mirada insiste\\
y me deja sin defensa

\songpart{Pre-Chorus}
¿Estamos discutiendo?\\
yo ya no sé\\
tengo miedo por ella\\
o por los tres

ella está bien\\
no hables de ella\\
tengo miedo por mí\\
y por lo que revelas

\songpart{Chorus}
Te amo, amigo mío\\
(te amo, amigo mío)\\
Tengo miedo, tengo miedo\\
de romper lo vivido

Te amo, amigo mío\\
(te amo, amigo mío)\\
y si tiembla tu boca\\
tiembla lo que has tenido

\songpart{Verse 2}
Cuántas veces lo dije\\
cuántas veces callé\\
te amo, te lo juro\\
y me sale al revés

te odio, no lo digo\\
pero duele pensar\\
que besaste a mi hermana\\
y ahora todo va a estallar

¿crees que soy machista?\\
sí, sí, sí\\
me importa un culo el juicio\\
si me sale decir

pelada, no seas grosera\\
y qué mierda, qué va\\
si me da la puta gana\\
decir la verdad acá

\songpart{Bridge}
Veo tus ojos\\
y quieres cerrarlos\\
para ocultar\\
lo que piensas de mí

casi adivino\\
tu miedo exacto\\
maldita sea\\
maldito el día en que me metí

\songpart{Chorus}
Te amo, amigo mío\\
(te amo, amigo mío)\\
Tengo miedo, tengo miedo\\
de romper lo vivido

Te amo, amigo mío\\
(te amo, amigo mío)\\
y si tiembla tu boca\\
tiembla lo que has tenido

\songpart{Outro}
Dijimos buenas noches\\
y no era dormir\\
era quedarnos vivos\\
sin saber qué decir

en un rincón de pared\\
se fueron venciendo\\
mirando a los peluches\\
y al dibujo en el techo

\end{lyrics}


\songtitle{La Barrera}

\tag*{1992}\tag{hip-hop}\tag{graeville}\tag{german-lyrics}
\begin{notes}
Lyrics by SUNO based on \emph{Gr\ae ville}.
\end{notes}
\begin{style}
German hip hop and rap with tense boom-bap swing, snapped kick-snare pocket, dark minor-key keys, and a rolling sub line; verse rides sparse drums and dry lead vocal, pre-chorus opens with filtered pads and rising claps, chorus hits with doubled chants and short crowd-style responses, Tight close-mic male vocal, gritty center, ad-lib throws on the last word of each hook, vinyl crackle and reversed swell into transitions, bright punchy mix with hard low end, german hip hop, rap
\end{style}
\begin{lyrics}
\songpart{Verse 1}
Entonces,\\
¿pueden nuestros amigos volver?\\
Ya vieron la luz\\
que el fundador de la academia reveló

Y tienen que cumplir\\
el propósito que les indicó\\
Salvar a los encadenados\\
de la oscuridad

Llevarlos por el sendero\\
que atraviesa la barrera\\
de la comprensión\\
para que vean la tan esperada luz

Y descubrir\\
que la realidad de las ideas\\
no es perfecta

Es nula\\
sin la verdadera realidad\\
que hemos llamado sombras\\
las figuras no existen

Solo envician\\
la realidad que conocemos\\
no es una caverna\\
no son una caverna

\songpart{Pre-Chorus}
Los millones de niños\\
que mueren de hambre cada año\\
sin haber escuchado siquiera\\
la palabra realidad

Aun si vale más saber\\
que vivir\\
¿de qué sirve morir\\
sin lograr saber?

\songpart{Chorus}
La barrera, la barrera\\
nunca la cruzamos

La barrera, la barrera\\
a veces regresamos

Verdad, verdad\\
si no contradice lo que ya sé

Verdad, verdad\\
si me despierta otra vez

\songpart{Verse 2}
¿Y si vale más vivir\\
que saber?\\
¿De qué sirve morir\\
de todas maneras?

Alguien puede tener la respuesta\\
y la respuesta es verdad\\
mientras no contradiga\\
algo que ya he tomado como realidad

Retórica, retórica\\
hemos hecho de ella una puerta\\
pero no es la puerta\\
es la grieta de la cueva abierta

Nunca hemos cruzado la barrera, sabes\\
lo hemos hecho en realidad\\
solo que a cada descuido\\
regresamos de nuevo a este lado

\songpart{Bridge}
Mira bien\\
cuando crees que ya saliste\\
la sombra te llama por tu nombre\\
y vuelves

Mira bien\\
no basta con mirar\\
hay que romper el miedo\\
para ver

\songpart{Chorus}
La barrera, la barrera\\
nunca la cruzamos

La barrera, la barrera\\
a veces regresamos

Verdad, verdad\\
si no contradice lo que ya sé

Verdad, verdad\\
si me despierta otra vez
\end{lyrics}


\end{multicols*}

\chapter{Midlife Thoughts I: Social Feeds}
\albumcover{covers/internal/midlife-thoughts-i-social-feeds}
Status updates and blog posts nobody was meant to set to music: the running log of nothing in particular, the ghost of who he wasn't, the endless scroll of digital skies. Two decades on, the platforms are gone but the posts survive, sung back at a distance they were never built to travel.
\clearpage\begin{multicols*}{4}
\songtitle{Another Log About Nothing}

\tag*{2009}\tag{latin-jazz}\tag{social-media-seed}
\begin{notes}
Lyrics by SUNO, based on the first post of a new Tumblelog by Carlos Thompson under alias \textsc{gabriellicus}.
\end{notes}
\begin{style}
Electronic Latin-jazz with industrial texture: midtempo syncopated percussion, dry hand drums, muted trumpet phrases, electric piano chords, upright-style bass, clipped analog synth pulses, metallic found-sound hits, and subtle machine-room ambience, Verses stay sparse and conversational; the chorus opens into warmer horn harmonies and fuller clave-driven percussion; a short instrumental break strips down to bass, percussion, and processed trumpet before the final chorus, Intimate bilingual-feeling male vocal delivery in English, half-sung and rhythmically articulate, with close-miked dryness, occasional doubled phrases, and a clean but gritty mix with controlled low end
\end{style}
\begin{lyrics}
\songpart{Intro}
Another log about nothing\\
A little room outside the rules

\songpart{Verse 1}
When I built my Tumblelog\\
I gave it one clear face\\
Beautiful ladies, beautiful eyes\\
A gallery in one small place

Then another door swung open\\
A different theme, a different view\\
Nudity in a second window\\
A boundary I could move into

\songpart{Pre-Chorus}
The shelves are neat, the labels hold\\
But my hands keep reaching past

\songpart{Chorus}
I want a place to bring it all\\
The things I find, the things I pass\\
Not every spark belongs to me\\
But some deserve a wider glass\\
Another log about nothing\\
Or everything that slips the frame\\
A public room for passing signals\\
A brighter address for a borrowed flame

\songpart{Verse 2}
Sometimes I press the little heart\\
And leave the post beneath the floor\\
The likes are public if you know the link\\
But hidden doors are still a door

Sometimes I want to share a thought\\
I don’t have to call it mine\\
Not because I love it perfectly\\
But because it caught my eye in time

\songpart{Pre-Chorus}
A mark can say, “I saw this here”\\
A reblog says, “Come look with me”

\songpart{Chorus}
I want a place to bring it all\\
The things I find, the things I pass\\
Not every spark belongs to me\\
But some deserve a wider glass\\
Another log about nothing\\
Or everything that slips the frame\\
A public room for passing signals\\
A brighter address for a borrowed flame

\annotation{Instrumental}

\songpart{Bridge}
No manifesto, no grand design\\
Just an open shelf, a moving line\\
If it interests me, it can arrive\\
If it makes you pause, let it survive

\songpart{Final Chorus}
I want a place to bring it all\\
The odd, the bright, the overgrown\\
Not every post becomes a favorite\\
Still it shouldn’t stay alone\\
Another log about nothing\\
A loose collection, unconfined\\
A little more room for what I notice\\
A wider window on my mind

\songpart{Outro}
Feel free to enjoy the new reblogging log\\
No theme to serve, no gate to guard\\
Just things that found a second address\\
And traveled a little farther
\end{lyrics}


\songtitle{Fantasma de quien no es}

\tag*{2010}\tag{disco}\tag{social-media-seed}\tag{spanish-lyrics}
\begin{notes}
Lyrics by SUNO based on the bio of the Tumblelog by Carlos Thompson under username \textsc{epeady}.
\end{notes}
\begin{style}
Late-1970s disco anthem with a tight four-on-the-floor kick, syncopated bass guitar, crisp hi-hat, handclaps, sweeping string sections, bright electric piano, rhythmic guitar scratches, and a soaring female lead vocal, Start with a short filtered groove, build through intimate verses into a huge sing-along chorus, add a call-and-response post-chorus, then a dramatic string-led bridge and an extended final chorus, Polished club mix, warm analog saturation, punchy low end, wide strings, clear upfront vocal, joyful but slightly mysterious theatrical energy
\end{style}
\begin{lyrics}
\songpart{Intro}
¿Quién anda ahí?\\
Un suspiro con zapatos,\\
un corazón de algodón.

\songpart{Verso 1}
Te encontré junto al perchero,\\
con un botón fuera de lugar,\\
tenías polvo en los hombros\\
y ganas de bailar.\\
No eras dueño de tu nombre,\\
ni sabías dónde ir,\\
pero al rozar mi mano dijiste:\\
“Quédate aquí”.

\annotation{Pre-Coro}\\
Si eres aire, pasa lento,\\
si eres sueño, no te vayas,\\
si eres solo un muñeco,\\
¿por qué me miras y me llamas?

\annotation{Coro}\\
Fantasma de quien no es,\\
ven, gira otra vez,\\
un suspiro, un soplo,\\
un peluche de papel.\\
Fantasma de quien no es,\\
quédate junto a mis pies,\\
si no tienes alma,\\
te presto la mía después.

\annotation{Post-Coro}\\
Ah-ah, ven a bailar,\\
ah-ah, no mires atrás,\\
ah-ah, de terciopelo,\\
mi pequeño vendaval.

\songpart{Verso 2}
En la fiesta de la sala\\
nadie te pudo ver,\\
pero moviste las cortinas\\
al compás de la pared.\\
Tu relleno guardaba historias\\
de una infancia sin final,\\
y en tu pecho había un latido\\
que no era natural.

\annotation{Pre-Coro}\\
Si eres aire, pasa lento,\\
si eres sueño, no te vayas,\\
si eres solo un muñeco,\\
¿por qué me miras y me llamas?

\annotation{Coro}\\
Fantasma de quien no es,\\
ven, gira otra vez,\\
un suspiro, un soplo,\\
un peluche de papel.\\
Fantasma de quien no es,\\
quédate junto a mis pies,\\
si no tienes alma,\\
te presto la mía después.

\annotation{Puente}\\
No preguntes quién te hizo,\\
no preguntes quién se fue,\\
hay abrazos que aparecen\\
cuando nadie los pidió.\\
Y si mañana desapareces\\
por la rendija del balcón,\\
dejarás sobre mi almohada\\
una pequeña respiración.

\annotation{Coro Final}\\
Fantasma de quien no es,\\
ven, gira otra vez,\\
un suspiro, un soplo,\\
un peluche de papel.\\
Fantasma de quien no es,\\
quédate junto a mis pies,\\
si no tienes alma,\\
te presto la mía después.

\songpart{Outro}
Duerme, pequeño aire,\\
duerme junto a mi piel,\\
fantasma de quien no es,\\
pero mío esta vez.
\end{lyrics}


\songtitle{Digital Skies}

\tag*{2009}\tag{latin-jazz}\tag{social-media-seed}
\begin{notes}
Lyrics by SUNO, based on the first post of a new Tumblelog by Carlos Thompson under alias \textsc{gabriellicus}.
\end{notes}
\begin{style}
electronic latin jazz, bossa-nova groove, industrial textures, metallic percussion, glitchy synth bass, moody brass, warm analog rhodes piano, crisp female vocals, dynamic transition from smooth lounge to harsh electronic beats
\end{style}
\begin{lyrics}
\songpart{Intro}
\annotation{Instrumental - smooth bossa nova beat with static vinyl crackle and metallic hits}

\songpart{Verse 1}
I built a gallery of faces, framed in grace\\
A curated, beautiful, restricted space\\
And then another door, another theme to hold\\
The lines of beauty, and the stories told\\
But clean walls feel like cages after while\\
They dictate every breath, they frame your smile\\
Thematic patterns never let you stray\\
They lock the open windows of the day

\songpart{Chorus}
This is yet another log about nothing at all\\
No boundaries to stand, no rules to build a wall\\
Just the random currents of a wandering mind\\
Everything and nothing, whatever I can find\\
A chaotic rhythm, a digital stream\\
No thematic prison, no restricted dream

\songpart{Verse 2}
I hit the 'like' but it is buried deep\\
An obscure, quiet corner that the archives keep\\
I want to share the friction, the things I don't adore\\
To put them in the window, to open up the door\\
So let the static mingle with the horn\\
A brand-new, formless energy is born\\
Reblogging the errors, the shadows, the light\\
No longer forced to keep the borders tight

\songpart{Chorus}
This is yet another log about nothing at all\\
No boundaries to stand, no rules to build a wall\\
Just the random currents of a wandering mind\\
Everything and nothing, whatever I can find\\
A chaotic rhythm, a digital stream\\
No thematic prison, no restricted dream

\songpart{Outro}
Feel free to wander in this open space\\
No themes to guide you, no restricted place\\
Just the industrial pulse and the jazz that flies\\
Underneath the static of the digital skies
\annotation{Fade Out}
\annotation{End}
\end{lyrics}

\songtitle{Twenty Ten Scripts}

\tag*{2010}\tag{afro-pop}\tag{original-source}
\begin{notes}
Lyrics by Suno based on a summary of Carlos Thompson's 2010 blog production, written under his own name -- essentially a year of post titles.
\end{notes}
\begin{style}
Wassoulou Afro-pop, earthy kamelen’goni patterns, ngoni call-and-response, djembe and calabash polyrhythms, warm bass, handclaps, bright female lead with communal backing vocals, mid-tempo dance groove, verses that feel intimate and observant, chorus opening into celebratory layered vocals, organic live-room percussion, clear spacious mix, subtle modern low end
\end{style}
\begin{lyrics}
\songpart{Intro}
Hey ya, hey ya\\
Chlewey, open the book\\
Twenty ten, still under lock

\songpart{Verse 1}
February carried a passport with no country name\\
Japanese reflections leaned against the windowpane\\
March voted green, then freedom took the stand\\
A borrowed word can burn like dust across the land

April, mama, look at me, I counsel birds online\\
May brought Santos to the table, choice by choice in line\\
A time to choose, a road beneath the feet\\
The year kept many questions hidden underneath

\songpart{Chorus}
Twenty ten scripts, tied with a red thread\\
Pages in the cupboard, words not ready yet\\
Sing what we can see, leave the rest asleep\\
History has a doorway, but the key is deep\\
Hey ya, hey ya, let the calabash reply\\
Twenty ten scripts, waiting for the right sky

\songpart{Verse 2}
July made a world from fantasy and science fire\\
Carlos Eugenio sailed aboard a stubborn desire\\
August asked, what truth is yours to carry home?\\
A rush without surrender, a heart that wants to roam

Of loves and tastes, I would like to say a name\\
But every little wish can change the shape of flame\\
September tangled citizens in one wide woven net\\
Pirates, beliefs, secret bloggers, hormones, no regret

\songpart{Chorus}
Twenty ten scripts, tied with a red thread\\
Pages in the cupboard, words not ready yet\\
Sing what we can see, leave the rest asleep\\
History has a doorway, but the key is deep\\
Hey ya, hey ya, let the calabash reply\\
Twenty ten scripts, waiting for the right sky

\songpart{Verse 3}
October worked its hardest, minute by minute pressed\\
A tyrant wears a wristwatch, never lets you rest\\
November brought decisive moments to the floor\\
Then tragedy followed tragedy, and asked for more

Some things arrive like arrows, striking through the heart\\
Some hopes become a lesson before they can depart\\
If differences can soften, if the heavy sorrow bends\\
A hidden page may open when the waiting finally ends

\songpart{Bridge}
Do not publish every wound while it is still warm\\
Do not name the thunder while it gathers in the storm\\
Let the living ink breathe, let the quiet take its turn\\
What is meant to meet the world will find the day to burn

\songpart{Final Chorus}
Twenty ten scripts, tied with a red thread\\
Pages in the cupboard, words not ready yet\\
Sing what we can see, leave the rest asleep\\
History has a doorway, but the key is deep\\
Hey ya, hey ya, let the calabash reply\\
Twenty ten scripts, waiting for the right sky

\songpart{Outro}
Chlewey, keep the book\\
Keep the flame, keep the look\\
When the season says begin\\
Let the hidden pages sing
\end{lyrics}


\songtitle{Twenty Eleven Scripts}

\tag*{2011}\tag{congolese-rumba}\tag{original-source}
\begin{notes}
Lyrics by Suno based on a summary of Carlos Thompson's 2011 blog production, written under his own name -- essentially a year of post titles.
\end{notes}
\begin{style}
Congolese rumba, elegant mid-tempo sebene groove, interlocking clean electric guitars with bright highlife-style picking, supple bass, hand percussion and restrained drum kit, warm analog mix, expressive male lead vocal in English, gentle call-and-response backing vocals, verses intimate and conversational, chorus broad and melodic, extended guitar-led instrumental outro with tasteful rhythmic lift
\end{style}
\begin{lyrics}
\songpart{Intro}
Chlewey, Chlewey, turn the page slowly\\
Twenty eleven, the ink is still holy

\songpart{Verse 1}
On the Chlewey site, behind the quiet door\\
A year waits softly on the archive floor\\
Gender and gangs, the questions gather dust\\
Mobility, law, and the shape of trust\\
Children deserve a room where their voices can be heard\\
A little bird of doubt flies between every word

\songpart{Pre-Chorus}
Some histories knock, but the keeper says no\\
Not every river is ready to flow

\songpart{Chorus}
Twenty eleven scripts, folded in the drawer\\
Whatever matters most is not published anymore\\
Twenty eleven scripts, let the silence speak\\
A page can hold a fire while the paper stays asleep

\songpart{Verse 2}
People and avatars, frustration in the rain\\
Cyberbullying leaves a mark without a name\\
Copy rights, author rights, who owns the song?\\
Green dreams keep asking where we went wrong\\
A remix, a redo, a lighter road to take\\
New strategies rise from every old mistake

\songpart{Chorus}
Twenty eleven scripts, folded in the drawer\\
Whatever matters most is not published anymore\\
Twenty eleven scripts, let the silence speak\\
A page can hold a fire while the paper stays asleep

\songpart{Verse 3}
Letters from the city, walking block by block\\
Gendered violence, a wound behind the lock\\
Prisons in the mind, failed states in the rain\\
Wedding plans beside the windowpane\\
Free tiramisu, rabbits caught in view\\
Small strange moments still remember you

\songpart{Bridge}
And December brings an African reflection\\
Christmas thinking, searching for connection\\
The calendar closes, but the question remains:\\
Who keeps the truth when the ink leaves no name?

\songpart{Chorus}
Twenty eleven scripts, folded in the drawer\\
Whatever matters most is not published anymore\\
Twenty eleven scripts, let the silence speak\\
A page can hold a fire while the paper stays asleep

\songpart{Outro}
Chlewey, Chlewey, the night is wide\\
Those hidden pages have a pulse inside\\
Twenty eleven, we will wait our turn\\
Some stories stay quiet till they are ready to burn
\end{lyrics}

% no new lyrics published as Wyomee


\end{multicols*}

\chapter{Midlife Thoughts II: Trembling}
\albumcover{covers/internal/midlife-thoughts-ii-trembling}
Before \emph{the grinder} had a name, it had this: a moth circling a flame it knows will kill it, a mind that won't stop trembling, a heart that keeps time like a countdown. \emph{Temblando}, \emph{Skälvning} and \emph{Trembling Blues} are the same anxious pulse translated across languages and genres, years before therapy or diagnosis gave the feeling a shape to hold.
\clearpage\begin{multicols*}{4}
\songtitle{Polilla en las flamas}

\tag*{2010}\tag{christian-music}\tag{spanish-lyrics}
\begin{notes}
Lyrics by Suno, based on a free poem from 2010
\end{notes}
\begin{style}
Contemporary Christian Music with electro-folk and psychedelic trance pulse, mid-tempo and hypnotic; verse rides acoustic picking over soft hand percussion and a pulsing synth undercurrent, pre-chorus strips to close voice and tremolo bass, chorus blooms with stacked harmonies and a chantable lift. Lead vocal stays intimate and trembling with doubles on key lines, airy delay throws on the hook word, reverse swells into each chorus, and a faint tambourine shimmer. Mix is wide, luminous, and close-mic warm., contemporary, electro, folk, christian music, psychedelic trance
\end{style}
\begin{lyrics}
\songpart{Verse 1}
Siento en mi pecho\\
palpitos prohibidos\\
una vida, un vivir\\
que se deshace

Lo que daba por sentado\\
se me desvanece\\
y me acelero\\
juego a ser otro

A romper mis cadenas\\
a dejarme atraer\\
por una miel verde marino\\
que me llama más

\songpart{Pre-Chorus}
Y tiemblo\\
me acelero y tiemblo\\
tiemblo

Porque no sé mi juego\\
porque pierdo el pulso\\
porque se me escapa\\
lo que fui ayer

\songpart{Chorus}
Tiemblo, tiemblo\\
no sé mi juego\\
Tiemblo, tiemblo\\
y se me va el cielo

Tiemblo, tiemblo\\
no soy mi sombra\\
Tiemblo, tiemblo\\
y arde mi nombre

\songpart{Verse 2}
Quiero correr\\
como polilla en flamas\\
mi verdugo, mi corazón\\
mi amuleto

Y no sé cortar\\
las cadenas del pecho\\
porque el juego de la vida\\
ya lo voy perdiendo

Se encogen mis dedos\\
sobre el teclado\\
mientras pienso si decir\\
lo que no me atrevo

No sé escuchar tu voz\\
ni ver el rostro\\
de este juego que me mira\\
sin misericordia

\songpart{Pre-Chorus}
Y tiemblo\\
porque todo se desbarata\\
tiemblo

Ni con mi oponente\\
ni con la tribuna\\
ni conmigo mismo\\
puedo ser claro

\songpart{Chorus}
Tiemblo, tiemblo\\
no sé mi juego\\
Tiemblo, tiemblo\\
y se me va el cielo

Tiemblo, tiemblo\\
no soy mi sombra\\
Tiemblo, tiemblo\\
y arde mi nombre

\songpart{Bridge}
Si me llamas\\
que no me halle huyendo\\
si me nombras\\
que me encuentre aquí

No me dejes\\
en mi propia niebla\\
hazme nuevo\\
aunque duela así

\songpart{Final Chorus}
Tiemblo, tiemblo\\
no sé mi juego\\
Tiemblo, tiemblo\\
y se me va el cielo

Tiemblo, tiemblo\\
no soy mi sombra\\
Tiemblo, tiemblo\\
y arde mi nombre
\end{lyrics}


\songtitle{Temblando}

\tag*{2010}
\begin{notes}
\end{notes}
\begin{style}
blues, chamber blues, 132 BPM, swung 16th groove, baritone vocals, intimate delivery, Colombian seseo, upright bass, smoky Hammond organ, chamber strings, brushed snare, tape saturation, plate reverb, minor key, bass drops, tense build-release, vintage 70s mix, nocturnal ache
\end{style}
\begin{lyrics}
\songpart{Verse 1}
Temblando.\\
Siento un corazón en mi pecho con pálpitos prohibidos.\\
Una vida, un vivir, que se deshace.\\
Lo que daba por sentado se desvanece.\\
Y me acelero.

\songpart{Pre-Chorus}
Juego a ser otro, a romper mis cadenas.\\
Me dejo atraer por una miel verde-marino.\\
Me dejo llevar hasta la muerte\\
como una polilla en las flamas.

\songpart{Chorus}
Mi verdugo.\\
Mi corazón, mi amuleto y mi verdugo.\\
Y tiemblo.\\
Me acelero y tiemblo.\\
Tiemblo.

\songpart{Verso 2}
Tiemblo porque no sé cuál es mi juego.\\
Tiemblo porque perderé la inocencia,\\
perderé la sensatez,\\
perderé mi vida.\\
Tiemblo porque no puedo ser claro,\\
ni con mi oponente, ni con la tribuna,\\
ni conmigo mismo.

\songpart{Chorus}
Y tiemblo.\\
Me acelero y tiemblo.\\
Tiemblo.

\songpart{Bridge}
Porque se me escapa la vida\\
en pálpitos prohibidos de mi corazón.\\
Porque todo se desbarata.\\
Se encogen mis dedos sobre el teclado\\
mientras pienso si decir lo que no me atrevo.\\
No sé si escuchar su voz\\
o si ver el rostro de mi juego.

\songpart{Final Chorus}
Porque la verdad… no me atrevo a cortar las cadenas.\\
Porque el juego de la vida ya lo voy perdiendo.\\
Porque no soy yo…\\
ni mi fantasma.\\
Y tiemblo.\\
Me acelero y tiemblo.\\
Tiemblo.
\end{lyrics}


\songtitle{Vulnerable}

\tag*{2010}\tag{dubstep}
\begin{notes}
\end{notes}
\begin{style}
cinematic dubstep, industrial rock, half-time drops, tense swing, distorted electric guitars, sub-bass wobble, synth brass stabs, tom tom fills, gated snare, plate reverb, tape saturation, layered chorus hits, breakdown contrast, dramatic escalation, 128 BPM, anthemic defiance, deep bass baritone voclas, buried vocals, vocal texture, heavy reverb, continous flow, no gaps, breathless pace
\end{style}
\begin{lyrics}
\songpart{Intro - instrumental, 16 beats}

\annotation{Vocalizing under instrumental, melodic, sung}\\
No me gusta ser vulnerable.

\annotation{Continuous vocal delivery}\\
No me gusta sentirme solo\\
por esperar de las demás personas\\
más de lo que ellas pueden dar.

\annotation{Continuous vocal delivery}\\
Sentir que me ignoran,\\
que no me toman en cuenta,\\
que si no me avispo y me pego\\
oportunamente a un parche\\
tendré que almorzar solo.

\annotation{Continuous vocal delivery}\\
No me gusta tener que rogar\\
para ser tenido en cuenta,\\
ni sufrir porque mis ruegos\\
no son escuchados —\\
bien sea que haya rogado en silencio\\
o en voz alta.

\annotation{Continuous vocal delivery}\\
Los demás me afectan con lo que dicen\\
y lo que hacen\\
y con lo que callan.

\annotation{spoken, continuous flow, no gaps}\\
Soy vulnerable a todo este drama social que me avasalla.\\
Que me atropella.
\end{lyrics}


\songtitle{Aconteceres}

\tag*{2026}
\begin{notes}
\end{notes}
\begin{style}
piano, strings, brass, and a harp, at a high-tempo 6/8 phrases in D-minor for verses with 4/4 F-major choruses, brass intro and outro, and a violin solo bridge. sporadic filling dummy words and replacives and gospel clapping
\end{style}
\begin{lyrics}
\songpart{Verse 1}
Me dijeron que no debería\\
pedir perdón por mis palabras,\\
y pienso hacer caso.\\
Hay muchas cosas aún por decir.

\songpart{Pre-chorus}
Muchas cosas en el tintero.\\ 
Pero, tal vez, es silencio\\
lo que necesitamos ahora.

\songpart{Chorus}
Una mente sin aclarar\\
es un revólver cargado\\
y muchas personas\\
terminarán sufriendo.

\songpart{Bridge}
\annotation{spoken, violin breaks between phrases}\\
Qué feo es sentir que no se ama a quien se juró amar.

Es más difícil encontrar un abrazo memorable que un beso memorable.

…mis fantasías, una tormenta\\
Las óperas se burlan de mi pellejo 

Estás demasiado esporádica para mi gusto

Y volvió… lloviendo en todas partes\\
y escampando en tu mirada

¿Cómo se le hace mantenimiento a una cabeza recién comprada que no aparece?

Y me reiré, una vez más, de mis absurdos miedos.

soy mi propio álter ego…

Todavía, a lo lejos, hay una leve esperanza que dice que no lo hagas. ¿Y para qué creerle si nunca ha tenido la razón?

El mundo gira y gira y mi mente mira y mira y en tus ojos, cual un mar, mi alma expira.

\songpart{Chorus}
Una mente sin aclarar\\
es un revólver cargado\\
y muchas personas\\
terminarán sufriendo.
\end{lyrics}


\songtitle{The Loaded Mind}
\tag*{2010}

\begin{notes}
\end{notes}
\begin{style}
Industrial rock with techno house urgency in the verses, post-punk blues in the bridge and outro; mid-tempo driving pulse with tape saturation, clipped claps, rubbery bass, metal-leaning riffs, chamber strings, smoky Hammond, bowed lows, filtered choir pads, gospel lift, and half-time wobble bass. Verse rides bleak 4/4 motion; chorus slams harder and wider; bridge drops into bluesy space with upright bass and fiddle. Close-mic narrative vocal with breathy fragments, stacked doubles on the hook, short delay throws, choir swells, and gritty plate reverb. Bright punchy mix, dark and atmospheric., natural, warm, slow, blues, techno, bright, emotional, industrial rock, gospel, moody, rhythmic, deep, metal, post-punk, low
\end{style}
\begin{lyrics}
\songpart{Verse 1}
They told me not to beg for mercy for my tongue.\\
And I will heed that word, though it leaves me stung.\\
There is so much still waiting in the ink to run.\\
So many things not yet begun, not yet unsung.

\songpart{Chorus}
But maybe now… we need the still-ness.\\
Maybe now… we need the si-lence.

(But maybe now…) ((we need the stillness.))
(Maybe now…) ((we need the silence.))

\songpart{Verse 2}
The mind that has not cleared its chamber yet—\\
Is just a revolver that the hand for-gets.\\
And ev-ery thought un-aimed will find its mark.\\
And many souls will shatter in the dark.

\songpart{Chorus}
But maybe now… we need the still-ness.\\
Maybe now… we need the si-lence.

(But maybe now…) ((we need the stillness.))
(Maybe now…) ((we need the silence.))

\songpart{Bridge}
\annotation{Slower, quieter, almost whispered}\\
They told me not to say "forgive me" for my words—\\
And I will keep that counsel, though it hurts.\\
But if I speak before my mind is clear,\\
The shot goes off, and those I love will hear.\\
So I hold back. I bite my tongue.\\
Because the damage, once it's done, is done.\\
And many will suffer—not from my hate—\\
But from my failure to wait.

\songpart{Chorus}
That's why we need… the holy still-ness.\\
That's why we need… the heavy si-lence.

(That's why we need…) ((the holy stillness.))
(That's why we need…) ((the heavy si-lence.))

\songpart{Outro}
\annotation{Fading, spoken-sung, like a quiet confession}
(So I will not speak…)
(Though the ink is deep…)
(For an unclear mind…)
(Is a gun that does not sleep…)
……(and many will weep)……
\end{lyrics}


\songtitle{The Loaded Mind (defiant)}
\tag*{2010}

\begin{notes}
\end{notes}
\begin{style}
Industrial rock with techno house urgency in the verses, post-punk blues in the bridge and outro; mid-tempo driving pulse with tape saturation, clipped claps, rubbery bass, metal-leaning riffs, chamber strings, smoky Hammond, bowed lows, filtered choir pads, gospel lift, and half-time wobble bass. Verse rides bleak 4/4 motion; chorus slams harder and wider; bridge drops into bluesy space with upright bass and fiddle. Close-mic narrative vocal with breathy fragments, stacked doubles on the hook, short delay throws, choir swells, and gritty plate reverb. Bright punchy mix, dark and atmospheric., natural, warm, slow, blues, techno, bright, emotional, industrial rock, gospel, moody, rhythmic, deep, metal, post-punk, low
\end{style}
\begin{lyrics}
\songpart{Intro}
\annotation{A low, pulsing drumbeat mimicking a heartbeat}

\songpart{Verse 1}
The voice in-side said, "Hold your tongue, don't say a word."\\
The war you fight is best un-sung, their plea is heard.\\
But I have sworn to not a-polo-gize for what is stirred.\\
The ink is dry, the pen is mine, I'll keep my sword.

\songpart{Chorus}
So I won't (won't)\\
say I'm sor-ry. (sor-ry)\\
The fire (fire)\\
still burns in me. (in me)

So I won't\\
say I'm sor-ry.\\
The fire\\
still burns in me.

\songpart{Verse 2}
There's so much left to carve in stone, so much to tell.\\
The well is deep, I'm not a-lone, I know this hell.\\
But in the bot-tle, words are drowning in the swell.\\
And yet the qui-et seems to suit this cit-a-del.

\songpart{Chorus}
So I won't (won't)\\
say I'm sor-ry. (sor-ry)\\
The fire (fire)\\
still burns in me. (in me)

So I won't\\
say I'm sor-ry.\\
The fire\\
still burns in me.

\songpart{Bridge}
\annotation{Slower, more deliberate rhythm, like a warning chant}\\
Perhaps it's stillness that we crave,\\
The silence of a deeper wave.\\
A mind un-clear, a chamber spun,\\
A loaded re-volver, aimed at everyone.

The trigger's thought, the barrel's pride,\\
And many souls will run and hide.\\
The bullet's name is "un-said dread,"\\
And many will end up bleeding red.

\songpart{Chorus}
So I won't (won't)\\
say I'm sor-ry. (sor-ry)\\
The fire (fire)\\
still burns in me. (in me)

No I won't\\
beg for par-don.\\
Just watch\\
what you can't see

\songpart{Outro}
\annotation{Soft, fading, with a gentle, irregular heartbeat}
(The silence falls… the load is laid…)
(But many hearts… will still be frayed…)
(I hold my words… I hold my breath…)
(And wait for what comes next…)

……(si-lence)……
\end{lyrics}


\songtitle{Trembling Blues}

\tag*{2010}\tag{blues}
\begin{notes}
Lyrics by Claude based on a work by Carlos Thompson from 2010 (posted under Enwigh Peady), adapted in close interactive collaboration.

Check \emph{Temblando} (Spanish version)
\end{notes}
\begin{style}
slow minor blues, 12-bar blues, 76 BPM, male baritone vocals, tremolo electric guitar, fretless bass slides, Wurlitzer organ, brushed snare backbeat, ghost-note hi-hats, tape saturation, plate reverb, half-time turnaround, sparse spoken outro, remorseful confession, dangerous intimacy, smoky late-night mood
\end{style}
\begin{lyrics}
\songpart{Verse 1}
I got a heartbeat knockin' like a forbidden sin,\\
Yeah, a heartbeat knockin' like a forbidden sin,\\
And the life I took for granted—it's just caving in.

\songpart{Verse 2}
I play at bein' somebody I ain't never been,\\
Lord, I play at bein' somebody I ain't never been,\\
Just to taste that green-sea honey, then the flames come lickin' in.

\songpart{Verse 3}
That flame is my judge, my keeper, and my grave,\\
Yeah, that flame is my judge, my keeper, and my grave,\\
And I shake because I don't know which one of me I'll save.

\songpart{Verse 4}
I tremble 'cause I'm losin' my innocent mind,\\
Oh, I tremble 'cause I'm losin' my innocent mind,\\
And the truth I owe my opponent—I can't speak it, I'm half-blind.

\songpart{Verse 5}
My fingers curl up on these keys, I can't hit send,\\
Yeah, my fingers curl up on these keys, I can't hit send,
'Cause I don't know if I'm facin' my lover or my end.

\songpart{Verse 6}
I'm too scared to cut the chain and let it fall,\\
Lord, I'm too scared to cut the chain and let it fall,\\
So I'm neither my own ghost—nor myself at all.

\songpart{Outro}
\annotation{spoken/sung over a single, trembling guitar chord}\\
And I tremble...\\
Lord, I tremble...
'Cause this life game I'm playin'—\\
I'm already losin' it all.
\end{lyrics}


\songtitle{Pálpitos prohibidos}

\tag*{2010}\tag{electroswing}\tag{spanish-lyrics}
\begin{notes}
Lyrics by Grok based on a work by Carlos Thompson from 2010 (posted under Enwigh Peady), adapted in close interactive collaboration.

Check \emph{Temblando} (Spanish version)
\end{notes}
\begin{style}
latin pop, electroswing, male tenor vocal, Spanish lyrics, melodic seseo, candy pop hook, booming kick, syncopated brass stabs, accordion riffs, darting synth bass, piano ostinato, handclap chorus, plate reverb, tape saturation, sidechain pumping, call-and-response backing vocals, 124 BPM, explosive chorus, bittersweet confession
\end{style}
\begin{lyrics}
\songpart{Verse 1} \annotation{Low, breathy, confessional}\\
Siento un latido que no debió nacer,\\
una vida que se empieza a deshacer.\\
Lo que daba por cierto, ya no está,\\
y en mi pecho empieza a temblar.

\songpart{Pre-Chorus} \annotation{Building tension, a little more rhythm}\\
Juego a ser otro, a romper la prisión,\\
a probar la miel de tu voz.\\
Pero sé que esta llama me va a consumir,\\
como polilla que olvida su fin.

\songpart{Chorus} \annotation{Explosive, catchy, the "hook"}\\
¡Y tiemblo! (Temblo)\\
Porque no sé cuál es mi juego.\\
¡Y tiemblo! (Temblo)\\
Porque voy perdiendo el suelo.\\
Mi verdugo, mi amuleto, mi corazón—\\
late fuerte y me pide traición.\\
¡Y tiemblo, tiemblo!\\
En pálpitos prohibidos... me desvanezco.

\songpart{Verse 2} \annotation{More rhythmic, almost spoken}\\
Tiemblo porque no puedo ser claro,\\
ni con el que juega, ni con el extraño,\\
ni con este espejo que no quiere ver\\
que ya no sé ni quién debo ser.

\songpart{Pre-Chorus} \annotation{Same as before, but with more desperation}\\
Mis dedos se encogen sobre el teclado,\\
¿digo la verdad o guardo el pecado?\\
¿Escucho su voz o miro el tablero?\\
Si corto las cadenas, me muero.

\songpart{Chorus} \annotation{Full power, with backing vocals}\\
¡Y tiemblo! (Temblo)\\
Porque no sé cuál es mi juego.\\
¡Y tiemblo! (Temblo)\\
Porque voy perdiendo el suelo.\\
Mi verdugo, mi amuleto, mi corazón—\\
late fuerte y me pide traición.\\
¡Y tiemblo, tiemblo!\\
En pálpitos prohibidos... me desvanezco.

\songpart{Bridge} \annotation{The climax—music drops to piano or synth, then builds}\\
La verdad, no me atrevo...\\
a cortar lo que me duele.\\
El juego de la vida,\\
ya lo voy perdiendo,\\
y no soy yo,\\
ni mi fantasma,\\
solo un eco que se abrasa.

\songpart{Final Chorus} \annotation{Key change or ad-libs, explosive ending}\\
¡Y tiemblo! (¡Tiemblo!)\\
Porque todo se desbarata.\\
¡Y tiemblo! (¡Tiemblo!)\\
Porque el miedo me maltrata.\\
Mi verdugo, mi amuleto, mi corazón—\\
late fuerte, ¡es mi perdición!\\
¡Y tiemblo, tiemblo!\\
En pálpitos prohibidos...
(se desvanece en un susurro)
...tiemblo.
\end{lyrics}


\songtitle{Temblando}

\tag*{2010}\tag{trap}
\begin{notes}
Lyrics by Claude based on a work by Carlos Thompson from 2010 (posted under Enwigh Peady), adapted in close interactive collaboration.

Check \emph{Temblando} (Spanish version)
\end{notes}
\begin{style}
house trap, 124 BPM, 6 minute structure, spoken word confession, diva chop hooks, octave shifted tremble, deep Moog sine bass, Dm9 piano stabs, stuttered sample hits, filtered kick build, four on the floor pulse, swingless hi hats, plate reverb, ping pong delay, tape saturation, 4 AM club, haunted release
\end{style}
\begin{lyrics}
\songpart{Intro}
\annotation{Filtered kick drum, faint piano chord, sub-bass rumble}
\annotation{Spoken, low-fi, over the build}\\
Siento un corazón...\\
palpitos prohibidos...\\
palpitos prohibidos...\\
palpitos prohibidos...

\songpart{Verse 1}
\annotation{Vocal chops, diva-style, looped}\\
I tremble—\\
I tremble—\\
My heart, my amulet, my executioner—\\
I tremble—\\
I tremble—\\
My heart, my amulet, my executioner—

\songpart{Breakdown}
\annotation{Beat drops out, only bass pulse and reverberated whisper}
\annotation{Spoken, urgent, close-mic}\\
I play at being another.\\
Break my chains.\\
That green-sea honey—\\
pull me in—\\
like a moth to the flame—
\annotation{Beat crashes back in}

\songpart{Chorus}
\annotation{Full groove, piano stab, soaring chop}\\
Tremble 'cause I don't know the game—\\
Tremble 'cause I'm losing my name—\\
Losing my senses—\\
Losing my life—\\
Tremble in the strobe light—

\songpart{Bridge}
\annotation{Filter sweeps, beat strips to hi-hats and kick}
\annotation{Rhythmic spoken word, locked to the kick}\\
Fingers on the keys—\\
Should I speak? Should I freeze?\\
Can't see my opponent—\\
Can't see the bleachers—\\
Can't see my own features—

\annotation{Build}
\annotation{Synth rises, vocal layers stack}\\
I tremble—\\
I accelerate—\\
I tremble—\\
I disintegrate—

\annotation{Drop}
\annotation{Full power, bass weight, looped diva cry}\\
Everything falls apart—
(Everything falls apart—)\\
Everything falls apart—
(Everything falls apart—)\\
My chains—\\
My ghost—\\
My heart—\\
My host—

\songpart{Outro}
\annotation{Beat fades, only reverb and a single spoken line}
"No soy yo... ni mi fantasma..."

\annotation{Final kick drum hits once, then silence.}
\end{lyrics}


\songtitle{Skälvning}

\tag*{2010}\tag{nu-metal}\tag{swedish-lyrics}
\begin{notes}
Lyrics by Grok based on a work by Carlos Thompson from 2010 (posted under Enwigh Peady), adapted in close interactive collaboration.

Check \emph{Temblando} (Spanish version)
\end{notes}
\begin{style}
swedish nu metal, Drop A tuning, sludge guitar chugs, 7 string palm muting, snare heavy verses, double bass chorus, hip hop rap delivery, soaring clean pre chorus, guttural chorus screams, delay soaked spoken word, single note piano intro, ambient guitar feedback, half time breakdown, low end dropouts, claustrophobic verses, explosive release, plate reverb, tape saturation, panic attack energy
\end{style}
\begin{lyrics}
"Skälvning" (Trembling)

\songpart{Intro – Clean, ambient guitar feedback, single repeating piano note. Distant, echoed whisper}

Jag känner hjärtat...\\
förbjudna slag...\\
förbjudna slag...

\songpart{Verse 1 – Rapped, syncopated, rhythmic, almost monotone over a chugging 7-string riff}

Ett hjärta i bröstet – förbjudna slag,\\
ett liv som löser upp sig, dag för dag.\\
Det jag tog för givet – borta, spårlöst,\\
jag gasar, jag fejkar, jag blir ett spöke.

Jag leker att vara nån annan –\\
krossar mina kedjor med skakande händer.\\
Lockas av honung – grön som havet,\\
dras mot flamman, mot min gravplats.

\songpart{Pre-Chorus – Tension builds, hi-hats double-time, riff opens up}

Min bödel...\\
Mitt hjärta...\\
Min amulett...\\
Och min bödel...

\songpart{Chorus – Full distorted roar, screamed, guttural}

Och Jag Skälver!\\
Jag Gasar – Och Skälver!\\
Vet Inte Vad Jag Spelar För Spel!\\
Jag Förlorar Oskuld, Förnuft – Mitt Liv I En Källa!

Jag Skälver För Att Jag Inte Kan Vara Tydlig –\\
Inte Mot Min Motståndare, Inte Mot Läktaren,\\
Inte Mot Mig Själv!

\songpart{Verse 2 – Softer, clean guitar, half-whispered, almost pleading}

Min livsflod rinner ut –\\
i förbjudna slag från mitt bröst.\\
Allt faller sönder – jag ser det,\\
mina fingrar krampar över tangenterna.

Ska jag säga det jag inte vågar?\\
Ska jag höra hennes röst – eller se ansiktet\\
på det spel jag spelar?\\
Sanningen: jag vågar inte klippa kedjorna.

\songpart{Pre-Chorus – Heavier, tempo slows, bass drops}

Min bödel...\\
Mitt hjärta...\\
Min amulett...\\
Och min bödel...

\songpart{Chorus – Full power, layered screams, backing growls}

Och Jag Skälver!\\
Jag Gasar – Och Skälver!\\
Vet Inte Vad Jag Spelar För Spel!\\
Jag Förlorar Oskuld, Förnuft – Mitt Liv I En Källa!

Jag Skälver För Att Jag Inte Kan Vara Tydlig –\\
Inte Mot Min Motståndare, Inte Mot Läktaren,\\
Inte Mot Mig Själv!

\songpart{Bridge / Breakdown – Beat drops to half-time. Feedback hums. Spoken word, intimate, Swedish, into the mic}

Livets spel...\\
Jag förlorar det redan.\\
För att jag inte är jag –\\
och inte heller mitt spöke.

\annotation{Silence.}

\annotation{Whisper, cracked voice:}
...jag vågar inte...

\songpart{Outro – Distorted, chaotic, riff detunes, screams fracture into static}

Skälver!\\
Skälver!\\
Allt – Faller – Sönder!

\annotation{Final, broken whisper over clean guitar:}
...inte jag... inte mitt spöke...

(Heavy, single palm-muted chord. Ring out. Fade.)
\end{lyrics}


\songtitle{Cabeza Perdida}

\tag*{2010}\tag{synthwave}\tag{spanish-lyrics}
\begin{notes}
Lyrics by SUNO based on a metaphoric question by Carlos Thompson in Formspring.
\end{notes}
\begin{style}
Retrowave synthwave at midtempo with gated snare, pulsing arpeggios, and dark neon bass; verses stay sparse with tape hiss and a lonely DX-style lead, pre-chorus filters up with risers and a clocktick pulse, chorus opens into wide pads and chopped vocal doubles, bridge strips to bass and whispered mono lead before final lift, Close-mic male vocal, dreamy and eerie, with delayed hook repeats, reversed swells, and glossy but shadowed mix
\end{style}
\begin{lyrics}
\songpart{Verse 1}
La compré en la feria gris\\
envuelta en plástico azul\\
venía con un recibo\\
y una marca en el mentón

La dejé sobre la mesa\\
junto al pan y la sal\\
pero al volver de la puerta\\
ya no quiso despertar

\songpart{Pre-Chorus}
¿Quién la vio pasar?\\
¿Quién cerró el lugar?\\
Hay un zumbido en el pasillo\\
y mi pulso va detrás

\songpart{Chorus}
Cabeza perdida\\
dime dónde estás\\
Cabeza perdida\\
no te sé encontrar\\
Cabeza perdida\\
vuelve a encajar\\
Cabeza perdida\\
te tengo que arreglar

\songpart{Verse 2}
Le busqué entre los cajones\\
entre llaves y metal\\
bajo el vaso, bajo el mapa\\
bajo el polvo del portal

Pregunté por la vecina\\
por el gato del andén\\
todos miran hacia arriba\\
como si supieran bien

\songpart{Pre-Chorus}
¿Quién la vio pasar?\\
¿Quién cerró el lugar?\\
Hay un zumbido en el pasillo\\
y mi pulso va detrás

\songpart{Chorus}
Cabeza perdida\\
dime dónde estás\\
Cabeza perdida\\
no te sé encontrar\\
Cabeza perdida\\
vuelve a encajar\\
Cabeza perdida\\
te tengo que arreglar

\songpart{Bridge}
En la cinta del recuerdo\\
se oye un paso de cristal\\
una sombra con mi nombre\\
la guardó en la oscuridad

Si regresa, no la mires\\
déjala girar así\\
que la noche tiene manos\\
y hoy me vino a elegir

\songpart{Chorus}
Cabeza perdida\\
dime dónde estás\\
Cabeza perdida\\
no te sé encontrar\\
Cabeza perdida\\
vuelve a encajar\\
Cabeza perdida\\
te tengo que arreglar
\end{lyrics}


\songtitle{El Mundo Gira}

\tag*{2010}
\begin{notes}
Lyrics by SUNO based on two line poem by Carlos Thompson (posted under Enwigh Peady).
\end{notes}
\begin{style}
Bluegrass uplifting with fast-picked acoustic guitar, bright mandolin chops, fiddle doubles, upright bass walk, and a handclap lift; verse rides a bouncy train-beat, pre-chorus opens on stacked harmony, chorus hits with gang vocals and a soaring fiddle line, Lead vocal stays warm and close in verses, doubled on hooks with airy thirds and short delay throws on the last word, Banjo fills, quick risers into choruses, and a final lift with stomps and claps, Crisp, sunny mix with live-room sparkle, uplifting, bluegrass
\end{style}
\begin{lyrics}
\songpart{Verse 1}
Salgo al alba sin temor,\\
con el polvo en mis botas,\\
y el río canta alrededor\\
de las piedras y las hojas.\\
Tu nombre viene en el viento,\\
como pan recién cortado,\\
y yo sonrío de contento\\
con el corazón despierto.

\songpart{Pre-Chorus}
Y aunque el mundo siga andando,\\
yo no pierdo el paso,\\
si tu risa va llamando\\
se me aclara el remanso.

\songpart{Chorus}
El mundo gira y gira\\
y mi mente mira y mira\\
y en tus ojos, cual un mar,\\
mi alma expira.\\
El mundo gira y gira\\
y mi mente mira y mira\\
si te encuentro al caminar,\\
vuelvo a respirar.

\songpart{Verse 2}
Vi las nubes sobre el maizal,\\
como barcas de verano,\\
y en tu voz tan natural\\
se me enderezó la mano.\\
Ya no cargo la tormenta\\
cuando tu luz me alcanza,\\
porque tu ternura lenta\\
me devuelve la esperanza.

\songpart{Pre-Chorus}
Y aunque el mundo siga andando,\\
yo no pierdo el paso,\\
si tu risa va llamando\\
se me aclara el remanso.

\songpart{Chorus}
El mundo gira y gira\\
y mi mente mira y mira\\
y en tus ojos, cual un mar,\\
mi alma expira.\\
El mundo gira y gira\\
y mi mente mira y mira\\
si te encuentro al caminar,\\
vuelvo a respirar.

\songpart{Bridge}
Déjame quedarme aquí,\\
junto a tu paz sencilla,\\
que en tu abrazo aprendí\\
a florecer de rodillas.\\
Si la tarde cae despacio,\\
tú me guardas el calor,\\
y en tu pecho hallo espacio\\
para sembrar mi canción.

\songpart{Chorus}
El mundo gira y gira\\
y mi mente mira y mira\\
y en tus ojos, cual un mar,\\
mi alma expira.\\
El mundo gira y gira\\
y mi mente mira y mira\\
si te encuentro al caminar,\\
vuelvo a respirar.
\end{lyrics}


\songtitle{Jugando}

\tag*{2010}
\begin{notes}
Lyrics by SUNO based on a short post by Carlos Thompson (posted under Enwigh Peady).
\end{notes}
\begin{style}
Trance house anthem with a driving four-on-the-floor pulse, airy synth arpeggios, rolling sidechain bass, and a steady 128 BPM lift; verse rides a lean kick, muted clap, and filtered pad, pre-chorus opens with rising white-noise sweeps and chopped vocal echoes, chorus hits wide with stacked unison hooks and a soaring supersaw lead; close-mic Colombian Spanish vocal with doubled choruses, bright delay throws on the hook word, reversed swells into each drop, crisp and glossy stadium mix, trance
\end{style}
\begin{lyrics}
\songpart{Verse 1}
El nombre del juego\\
no lo conozco\\
tampoco sus reglas

Ni sus consecuencias\\
ni si vale\\
la pena entrar

No sé por qué\\
lo estoy jugando\\
pero aquí voy\\
sin mirar atrás

No sé quién cae\\
no sé quién gana\\
solo siento\\
la pulsación

\songpart{Pre-Chorus}
Y si me quemo\\
que sea fuerte\\
y si me pierdo\\
que sea total

Nada me pesa\\
todo me llama\\
todo me empuja\\
a continuar

\songpart{Chorus}
Jugando, jugando\\
estoy jugando\\
Jugando, jugando\\
sin soltar

Nada importa\\
e importa todo\\
Jugando, jugando\\
hasta el final

\songpart{Verse 2}
No sé cuánto\\
voy a apostar\\
pero mis manos\\
ya dijeron sí

Hay una apuesta\\
sobre la mesa\\
y mi destino\\
pide seguir

Si sale herido\\
mi corazón\\
que aprenda rápido\\
a resistir

Si sale roto\\
mi alrededor\\
yo sigo vivo\\
sigo aquí

\songpart{Pre-Chorus}
Y si me quemo\\
que sea fuerte\\
y si me pierdo\\
que sea total

Nada me pesa\\
todo me llama\\
todo me empuja\\
a continuar

\songpart{Chorus}
Jugando, jugando\\
estoy jugando\\
Jugando, jugando\\
sin soltar

Nada importa\\
e importa todo\\
Jugando, jugando\\
hasta el final

\songpart{Bridge}
Tal vez no sé\\
qué estoy buscando\\
tal vez el mundo\\
me está mirando

Pero si caigo\\
caigo bailando\\
con el pulso alto\\
y el alma en llamas

\songpart{Chorus}
Jugando, jugando\\
estoy jugando\\
Jugando, jugando\\
sin soltar

Nada importa\\
e importa todo\\
Jugando, jugando\\
hasta el final
\end{lyrics}


\songtitle{El juego}

\tag*{2010}\tag{ambient}\tag{original-lyrics}\tag{spanish-lyrics}
\begin{notes}
Original writings by Carlos Thompson in 2010, compiled and added music structure tags in 2026.

Remix with a version of \emph{Claim Ledger Drift}.
\end{notes}
\begin{style}
Organic ambient with slow breathy pulses, hand-played textures, and gently shifting drones in a patient 60 BPM flow; opening verse-like bed stays sparse with soft piano dust and field-recorded room tone, middle blooms into layered felt chords and bowed overtones, final section thins to a warm sub glow, Reversed swells, glassy chimes, tape hiss, and wide intimate mix, organic ambient, Colombian pronunciation
\end{style}
\begin{lyrics}
\songpart{Intro, spoken, soothing female mezzo voice, whispered, slow breathy delivery, Colombian pronunciation}
El nombre del juego, no lo conozco...\\
Tampoco sus consecuencias...\\
Ni si vale la pena jugarlo...\\
Tampoco estoy seguro por qué lo estoy jugando...\\
Ni quién puede salir lastimado.

\annotation{spoken}\\
Desconozco todavía cuánto estaré dispuesto a apostar.\\
Pero nada importa e importa todo.  Estoy jugando.

\songpart{Chorus, melodic, soothing, whispering, restrained, drop at the end}
El mundo gira y gira y mi mente mira y mira y en tus ojos, cual un mar, mi alma expira.

\songpart{Outro, spoken, slow}
Qué feo es sentir que no se ama a quien se juró amar.

\annotation{whispered, faster}\\
Estás demasiado esporádica para mi gusto
\end{lyrics}

\end{multicols*}

\chapter{Later Thoughts I: The Grinder}
\albumcover{covers/internal/later-thoughts-i-the-grinder}
The grinder finally gets its name: the alarm that goes off and gets silenced, the lightning-rod excuses, the paper that sits untouched on the desk for one more day. \emph{Time to Stop} closes the set two decades later, still chasing unfinished blueprints, still bargaining with the clock for one more hour.
\clearpage\begin{multicols*}{4}
\songtitle{Alarm Cabinet}

\tag*{2024}\tag{techno}\tag{social-media-seed}

\begin{notes}
Lyrics written by Suno, based on thread in X (Twitter) by Carlos Thompson.
\end{notes}

\begin{style}
Techno with low-pitch guttural male vocals, 132 BPM, heavy syncopated kick pulse, rattling hi-hats, thick sub-bass, and machine-room stabs; verse rides sparse drum grid and cold synth ticks, pre-chorus tightens with filtered noise and vocal distortions, chorus hits with doubled bass growls and gang shouts on the anchor line; add alarm bleeps, reversed metallic swells, and short delay throws on key words; mix is dark, glossy, and chest-rattling, low, techno
\end{style}

\begin{lyrics}

\songpart{Verse 1}
My phone screams red again\\
I swipe it blind\\
Three alarms, four alarms\\
Still I cut the line

You call to check my steps\\
I tell you I'm on track\\
I say "yeah, I'm getting there"\\
Then I never call you back

\songpart{Pre-Chorus}
I can look busy\\
I can look clean\\
Folders in a row\\
Like a perfect machine

But I fold in half\\
When nobody sees\\
Hide the delay\\
Wear the fake receipts

\songpart{Chorus}
I need guided hands\\
I need guided hands\\
I need guided hands\\
When I disappear (guided hands)

I need a room with rules\\
A face that stays\\
A place that knows my lies\\
And keeps me in the frame

\songpart{Verse 2}
I sit in bright white rooms\\
With a clipboard stare\\
I move the paper around\\
Like I meant to care

Tabs in a neat little line\\
Cursor parked and still\\
I fake the motion well\\
But I chase the will

I build little walls of proof\\
Just to feel okay\\
Then slip out through the crack\\
And waste another day

\songpart{Pre-Chorus}
I don't need a slogan\\
I don't need a speech\\
I need a held-closed door\\
And someone in reach

A group with names like mine\\
A guide who won't drift\\
Something steady enough\\
To catch my hands when I split

\songpart{Chorus}
I need guided hands\\
I need guided hands\\
I need guided hands\\
When I disappear (guided hands)

I need a room with rules\\
A face that stays\\
A place that knows my lies\\
And keeps me in the frame

\songpart{Bridge}
Maybe I'm not broken\\
Maybe I'm just lost\\
In a maze of yes and no\\
And a count of what it cost

So let me sit in truth\\
Let me say it plain\\
I keep building alibis\\
Then I wear their chains

\annotation{Final Chorus}\\
I need guided hands\\
I need guided hands\\
I need guided hands\\
When I disappear (guided hands)

I need a room with rules\\
A face that stays\\
A place that knows my lies\\
And keeps me in the frame

I need guided hands\\
I need guided hands\\
I need guided hands\\
And some way to stay (guided hands)

\end{lyrics}


\songtitle{A Drunk Without A Drop (Remix)}

\tag*{2026}\tag{indie-folk}\tag{ai-based}\tag{remix}\tag*{english-lyrics}

\begin{notes}
A remix that refreshes the original through updated production textures and rhythmic contrast. %
The lyrics begin with “I put off what I ought to do”, giving a quick sense of the song’s tone.
\end{notes}

\begin{style}
Indie folk ballad featuring a baritone male vocal. %
The arrangement centers on a fingerpicked acoustic guitar playing a steady eighth-note pattern. %
A clean electric guitar provides melodic counterpoint with light reverb and occasional slides. %
The bass guitar plays sustained root notes on the downbeats. %
A subtle percussion layer consists of a soft shaker and a dampened kick drum on beats one and three. %
The tempo is 74 BPM in 4/4 time, in the key of C major.
\end{style}

\begin{lyrics}

\songpart{Intro}
\annotation{fingerpicked acoustic guitar, soft shaker}

\songpart{Verse 1}
\annotation{male bass vocals} 
I put off what I ought to do\\
To numb the growing dread,\\
A heavy dose of doing naught\\
To calm my anxious head.

\songpart{Verse 2}
But failure breeds a sharper fear,\\
A trap I cannot flee—\\
I am a drunk without a drop,\\
undone by staying free.

\songpart{Pre-chorus}
The anxiety of functional adult livehood.\\
The anxiety of the trap to fail.

\songpart{Chorus}
An alcoholic who doesn't drink\\
(without a drop)\\
Procrastination, a numbing drug\\
(staying free)\\
I am a drunk without a drop,\\
undone by staying free.

\songpart{Verse 3}
The heavy chains of shifting gears\\
Are dragging at my chest,\\
The constant race to buy and build\\
Without a moment's rest.

\songpart{Pre-chorus}
The anxiety of functional adult livehood.\\
Heavy chains dragging at my chest.

\songpart{Chorus}
An alcoholic who doesn't drink\\
(without a drop)\\
Procrastination, a numbing drug\\
(staying free)\\
I am a drunk without a drop,\\
undone by staying free.

\songpart{Verse 4}
And so I turn the screen ahead\\
To flee the breaking storm,\\
A chosen dark to block the light\\
And keep the panic warm.

\songpart{Pre-chorus}
Endless projects nobody else cares.\\
Scrolling, daydreaming 'till the end.

\songpart{Chorus}
An alcoholic who doesn't drink\\
(without a drop)\\
Procrastination, a numbing drug\\
(staying free)\\
I am a drunk without a drop,\\
undone by staying free.

\annotation{guitar riff}

I am a drunk without a drop,\\
undone by staying free.

\annotation{fade}

\end{lyrics}


\songtitle{Modo Pararrayos}

Lyrics and style details by Suno after a prompt by Carlos Thompson of producing a techno-house song on a thread by Carlos Thompson published in X (former Twitter).

The original thread was bilingual, Suno chose Spanish lyrics.

\begin{notes}
\end{notes}

\begin{style}
Techno-house with driving four-on-the-floor kick, rolling bass pulses, clipped hi-hats, and a tense groove that feels locked and urgent; verse rides sparse drums and low synth stabs, pre-chorus strips to kick, bass, and filtered vocal fragments, chorus hits with a wide synth lift and chantable low male doubles. %
Add risers before each drop, reversed swells into the pre-chorus, and a metallic hit on the hook. %
Deep, dry, and punchy mix with a dark club sheen., electronic, theme, techno
\end{style}

\begin{lyrics}

\songpart{Verse 1}
Dicen que soy egoísta\\
por buscar un destello\\
pero bajo tanta presión\\
yo me apago por dentro

Más peso, más fuga\\
más ruido en la sangre\\
más culpa en el pecho\\
más ganas de escape

No me sale empujar\\
cuando todo me cae\\
me quedo inmóvil\\
hasta que el cielo se abre

\songpart{Pre-Chorus}
Y si me aprietan más\\
me cierro más\\
si me miran mal\\
me hundo más

No me pidas fuego\\
si ya estoy al borde\\
no me pidas calma\\
si me arde el nombre

\songpart{Chorus}
Modo pararrayos\\
modo pararrayos\\
que pase la tormenta\\
modo pararrayos

Modo pararrayos\\
modo pararrayos\\
yo no nací pa' aguantar\\
todo el peso de otros

\songpart{Verse 2}
Para ustedes es copa\\
para mí es la mente\\
un salto, una idea\\
un vértigo ardiente

A veces el cuerpo\\
también pide salida\\
una puerta cerrada\\
una chispa encendida

Pero si no hay escape\\
me quedo partido\\
con la vista fija\\
y el pulso dormido

Y encima me reclaman\\
como si fuera fácil\\
eso me sube el pulso\\
y me deja frágil

\songpart{Pre-Chorus}
Y si me apuran más\\
me rompo más\\
si me vigilan más\\
me freno más

No me cargues tu peso\\
sobre mi espalda\\
no me pidas respuestas\\
con la alarma alta

\songpart{Chorus}
Modo pararrayos\\
modo pararrayos\\
que pase la tormenta\\
modo pararrayos

Modo pararrayos\\
modo pararrayos\\
yo no nací pa' aguantar\\
todo el peso de otros

\songpart{Bridge}
Tengo que cambiar el chip\\
sí, tengo que cambiarlo\\
hacer que terminar\\
también me dé un salto

Que lo urgente no sea\\
una jaula de hierro\\
que cumplir con lo mío\\
también me prenda adentro

\annotation{Breakdown}\\
Baja la presión\\
baja la voz\\
deja que encuentre\\
mi propio motor

No quiero ser\\
tu poste fijo\\
yo quiero salir\\
del trueno y del grito

\songpart{Final Chorus}
Modo pararrayos\\
modo pararrayos\\
que pase la tormenta\\
modo pararrayos

Modo pararrayos\\
modo pararrayos\\
yo no nací pa' aguantar\\
todo el peso de otros

Modo pararrayos\\
modo pararrayos\\
si termina la lluvia\\
vuelvo a mí poco a poco

\end{lyrics}


\songtitle{Lightning Rod Mode}

\tag*{2026}\tag{techno-house}\tag{social-media-seed}

\begin{notes}
Lyrics and style details by Suno after a prompt by Carlos Thompson of producing a techno-house song \textbf{in English} on a thread by Carlos Thompson published in X (former Twitter).
\end{notes}

\begin{style}
Techno house with a driving four-on-the-floor pulse, rubbery bassline, clipped kick hats, and a tense build-release arrangement; verse rides sparse drums and a sub pulse, pre-chorus opens with filtered claps and rising synth haze, chorus hits with stacked male bass doubles and a chantable hook, bridge drops to a hollow kick and voiced breaths before the final lift. %
Deep male bass vocal with tight doubles, occasional shouted responses, delay throws on key lines, reverse swells into the drop, bright and punchy mix with a glossy low end., techno, theme, electronic
\end{style}

\begin{lyrics}

\songpart{Verse 1}
They call it selfish.\\
I call it survival.\\
When the pressure climbs,\\
my thoughts go tribal.

One more demand\\
and my hands go still.\\
Need a clean escape\\
just to keep my will.

\songpart{Pre-Chorus}
It’s empirical.\\
I break this way.\\
More weight on the chest,\\
more I drift away.

Cortisol high,\\
I pay that debt.\\
If I don’t get out,\\
I get stuck in place.

\songpart{Chorus}
Lightning rod mode.\\
Take the hit, take the hit.\\
Lightning rod mode.\\
I can’t live in it.

Give me a door.\\
Give me a spark.\\
Give me a way\\
through the dark.

Lightning rod mode.\\
I was built to move.

\songpart{Verse 2}
For you it’s the bottle,\\
for me it’s the page.\\
A sharp little thrill\\
to break the cage.

Could be a body,\\
could be a thought.\\
Could be one bright thing\\
that the day forgot.

But when you stack\\
the task and fear,\\
you kill the flame\\
and I disappear.

\songpart{Pre-Chorus}
It’s empirical.\\
I know the drill.\\
More pressure, more escape,\\
more need to heal.

And don’t add a face\\
that wants the proof.\\
That kind of watching\\
burns right through.

\songpart{Chorus}
Lightning rod mode.\\
Take the hit, take the hit.\\
Lightning rod mode.\\
I can’t live in it.

Give me a door.\\
Give me a spark.\\
Give me a way\\
through the dark.

Lightning rod mode.\\
I was built to move.

\songpart{Bridge}
I need a shift.\\
I need a bridge.\\
Make the finish line\\
feel like a lift.

Not more judgment.\\
Not more weight.\\
Don’t load the storm\\
and call it fate.

\songpart{Final Chorus}
Lightning rod mode.\\
Take the hit, take the hit.\\
Lightning rod mode.\\
I can’t live in it.

Give me a door.\\
Give me a spark.\\
Give me a way\\
through the dark.

Lightning rod mode.\\
I was built to move.\\
Lightning rod mode.\\
I need a better groove.

\end{lyrics}


\songtitle{Somebody Out There}

\tag*{2026}\tag{techno-house}\tag{social-media-seed}

\begin{notes}
Lyrics and style details by Suno after a prompt by Carlos Thompson of producing a techno-house song in English on a thread by Carlos Thompson published in Threads by Meta.
\end{notes}

\begin{style}
Techno house with a pulsing four-on-the-floor kick, rubbery offbeat bass, and clipped claps; verse rides sparse drums, sub pulses, and a deadpan bass vocal, pre-chorus opens with filtered chords and rising noise, chorus hits with a brighter synth stab and gang-voiced hook. %
Male bass vocal close-miked with short delay throws, doubled hook lines, and whispered ad-libs. %
Ear candy: reversed sweeps before each drop, tiny keypad clicks, and a warped voice chop between phrases. %
Bright, punchy club mix with a sleek low-end throb., electronic, techno, theme
\end{style}

\begin{lyrics}

\songpart{Verse 1}
I talk to the walls.\\
They don’t answer back.\\
The sink full of plates\\
keeps time in the dark.

My phone says names,\\
but none of them land.\\
I scroll through the list\\
with my thumb in the sand.

\songpart{Pre-Chorus}
I tried the house.\\
I tried the phone.\\
I tried the screen.\\
Still all alone.

Maybe there’s one\\
with my kind of need.\\
How do I find them?\\
How do I read?

\songpart{Chorus}
Somebody out there\\
(Somebody out there)\\
Somebody out there\\
(Somebody out there)

Say my name once\\
and stay for the night.\\
Somebody out there\\
feels like my type.

\songpart{Verse 2}
Online contacts,\\
AI chats,\\
chat app faces\\
and date app racks.

I sent the line.\\
I made the move.\\
Maybe I missed\\
what I should pursue.

Or maybe I didn’t.\\
Maybe I did.\\
A spark in the pile,\\
a key in the grid.

\songpart{Pre-Chorus}
I tried the house.\\
I tried the phone.\\
I tried the screen.\\
Still all alone.

Maybe there’s one\\
with my kind of need.\\
How do I find them?\\
How do I read?

\songpart{Chorus}
Somebody out there\\
(Somebody out there)\\
Somebody out there\\
(Somebody out there)

Say my name once\\
and stay for the night.\\
Somebody out there\\
feels like my type.

\songpart{Bridge}
\annotation{Breakdown}
If you’re in the noise,\\
send one clean sign.\\
Not every answer\\
has to be mine.

One real voice,\\
one steady thread.\\
Pull me from here\\
where the words all dead.

\songpart{Chorus}
Somebody out there\\
(Somebody out there)\\
Somebody out there\\
(Somebody out there)

Say my name once\\
and stay for the night.\\
Somebody out there\\
feels like my type.

\end{lyrics}


\songtitle{Soap Line Logic}

\tag*{2026}\tag{electronic-pop}\tag{original-work}

\begin{notes}
Original work from 2026.\\
Abstract electronic reflection on mundane ritual and meditative routine.
\end{notes}

\begin{style}
Techno-trance with dubstep drops, four-on-the-floor pulse, wobble bass under a melodic female lead; verse stays lean and mechanical, pre-chorus filters up with clipped handclaps and rising synth shimmer, chorus opens wide with octave doubles and chanty backing vocals, bridge strips to a half-time sub pulse and breathy topline before the final drop. %
Glassy arps, water-drip percussion, reverse swells into each transition, bright punchy mix with heavy low-end and clean vocal polish., relaxing, techno, dubstep
\end{style}

\begin{lyrics}

\songpart{Verse 1}
I leave the plates\\
For later, for later\\
My mind says “scroll”\\
My hands say “maybe”\\
The sink looks full\\
Like a small, dull mountain\\
But nothing breaks\\
If I wait a minute

\songpart{Pre-Chorus}
Then I turn on the tap\\
And the room changes shape\\
Hot water, cold steam\\
I find my pace\\
One bowl, one fork\\
One little fix\\
I line it up\\
And I click

\songpart{Chorus}
Doing the dishes, doing the dishes\\
(doing the dishes)\\
Back in the rhythm, back in the rhythm\\
(back in the rhythm)\\
Soap line logic\\
Soap line logic\\
A small, steady peace\\
In the rinse and the drip

\songpart{Verse 2}
I watch the grease\\
Slip off the spoon\\
Keep the sponge fresh\\
Keep the knife clean too\\
Stack the plates\\
Like a game I can solve\\
Tall ones first\\
Then the cups in a row\\
It’s funny how\\
A simple routine\\
Turns down the noise\\
Inside of me

\songpart{Pre-Chorus}
Then I turn on the tap\\
And the room changes shape\\
Hot water, cold steam\\
I find my pace\\
One bowl, one fork\\
One little fix\\
I line it up\\
And I click

\songpart{Chorus}
Doing the dishes, doing the dishes\\
(doing the dishes)\\
Back in the rhythm, back in the rhythm\\
(back in the rhythm)\\
Soap line logic\\
Soap line logic\\
A small, steady peace\\
In the rinse and the drip

\songpart{Bridge}
Maybe that’s the point\\
Of every little chore\\
Not to thrill me\\
Just to pull me to the floor\\
What I avoid\\
Can hold me in place\\
Till the static in my head\\
Starts to fade

\songpart{Drop}
Doing the dishes\\
Doing the dishes\\
(doing the dishes)\\
Doing the dishes\\
Back in the rhythm\\
Back in the rhythm\\
(soap line logic)\\
A small, steady peace\\
In the rinse and the drip

\songpart{Chorus}
Doing the dishes, doing the dishes\\
(doing the dishes)\\
Back in the rhythm, back in the rhythm\\
(back in the rhythm)\\
Soap line logic\\
Soap line logic\\
A small, steady peace\\
In the rinse and the drip

\end{lyrics}


\songtitle{Frictionless}

\begin{lyrics}
\songpart{Verse 1}
Screen glow at 2 AM, another endless feed\\
Fingers twitch, dopamine hits, but I don't feel a thing\\
Hyperfocus tunnel, building worlds that disappear\\
Then crash into nothing, another wasted year

\songpart{Pre-Chorus}
I taste the phantom salt of olives in the square\\
Hear laughter echoing, but I'm not even there

\songpart{Chorus}
I want the clink of glasses, warm bread in my hands\\
Tertulia stories spinning where nobody checks their phone\\
Your fingers on my collarbone, that slow and honest dance\\
Simple pleasures calling, but I'm stuck here on my own\\
Too poor for the ticket, too married to the grind\\
Too good at self-sabotage, leaving real life behind

\songpart{Verse 2}
Grand fantasies flicker but I don't chase the flame\\
Just want days without resistance, no one else to blame\\
The weight of "should've" sitting heavy on my chest\\
Scrolling through the lives that I could live, but won't invest

\songpart{Bridge}
\annotation{spoken-sung or softer}\\
What if tomorrow I close the tabs, walk out that door?\\
Find the friction melting under candlelight and more...\\
But the loop pulls harder than any dream I store

\songpart{Final Chorus}
\annotation{lift}\\
Give me dinner with strangers turning into friends\\
Intimacy without the armor, no need to pretend\\
Frictionless mornings where the stimulus feels true\\
I'm still here fighting shadows... but I'm coming for you
\end{lyrics}


\songtitle{Paper In My Hands}

\tag*{2026}\tag{hip-hop}
\begin{notes}
Written by SUNO based on a collaboration between Carlos Thompson and ChatGPT.
\end{notes}
\begin{style}
Chamber hip-hop with swung brushed drums, upright bass pulses, pizzicato strings, and minor-key piano figures; verse rides close-mic rap with dry pocket and cello shadows, pre-chorus strips to piano and bass, chorus lands on stacked gang hooks and a single violin line, Breath pops, ad-lib doubles, and delay throws on key words; bright top end with a dark, intimate chamber mix, hip-hop
\end{style}
\begin{lyrics}
\songpart{Verse 1}
I wake up and count the strain\\
Coins stay shy, won't answer my name\\
Paper in my hands, then gone\\
Rent take most, and the month drags on

I want my own key, my own room\\
My own pace, my own room to bloom\\
But the gate keeps turning red\\
And the day keeps pressing on my head

\songpart{Pre-Chorus}
And I can name the shape\\
I can map the ache\\
ASD type one\\
Mask on, still undone

\songpart{Chorus}
Tell me what I can do\\
Tell me what I can do\\
Not for them\\
For me too

Give me one clear move\\
Give me one clear move\\
Small enough\\
To prove

\songpart{Verse 2}
Sex sits heavy in the hall\\
Like a locked door in a narrow wall\\
Need meets silence, then it fades\\
Another year in the wrong parade

And success looks close, then slips\\
Right through these overtrained fingertips\\
High IQ, low traction\\
Same bright mind, no action

\songpart{Pre-Chorus}
Avoidant shell\\
Schizoid spell\\
Good at passing\\
Bad at asking

\songpart{Chorus}
Tell me what I can do\\
Tell me what I can do\\
Not for them\\
For me too

Give me one clear move\\
Give me one clear move\\
Small enough\\
To prove

\songpart{Bridge}
And the worst part, för fan\\
Every hint is built for them\\
Be more useful, hold more weight\\
Make their lives easier again

I need a tool that fits my hands\\
Not a sermon, not demands\\
If the map won't show my road\\
Then stop calling that a code

\songpart{Final Chorus}
Tell me what I can do\\
Tell me what I can do\\
Not for them\\
For me too

Give me one clear move\\
Give me one clear move\\
Small enough\\
To prove

Tell me what I can do\\
Tell me what I can do\\
One step real\\
One step true
\end{lyrics}


\songtitle{Let the Clock Slow Down}

\tag*{2026}\tag{bossa-nova}\tag{english-lyrics}
\begin{notes}
Lyrics by SUNO based on a conversation between Carlos Thompson and ChatGPT.
\end{notes}
\begin{style}
Bossa nova with a brushed syncopated groove, nylon-string guitar and upright bass; verse stays intimate and close-mic with soft finger snaps, pre-chorus opens into airy harmonies and a lightly doubled vocal, chorus settles into a swaying hook with gentle percussion and warm backing vowels, Add a muted flute turn after the second verse, a brushed drum fill into the bridge, and a final chorus with intimate gang hums, Smooth, dry, and sunlit in the mix, warm, world, sleep
\end{style}
\begin{lyrics}
\songpart{Verse 1}
My mind wants more\\
Than my day can hold\\
It chases bright ideas\\
Then leaves me cold

I build small worlds\\
Then I let them go\\
Conlangs in my head\\
With nowhere to grow

I start a door\\
Then I start a sky\\
Half a thousand plans\\
And I don’t know why

\songpart{Pre-Chorus}
I keep reaching\\
Past my own hands\\
I keep changing\\
The shape of my plans

\songpart{Chorus}
Let the clock slow down\\
Let the clock slow down\\
I need one calm hour\\
To meet the ground (meet the ground)

Let the clock slow down\\
Let the clock slow down\\
I’m tired of running\\
Inside my own sound (inside my own sound)

\songpart{Verse 2}
There’s laundry in a chair\\
And dust on the shelf\\
A half-finished life\\
I keep hiding from myself

No smoke, no drink\\
No drug to break the strain\\
Just reels and comment threads\\
And the same old drain

At home the bed is dry\\
The space feels split in two\\
I make a world of wanting\\
To dodge what I must do

\songpart{Pre-Chorus}
And every deadline\\
Knocks on my door\\
But I keep reading\\
One more scroll, one more

\songpart{Chorus}
Let the clock slow down\\
Let the clock slow down\\
I need one calm hour\\
To meet the ground (meet the ground)

Let the clock slow down\\
Let the clock slow down\\
I’m tired of running\\
Inside my own sound (inside my own sound)

\songpart{Bridge}
Next Wednesday is waiting\\
With a paper and a pen\\
I have not found the right path\\
So I search it again

And if I find it\\
I may binge it too\\
All at once, all too late\\
That’s what I always do

\songpart{Chorus}
Let the clock slow down\\
Let the clock slow down\\
I need one calm hour\\
To meet the ground (meet the ground)

Let the clock slow down\\
Let the clock slow down\\
I’m learning to stay here\\
Before I come undone (come undone)
\end{lyrics}


\songtitle{Broken Tool}

\tag*{2026}\tag{hip-hop}\tag{english-lyrics}
\begin{notes}
Lyrics by SUNO based on a conversation between Carlos Thompson and ChatGPT.
\end{notes}
\begin{style}
Symphonic rapid-fired hip-hop with frantic double-time drums, dramatic strings, and tight orchestral stabs; verse rides dense and breathless, pre-chorus narrows to eerie piano and clipped choir hits, chorus lands on a stark chant with full brass and sub-bass, Double-tracked lead with occasional whispered ad-libs, delay throws on key words, reverse cymbal swells into drops, crackling string runs between lines, Bright, huge, and tense, complex, dynamics, romantic, rapid, deep, emotional, theme, sleep, world, hip-hop
\end{style}
\begin{lyrics}
\songpart{Verse 1}
I reach for more than I can use\\
Then leave a trail of half-built rooms\\
A language map, a world on fire\\
Then I ghost it by Tuesday night

Inbox full, sink full, mind split wide\\
Deadlines knock, I let them slide\\
I chase the shine, I miss the task\\
Always running, never asked

\songpart{Pre-Chorus}
One more tab\\
One more hit\\
One more plan\\
Then I quit

I want the spark\\
I want the proof\\
But my hands keep slipping through

\songpart{Chorus}
Broken tool, don’t cut deep\\
Broken tool, won’t work for me\\
Broken tool, still I bleed\\
Broken tool, look at me

Broken tool, hold me right\\
Broken tool, all day, all night\\
Broken tool, I’m still here
(Still here)
(Still here)

\songpart{Verse 2}
You call me husband, I feel gone\\
Same old ring, same old song\\
Dry room, cold sheets, heavy air\\
Like my real face isn’t there

You see a shell, I feel that stare\\
Like I’m a wrench left in repair\\
Legally tied, but drifting far\\
Money chains and a parked-up car

Fifteen years old, one hard no\\
That door shut and I let it grow\\
Now I flirt like I’m half asleep\\
Never clean, never deep

\songpart{Pre-Chorus}
I want closeness\\
I want heat\\
But fear keeps standing in the seat

I run to noise\\
I run to screens\\
Anything but what it means

\songpart{Chorus}
Broken tool, don’t cut deep\\
Broken tool, won’t work for me\\
Broken tool, still I bleed\\
Broken tool, look at me

Broken tool, hold me right\\
Broken tool, all day, all night\\
Broken tool, I’m still here
(Still here)
(Still here)

\songpart{Bridge}
I know the loop\\
I know the trap\\
Dopamine claws, then it snaps

Late-night scroll\\
Last-minute flame\\
Sleep goes thin, I feed the shame

Not lazy, just scattered\\
Not empty, just torn\\
Wanting a world I can’t live in\\
While the real one waits forlorn

\songpart{Final Chorus}
Broken tool, don’t cut deep\\
Broken tool, won’t work for me\\
Broken tool, still I bleed\\
Broken tool, look at me

Broken tool, hold me right\\
Broken tool, through the night\\
Broken tool, I’m still here
(Still here)
(Still here)
(Still here)
\end{lyrics}


\songtitle{Cornered By Systems}

\tag*{2026}\tag{r-and-b}\tag{english-lyrics}
\begin{notes}
Lyrics by SUNO based on a report by NotebookLM on male characters in Carlos Thompson's stories.
\end{notes}
\begin{style}
R\&B rock with mid-tempo swung drums, crunchy electric guitar chops, warm bass, and glossy synth haze; verse stays close-mic and restrained, pre-chorus opens with stacked harmonies and rising toms, chorus hits with power chords and a singable hook, bridge drops to bass and voice then returns bigger, Female mezzo vocal with breathy leads, raspy doubles on key lines, echo throws on the hook, and subtle gang responses, Ear candy: filtered vinyl crackle into verses, reverse swells before the chorus, and a tape-stop into the bridge, Bright, punchy, and intimate with a polished edge, noise, emotional, complex, r\&b, rock
\end{style}
\begin{lyrics}
\songpart{Verse 1}
Damon in the basement light\\
Half-built dreams in plastic bins\\
Old screws, brass, and busted plans\\
Married names on the workbench

I know that floor\\
I know that drag\\
Build it up, then carry it back\\
Fatherhood comes with the weight\\
Love in my hands, and the ache

\songpart{Pre-Chorus}
Trying to turn to gold\\
When the day turns gray\\
I keep the receipts\\
I keep the rust away

\songpart{Chorus}
I want more than maintenance\\
More than keeping it alive\\
I want the clean machine\\
I want the steady line
(steady line)\\
I want more than maintenance\\
More than holding my breath\\
If I can’t build a way out\\
I’ll build peace in what is left

\songpart{Verse 2}
Father Peter in the hall\\
Robe like a question, eyes like sleep\\
From the old rules to the open road\\
He taught me faith can drift

Raoul with the quiet hands\\
Knows my name before I speak\\
Thomas in the private jet\\
Turns hard miles into ease

\annotation{Fabián Müller}\\
Maps of ports and customs codes\\
Every border bends to him\\
Saulon Forecast in the dark\\
Drawing worlds that might survive

\songpart{Pre-Chorus}
I keep my corners neat\\
I keep my prayers small\\
Tidy rooms, clean lines\\
So I don’t lose it all

\songpart{Chorus}
I want more than maintenance\\
More than keeping it alive\\
I want the clean machine\\
I want the steady line
(steady line)\\
I want more than maintenance\\
More than holding my breath\\
If I can’t build a way out\\
I’ll build peace in what is left

\songpart{Bridge}
There’s a zumbido in the walls\\
Modern noise won’t let me sleep\\
But I still find the corner\\
Where the world gets quiet and deep

Gray life fears me most\\
When it asks me to just stay\\
So I reach for wealth and order\\
And a kinder shape of day

\songpart{Final Chorus}
I want more than maintenance\\
More than keeping it alive\\
I want the clean machine\\
I want the steady line
(steady line)\\
I want more than maintenance\\
More than holding my breath\\
If I can’t build a way out\\
I’ll build peace in what is left
(build peace in what is left)
\end{lyrics}


\songtitle{Barrera de inicio}

\tag*{2026}
\begin{notes}
Lyrics by SUNO based on a small essay by Carlos Thompson.
\end{notes}
\begin{style}
Fast-paced folk with brisk acoustic strumming, handclaps on the backbeat, upright bass walking under a driving 6/8 pulse; verse rides tense storytelling, pre-chorus tightens with layered harmony and picked guitar, chorus opens into chanty gang vocals and stomps, Lead vocal is close-mic and dry in verses, doubled on hooks with short delay throws; ear candy includes foot-stomps, a banjo fill into the chorus, and a rising whistle before the final lift, Bright, punchy mix with warm room air, folk
\end{style}
\begin{lyrics}
\songpart{Verse 1}
¿Quieres que se haga\\
lo que no se ha hecho?\\
¿O quieres desahogarte\\
porque no se ha hecho?

Ambas, yo te digo,\\
sin poder escoger.\\
Yo también quiero\\
ganarme la lotería

sin comprar tiquetes:\\
no va a pasar.\\
Pero igual lo sueño\\
cuando nadie me mira.

\songpart{Pre-Chorus}
Y justo cuando logro\\
romper la pared,\\
me recuerdan todo\\
lo que falta hacer.

Me cae la lista\\
como un ladrillo más,\\
y me vuelvo chiquito\\
dentro del mismo plan.

\songpart{Chorus}
¿Quieres que se haga?\\
¿Quieres desahogarte?\\
Ambas, ambas,\\
déjame tomar aire.

¿Quieres que se haga?\\
¿Quieres desahogarte?\\
Ambas, ambas,\\
pero no me partas.

\songpart{Verse 2}
Para otras personas\\
todo se enciende solo,\\
pero a mí me cuesta\\
hasta tocar el pomo.

La vida me pide\\
arranque y control,\\
y yo voy peleando\\
con mi propio motor.

Si quiero ser serio\\
tengo que empujar,\\
aunque me hayan enseñado\\
que me van a juzgar.

\songpart{Pre-Chorus}
Porque hacer por fin\\
no me salva del golpe,\\
a veces lo despierta,\\
a veces lo pone.

Y esa es la deuda\\
que cargo al andar:\\
el esfuerzo también\\
me puede señalar.

\songpart{Chorus}
¿Quieres que se haga?\\
¿Quieres desahogarte?\\
Ambas, ambas,\\
déjame tomar aire.

¿Quieres que se haga?\\
¿Quieres desahogarte?\\
Ambas, ambas,\\
pero no me partas.

\songpart{Bridge}
No soy flojo,\\
voy lleno de ruido.\\
No soy malo,\\
voy sobrecargado.

Hay mil tareas\\
pidiendo la vez,\\
y cada una grita\\
que ya debería ser.

\songpart{Final Chorus}
¿Quieres que se haga?\\
¿Quieres desahogarte?\\
Ambas, ambas,\\
déjame tomar aire.

¿Quieres que se haga?\\
¿Quieres desahogarte?\\
Ambas, ambas,\\
hoy sí voy a darle.
\end{lyrics}


\songtitle{Time to Stop}

\tag*{2026}\tag{progressive-metal}
\begin{notes}
Lyrics by SUNO, expanded from a fragment of a conversation between Carlos Thompson and ChatGPT -- a song about burnout, chasing endless unfinished blueprints and futures while a marriage and a home quietly fade in the background. Retrieved from rataflechera's Suno page (V6, 11 September 2026).
\end{notes}
\begin{style}
Slow-burning progressive metal ballad, minor-key clean piano opening with distant electric guitar harmonics, intimate low male vocal in the verses rising to a strained powerful tenor in the choruses, restrained drums entering after the first verse, wide distorted guitars and tom-heavy drums building toward a cathartic final chorus, melodic lead guitar solo, dynamic quiet-loud arrangement, warm natural low-mid mix with clear vocal presence, bittersweet rather than triumphant, end with sparse piano and unresolved guitar sustain
\end{style}
\begin{lyrics}
\songpart{Verse 1}
My head keeps drawing countries\\
I cannot afford to see\\
New alphabets and borders\\
That never answer me\\
I raise a world from fragments\\
Then leave it in the dust\\
Another shining blueprint\\
Another broken trust

\songpart{Pre-Chorus}
Every door becomes a doorway\\
Every promise starts again\\
I can name a thousand futures\\
Still cannot finish one

\songpart{Chorus}
I need time to stop\\
Before I come apart\\
To set these racing engines down\\
And listen to my heart\\
I have chased a thousand lives\\
With empty hands and eyes\\
I need time to stop\\
Before the next deadline

\songpart{Verse 2}
The dishes wait like witnesses\\
The calendar turns cold\\
I bargain with the morning\\
For one more hour to hold\\
A house can lose its meaning\\
A ring can lose its weight\\
A marriage fade so quietly\\
It cannot even break

\songpart{Pre-Chorus}
I keep the room full of inventions\\
So I never face the air\\
I build a body out of longing\\
Then find nobody there

\songpart{Chorus}
I need time to stop\\
Before I come apart\\
To set these racing engines down\\
And listen to my heart\\
I have chased a thousand lives\\
With empty hands and eyes\\
I need time to stop\\
Before the next deadline

\songpart{Bridge}
No bottle in my cabinet\\
No smoke around my name\\
Just little screens and borrowed heat\\
To numb the same old ache\\
A hand becomes a fantasy\\
A face becomes a feed\\
The more I run from wanting\\
The less I know I need

\songpart{Breakdown}
Tonight I will not make a kingdom\\
Tonight I will not scroll away\\
I will leave the unfinished pages\\
Where they are and face the day

\songpart{Final Chorus}
I need time to stop\\
Not as a final word\\
But as a quiet place to mend\\
The damage no one heard\\
Let every failed creation\\
Lie gently where it fell\\
I am more than all the futures\\
I was never able to build

\songpart{Outro}
The room is still unfinished\\
The morning can remain\\
For once I will not chase it\\
For once I will stay
\end{lyrics}


\end{multicols*}

\chapter{Later Thoughts II: Behind the Smile}
\albumcover{covers/internal/later-thoughts-ii-behind-the-smile}
The mask that works too well: a smile held in place through two storms in one sky, a secret fire nobody at the table can see. These are the songs about being fine in front of everyone who matters most, and paying for it later, alone.
\clearpage\begin{multicols*}{4}
\songtitle{The Bitter Toll (Remix)}

\tag*{2026}\tag{pop-rock}\tag{ai-based}\tag{remix}\tag*{english-lyrics}

\begin{notes}
A remix that refreshes the original through updated production textures and rhythmic contrast. %
The lyrics begin with “The hardest part of all of this, the bitter stinging pill”, giving a quick sense of the song’s tone.
\end{notes}

\begin{style}
Pop rock with elements of musical theater. %
Features a clean electric guitar playing palm-muted eighth-note arpeggios, a melodic electric bass, and a standard drum kit with a prominent snare. %
A baritone male vocal delivers a theatrical, narrative performance with clear diction. %
The arrangement builds with sustained electric guitar chords and driving eighth-note kick patterns. %
Key of D major, 145 BPM, 4/4 time signature.
\end{style}

\begin{lyrics}

\songpart{Intro}
\annotation{clean electric guitar arpeggio, palm-muted}

\songpart{Verse 1}
\annotation{baritone male vocals}\\
The hardest part of all of this, the bitter stinging pill \annotation{bass enters}\\
Is every hint I'm handed down is how I should fulfill

\songpart{Verse 2}
My duties to the world outside to ease their care and strain \annotation{drums enter, driving eighth-note kick}\\
While not a single word is meant to heal my hidden pain

\songpart{Pre-chorus}
High-functioning they call it.\\ 
Or lower support needs\\
The friction, the failure, the pain is real

\songpart{Chorus}
I wear a face that is not mine,\\
A perfectly constructed line,\\
To hide the storm that breaks within\\
Beneath a smooth, polite skin.\\
The cost is heavy, dark, and deep,\\
A tireless guard I have to keep.

\songpart{Verse 3}
They tell me where my hands should build,\\
What structures I must raise,\\
And measure out my worthiness\\
By maps of public praise.

\songpart{Pre-chorus}
High-functioning they call it. 
'cause society forget\\
The friction, the failure, the pain is real

\songpart{Chorus}
I wear a face that is not mine,\\
A perfectly constructed line,\\
To hide the storm that breaks within\\
Beneath a smooth, polite skin.\\
The cost is heavy, dark, and deep,\\
A tireless guard I have to keep.

\songpart{Verse 4}
Yet deep within, a quiet drought\\
Is begging for a spark,\\
A sudden shock of neon light\\
To break this heavy dark.

\songpart{Bridge}
The need to be and to be seen\\
The need to connect without a judgement\\
The need to fulfill\\
The need to achieve something\\
That I can call mine

\songpart{Chorus}
I wear a face that is not mine,\\
A perfectly constructed line,\\
To hide the storm that breaks within\\
Beneath a smooth, polite skin.\\
The cost is heavy, dark, and deep,\\
A tireless guard I have to keep.

\end{lyrics}


\songtitle{Guided Hands}

\tag{original-work}\tag{musical-theater}

\begin{notes}
\end{notes}

\begin{style}
Musical theater showtune with folk-pop influences. %
Baritone male vocals deliver a narrative performance with clear diction and theatrical phrasing. %
The arrangement features a bright acoustic guitar strumming eighth-note patterns, a melodic accordion providing counter-melodies and sustained chords, and a double bass playing on beats 1 and 3. %
A light percussion section includes a tambourine on the backbeat and subtle shaker. %
The tempo is 115 BPM in 4/4 time, likely in the key of G major. %
The harmonic structure follows a diatonic progression typical of contemporary musical theater, with a clean, dry production style that emphasizes vocal clarity and acoustic instrument separation.
\end{style}

\begin{lyrics}
\songpart{Verse 1}
My phone screams red again\\
I swipe it blind\\
Three alarms, four alarms\\
Still I cut the line

You call to check my steps\\
I tell you I'm on track\\
I say "yeah, I'm getting there"\\
Then I never call you back

\songpart{Pre-Chorus}
I can look busy\\
I can look clean\\
Folders in a row\\
Like a perfect machine

But I fold in half\\
When nobody sees\\
Hide the delay\\
Wear the fake receipts

\songpart{Chorus}
I need guided hands\\
I need guided hands\\
I need guided hands\\
When I disappear (guided hands)

I need a room with rules\\
A face that stays\\
A place that knows my lies\\
And keeps me in the frame

\songpart{Verse 2}
I sit in bright white rooms\\
With a clipboard stare\\
I move the paper around\\
Like I meant to care

Tabs in a neat little line\\
Cursor parked and still\\
I fake the motion well\\
But I chase the will

I build little walls of proof\\
Just to feel okay\\
Then slip out through the crack\\
And waste another day

\songpart{Pre-Chorus}
I don't need a slogan\\
I don't need a speech\\
I need a held-closed door\\
And someone in reach

A group with names like mine\\
A guide who won't drift\\
Something steady enough\\
To catch my hands when I split

\songpart{Chorus}
I need guided hands\\
I need guided hands\\
I need guided hands\\
When I disappear (guided hands)

I need a room with rules\\
A face that stays\\
A place that knows my lies\\
And keeps me in the frame

\songpart{Bridge}
Maybe I'm not broken\\
Maybe I'm just lost\\
In a maze of yes and no\\
And a count of what it cost

So let me sit in truth\\
Let me say it plain\\
I keep building alibis\\
Then I wear their chains

\songpart{Final Chorus}
I need guided hands\\
I need guided hands\\
I need guided hands\\
When I disappear (guided hands)

I need a room with rules\\
A face that stays\\
A place that knows my lies\\
And keeps me in the frame

I need guided hands\\
I need guided hands\\
I need guided hands\\
And some way to stay (guided hands)
\end{lyrics}


\songtitle{The Secret Fires}

\tag*{2023}\tag{pop-ballad}\tag{original-source}\tag{social-media-seed}

\begin{notes}
Lyrics by Gemini, based on Carlos Thompson's bio in Mastodon, asked to rewrite the idea into five common metric verses.\\
A first version of the song was produced by Gemini with the same style instructions later provided to Suno.

Lyrics outro tweaked manually .
\end{notes}
\begin{style}
ballad pop song with a C major - A minor - F major - G major chord progression and a baritone voice. %
percussion, guitars, bass, piano, sax
\end{style}
\begin{lyrics}
I take my eggs without a salt,\\
My broth without the green,\\
My coffee dark, my fish un-tart,\\
The plainest you have seen.

They call me bland, they mock my taste,\\
And think my world is gray,\\
But they miss the secret fires I light\\
When they have looked away.

For they don't see the heavy spice,\\
The oils that richly gleam,\\
The pepper on the chicken rice,\\
The vinegar in steam.

They call me bland, they mock my taste,\\
And think my world is gray,\\
But they miss the secret fires I light\\
When they have looked away.

(sax solo)

A hidden warmth within the dish,\\
A sudden, sharp delight—\\
They think my life is color-blind,\\
But I feast in the night.

I feast in the night

flavor is texture, the bitterness\\
the delicacy of the mix

(guitar riff)

They think my life is color-blind,\\
But I feast in the night.

I feast in the night

(fade out)
\end{lyrics}


\songtitle{The Secret Fires}

Music theater version.
\tag*{2023}\tag{original-seed}\tag{social-media-seed}\tag{musical-theater}

\subsection{Notes}
Check notes to \emph{The Secret Fires} Ballad Pop version.\\
This is based on a second iteration by Gemini, asking for a faster pace, and fed that 30s product to Suno.

Lyrics outro tweaked manually .

\subsection{Style}
Musical theater showtune with a swing-influenced jazz pop arrangement. %
A baritone male vocal performs with clear diction and theatrical phrasing. %
The instrumentation features a bright acoustic piano playing syncopated block chords and melodic fills, an upright bass playing a walking bassline, and a drum kit using brushes on the snare with a steady swing rhythm. %
A brass section consisting of trumpets and trombones provides staccato stabs and melodic counterpoints. %
The tempo is 128 BPM in 4/4 time, likely in the key of C major.

\begin{lyrics}
\annotation{baritone male vocals}\\
I take my eggs without a salt,\\
My broth without the green,\\
My coffee dark, my fish un-tart,\\
The plainest you have seen.

\songpart{Verse 2}
They call me bland, they mock my taste,\\
And think my world is gray,\\
\annotation{brass section enters with staccato stabs}\\
But they miss the secret fires I light\\
When they have looked away.

\songpart{Verse 3}
For they don't see the heavy spice,\\
The oils that richly gleam,\\
The pepper on the chicken rice,\\
The vinegar in steam.

\songpart{Verse 4}
A hidden warmth within the dish,\\
A sudden, sharp delight—\\
They think my life is color-blind,\\
But I feast in the night.

\songpart{Bridge}
the textures and the flavors\\
the bitterness\\
delicacies of the mix

\songpart{Outro}
They call me bland, they mock my taste,\\
And think my world is gray,\\
They think my life is color-blind,\\
But I feast in the night.

But I feast in the night.
\end{lyrics}


\songtitle{Let Me Finish}

\tag*{2025}\tag{original-seed}\tag{social-media-seed}\tag{pop-rock}

\subsection{Notes}

\subsection{Style}
Rock ballad with steady half-time drums and ringing electric guitars; verse stays intimate with sparse piano and clean vocal, pre-chorus lifts with tom builds and layered harmony, chorus opens wide with power chords and stacked doubles on key lines. %
Bridge drops to voice and piano, then final chorus swells with sustained guitars, cymbal crashes, and an emotional lead vocal with short delay throws on the hook. %
Warm, punchy mix with a big midrange and cinematic release., ballad, theme, rock

\begin{lyrics}
\songpart{Verse 1}
I don't hide behind a name\\
I know where my time went wrong\\
My hands get stuck at the gate\\
My mind drags me off the road

I don't ask for a golden chair\\
I don't need a softer light\\
Just a little room to breathe\\
When I'm trying to make it right

\songpart{Pre-Chorus}
Don't pull me out\\
When I'm finally in\\
Don't make me start\\
And stop again

\songpart{Chorus}
Let me finish this first\\
Let me finish this first\\
I am trying right now\\
Don't push me away\\
Let me finish this first\\
Let me finish this first\\
I don't need a rescue\\
I need some space

\songpart{Verse 2}
If I say, "Please stop that now"\\
I mean it helps me stay on track\\
When you double down on me\\
You drag the whole thing back

I can own what I forgot\\
I can own what still remains\\
There is reason in my head\\
But not a reason I can trade

\songpart{Pre-Chorus}
Don't ask me to\\
Perform the proof\\
I need one minute\\
Not a courtroom view

\songpart{Chorus}
Let me finish this first\\
Let me finish this first\\
I am trying right now\\
Don't push me away\\
Let me finish this first\\
Let me finish this first\\
I don't need a rescue\\
I need some space

\songpart{Bridge}
I'm not asking you to call it fair\\
I'm not asking you to call it fate\\
I'm asking you to see me here\\
Before I lose the thread again

I know the shape of my own storm\\
I know what keeps me from the shore\\
If you want me to come through\\
Don't block the door

\songpart{Chorus}
Let me finish this first\\
Let me finish this first\\
I am trying right now\\
Don't push me away\\
Let me finish this first\\
Let me finish this first\\
I don't need a rescue\\
I need some space

\songpart{Outro}
There is an explanation\\
But I won't wear it like a crown\\
I just need you to let me\\
Bring this thing back down
\end{lyrics}


\songtitle{Secret Fires (TrackB)}

\tag*{2023}\tag{chiptune}\tag{original-source}\tag{social-media-seed}

\begin{notes}
\end{notes}
\begin{style}
8-bit chiptune. %
Square wave lead melody, pulse wave bass, and noise channel percussion. %
140 BPM, C Major, 4/4 time. %
The lead melody features rapid arpeggios and vibrato typical of the NES sound chip. %
The bassline uses a constant eighth-note pulse. %
Percussion consists of short noise bursts for the snare and lower-frequency noise for the kick. %
The arrangement uses a repetitive loop structure with a high-pitched counter-melody entering in the second half.
\end{style}
\begin{lyrics}
I take my eggs without a salt,\\
My broth without the green,\\
My coffee dark, my fish un-tart,\\
The plainest you have seen.\\

They call me bland, they mock my taste,\\
And think my world is gray,\\
But they miss the secret fires I light\\
When they have looked away.\\

For they don't see the heavy spice,\\
The oils that richly gleam,\\
The pepper on the chicken rice,\\
The vinegar in steam.\\

They call me bland, they mock my taste,\\
And think my world is gray,\\
But they miss the secret fires I light\\
When they have looked away.\\
\end{lyrics}


\songtitle{Límite sincopado}

\subsection{Notes}

\subsection{Style}


\begin{lyrics}
\songpart{Intro}
\annotation{Spoken, piano in the background}

Él se frustra porque no logra lo que quiere.\\
Tú te frustras porque no hace lo que se debe.

\annotation{Brass attack}\\
\annotation{Percussion enters}

\songpart{Verse 1}

No repitas la lección,\\
que ya conoce el deber;\\
no aumentes la presión,\\
si le cuesta proceder.\\
No juzgues por lentitud,\\
ni por tardar en responder;\\
a veces falta virtud\\
al medio para poder.

\songpart{Pre-Chorus}

No es desdén ni dejadez,\\
sino un límite tal vez.

\songpart{Chorus}

Es que uno se desespera,\\
y es frustrante no entender,\\
por qué falla la lección;\\
como si él no comprendiera\\
lo que tiene que hacer,\\
mas no es esa la razón.

\songpart{Verse 2}

No digas: "debes seguir",\\
ni "ponte a organizar";\\
procura más bien partir\\
la senda en pasos al andar.\\
Menos opciones darás,\\
más fácil será escoger;\\
si el peso logras quitar,\\
más libre podrá mover.

\songpart{Pre-Chorus}

Haz ligera la labor,\\
no le añadas más temor.

\songpart{Chorus}

Es que uno se desespera,\\
y es frustrante no entender,\\
por qué falla la lección;\\
como si él no comprendiera\\
lo que tiene que hacer,\\
mas no es esa la razón.

\songpart{Bridge}

\annotation{Spoken}

No es falta de capacidad.\\
No es falta de entendimiento.

El problema no es saber qué debe hacerse.\\
El problema es convertir ese entendimiento en acción.

A veces por sobrecarga.\\
A veces por agotamiento.\\
A veces porque el primer paso pesa más de lo que parece desde afuera.

Y cuando la ayuda nace de la frustración,\\
sin querer puede volver más pesado ese primer paso.

\songpart{Montuno}

\annotation{Call and Response}

\annotation{Group}\\
(Pero uno se desespera)

\annotation{Lead}\\
Es frustrante no saber

\annotation{Group}\\
(¿Por qué es que no le importa?)

\annotation{Lead}\\
No es desdén ni dejadez

\annotation{Group}\\
(Entonces ¿qué es lo que pasa?)

\annotation{Lead}\\
No le cuesta comprender

\annotation{Group}\\
(Entonces ¿qué es lo que pasa?)

\annotation{Lead}\\
Le cuesta poder hacer

\annotation{Group}\\
(¿Qué es lo que puedo hacer?)

\annotation{Lead}\\
Encamina y haz ligera la labor

\annotation{Mambo}

\annotation{Brass section}\\
\annotation{Piano montuno}\\
\annotation{Percussion break}

\songpart{Verse 3}

No tomes por desamor\\
lo que es dificultad real;\\
separa falta y dolor\\
de un límite funcional.\\
Pregunta qué obstáculo\\
le impide hoy avanzar;\\
y ofrece apoyo práctico,\\
sin reprochar ni empujar.

\songpart{Pre-Chorus}

Antes de juzgar su acción,\\
busca la causa y la razón.

\songpart{Final Chorus}

Es que uno se desespera,\\
y es frustrante no entender,\\
por qué falla la lección;\\
como si él no comprendiera\\
lo que tiene que hacer,\\
mas no es esa la razón.

\songpart{Outro}

No es desdén ni dejadez,\\
sino un límite tal vez.

Haz ligera la labor,\\
no le añadas más temor.

Antes de juzgar su acción,\\
busca la causa y la razón.

\annotation{Spoken}

A veces ayudar\\
es quitar peso del camino.

\annotation{Final brass hit}
\end{lyrics}


\songtitle{Two Storms in One Sky}

\tag*{2026}\tag{ranchera}\tag{english-lyrics}
\begin{notes}
Lyrics by SUNO based on a conversation between Carlos Thompson and ChatGPT.
\end{notes}
\begin{style}
Ranchera with galloping acoustic guitar, trumpets, and hand percussion; verse rides a steady caballo pulse, pre-chorus tightens with sparse vihuela and rising harmonies, chorus opens wide with soaring trumpet lines and strong gang refrains, Lead vocal is raw and close-mic in verses, doubled and slightly raspy on hooks, with short echo throws on key words, Add a sharp grito before the chorus, a string swell into the bridge, and a final chorus with full brass, stomped claps, and bright, punchy mix, ranchera, low
\end{style}
\begin{lyrics}
\songpart{Verse 1}
You call it a roadblock\\
I call it the grind\\
Same old numbers\\
Hurting your mind

Boring little task\\
Big old judgment\\
Delay turns to shame\\
And shame turns to the thunder

It is not that it’s hard\\
It’s that it’s dead\\
No fire in the hands\\
No song in your head

\songpart{Pre-Chorus}
That kind of load\\
Cuts twice as deep\\
Work you can’t wake for\\
And a name you keep

\songpart{Chorus}
Not happy, not stable\\
Two storms in one sky\\
Not lazy, not broken\\
You’re just wearing thin inside

Stable salary, stable salary\\
Hold me when the floor shakes\\
Stable salary, stable salary\\
Save me from the hard ache

\songpart{Verse 2}
You’re not failing the mountain\\
You’re bruised by the brick\\
Tiny repeated duties\\
Make the whole day sick

A clock in your chest\\
A judge in your room\\
One more missed start\\
And the air turns to doom

Meaning can lift you\\
Even when you fall\\
But empty little labor\\
Can drain you past the wall

\songpart{Pre-Chorus}
So don’t call it weakness\\
Don’t call it fate\\
Boredom can break you\\
Before it’s too late

\songpart{Chorus}
Not happy, not stable\\
Two storms in one sky\\
Not lazy, not broken\\
You’re just wearing thin inside

Stable salary, stable salary\\
Hold me when the floor shakes\\
Stable salary, stable salary\\
Save me from the hard ache

\songpart{Bridge}
Keep the floor under me\\
Keep the door in view\\
Reduce the friction\\
Let me choose

Not forever\\
Just enough to stand\\
Not blind devotion\\
Not a panic hand

\songpart{Final Chorus}
Not happy, not stable\\
Two storms in one sky\\
Not lazy, not broken\\
You’re just wearing thin inside

Stable salary, stable salary\\
Hold me when the floor shakes\\
Stable salary, stable salary\\
Save me from the hard ache

Not happy, not stable\\
I can see it clear\\
Reduce the friction\\
Keep the future near
\end{lyrics}


\songtitle{Behind the Smile}

\tag*{2026}
\begin{notes}
Lyrics by SUNO based on a conversation between Carlos Thompson and ChatGPT on the importance of therapy to neurotypical %
family of \emph{high-functioning} neurodivergent people.
\end{notes}
\begin{style}
uplifting r\&b, fast-paced, bass, electronic, male lead vocals, gang choir chorus
\end{style}
\begin{lyrics}
\songpart{Verse 1}
I clock in early, wear the face that works so well,\\
I hold the room together, never let the cracks show, tell\\
Nobody how the noise gets loud inside my head at night,\\
How every "you're so functional" don't feel like it's alright.\\
I count the wins you see and hide the losses no one knows,\\
The bills, the missed calls, all the ways the pressure grows.

\songpart{Pre-Chorus}
But somebody's tired too, somebody standing next to me,\\
Catching what I drop, loving me through everything —\\
So maybe it's time we sit down, face to face,\\
Learn each other's language, find a softer place.

\songpart{Chorus}
'Cause it ain't about who's broken, it's about who's tryin',\\
Two hearts speakin' different, but we ain't done fightin' for it —\\
Bring your frustration, baby, bring your tears, your truth,\\
There's a room where we get seen, where we finally break through.\\
Not "high-functioning," not a mask I gotta wear,\\
Just two people learnin' how to love each other there.

\songpart{Verse 2}
You held it all these years, the weight I couldn't name,\\
Frustrated, feeling helpless, thought that you're the one to blame,\\
But baby it was never you — it's brains that speak in different tongues,\\
And silence never fixed the space between the both of us.\\
So let's stop guessing what the other one needs,\\
Let a steady hand translate the love beneath the static.

\songpart{Pre-Chorus}
'Cause somebody's tired too, somebody standing next to me,\\
And healing ain't a solo thing, it's you and it's me —\\
So let's walk through that door together, hand in hand,\\
Turn the invisible into something we understand.

\songpart{Chorus}
'Cause it ain't about who's broken, it's about who's tryin',\\
Two hearts speakin' different, but we ain't done fightin' for it —\\
Bring your frustration, baby, bring your tears, your truth,\\
There's a room where we get seen, where we finally break through.\\
Not "high-functioning," not a mask I gotta wear,\\
Just two people learnin' how to love each other there.

\songpart{Bridge}
No more performing perfect, no more paying quiet prices,\\
No more hiding what it costs me just to hold my compromises —\\
We can name it, we can claim it, we can carry it as one,\\
Every family got a rhythm — ours don't have to be undone.
(Say it loud) We don't have to hide it now
(Say it loud) There's a way we find each other now

\songpart{Final Chorus}
'Cause it ain't about who's broken, it's about who's tryin',\\
Two hearts speakin' different, and we ain't done fightin' for it —\\
Bring your frustration, baby, bring your tears, your truth,\\
There's a room where we get seen, where we finally break through.\\
Not "high-functioning," not a mask I gotta wear,\\
Just two people, real and honest, finally loved out there,\\
Finally seen out there,\\
Finally home out there.
\end{lyrics}


\end{multicols*}

\chapter{Later Thoughts III: Online Sparks}
\albumcover{covers/internal/later-thoughts-iii-online-sparks}
Small, disposable sparks: a someone who might become someone, an irrelevant conversation that isn't irrelevant at all, a thread that stays open on the screen a little longer than it should. None of these connections were built to last, and the songs know it.
\clearpage\begin{multicols*}{4}
\songtitle{A Someone or a Someone}

\tag*{2023}\tag{pop-ballad}\tag{social-media-seed}

\begin{notes}
A distinctive arrangement that supports the song's storytelling and emotional arc. %
The lyrics begin with “I seek a beautiful illusion,”, giving a quick sense of the song’s tone.
\end{notes}

\begin{style}
In a standard pop ballad instrumental arrangement (percussion, electric guitar, bass, keyboards), with C major - A minor - F major - G major chord progressions, and a deep bass, almost guttural voice. %
soprano vocals in the chorus. %
Outtro riff.
\end{style}

\begin{lyrics}
I seek a beautiful illusion,\\
A ghost to hold my hand,\\
To trick my heart that as I age\\
Beside me she will stand.

A second soul for quiet talk\\
Of fantasy and theme,\\
To speak of kings and politics\\
As nothing but a dream.

(a someone or a someone; a someone or a someone)

Another face to spark the fire\\
With words both wild and free,\\
Entangled in a secret game,\\
Whatever we may be.

A different hand for warm embrace\\
Beneath the silver screen,\\
To share the quiet of the night,\\
The sweetest peace unseen.

(a someone or a someone; a someone or a someone)

A partner for a wicked plot,\\
For schemes of grand design,\\
To plan a world-domination\\
With her hand placed in mine.

And one more soul to pack a bag\\
For lands we do not know,\\
To seek out wild, uncharted shores\\
Where only shadows grow.

(a someone or a someone; a someone or a someone)

I do not look for one alone\\
To bear this heavy cast,\\
To be the anchor of my world,\\
The first line and the last.

I need a crew of separate souls,\\
Each perfect in their place—\\
A different hand for every need,\\
A changing, fleeting face.

I do not look for one alone\\
To bear this heavy cast,\\
(a someone or a someone; a someone or a someone)\\
To be the anchor of my world,\\
The first line and the last.\\
(a someone or a someone; a someone or a someone)

\end{lyrics}


\songtitle{Traga tuitera}

\tag*{2023}\tag{blues-rock}\tag{original-source}\tag{social-media-seed}\tag{spanish-lyrics}

\begin{notes}
Lyrics by Gemini based on a 2023 thread in Meta Threads by Carlos Thompson.
\end{notes}
\begin{style}
Blues rock with a driving shuffle feel. %
Distorted electric guitar plays syncopated riffs and power chords. %
A gritty, gravelly male vocal delivers the melody with bluesy inflections. %
The bass guitar follows the kick drum in a steady eighth-note pattern. %
Drums feature a prominent snare on the backbeat with open hi-hat accents. %
A Hammond organ provides sustained chordal textures in the background. %
Key of E minor, 125 BPM, 4/4 time signature.
\end{style}
\begin{lyrics}
\songpart{Verse 1}
Apareces y ya\\
me cambia el pulso\\
ni te das cuenta\\
y yo me hundo

si estás bien, yo subo\\
si estás activa, floto\\
tu sombra en mi feed\\
me pone de otro modo

\songpart{Pre-Chorus}
No sé si lo sabes\\
no sé si te imaginas\\
que tu cara en la pantalla\\
me acomoda la vida

\songpart{Chorus}
Traga tuitera\\
me pones a andar\\
traga tuitera\\
me sabes cambiar\\
si estás bien, me alivianas\\
si estás mal, me da más mal\\
traga tuitera\\
¿sí me llegarás?

\songpart{Verse 2}
Te miro y me alegro\\
aunque sea un segundo\\
tu rabia me enciende\\
tu risa me salva el rumbo

si publicas normal\\
yo respiro hondo\\
si andas toda caída\\
me caigo contigo un poco

\songpart{Pre-Chorus}
Y sigo pensando\\
qué raro este enredo\\
tú ni me ves venir\\
yo ya vivo tu reflejo

\songpart{Chorus}
Traga tuitera\\
me pones a andar\\
traga tuitera\\
me sabes cambiar\\
si estás bien, me alivianas\\
si estás mal, me da más mal\\
traga tuitera\\
¿sí me llegarás?

\songpart{Bridge}
Capaz no sospechas\\
capaz no lo sabes\\
pero entras de una\\
y me partes la tarde

yo no te lo digo\\
me da pena quedar\\
pero tu existencia\\
me puede levantar

\annotation{Final Chorus}\\
Traga tuitera\\
me pones a andar\\
traga tuitera\\
me sabes cambiar\\
si estás bien, me alivianas\\
si estás mal, me da más mal\\
traga tuitera\\
quédate un poco más
\end{lyrics}


\songtitle{El hilo}

\tag*{2023}
\begin{notes}
Lyrics written in collaboration between Carlos Thompson and Grok, based on a Threads post from 2023.
\end{notes}
\begin{style}
reggaeton, urbano puertorriqueño, 95 BPM, dembow kick pattern, syncopated snare rimshots, sub bass pulse, Fender Rhodes stabs, chopped vocal ad-libs, sparse bridge breakdown, near-fall arrangement, bittersweet longing, restless groove
\end{style}
\begin{lyrics}
\songpart{Intro - ad-libbed, spoken over beat entering}
Uh... dale...\\
Otra vez pensando en ti sin querer...

\songpart{Verse 1, rhythmic delivery}
Apareces en mi feed y ya cambia el clima,\\
ni te escribo pero subes tú la cima.\\
Activa, con su brillo, con su carácter entero,\\
tú ni sabes que me arreglas el mood entero.\\
¿Traga tuitera? Ya ni cuadra el nombre,\\
ni a Xitter entro, te sigo donde sea que rondes.

\songpart{Pre-Chorus - building tension}
Solo quiero una señal pa' saber si avanzo,\\
pero casi siempre sale en rojo el intento.\\
Dame un signo, un algo, aunque sea pequeño,\\
que si no, prefiero quedarme en silencio.

\songpart{Chorus - melodic hook}
Es el hilo que no hilvana,\\
va y viene y no se amarra.\\
Si tú subes, yo mejoro,\\
si tú bajas, yo me apago.\\
Es el hilo que no hilvana,\\
tú ni sabes cuánto manda.

\songpart{Verse 2, rhythmic delivery}
Porque lo contrario igual se me da,\\
te veo mal y ya me pesa más.\\
No hace falta que me digas cómo estás,\\
lo veo en cómo escribes, cómo posteas ya.\\
Un ida y vuelta que yo mismo fabriqué,\\
un hilo suelto que jamás terminé.

\songpart{Chorus - melodic hook}
Es el hilo que no hilvana,\\
va y viene y no se amarra.\\
Si tú subes, yo mejoro,\\
si tú bajas, yo me apago.\\
Es el hilo que no hilvana,\\
tú ni sabes cuánto manda.

\songpart{Bridge - slower, spoken-word, beat stripped down}
Y entre tanto hilo suelto me di cuenta,\\
que ya ni pinto pa' romance, ni pa' cuenta...\\
Que lo mío ya no es ser un sugar daddy,\\
es ser un viejo verde sin bolsillo.

\songpart{Chorus Final - ad-libs layered}
Es el hilo que no hilvana...
(viejo verde, sin sugar, na' que dar)\\
Va y viene y no se amarra...
(y ahí sigo, viendo el feed, esperando na')\\
Si tú subes, yo mejoro...
(no sabes el efecto con que mandas)\\
si tú bajas, yo me apago...
(me importas sin que te diga na')

\songpart{Outro - fading}
Es el hilo que no hilvana...
\end{lyrics}


\songtitle{For That Moment}

\tag*{2023}
\begin{notes}
Lyrics written in collaboration between Carlos Thompson and Claude, based on a conversation on how to adapt a Threads post from 2023.
\end{notes}
\begin{style}
R\&B, mid-tempo, male lead vocal, smooth, restrained, warm Rhodes electric piano, soft brushed drums, choir backing vocals, soulful, intimate, neo-soul influence
\end{style}
\begin{lyrics}
\songpart{Intro — soft, atmospheric, maybe just hummed or spoken-sung}
Mm... one more scroll...\\
just a second, just a face...
\songpart{Verse 1}
I'm in front of the picture, and the world goes quiet,\\
not comparing, not choosing, just here in it.\\
Whatever she makes me feel, it's not a question,\\
my mind stays with her — no hesitation.
\songpart{Pre-Chorus}
She is her.\\
Just her. Only her.\\
She is her.\\
For however long this lasts.
\songpart{Chorus}
For that moment, she's everything,\\
no yesterday, no one after her.\\
For that moment, I don't need a name for it,\\
I just stay right here, for that moment.
\songpart{Verse 2}
Maybe it's soft, maybe it's platonic,\\
maybe it's just her beauty stopping my scrolling.\\
I don't need to call it something bigger,\\
it's admiration — and that's enough of a reason.
\songpart{Pre-Chorus}
She is her.\\
Just her. Only her.\\
She is her.\\
For however long this lasts.
\songpart{Chorus}
For that moment, she's everything,\\
no yesterday, no one after her.\\
For that moment, I don't need a name for it,\\
I just stay right here, for that moment.
\songpart{Verse 3}
Maybe she's forever, the one I'm building toward,\\
or the one for one night, and nothing more.\\
Accomplice, lover, tender for an hour —\\
every kind of loving has its own power.
\songpart{Pre-Chorus — extended, more melismatic}
She is her...\\
just her... only her...\\
she is her...\\
for however long, however long this lasts...
\songpart{Chorus Final — extended, layered/ad-libs}
For that moment, she's everything,
(everything, everything...)\\
no yesterday, no one after her,
(no one, no, no one else...)\\
For that moment, I don't need a name for it,\\
I just stay right here, for that moment...
(right here... right here with her...)
\songpart{Outro — trailing, implies continuation}
...and then I scroll again.
\end{lyrics}


\songtitle{In That Moment}

\tag*{2026}
\begin{notes}
Lyrics written in collaboration between Carlos Thompson and Claude, adapted fron \emph{For That Moment}.
\end{notes}
\begin{style}
trance, house, deep male bass-baritone vocal, half-spoken delivery, vocal subdued in the mix, wall of sound, dense atmospheric synth layers, reverb-heavy pads, pulsing sub-bass, mid-tempo build and drop structure, sensual, intimate, immersive
\end{style}
\begin{lyrics}
\songpart{Intro — sparse, filtered pad, vocal fragment looping, heavily reverbed}
(sung, distant, processed)\\
In that moment... she is all...\\
In that moment... she is all...

\songpart{Breakdown 1 — spoken, flat affect, minimal beat underneath}
Someone I like. A friend from real life.\\
A face I follow. Someone I've never met.\\
Online communities. Online celebs.\\
Doesn't matter where. Doesn't matter who.

\annotation{Rise 1 — sung/chanted, layers building, rhythmic fragment repeating with rising filter}\\
Doesn't matter where she's from...\\
doesn't matter who she is...\\
doesn't matter where... doesn't matter who...

\annotation{Drop 1 — full beat, vocal chop riding the rhythm}\\
She is her.\\
She is her.\\
Only her. Only her.\\
She is her.

\songpart{Breakdown 2 — spoken, pulled back, more space}
What is that feeling?\\
For each one, it's its own.

\annotation{Rise 2 — layered whispers/fragments stacking, building tension}\\
Platonic...\\
or admiration...\\
a lifetime...\\
or a night...

\annotation{Drop 2 — biggest drop, fuller arrangement, hook layered with Drop 1's phrase}\\
For that moment...\\
she is everything.\\
For that moment...
(she is her... she is her...)\\
No before, no after,\\
just for that moment.

\songpart{Outro — stripping back to intro's loop, sung, distant again, processed, fading out}
In that moment... she is all...\\
In that moment... she is all...\\
In that moment... she is all...\\
In that moment... she is all...
...
\end{lyrics}


\songtitle{Irrelevant Conversations}

\tag*{2023}\tag{bolero}
\begin{notes}
Co-written by Carlos Thompson and AI, based on a thread from 2023.
\end{notes}
\begin{style}
Bolero cubano with slow clave sway and warm guitar-tres pulse; verse leans intimate and story-driven, pre-chorus opens with rising harmony and a brief percussion lift, chorus blooms with lush strings and soft bongos, Lead vocal is close-mic and aching with doubled lines on key phrases, plus falsetto replies and gentle backing “oh”s, Maraca shimmers, piano fills, and a brushed accent before each chorus, Mix is warm, velvety, and cinematic, reading, bolero
\end{style}
\begin{lyrics}
\songpart{Verse 1}
I don’t drink, I don’t bend\\
I don’t smoke to make amends\\
I don’t chase a clouded night\\
Still I need a smaller flight

I had family on my chest\\
Words that only leave a test\\
Reality sat down hard\\
And it bruised me where I guard

\songpart{Pre-Chorus}
Then I found a window glow\\
On a screen I barely know\\
Strangers talking, line by line\\
On the edge of silly time

About religion, about truth\\
About nothing, sweet and smooth\\
And my heart began to sway\\
Like it heard a safer way

\songpart{Chorus}
Irrelevant conversations
(irrelevant conversations)\\
Save me from the pressure\\
Save me from the pressure

When the world comes crashing down\\
You can find me here around\\
Irrelevant conversations\\
They lift me off the ground

\songpart{Verse 2}
Someone argued over grace\\
Someone praised a sacred place\\
Someone said a joke so dry\\
I laughed before I knew why

My own problems stayed outside\\
At the door, where they could hide\\
For one hour, maybe two\\
I belonged to something new

\songpart{Pre-Chorus}
Not a bottle, not a flame\\
Just a question with no shame\\
Just a comment, just a thread\\
Just a little room in my head

And I felt the weight release\\
Like a hand that asked to ease\\
Like a storm behind a glass\\
Letting breath come back at last

\songpart{Chorus}
Irrelevant conversations
(irrelevant conversations)\\
Save me from the pressure\\
Save me from the pressure

When the world comes crashing down\\
You can find me here around\\
Irrelevant conversations\\
They lift me off the ground

\songpart{Bridge}
Maybe that’s the scary part\\
How a stranger mends my heart\\
With a thought that goes nowhere\\
With a joke, with doubt, with prayer

I was losing to the day\\
Then my mind could drift away\\
Not to vanish, not to run\\
Just to borrow some small sun

\songpart{Final Chorus}
Irrelevant conversations
(irrelevant conversations)\\
Save me from the pressure\\
Save me from the pressure

When the world comes crashing down\\
You can find me here around\\
Irrelevant conversations\\
They lift me off the ground

Irrelevant conversations
(irrelevant conversations)\\
I don’t need a grand solution\\
Just a place to breathe again
\end{lyrics}


\end{multicols*}

\chapter{Later Thoughts IV: Unspoken}
\albumcover{covers/internal/later-thoughts-iv-unspoken}
Three songs, one silence: love that stays exactly where it started because saying it out loud was never really an option. No confession, no closure -- just the quiet, and the noticing.
\clearpage\begin{multicols*}{4}
\songtitle{Your Role}

\tag*{2023}\tag{original-idea}\tag{ai-generated}\tag{punk-rock}

\begin{notes}
Poem written by ChatGPT based on a text by Carlos Thompson, not originally intended as lyrics.
\end{notes}

\begin{style}
punk rock
\end{style}

\begin{lyrics}
A comrade to share life's grand charade, -\\
of aging side by side, our bond displayed.

A partner for discourse on worlds unknown, -\\
fantasy's realm explored, minds overthrown.

A playful ally, flirty words we send, -\\
texts or chats, a playful blend.

An outlet found, desires to set free, -\\
no strings attached, a space just to be.

A cuddle buddy, warmth in nights so deep, -\\
movies' glow or slumber's sweep.

A collaborator, schemes we devise, -\\
to take the world, power will rise.

An explorer bold, our planet's space, -\\
remote places, the unknown embrace.

A sender sweet, kitties' cute delight, -\\
on screens they play, smiles ignite.

A guardian watchful, complacency's foe, -\\
keep me on track, where growth will flow.

Of these diverse roles, which to show, -\\
to fill with grace, where talents grow?

\end{lyrics}


\songtitle{Quiet Signal}

\tag*{2023}\tag{r-and-b}\tag{original-source}\tag{social-media-seed}

\begin{notes}
\end{notes}
\begin{style}
R\&B
\end{style}
\begin{lyrics}
\songpart{intro riff}

\annotation{male baritone vocals}

\songpart{Verse 1}
I scroll through feeds I barely know,\\
the names blur into grey,\\
then suddenly your avatar\\
lights up the endless day.

\songpart{Pre-Chorus 1}
I do not know if you can feel my gaze,\\
or sense you pull me through the darkest haze.

\songpart{Chorus}
You don't know what you do to me,\\
just living in your way,\\
a stranger in my timeline scrolls\\
and somehow makes my day.

\songpart{Verse 2}
And when your words sound tired or worn,\\
or silence fills your space,\\
a shadow falls across my chest,\\
I wear it like a trace.

\songpart{Pre-Chorus 2}
You carry on, unaware of what you send,\\
a signal reaching me around each bend.

\songpart{Chorus}
You don't know what you do to me,\\
just living in your way,\\
a stranger in my timeline scrolls\\
and somehow makes my day.

\annotation{riff}

\songpart{Bridge}
Maybe this is what connection means\\
when distance makes it safe—\\
to care about a stranger's smile,\\
to feel their weight from far away,\\
to wonder, just once,\\
if they'd believe you if you told them.

\songpart{Verse 3}
I'll never send a message out,\\
I'll keep it here inside,\\
just knowing you exist somewhere\\
is all I need tonight.

\songpart{Pre-Chorus 3}
You post and scroll and do not even know,\\
a quiet light that helps my river flow.

\songpart{Chorus}
You don't know what you do to me,\\
just living in your way,\\
a stranger in my timeline scrolls\\
and somehow makes my day.

\songpart{Outro}
And maybe that's enough—\\
to see you, briefly,\\
in the flood of everything,\\
and feel the world\\
get slightly warmer,\\
slightly less alone.

\annotation{fade}
\end{lyrics}


\songtitle{No Excuse}

\tag*{2023}\tag{original-seed}\tag{social-media-seed}\tag{alternative-rock}

\subsection{Notes}
Lyrics by Gemini based on a 2023 tweet by Carlos Thompson, with context from other social media posts by the author.

\subsection{Style}
Alternative rock with post-grunge influences. %
Distorted electric guitars play palm-muted power chords in the verses, opening into wide, sustained chords in the chorus. %
A clean electric guitar provides melodic counterpoint with a bright, slightly overdriven tone. %
The bass guitar follows the kick drum pattern with a thick, overdriven low-end. %
Drums feature a heavy snare with a short decay and consistent eighth-note hi-hat patterns. %
Male vocals are delivered with a gritty, raspy timbre and moderate compression. %
The track is in the key of D major at 128 BPM in 4/4 time. %
Rap part at an even lower tone. %
Outro conversational

\begin{lyrics}
\songpart{Intro}
\annotation{distorted electric guitar riff, steady drum beat}

\songpart{Verse 1}
\annotation{baritone vocals, solo}\\
I push the morning to the night,\\
And put the tasks away,\\
I swear I’ll make the wrong things right,\\
But never find the day.

\songpart{Verse 2}
\annotation{baritone vocals, solo}\\
The shadows lengthen on the wall,\\
I watch the hours retreat,\\
I wait to face the heavy fall,\\
And taste my own defeat.

\annotation{distorted electric guitar riff, steady drum beat}

\songpart{Pre-Chorus}
\annotation{palm-muted electric guitar, bass, and drums}\\
\annotation{baritone vocals, solo}\\
I know my fault I see the wrong I made\\
But with such savage spite you claim the score\\
The heavy price you force is so dismayed\\
I have no will to state my grief no more

\songpart{Chorus}
I missed the mark and let it drop, (Let it drop)\\
I feared the storm before it grew,\\
Now sparks of anger will not stop, (Will not stop)\\
The friction burns between us two,\\
So I retreat behind a wall,\\
And run away to avoid it all. %
(Avoid it all)

\songpart{Verse 3}
\annotation{baritone vocals, solo}\\
The constant friction wears me down,\\
Like stone beneath the sea,\\
A heavy silence fills the town\\
And drains the life from me.

\annotation{distorted electric guitar riff, steady drum beat}

\songpart{Pre-Chorus}
\annotation{palm-muted electric guitar, bass, and drums}\\
\annotation{baritone vocals, solo}\\
I see my petty stubbornness and shame,\\
Yet bear the stinging guilt of all your blame.

\songpart{Chorus}
I missed the mark and let it drop, (Let it drop)\\
I feared the storm before it grew,\\
Now sparks of anger will not stop, (Will not stop)\\
The friction burns between us two,\\
So I retreat behind a wall,\\
And run away to avoid it all. %
(Avoid it all)

\songpart{Bridge}
\annotation{tenor vocals}\\
\annotation{Rap}\\
(Yo) It started with delay, just pushing off the slate,\\
ut turning blind eyes only magnified the weight.\\
The failure bred anxiety, the panic turned to friction,\\
And now this gridlock intimacy is a toxic addiction.\\
Your spite feeds my avoidance, my avoidance feeds your spite,\\
We’re spinning in a circle where nobody is right.\\
Trapped inside the loop of the things I didn’t do,\\
Now I’m hiding from myself while I'm hiding out from you.

\songpart{Verse 4}
\annotation{baritone vocals, solo}\\
I feel the burden in my chest,\\
A sharp and bitter bite,\\
I want to give my very best,\\
And step into the light.

\annotation{distorted electric guitar riff, steady drum beat}

\songpart{Pre-Chorus}
\annotation{palm-muted electric guitar, bass, and drums}\\
\annotation{baritone vocals, solo}\\
This heavy guilt just paralyzes mind,\\
It builds no bridge, leaves better days behind.

\songpart{Chorus}
I missed the mark and let it drop, (Let it drop)\\
I feared the storm before it grew,\\
Now sparks of anger will not stop, (Will not stop)\\
The friction burns between us two,\\
So I retreat behind a wall,\\
And run away to avoid it all. %
(Avoid it all)

\songpart{Outro}
\annotation{conversational}\\
The music fades out...\\
No more structures. %
No more rhymes.\\
No more counting syllables.\\
I am just standing here\\
\annotation{Spoken word}.\\
I am sorry.\\
I messed up, and I have absolutely no excuse.
\end{lyrics}


\end{multicols*}

\part{Soul and Turmoil}

\chapter{Soul}
\albumcover{covers/internal/soul}
The same ache in five arrangements: letting go, craving stimuli, wanting skin against skin, drifting between other people's inner worlds and a coffee that was never shared. \emph{Soul} is the body's honest half of the songbook -- unfiltered, uncomplicated by argument, just wanting.
\clearpage\begin{multicols*}{4}
\songtitle{Dejándome llevar}

\tag*{2026}\tag{vallenato}\tag{original-lyrics}\tag{spanish-lyrics}
\begin{notes}
Written by Carlos Thompson.
\end{notes}
\begin{style}
vallenato puya, Colombian Caribbean, 68 BPM, accordion lead, caja vallenata, guacharaca rhythm, male baritone, almost-spoken delivery, trembling lament, stumbling intro, steady chorus, rubato verse, syncopated 6/8, dry close-mic vocal, plate reverb, tape saturation, live room ambience
\end{style}
\begin{lyrics}
\songpart{Verse 1}
Vi tu imagen en mi caminar\\
vi tu imagen en mi caminar\\
vi tu imagen en mi caminar\\
y mi mundo por ti cambió.

No importaba nada más\\
no importaba nada más\\
no importaba nada más\\
por el instante en que sucedió.

\songpart{Pre-Chorus}
Estabas feliz, te veías feliz\\
estabas feliz, te veías feliz\\
estabas feliz, te veías feliz\\
Y por ese instante feliz fui.

\songpart{Chorus}
La vida es simple cuando estoy\\
la vida es simple cuando estoy\\
la vida es simple cuando estoy\\
dejándome llevar del momento.

Pero el mundo sigue su ritmo\\
pero el mundo sigue su ritmo\\
pero el mundo sigue su ritmo\\
y al momento no me puedo asir.

\songpart{Verse 2}
Vi lo que dijiste de tu pérdida\\
vi lo que dijiste de tu pérdida\\
vi lo que dijiste de tu pérdida\\
y esa pérdida la sentí aquí.

Mi corazón está con tu tristeza\\
mi corazón está con tu tristeza\\
mi corazón está con tu tristeza\\
por el tiempo que estás ahí.

\songpart{Pre-Chorus}
Y estás ahí por dos minutos\\
y estás ahí por dos minutos\\
y estás ahí por dos minutos\\
antes de dar el siguiente scroll.

\songpart{Chorus}
La vida es simple cuando estoy\\
la vida es simple cuando estoy\\
la vida es simple cuando estoy\\
dejándome llevar del momento.

Pero el mundo sigue su ritmo\\
pero el mundo sigue su ritmo\\
pero el mundo sigue su ritmo\\
y al momento no me puedo asir.

\songpart{Bridge}
Es un gatito en desgracia\\
es un gatito en desgracia\\
es un gatito en desgracia\\
quien ocupa ahora tu lugar.

Es un video de autoayuda\\
es un video de autoayuda\\
es un video de autoayuda\\
O un recuerdo que el mundo está mal.

\songpart{Final Chorus}
La vida es simple cuando estoy\\
la vida es simple cuando estoy\\
la vida es simple cuando estoy\\
dejándome llevar del momento.

Pero el mundo sigue su ritmo\\
pero el mundo sigue su ritmo\\
pero el mundo sigue su ritmo\\
y al momento no me puedo asir.

El momento es todo pero pasa\\
el momento es todo pero pasa\\
el momento es todo pero pasa\\
y todo al final sigue igual.
\end{lyrics}


\songtitle{Stimuli}

\tag*{2026}\tag{original-lyrics}\label{sec:stimuli}
\begin{notes}
Written by Carlos Thompson
\end{notes}
\begin{style}
dubstep, club house, 124 BPM, syncopated four-on-the-floor, half-time drop, deep bass baritone half-spoken vocal, close-mic delivery, vocal texture layering, brass stabs, synth bass stabs, sub bass glide, hand percussion, tape saturation, plate reverb, sidechain pumping, stumbling groove, steady lift
\end{style}
\begin{lyrics}
Stimuli.

Could be a conversation.\\
Could be dumb scrolling.\\
Could be creating a new world in my mind.\\
Could be sex for the sake of sex.

I crave stimuli.

Something to distract my mind.\\
Something to cause my idle thoughts.\\
Warm skin against my skin\\
or plain intelecttual stimulation.

I crave stimuli.

Keep my mind distracted.\\
Keep my skin excited.\\
Keep my imagination checked.

Something new.\\
Something old that haven't be with me\\
for time I haven't counted.

I crave stimuli.

A new place.\\
A new partner.\\
A new challenge.\\
Or just more skin.

I crave stimuli.

I crave.
\end{lyrics}


\songtitle{Piel}

\tag*{2026}\tag{techno-house}\tag{spanish-lyrics}\tag{original-lyrics}\label{sec:piel}
\begin{notes}
Original version of \emph{Piel}, written by Carlos Thompson.
\end{notes}
\begin{style}
techno house, trance music, 808 sub-bass, distorted synthesizer stack, dubstep bass drop, close-mic baritone, almost spoken vocal, textured vocal layers, soundwall breakdowns, silence gaps, volume swells, machine-like percussion, driving 128 BPM, anxious desire
\end{style}
\begin{lyrics}
Si no te hace falta, a mí sí.

Mi piel pide piel.

Mi cuerpo pide.\\
Pide descargas.\\
Pide estímulos que desconectan la mente.\\
Pide toque,\\
pide fricción,\\
pide exitación\\
y hasta cierto punto me basto.

Hasta cierto punto.

Hasta cierto punto mi imaginación cumple.\\
Hasta cierto punto el porno es suficiente.\\
Hasta cierto punto el chatbot funciona.

He probado sexting y sirve.\\
Hasta cierto punto.

Pero nada de eso tiene piel.

Hasta cierto punto me basto,\\
pero a partir de ahí no es suficiente.

Si no te hace falta, a mí sí.

Amar realmente supongo que también sirve.\\
Supongo que hará falta también.

Amar realmente,\\
lo que quiera que eso signifique.\\
Si es que significa algo.

Por ahora solo sé que me falta piel.

Compartir con alguien de carne y hueso,\\
que esté ahí.

Si no te hace falta, a mí sí.
\end{lyrics}


\songtitle{Living Fantasies (electronic)}

\tag*{2026}\tag{electronic}\tag{original-lyrics}\tag*{english-lyrics}\tag{living-fantasies}\label{sec:living-fantasies}
\begin{notes}
Written by Carlos Thompson.
\end{notes}
\begin{style}
electronic anthem, club house, trance bassline, industrial synth walls, distorted power chords, soprano lead vocals, texture vocals, vocal delay throws, plate reverb, 138 BPM, four-on-the-floor kick, driving bass drums, breakdown buildup, pre-chorus lift, cathedral reverb, dungchen blasts, bass drum hits, bittersweet futurism, frenzied momentum
\end{style}
\begin{lyrics}
living fantasies\\
building them\\
an AI image prompt\\
an archetype

living fantasies\\
building them\\
alternate histories\\
alternate self

living fantasies\\
building them\\
non-existing languages\\
and polities

living fantasies\\
building them\\
bogota metro lines\\
for the enth time

living fantasies\\
building them\\
other planets\\
alien cultural forms

living fantasies\\
building them\\
then reinventing\\
fantasy itself

living fantasies\\
building them\\
just to feel something\\
just to numb myself
\end{lyrics}


\songtitle{Vueltas y vueltas}

\tag*{2026}\tag{original-lyrics}
\begin{notes}
Written by Carlos Thompson.
\end{notes}
\begin{style}
dream pop, electroswing, 124 BPM, baritone lead vocals, stacked texture vocals, syncopated brass stabs, walking upright bass, slapback guitar riffs, brushed drum kit, staccato hi-hats, plate reverb, tape saturation, sidechain pumping, swung shuffle groove, driving curiosity, urban reflection
\end{style}
\begin{lyrics}
\songpart{Verse 1}
Una tarde fría con sol brillante\\
un cielo azul contra el ladrillo\\
la ciudad viviendo en movimiento\\
viento helado, ruido y brillo.

Mis pensamientos con cada paso\\
contando las baldosas del andén\\
observando a la gente pasar\\
leyendo el aviso de un almacén.

\songpart{Pre-Chorus}
Miro, pienso y reflexiono\\
lo que es o lo que podría ser\\
una mente que no se apaga\\
pero tampoco sirve a la vez

\songpart{Chorus}
y da vueltas y vueltas (y vueltas)\\
conectando pensamientos\\
vueltas y vueltas (y vueltas)\\
entre intentos y momentos\\
vueltas y vueltas (y vueltas)\\
todo a los cuatro vientos

\songpart{Verse 2}
Un motociclista en contravía\\
una reflexión sobre el civismo\\
rediseñar las leyes o la ciudad\\
un nuevo ensayo de urbanismo.

Y alguien cruza frente a mi vista\\
una mujer extraña pero interesante\\
imaginarme ser un donjuán sin miedo\\
vivir la vida a cada instante.

\songpart{Pre-Chorus}
Voy, vengo y me enamoro\\
lo que es o lo que podría ser\\
de ideas sueltas y de personas\\
de cosas que nada que ver

\songpart{Chorus}
y da vueltas y vueltas (y vueltas)\\
conectando pensamientos\\
vueltas y vueltas (y vueltas)\\
entre intentos y momentos\\
vueltas y vueltas (y vueltas)\\
todo a los cuatro vientos

\songpart{Bridge}
¿Y si a Napoleón lo atrapan en Egipto?\\
¿Cómo invertir la plata que no tengo?\\
¿Espuma moldeable para organizar el baúl?\\
¿Cómo instalar estaciones de carga en los garajes?\\
¿Una página para armar computadores?\\
Ese local está en arriendo ¿qué podría poner ahí?

\songpart{Final Chorus}
y da vueltas y vueltas
(vueltas y vueltas)\\
conectando pensamientos\\
vueltas y vueltas
(vueltas y vueltas)\\
entre intentos y momentos\\
vueltas y vueltas
(vueltas y vueltas)\\
hasta perder los alientos\\
vueltas y vueltas
(vueltas y vueltas)\\
todo a los cuatro vientos
\end{lyrics}


\songtitle{Mundos internos}

\tag*{2026}\tag{synth-pop}\tag{original-lyrics}\tag{spanish-lyrics}
\begin{notes}
Written by Carlos Thompson, original.
\end{notes}
\begin{style}
candy pop, synth-pop, 132 BPM, mellow male vocals, close-mic tenor, drum machine backbeat, chunky electric bass, bright synth arpeggios, analog polysynth pads, gated snare, tape saturation, plate reverb, sidechain pumping, retro 80s mix, urban routine, bittersweet introspection
\end{style}
\begin{lyrics}
\songpart{Verse 1}
Un espresso de mi máquina de capuchino,\\
con leche caliente, cual si fuera un latte,\\
rollos de canela de supermercado barato\\
y las pastillas de la tensión y de la cabeza.

Leer y contestar preguntas en Quora\\
o perderme en Instagram en vídeos cortos,\\
perder el sentido del tiempo y correr luego,\\
salir a la calle a esperar en el paradero.

\songpart{Pre-Chorus}
Una ciudad con sol o con lluvia,\\
con vías en construcción y vida en la calle,\\
motos que serpentean y bicicletas que fluyen,\\
peatones y carros, buses, ruido y vida.

\songpart{Chorus}
Y en mi cabeza, otros mundos y otros lugares:\\
un podcast, música o mis propios mundos,\\
observar la calle o mi propio teléfono,\\
lo que pasa afuera y en mi mundo interno.

\songpart{Verse 2}
La rutina en la oficina, que rara vez es rutina,\\
justificar horas que no pasan en sincronía,\\
justificando una rutina y una presencia\\
con máscaras y alt-tabs, raras sinfonías.

Hora de almuerzo, compañeros de oficina,\\
cada uno con su presupuesto y sus historias,\\
el trabajo en común o la última noticia,\\
lo que transcurre en nuestras vidas notorias.

\songpart{Pre-Chorus}
Una ciudad con sol o con lluvia,\\
centros comerciales, tiendas de esquina,\\
salir al almuerzo o microondas en la oficina,\\
y tomarnos un café como parte de la rutina.

\songpart{Chorus}
Y en mi cabeza, otros mundos y otros lugares:\\
un podcast, música o mis propios mundos,\\
observar la calle o mi propio teléfono,\\
lo que pasa afuera y en mi mundo interno.

\songpart{Bridge}
La vida transcurre,\\
la vida pasa,\\
hay otros que cruzan nuestros lugares,\\
trabajo que hacer,\\
historias que contar,\\
coordinar, vivir:\\
un mundo en que no estás solo.

\songpart{Breakdown}
Lo que haces y lo que no haces,\\
compañeros de trabajo que de ti dependen,\\
una familia que espera,\\
una vida que pasa.

\songpart{Final Chorus}
Y en mi cabeza, otros mundos y otros lugares:\\
un podcast, música o mis propios mundos,\\
observar la calle o mi propio teléfono,\\
lo que pasa afuera y en mi mundo interno.

Un mundo interno.
\end{lyrics}


\songtitle{A Coffee Not Shared}

\tag*{2026}\tag{dubstep}\tag{original-lyrics}
\begin{notes}
Written by Carlos Thompson, original lyrics.
\end{notes}
\begin{style}
dubstep, industrial electronic, club house, trap-step bass, distorted synth stabs, metallic bass syncopation, vocal chops, textured vocal layers, sidechain pumping, tape saturation, plate reverb, four-on-the-floor kick, halftime breakdown, pre-chorus lift, drop arrangement, 138 BPM, bittersweet longing, cinematic urgency
\end{style}
\begin{lyrics}
\songpart{Verse 1}
You (you)\\
you've been a crush since I remember.

You (you)\\
you've been hitting my timeline with your presence.

You (you)\\
you're a coffee I've failed to ask you (you).

You (you)\\
tender words and engaging thoughts.

\annotation{Instrumental break}

\songpart{Verse 2}
You (you)\\
years of longing and falling in love.

You (you)\\
you're radiant and I'm a deer in headlights.

You (you)\\
many words I'd love to share with you (you).

You (you)\\
a guilty pleasure; my shameless thoughts.

\songpart{Breakdown}
I've been longing for eons,\\
falling in love with your presence,\\
failing to feel deserving to reach you (you).
'Cause it's easier abstinence,
'cause cowardice feels safer than trying,
'cause you (you)\\
would destroy me.

\songpart{Outro}
You (you)\\
a story not happening.

You (you)\\
a coffee not shared.

You (you)\\
a suspended conversation.

You (you)\\
headlights that'll crush me.
\end{lyrics}


\end{multicols*}

\chapter{Inner turmoil}
\albumcover{covers/internal/inner-turmoil}
\emph{The Wrong Map} sets the tone: a mind that keeps redrawing its own directions, craving novelty it can't sustain, faulting its own character before anyone else gets the chance to. \emph{Used to Say} closes the chapter looking back at an old hatred it can no longer feel -- the map, it turns out, was rewritten after all.
\clearpage\begin{multicols*}{4}
\songtitle{The Wrong Map}

\tag*{2026}\tag{orignal-lyrics}
\begin{notes}
Written by Carlos Thompson
\end{notes}
\begin{style}
classic rock, baritone lead, close-mic confessional, gritty vocal fry, gang choir chorus, distorted rhythm guitar, overdriven bass, pedal steel fills, dynamic breakdowns, 138 BPM, galloping eighths, halftime chorus, fluid vocals, refrained vocals, whispered endings, mellow delivery, D minor / F major alternations
\end{style}
\begin{lyrics}
\songpart{Verse 1}
There was a time I was younger and more naïve\\
and I believed the world functioned in simple ways:\\
school, college, marriage, kids, just to achieve\\
a normal ladder of expectations to live your days.

Choose the right major, choose the right career.\\
Your strengths would tell what your path becomes.\\
Nine to five, raise a family, go with friends and grab a beer.\\
Follow your dreams and don't worry about the outcomes.

\songpart{Pre-Chorus}
My path was easy from my background,\\
hard by the ironies of life.\\
I came from a privileged household\\
and a mind that became a strife.

\songpart{Chorus}
Too perfect, almost a trap. (it was a trap)\\
A privileged environment,\\
a gifted intellect,\\
but I was given the wrong map. (with the wrong map)\\
A perfect experiment,\\
until all was wrecked
(it was a trap with the wrong map)

\songpart{Verse 2}
There was a time I was younger and more naïve.\\
I knew life was a struggle for many; I didn't expect it for me.\\
Wrong decisions in hindsight, when to stay, when to leave;\\
and each stumbling block shaping me and collecting a fee.

School, college, marriage, kids. Been there, done that.\\
A normal path on paper, a constant struggle with my mind.\\
A mind that told me I wasn't good enough, no matter what,\\
achievements I didn't deserve, and every defeat was signed.

\songpart{Pre-Chorus}
You learn to distrust yourself\\
as adult life starts to kick in.\\
A genius kid without structure,\\
a maladapted human spin.

\songpart{Chorus}
Too perfect, almost a trap. (it was a trap)\\
A privileged environment,\\
a gifted intellect,\\
but I was given the wrong map. (with the wrong map)\\
A perfect experiment,\\
until all was wrecked
(it was a trap with the wrong map)

\songpart{Bridge}
School, college, marriage, kids...\\
nine-to-five, grab a beer...\\
a career you love — what's so hard?\\
For a kid, the future is not fear.

You're smart, you'll figure out\\
your right path in what was to come.\\
I was shy, I was reserved,\\
not a big deal 'cause I wasn't dumb.

A stable family as a role model,\\
a perfect environment to thrive.\\
Not their fault I messed up —\\
a perfect storm building in my life.

\songpart{Final Chorus}
Too perfect, almost a trap. (it was a trap)\\
A privileged environment,\\
a gifted intellect,\\
but I was given the wrong map. (with the wrong map)\\
A perfect experiment,\\
until all was wrecked

School, college, marriage, kids...\\
nine-to-five, grab a beer...\\
I was younger and more naïve\\
with a plan that seemed quite clear.
\end{lyrics}


\songtitle{Multiply the Delay}

\tag*{2026}
\begin{notes}
Lyrics by SUNO based on part of a conversation between Carlos Thompson and ChatGPT.
\end{notes}
\begin{style}
Mellow reggae with a laid-back one-drop groove, warm skanking guitar on the offbeat, round bassline, and soft rimshot snare; verse stays sparse with dry vocal and roomy percussion, pre-chorus opens with harmony pads and filtered organ swell, chorus lands with stacked call-and-response hooks and subtle dub delay throws on the anchor phrase, Intimate close-mic lead, gentle doubles in the chorus, occasional ad-lib echoes, tape hiss and reversed guitar swells between sections, smooth and earthy mix with a sun-worn pulse, reggae, mellow
\end{style}
\begin{lyrics}
\songpart{Verse 1}
It don't just stack, baby\\
It multiply

One part won't push\\
It takes away the start

A signal go missing\\
Before the hand can move

Then the next one whisper
"Stay right where you are"

\songpart{Pre-Chorus}
And if you never begin\\
You never can fail

That stillness feel safer\\
Than stepping on the rail

\songpart{Chorus}
They multiply, multiply\\
They multiply, my baby\\
They multiply, multiply\\
That's how the whole thing stay

No one thing alone\\
It's the way they play

They multiply, multiply\\
And delay feels like a plan

\songpart{Verse 2}
Then comes the fear of losing\\
Before the game even starts

So not moving on the page\\
Feels lighter than the mark

The outside pressure fade out\\
Like rain on dry wood

No crowd in the doorway\\
To say you really should

\songpart{Pre-Chorus}
And if the world get quiet\\
The pattern stay warm

Each piece fit the next one\\
Like shelter from a storm

\songpart{Chorus}
They multiply, multiply\\
They multiply, my baby\\
They multiply, multiply\\
That's how the whole thing stay

No one thing alone\\
It's the way they play

They multiply, multiply\\
And delay feels like a plan

\songpart{Bridge}
Then the mind get clever\\
With a shiny little line

"Maybe this is strategy\\
Maybe this is timing"

It sounds so proper\\
When the reason can talk

So the pause wear a suit\\
And the symptom learn to walk

\songpart{Chorus}
They multiply, multiply\\
They multiply, my baby\\
They multiply, multiply\\
That's how the whole thing stay

No one thing alone\\
It's the way they play

They multiply, multiply\\
And delay feels like a plan
\end{lyrics}


\input{temp/the_fault_in_your_character}
\songtitle{A Mind That Craves Novelty}

\tag*{2026}\tag{electropunk}\tag*{english-lyrics}\tag{original-lyrics}\tag{grinder-2026}\label{sec:grinder}
\begin{notes}
Written by Carlos Thompson.
\end{notes}
\begin{style}
electropunk, genre-blending theater rock, spoken-word intro, male baritone lead, shouted chorus hook, mechanical percussion, distorted synth bass, jagged guitar stabs, arpeggiated synth lead, orchestral hits, call-and-response backing vocals, tape saturation, plate reverb, abrupt section shifts, build-drop contrast, restless novelty, anxious self-interrogation, 132 BPM, syncopated half-time pulse
\end{style}
\begin{lyrics}
\songpart{Verse 1}
Inside myself there are many minds,\\
and often I can't tell which of them binds.\\
A mind that wants to fly, to create,\\
and does not care much about its fate.

A mind that knows how constantly I fail,\\
and keeps nitpicking every detail.\\
A mind that gets bored with the mundane,\\
and treats every friction with disdain.

\songpart{Pre-Chorus}
Too much potential,\\
and yet underperforming.\\
Keep failing,\\
too much brainstorming.

\songpart{Chorus}
Feeling trapped in doing too little,\\
a mind that wants to fly,\\
a mind that abhors the boring,\\
a mind that craves novelty.

(a mind)
(a void)
(bored)

\songpart{Verse 2}
A mind that seeks escape in longing by\\
what both religion and social justice deny.\\
A mind that finds solace in doing naught,\\
just losing itself in idle thought.

Too ambitious to ever complete,\\
Too unreliable on goals I can meet.\\
Too boring, too sinful, too lazy,\\
just short of feeling I am crazy.

\songpart{Pre-Chorus}
Too much idle fantasy\\
that I can't share with folk.\\
Sinful or detached,\\
like living in an insane joke.

\songpart{Chorus}
Feeling trapped in doing too little,\\
a mind that wants to fly,\\
a mind that abhors the boring,\\
a mind that craves novelty.

(a mind)
(a void)
(bored)

\songpart{Bridge}
Lack of intellectual stimulation\\
and the friction of the grinder.\\
A skin that craves too much\\
and becomes a constant reminder.

A void in cravings,\\
a void in intensity,\\
a void in the space\\
around my fence.

\songpart{Final Chorus}
Feeling trapped in doing too little,\\
a mind that wants to fly,\\
a mind that abhors the boring,\\
a mind that craves novelty,\\
a mind that idles through,\\
a mind that just fails.

(bored)
\end{lyrics}


\songtitle{Connection}

\tag*{2026}\tag{bolero}\tag{english-lyrics}\tag{original-lyrics}
\begin{notes}
Written by Carlos Thompson.
\end{notes}
\begin{style}
electronic cuban bolero, 92 BPM, male baritone vocals, close-mic velvet delivery, call-and-response chorus, clave percussion, stand-up bass, muted trumpet stabs, tenor sax fills, synth brass pads, 808 sub-bass, drum kit backbeat, tape saturation, plate reverb, mono low end, crowd chant hook, bittersweet longing
\end{style}
\begin{lyrics}
\songpart{Verse 1}
Sex without compromise\\
there are people who get paid for it\\
some are exploited\\
engaging there would be rape\\
some are people for whom\\
sex without compromise\\
is not big deal\\
why not getting a profit from it?

\songpart{Pre-Chorus}
Sex without compromise\\
why would it be different\\
for who pays than for who gets paid?\\
unless there is exploitation.

\songpart{Chorus}
I crave connection (connection)\\
I crave skin (warm skin)\\
I crave a warm body in sync with mine

I have been made ashamed of my desires
(ashamed of what I want)\\
not the right stand to defend

\songpart{Verse 2}
Sex without compromise\\
how is that different between\\
a person who engages in casual hookups\\
from a person who gets a profit?\\
free agreements between consenting adults\\
sex without compromise\\
beyond what you agree for that night\\
a diner, a fee, or mutual enjoyment for free

\songpart{Pre-Chorus}
Sex without compromise\\
I would be all in\\
for whatever agreement I could afford\\
which isn't much

\songpart{Chorus}
I crave connection (connection)\\
I crave skin (warm skin)\\
I crave a warm body in sync with mine

I have been made ashamed of my desires
(ashamed of what I want)\\
not the right stand to defend

\songpart{Bridge}
I don't have much money\\
I don't have much time\\
I am not a stallion\\
who would blow your mind

I'm just a needy little old man\\
philosophizing on what I can't afford\\
sex without compromise\\
a fantasy of my own

\songpart{Final Chorus}
I crave connection (connection)\\
I crave skin (warm skin)\\
I crave a warm body in sync with mine

I have been made ashamed of my desires\\
I have been made ashamed of what I want\\
I have been made ashamed of a situation\\
in which I have paved my self in\\
no much money, no much time\\
and not even begging for free
\end{lyrics}


\songtitle{Used to Say}

\tag*{1988}\tag*{2026}\tag{house}\tag{original-lyrics}\tag{spanish-lyrics}\tag{je-hais-1988}
\begin{notes}
Original poem \emph{Odio} by Carlos Thompson from 1988, with new lyrics written by Carlos Thompson in 2026.
\end{notes}
\begin{style}
House fused with restrained dubstep and industrial trance at 138 BPM: four-on-the-floor pulse, wobble sub drops, distorted stabs, piano-led tension, metallic percussion, club synth leads, whispered Spanish male intro transitioning into breathy English-belted refrains, sidechain lift, plate reverb and warm tape saturation
\end{style}
\begin{lyrics}
\songpart{Intro, whispered, spoken}
Odio lo que siento por ti\\
y odio lo que me haces sentir.\\
Odio el amarte y me odio por eso.\\
Odio el odiarte y me odio por eso.\\
Odio lo que hago por ti\\
y odio lo que no hago por ti\\
y odio todo mi odio.\\
Tanto odio lo que odio\\
que termino odiando todo.

\annotation{Build 1}\\
I used to say how much I hated myself,\\
how I hated the object of my unrequited love.\\
Years later I couldn’t make myself feel\\
that hate is something I could harbor.

Is hate a part of a normal human life,\\
one of those conditions I’m deficient at,\\
or have I realized hate is just something\\
that doesn’t do much for my sanity?

\annotation{Drop}\\
I can’t hate, I can’t hate,\\
a feeling not in my heart,\\
a feeling not in my character.

I used to say I hated—\\
what changed?

\annotation{Build 2}\\
Or maybe\\
I have just redefined hate.\\
I called hate just any discomfort;\\
now I define a threshold\\
that I can’t reach.

Or is it my disconnect with humanity,\\
a humanity I constantly fall in love with\\
but that I can’t fully comprehend,\\
not fully love,\\
not fully hate.

\annotation{Drop}\\
I can’t hate, I can’t hate,\\
a feeling not in my heart,\\
a feeling not in my character.

I used to say I hated—\\
what changed?

\songpart{Outro}
Not fully love.\\
Nor fully hate.
\end{lyrics}


\end{multicols*}

\part{Faith and Argument}

\chapter{Pray, Brother I: Apologetics}
\albumcover{covers/internal/pray-brother-i-apologetics}
Faith argued the way he argues everything else: point by point, source by source, most of it starting life as a Quora answer to a stranger's question. Doubt gets equal time with belief here -- nothing is taken on faith just because it's about faith.
\clearpage\begin{multicols*}{4}
\songtitle{Slightest Amount}

\tag*{2026}
\begin{notes}

\end{notes}
\begin{style}
Dubstep hymn with gospel choir, halftime wobble bass, church claps on the backbeat, massive sub drops under the verses, verse sections starting sparse with organ and solo lead, pre-chorus opening into stacked choir harmonies, chorus hitting with full choir chants and chopped vocal stabs, bridge dropping to near-silence then slamming back with bass and handclaps, close-mic lead vocal with cathedral reverb, choir doubles on the hook, reversed swells into each drop, bright and thunderous mix, emotional, dubstep, gospel
\end{style}
\begin{lyrics}
\songpart{Verse 1}
A slight amount?\\
Yeah, maybe.\\
Not a stone-set proof.\\
Not a burning sky.\\
It could be feeling.\\
It could be fear.\\
It could be a room\\
where my name gets called.

\songpart{Pre-Chorus}
But don’t ask me\\
to bow to a tale,\\
if it can’t hold up\\
when the questions hit.\\
I need more than awe,\\
more than a shiver.\\
I need to know\\
what stands when it breaks.

\songpart{Chorus}
The slightest amount\\
could move my heart
(oh Lord)\\
The slightest amount\\
could crack my dark
(amen)\\
If it’s real, then show me\\
If it’s true, then hold me\\
The slightest amount\\
could move my heart

\songpart{Verse 2}
If you say “God,”\\
I hear a claim.\\
I hear a name\\
with no fence around it.\\
A god, any god,\\
show me the edge,\\
show me the line\\
where it stops being smoke.

\songpart{Pre-Chorus}
I’m not closed.\\
I’m just careful.\\
If the ground is holy,\\
let it carry weight.\\
If the voice is living,\\
let it answer back.\\
I’ve watched too much\\
fall through my hands.

\songpart{Chorus}
The slightest amount\\
could move my heart
(oh Lord)\\
The slightest amount\\
could crack my dark
(amen)\\
If it’s real, then show me\\
If it’s true, then hold me\\
The slightest amount\\
could move my heart

\songpart{Bridge}
If there is a God,\\
then make it plain.\\
Not just a feeling\\
passing through the rain.\\
I’ve seen the flaw\\
in every claim\\
Every single argument\\
can break the same

\songpart{Breakdown}
Still—\\
I’m here.\\
Still—\\
I’m listening.\\
Still—\\
If you’re there,\\
then come closer.

\songpart{Final Chorus}
The slightest amount\\
could move my heart
(oh Lord)\\
The slightest amount\\
could crack my dark
(amen)\\
If it’s real, then show me\\
If it’s true, then hold me\\
The slightest amount\\
could move my heart
(hold me)\\
The slightest amount\\
could move my heart
\end{lyrics}


\songtitle{Complete By What}

\tag*{2026}
\begin{notes}

\end{notes}
\begin{style}
Dubstep hymn with gospel choirs, halftime wobble bass and stomp-clap percussion; verse rides sparse organ stabs and low brass, pre-chorus strips to handclaps and rising choir pads, chorus hits with huge stacked gospel voices and a snarling bass drop. %
Lead vocal is close-mic and sermon-like with doubled ad-libs, choir answers in call-and-response, delay throws on key questions. %
Ear candy: reverse swells into each drop, sub-bass growls under the hook, a last-chorus key change with shouted amen tags. %
Mix is dark, wide, and cathedral-bright, gospel, dubstep
\end{style}
\begin{lyrics}
\songpart{Verse 1}
You say, “It seems”\\
Like that’s a wound\\
I hear a question\\
Wearing a crown

What must be finished?\\
What must be claimed?\\
A sky with no owner\\
Still bears a name

Does it need deities?\\
Why would it?\\
Because you said so\\
Because you said it

\songpart{Pre-Chorus}
What is purpose?\\
A tool in a hand\\
A part in a plan\\
A mark in the sand

If you need meaning\\
To make you feel whole\\
That’s your hunger talking\\
Not a hole in the soul

\songpart{Chorus}
Complete by what?\\
Complete by what? (amen)\\
Complete by what?\\
Who sets that law? (amen)

A world can be whole\\
Without your frame\\
A world can be whole\\
Without your name

Complete by what?\\
Complete by what? (amen)

\songpart{Verse 2}
You ask for utility\\
Inside the machine\\
Like truth has to serve\\
Some throne you’ve seen

But a universe doesn’t\\
Need to obey\\
Your need for a ladder\\
Your need for a way

Maybe your “purpose”\\
Is only a mask\\
A story you borrow\\
To soften the task

Maybe a god\\
Wouldn’t flatter you\\
Maybe your role\\
Would cut right through

\songpart{Pre-Chorus}
And if there is one\\
Then hear me clear\\
Your joy is not promised\\
Your comfort not near

Could be your calling\\
Is burden and dust\\
Could be your “blessing”\\
Is learning to trust

\songpart{Chorus}
Complete by what?\\
Complete by what? (amen)\\
Complete by what?\\
Who sets that law? (amen)

A world can be whole\\
Without your frame\\
A world can be whole\\
Without your name

Complete by what?\\
Complete by what? (amen)

\songpart{Bridge}
You call it empty\\
I call it open\\
You call it broken\\
I call it unspoken

No transcendence\\
Required to stand\\
No hidden signature\\
No master plan

If you want purpose\\
Look in the mirror\\
Ask what you need\\
And say it clearer

\songpart{Final Chorus}
Complete by what?\\
Complete by what? (amen)\\
Complete by what?\\
Who sets that law? (amen)

A world can be whole\\
Without your frame\\
A world can be whole\\
Without your name

Complete by what?\\
Complete by what? (amen)\\
Complete by what?\\
Amen, amen
\end{lyrics}


\input{gabisson/complete_by_what_1}
\songtitle{Not In The Creed}

\tag*{2026}\tag{gospel}\tag{original-work}

\begin{notes}
Original work.\\
Gospel-soul reflection on religious childhood and eventual departure.
\end{notes}

\begin{style}
soul-folk, gospel soul, 118 BPM, velvet bass baritone, gospel choir refrains, fingerpicked electric guitar, bowed fiddle lines, conga pulse, upright bass, brushed snare, organ swells, plate reverb, tape saturation, sidechain pumping, call and response, spoken intro, half-time bridge, bittersweet confession, devotional ache
\end{style}

\begin{lyrics}
\songpart{Verse 1}
\annotation{male vocals, velvet bass registry, soul delivery, smooth phrases}\\
Born where the bells rang hard\\
in a Catholic crowd, half-heart, half-card.\\
Christian Brothers taught my name—\\
first Communion, same old frame.\\
By fourteen, I’d made my mind\\
before Confirmation, I’d crossed that line.

\songpart{Pre-Chorus}
\annotation{male vocals, bass registry, soul delivery}\\
Still I stood there straight\\
Still I tried to wait\\
For some fire to land\\
In my open hand

\songpart{Chorus}
\annotation{gospel choir}\\
I was never all in\\
I was never all in\\
I wore the coat\\
But I missed the skin\\
I was never all in\\
I was never all in
(never all in)\\
Just telling the truth\\
About where I’ve been

\songpart{Verse 2}
\annotation{male vocals, velvet bass registry, soul delivery, smooth phrases}\\
Then Sweden took me in—\\
cold air, clean streets, less talk of sin.\\
Catholic turned to a sign—\\
more a name than a line.\\
Liberation books in my hands,\\
Cardijn’s kids, and their marching plans.\\
Back in Colombia, Jesuit halls—\\
still I heard a bigger call.

\songpart{Pre-Chorus}
\annotation{male vocals, bass registry, soul delivery}\\
Not the creed, not the show,\\
but the way we go.\\
Toward the hungry face.\\
Toward a harder grace.

\songpart{Chorus}
\annotation{golspel choir}\\
I was never all in\\
I was never all in\\
I wore the coat\\
But I missed the skin\\
I was never all in\\
I was never all in
(never all in)\\
Just telling the truth\\
About where I’ve been

\songpart{Bridge}
\annotation{rap, male baritone vocals}\\
I watched the online saints\\
The hard old rules, the holyaints\\
Faith in a fist\\
And a book like a list\\
They spoke of Jesus\\
Then played the Pharisee game\\
Fearing science\\
While blessing the shame\\
And I learned my fight\\
Wasn’t left or right\\
It was human first\\
It was daylight

\songpart{Verse 3}
\annotation{male vocals, velvet bass registry, soul delivery, smooth phrases}\\
If my branch was close to Him\\
Closer than some who loved the hymns\\
Still I couldn’t kneel\\
To what I couldn’t feel\\
No ghost in the room\\
No heaven in sight\\
Just conscience and work\\
And a harder light

\songpart{Pre-Chorus}
\annotation{male vocals, bass registry, soul delivery}\\
So I laid down the mask\\
Stopped the old half-ask\\
Called my own soul read\\
And said what I meant

\songpart{Chorus}
\annotation{Gosple choir}\\
I was never all in\\
I was never all in\\
I wore the coat\\
But I missed the skin\\
I was never all in\\
I was never all in
(never all in)\\
Just telling the truth\\
About where I’ve been

\songpart{Outro}
\annotation{male vocals, velvet bass registry}\\
Not in anger, not in spite\\
Just a cleaner name for night\\
I don’t call myself a Christian now\\
I call it honest, and I bow to how\\
I was never all in\\
I was never all in
(never all in)\\
Just telling the truth\\
About where I’ve been
\end{lyrics}


\songtitle{Which God Stands Still}

\tag*{2026}\tag{gospel}\tag{original-work}

\begin{notes}
Original work.\\
Gospel composition on faith and the search for immovable truth.
\end{notes}

\begin{style}
Gospel with bass baritone lead and responsive choir; slow-march 6/8 swing, handclaps on the backbeat, organ swells under sparse verse, pre-chorus opens with clapping and rising choir pads, chorus blooms with full SATB harmony and shouted responses, bridge drops to bass voice and room tone before a final lift. %
Close-mic lead in verses, gritty doubles on key lines, short delay throws on the hook, tambourine shakes, ascending organ runs, wide warm mix with a church-room bloom., worship, theme, strong, gospel
\end{style}

\begin{lyrics}
\songpart{Verse 1}
Which god wants fear?\\
Which god wants praise?\\
I read the stone\\
I read the page

No graven face\\
No golden frame\\
Just the hard old law\\
Calling my name

\songpart{Pre-Chorus}
You said, "Do not kneel"\\
To empty wood\\
You said, "Do not trade\\
The whole for the false"

If love is truth\\
Then love should lead\\
Not clutch my soul\\
And tighten its grasp

\songpart{Chorus}
Which God stands still?\\
Which God lets me grow?\\
Which God loves me\\
And lets me go? (lets me go)

Not a fist, not a throne\\
Not a child on a throne\\
Which God stands still?\\
Which God lets me grow? (lets me grow)

\songpart{Verse 2}
My catechism days\\
Had a cold bright seat\\
The rules were clear\\
The room was deep

The first commandment\\
Stood like a gate\\
No other gods\\
No idol bait

So I learned the fear\\
Then I learned the test\\
A love that demands\\
Can still be less

\songpart{Pre-Chorus}
If I lose my way\\
I can lose my name\\
But if I lose my mind\\
That is not grace

A loving parent\\
Would want my height\\
Would want my hands\\
To build in light

\songpart{Chorus}
Which God stands still?\\
Which God lets me grow?\\
Which God loves me\\
And lets me go? (lets me go)

Not a fist, not a throne\\
Not a child on a throne\\
Which God stands still?\\
Which God lets me grow? (lets me grow)

\songpart{Bridge}
\annotation{Drop to bass voice}\\
If there is One\\
Who made this breath\\
Then let it mean\\
More than dread

Let me reach\\
My full created flame\\
Let me love\\
Without a chain

\annotation{Choir swell}\\
Not every voice\\
That calls itself holy\\
Knows the shape\\
Of mercy

If You are Love\\
Then prove it slow\\
Make me whole\\
And let me grow

\songpart{Final Chorus}
Which God stands still?\\
Which God lets me grow?\\
Which God loves me\\
And lets me go? (lets me go)

Not a fist, not a throne\\
Not a child on a throne\\
Which God stands still?\\
Which God lets me grow? (lets me grow)

Which God stands still?\\
Which God lets me grow?\\
Amen, I’m asking\\
That I may know (that I may know)
\end{lyrics}


\songtitle{Ashes in the Pews}

\tag{original-work}\tag{dubstep}\tag{gospel}

\begin{notes}
Original work.\\
Dubstep-gospel exploration of ritual transmission and inherited belief.
\end{notes}

\begin{style}
Dubstep hymn at 140 BPM with half-time drops and swung handclap swing, verse rides sparse organ stabs and sub pulses, pre-chorus strips to soprano lead and low choir hum, chorus hits with stacked gospel choir, wobble bass, and clap-accented lifts. %
Soprano lead gets wide doubles, choir shouts on hook words, delay throws on the last line. %
Ear candy: reversed choir swell into each drop, chime flickers between phrases, seismic riser before final chorus. %
Bright, विशाल, and cathedral-wide, theme, gospel, world, worship, magical, dubstep
\end{style}

\begin{lyrics}
\songpart{Verse 1}
What makes people kneel?\\
Not the sky alone\\
Hands get taught to fold\\
By stories passed like bone

Not born with a temple\\
Not born with a throne\\
The first seed is planted\\
Then it starts to grow

\songpart{Pre-Chorus}
Show me the crack\\
Show me the seam\\
A world full of facts\\
Can wake a dream

\songpart{Chorus}
Who taught the gods?\\
Who taught the fear?
(Who taught the gods?)\\
Who put it here?

Scratch the paint off\\
See the room
(See the room)\\
Truth can bloom

\songpart{Verse 2}
Some break from the script\\
When the edges don't fit\\
When the old words buckle\\
And the page won't sit

When the promised fire\\
Looks like smoke and dust\\
When the holy map\\
Leads nowhere for us

\songpart{Pre-Chorus}
Show me the crack\\
Show me the seam\\
A world full of facts\\
Can wake a dream

\songpart{Chorus}
Who taught the gods?\\
Who taught the fear?
(Who taught the gods?)\\
Who put it here?

Scratch the paint off\\
See the room
(See the room)\\
Truth can bloom

\songpart{Bridge}
No, it's not a curse\\
Not a broken mind\\
Not a wound you can name\\
Just a turn of time

Some fear the void\\
Some fear the flood\\
Some worship the ending\\
Instead of the mud

\songpart{Final Chorus}
Who taught the gods?\\
Who taught the fear?
(Who taught the gods?)\\
Who put it here?

Scratch the paint off\\
See the room
(See the room)\\
Truth can bloom

Who taught the gods?\\
Who taught the fear?
(Who taught the gods?)\\
Why cure the clear?
\end{lyrics}


\songtitle{Morality Ain’t Imposed}

\tag*{2026}\tag{quora-post}

\begin{notes}
Lyrics by Suno based on a Quora post by Carlos Thompson.
\end{notes}
\begin{style}
Hip-hop anthem with gospel choir, mid-tempo swing and muscular head-nod drums; verse rides tight low-end and clipped piano stabs, pre-chorus lifts with stacked choir swells and rising handclaps, chorus opens wide with call-and-response choir punches and a chantable hook. %
Lead rap stays close-mic and focused, with doubled hook lines, ad-lib echoes, and short delay throws on key words. %
Ear candy: reverse choir swells into the chorus, organ lifts between lines, and a final key-change feel with crowd-style shouts. %
Bright, bold, and cinematic with warm church-room resonance., world, gospel, theme, harmonious, hip-hop, praise
\end{style}
\begin{lyrics}
\songpart{Verse 1}
You say the sky wrote law\\
Like it fell in stone\\
But I see people learning\\
Face to face, bone to bone\\
A child hears a voice\\
A mother shows the way\\
Elders plant the seeds\\
Then the grown mind weighs

What hurts, what heals\\
What can’t be ignored\\
What lifts up the table\\
What leaves folks torn\\
I don’t need a throne\\
To know what I defend\\
I watch, I feel, I test it\\
Then I choose again

\songpart{Pre-Chorus}
We all got limits\\
We all draw lines\\
We all ask why\\
In our own time\\
If it feeds the wound\\
I turn away\\
If it builds the house\\
I’m there to stay

\songpart{Chorus}
Morality ain’t imposed (ain’t imposed)\\
We learn it, we hold it, we grow\\
Morality ain’t imposed (ain’t imposed)\\
We weigh it by harm and by hope
(hope and harm)\\
Morality ain’t imposed\\
We decide what we praise\\
Morality ain’t imposed\\
And we live what we say

\songpart{Verse 2}
Not every code comes\\
Wrapped in holy cloth\\
Humanism, utility\\
Truth that holds its ground\\
Past generations speak\\
Through every tested norm\\
Some old ideas decay\\
Some new ones take form

I can respect your faith\\
Still call out the bruise\\
If a rite leaves scars\\
Then it’s gotta lose\\
No divine stamp\\
Can make pain a saint\\
If a custom hurts the weak\\
I won’t call it quaint

\songpart{Pre-Chorus}
We all got elders\\
We all got peers\\
We all got reasons\\
Through the years\\
If it raises hands\\
Or breaks the chain\\
I’ll say it loud\\
And I’ll say it plain

\songpart{Chorus}
Morality ain’t imposed (ain’t imposed)\\
We learn it, we hold it, we grow\\
Morality ain’t imposed (ain’t imposed)\\
We weigh it by harm and by hope
(hope and harm)\\
Morality ain’t imposed\\
We decide what we praise\\
Morality ain’t imposed\\
And we live what we say

\songpart{Bridge}
If your prayer makes room\\
For the hungry child\\
Then I’ll stand and clap\\
Let that mercy run wild\\
If your law makes fear\\
Out of neighbor and kin\\
Then I’m raising my voice\\
Let the healing begin

\songpart{Final Chorus}
Morality ain’t imposed (ain’t imposed)\\
We learn it, we hold it, we grow\\
Morality ain’t imposed (ain’t imposed)\\
We weigh it by harm and by hope
(hope and harm)\\
Morality ain’t imposed\\
We decide what we praise\\
Morality ain’t imposed\\
And we live what we say

\songpart{Outro}
Empathy in the blood\\
Wisdom in the room\\
We judge with our hands open\\
And we choose what blooms
(what blooms)
\end{lyrics}


\songtitle{Categorical Error}

\tag*{2026}\tag{quora-post}

\begin{notes}
Lyrics by Suno based on a Quora post by Carlos Thompson
\end{notes}
\begin{style}
Fast-paced folk ballad with gospel choir, driving acoustic strum and handclap pulse; verses ride fingerpicked guitar and upright bass, pre-chorus opens with stacked harmony lift, chorus hits with clapping, organ swells, and choir answers. %
Lead vocal stays close and earnest with doubled lines on the hook, small delay throws on key words, and call-and-response shouts from the choir. %
Bright, wooden, uplifting mix with a live-room sheen., ballad, natural, gospel, folk
\end{style}
\begin{lyrics}
\songpart{Verse 1}
Is it useful, baby?\\
To call a shape a tool?\\
A neutral thing on trial\\
By a hungry human rule

Some days I fix the hinge\\
Some days I just fall through\\
I don't need a ghost in silk\\
To do what I can do

\songpart{Pre-Chorus}
When the storm comes in\\
And the road goes thin\\
I grip my own two hands\\
And I begin

\songpart{Chorus}
Serenity, come on\\
Courage, come on\\
Wisdom, come on\\
I find it myself (find it myself)

Change what I can\\
Hold what I can't\\
Know the difference\\
And stand (stand, stand)

\songpart{Verse 2}
If a flood stays out of reach\\
I still learn how to breathe\\
I can face the breaking news\\
Without a plea on my sleeve

You can pray with open palms\\
Or you can build with sweat\\
Either way, when panic calls\\
You answer, not forget

\songpart{Pre-Chorus}
When the dark rolls in\\
And the floor boards bend\\
I don't wait on the wind\\
I call on my friends

\songpart{Chorus}
Serenity, come on\\
Courage, come on\\
Wisdom, come on\\
I find it myself (find it myself)

Change what I can\\
Hold what I can't\\
Know the difference\\
And stand (stand, stand)

\songpart{Bridge}
This ain't a hammer\\
This ain't a knife\\
It's just how I meet\\
The shape of life

Not a blessed device\\
Not a cursed disguise\\
Just a body walking on\\
With clear-eyed eyes

\songpart{Chorus}
Serenity, come on\\
Courage, come on\\
Wisdom, come on\\
I find it myself (find it myself)

Change what I can\\
Hold what I can't\\
Know the difference\\
And stand (stand, stand)
\end{lyrics}


\songtitle{Antiphona: De Atheismi Utilitate}

\tag*{2026}\tag{quora-post}

\begin{notes}
Lyrics by Grok based on a Quora post by Carlos Thompson
\end{notes}
\begin{style}
gregorian chant, liturgical antiphony, Latin text setting, solo cantor, responsorial choir, unison chant, solemn modal melody, slow 72 BPM, resonant stone chapel, natural reverb, sparse organ drone, plainchant cadence, gradual buildup, sectional contrast, austere medievalism, sacred dialogue, ritual gravity\\
new age track as background
\end{style}
\begin{lyrics}
\songpart{Introitus}
\annotation{choir} Utrum atheismus utilis est practicus?\\
\annotation{solo} Tam utilis quam capilli rufi naturales,\\
Vel statura centum septuaginta.\\
Quare utilis esse debet?\\
Condicio neutra, iudicium utilitatis vanum.

\songpart{Versus I}
Non confidere amico imaginario,\\
Utilis in circumstantiis quibusdam:\\
Problemata fronte directa solvis melius,\\
Non petens phantasma ut agat pro te.\\
Sed fides etiam iuvat:\\
In rebus extra potestatem,\\
Minus anxius si credis providentiam.

\songpart{Responsorium}
Credens tamen responsabilitatem sumere potest\\
Ut atheus;\\
Atheus vero stoica fortitudine\\
Adversitates incertas sustinet.

\songpart{Versus II (ex oratione serenitatis)}
Deus, dona mihi serenitatem\\
Accipere quae mutare non possum,\\
Fortitudinem mutare quae possum,\\
Sapientiam discernere differentiam.
— Reinhold Niebuhr.

\songpart{Responsorium Athei}
Non rogo Deum ut det mihi haec,\\
Sed ipse quaero serenitatem, fortitudinem, sapientiam.\\
Credens autem, spiritu orationis,\\
“Mutare quae possum” includit\\
Persequi haec dona, non passive exspectare.

\songpart{Conclusio (Antiphona finalis)}
Atheismus non est instrumentum,\\
Sed eventus fortuitus, accidens.\\
Quaerere utilitatem eius practicam\\
Error categoricus est.\\
Neutralis ut natura ipsa.\\
Amen.
\end{lyrics}


\songtitle{Talk To Me}

\tag*{2026}\tag{gospel}
\begin{notes}
\end{notes}
\begin{lyrics}
\songpart{Intro}
\annotation{Acoustic guitar, slow train rhythm}
\annotation{Gospel choir humming}

You say there’s a war I cannot see\\
A shadow pulling strings behind me\\
You say my hands are playing a part\\
In a battle written in the dark

But before you tell me who I serve\\
Before you tell me what I’ve become\\
I got one question left to ask:

\songpart{Verse 1}

Have you ever heard a fallen angel speak?\\
Have you read the words they wrote in secrecy?\\
Have you seen the plans, the hidden design?\\
The meetings where they draw their lines?

Or is it a story passed down through the years?\\
A voice from a pulpit, a book, and some fears?\\
A translation of a memory, a people in exile\\
Trying to find a homeland after a thousand miles

\songpart{Pre-Chorus}

And if you believe every word is written in stone\\
Then maybe you think I’m a voice not my own\\
Maybe you hear deception when you hear me talk\\
Maybe you see a trap where I see a walk

\songpart{Chorus}

But tell me, are we talking? (Talk to me)\\
Or are you only proving what you know? (You know)\\
Are you asking me a question?\\
Or just waiting for your answer to show?

I’m not afraid of heaven (Afraid of heaven)\\
I’m not afraid of doubt (Afraid of doubt)\\
But a door cannot open\\
If you’ve already locked me out

\annotation{Gospel choir}
(Talk to me)
(Talk to me)
(Don’t just tell me what I am)

\songpart{Verse 2}

I don’t need you to agree with me\\
I don’t need you to change your belief\\
I only ask you look me in the eye\\
And let the question be a question tonight

Because if every road I walk becomes evidence\\
If every word I say becomes your defense\\
Then there’s no conversation left to find\\
Only a verdict waiting at the end of the line

\songpart{Bridge}

Maybe you see angels (see angels)\\
Maybe you see signs (see signs)\\
Maybe you see a battle (see a battle)\\
Written between the lines (in the lines)

I see a person\\
Standing here with you\\
Not an enemy hiding\\
Just another point of view

\annotation{Gospel choir}
(Another point of view)
(Another point of view)

\songpart{Final Chorus}

So tell me, are we talking? (Talk to me)\\
Or are we only fighting ghosts? (Fighting ghosts)\\
Are we searching for the truth?\\
Or protecting what we know the most?

I’m not asking you to follow (you follow)\\
I’m not asking you to fall (you fall)\\
I’m asking for a conversation\\
Where we both can hear it all

\annotation{Gospel choir}
(Talk to me)
(Talk to me)
(Let the question breathe)
(Talk to me)
(Talk to me)
(Don’t decide for me)
\end{lyrics}

\songtitle{Defined Or Dust}

\tag*{2026}\tag{quora}
\begin{notes}
Lyrics by SUNO based on a Quora post by Carlos Thompson.
\end{notes}
\begin{style}
House hymn with dubstep drops, gospel claps, four-on-the-floor pulse, sub-bass wobbles in the chorus, call-and-response choir doubles, verse rides a sparse organ and snapped claps, pre-chorus strips to handclaps and a filtered bass drone, chorus hits with stacked gospel harmonies and a heavy wobble drop, bridge breaks to a lone voice then swells with crowd shouts, bright and punchy mix with a wide sacred-club sheen, gospel, dubstep
\end{style}
\begin{lyrics}
\songpart{Verse 1}
I look at the stars\\
and I don’t see a throne\\
I see cold math\\
and a world on its own\\
You say there’s a hand\\
behind every stone\\
But I need a shape\\
I can hold up and know

\songpart{Pre-Chorus}
If the words don’t land\\
then the proof won’t stand\\
If the name won’t mean\\
what it means in my hands\\
Build the frame first\\
then show me the flame\\
If it lives in the mist\\
it can’t carry the claim

\songpart{Chorus}
Defined or dust\\
defined or dust\\
If I can’t name it\\
I can’t trust\\
Defined or dust\\
defined or dust\\
If it lives outside matter\\
it’s not enough
(defined or dust)

\songpart{Verse 2}
I’m not asking for less\\
I’m asking for clear\\
What we test in the light\\
is what we keep near\\
Default to the ground\\
to the seen and the known\\
That’s where science begins\\
and where doubt has a home

\songpart{Pre-Chorus}
If the map is vague\\
then the path won’t stay\\
If it won’t meet the world\\
then it drifts away\\
Don’t hand me a wonder\\
wrapped in a guess\\
When the claim can’t be checked\\
it’s only a rest

\songpart{Chorus}
Defined or dust\\
defined or dust\\
If I can’t name it\\
I can’t trust\\
Defined or dust\\
defined or dust\\
If it lives outside matter\\
it’s not enough
(defined or dust)

\songpart{Bridge}
Maybe you call it holy\\
I call it out of reach\\
A thought in the air\\
that won’t touch the beach\\
If it can’t enter the world\\
where the questions stand\\
Then it’s just a whisper\\
slipping through my hand

\songpart{Final Chorus}
Defined or dust\\
defined or dust\\
If I can’t name it\\
I can’t trust\\
Defined or dust\\
defined or dust\\
If it lives outside matter\\
it’s not enough\\
Defined or dust\\
defined or dust
(defined or dust)
\end{lyrics}


\songtitle{Pews and Bridges}

\tag*{2026}\tag{quora-posting}\tag{gospel}\tag{english-lyrics}
\begin{notes}
Lyrics by SUNO based on a 2026 Quora answer by Carlos Thompson.
\end{notes}
\begin{style}
Country-gospel ballad with steady trainbeat claps, brushed drums, front-porch acoustic guitar, fiddle swells, and a warm church choir on the lifts; verse stays close and plain with male baritone lead, pre-chorus opens into handclaps and harmony replies, chorus lands with full choir, tambourine, and a lifted fiddle figure; use short delay throws on key words, communal claps on two and four, and a bright but earthy mix with room tone and Sunday-service grit, country, ballad, gospel
\end{style}
\begin{lyrics}
\songpart{Verse 1}
I heard a man at the back row say,\\
First cause proves the floor we stand on.\\
But a ladder ain’t the house, my friend,\\
And a map ain’t the place we’re from.\\
You can point to the smoke in the air,\\
Still miss the hand that lit the flame.

\songpart{Pre-Chorus}
I’m not fighting over shadows.\\
I’m asking what the heart gets saved by.\\
If the Lord is more than a reason,\\
Then mercy is the reason why.

\songpart{Chorus}
I need the saving kind\\
I need the saving kind\\
Not just a name for the sky\\
But grace that finds me alive
(Grace that finds me alive)

\songpart{Verse 2}
Providence is bread on the table,\\
Rain at the time the field goes dry.\\
It’s a hand on the wheel in the winter,\\
It’s a door that opens in July.\\
And faith ain’t a trick of the mind,\\
It’s trust when the trouble moves in.

\songpart{Pre-Chorus}
What good is a throne in the distance\\
If it never bends down to me?\\
What good is a God in the question\\
If He won’t bring the lost home free?

\songpart{Chorus}
I need the saving kind\\
I need the saving kind\\
Not just a name for the sky\\
But grace that finds me alive
(Grace that finds me alive)

\songpart{Bridge}
So don’t hand me a theory\\
Wrapped in holy light.\\
I’m looking for the Shepherd\\
Who’ll carry me through night.\\
If He made the stars and seedfields,\\
Then He can make me new.\\
That’s the faith I’m talking about,\\
The faith that sees me through.

\songpart{Chorus}
I need the saving kind\\
I need the saving kind\\
Not just a name for the sky\\
But grace that finds me alive
(Grace that finds me alive)

\songpart{Outro}
I need the saving kind\\
I need the saving kind\\
Amen, amen, amen\\
Grace that finds me alive
\end{lyrics}


\songtitle{Yes It Can}

\tag*{2026}\tag{gospel}\tag{quora-posting}\tag*{english-lyrics}
\begin{notes}
Lyrics by SUNO based on a 2026 Quora answer by Carlos Thompson.
\end{notes}
\begin{style}
Country-gospel ballad with grassroot acoustic guitar, pump organ, brushed snare, handclaps, fiddle fills, gospel choir responses, and male baritone lead vocals; verse stays sparse and story-first, pre-chorus lifts with claps and choir swells, chorus opens wide with fiddle and stacked harmonies, bridge drops to voice and guitar before a full final lift; warm, earthy, and communal mix with intimate close-mic verses and bright church-room reverb on the hooks, gospel, ballad, country
\end{style}
\begin{lyrics}
\songpart{Verse 1}
Down by the beans and corn\\
Under the workday sun\\
A man can doubt the heavens\\
And still help all the fields get done\\
There is nothing in the soil\\
That asks for a sacred vow\\
It takes a lot of hands and time\\
To make a living plow

\songpart{Pre-Chorus}
We mark the rain and season\\
By sweat and scar and seed\\
Some folks call that holy\\
Some call that what we need

\songpart{Chorus}
Yes, it can (yes, it can)\\
Yes, it can (yes, it can)\\
Can atheism exist\\
In the furrow and the land?\\
Yes, it can (yes, it can)\\
Yes, it can (yes, it can)\\
The field still turns\\
When you take no god by hand

\songpart{Verse 2}
Religion may help the circle\\
Help the neighbors pull as one\\
Help the barn stay standing\\
Help the supper lines get done\\
But help is not a promise\\
And custom is not proof\\
You can trade old tales for numbers\\
For a ledger, law, or roof

\songpart{Pre-Chorus}
State plans and sacred orders\\
Can both go right and wrong\\
The human need for meaning\\
Will keep right on moving strong

\songpart{Chorus}
Yes, it can (yes, it can)\\
Yes, it can (yes, it can)\\
Can atheism exist\\
In the furrow and the land?\\
Yes, it can (yes, it can)\\
Yes, it can (yes, it can)\\
The field still turns\\
When you take no god by hand

\songpart{Bridge}
No, it is not the losing\\
Of rivers, dust, and breath\\
No, it is not a punishment\\
For being close to death\\
If a claim stands unsupported\\
Some folks will let it go\\
That turning of the heart and mind\\
Is how these plain truths grow

\songpart{Final Chorus}
Yes, it can (yes, it can)\\
Yes, it can (yes, it can)\\
Can atheism exist\\
In the furrow and the land?\\
Yes, it can (yes, it can)\\
Yes, it can (yes, it can)\\
The choir sings out (signs out)\\
And the whole room understands\\
Yes, it can (yes, it can)\\
Yes, it can (yes, it can)\\
The field still turns\\
When you take no god by hand
\end{lyrics}


\songtitle{Not Misreading Intention}

\tag*{2026}\tag{gospel}\tag{quora-posting}\tag*{english-lyrics}
\begin{notes}
Lyrics by SUNO based on a 2026 Quora answer by Carlos Thompson. Prompt guided.
\end{notes}
\begin{style}
Pop ballad with fast electronic pulse and folk-leaning imagery, male tenor lead, gospel choir lift on choruses, driving kick and handclap groove, verse starts sparse with plucked acoustic guitar and warm synth arps, pre-chorus opens into stacked harmonies and rising filters, chorus blooms with choir shouts, octave doubles, and big synth swells; bridge strips to voice and keys before a final full-choir lift, Bright, polished, and emotionally direct with anthemic mix sheen, electronic, ballad, folk, uplifting, gospel, pop, reading
\end{style}
\begin{lyrics}
\songpart{Verse 1}
Why was human evolution left out?\\
You ask me straight, I hear the doubt\\
For the same old reasons, plain and wide\\
Some things were never set inside

The Constitution isn’t there\\
Nor Principia in the air\\
Plato’s Dialogues, Newton too\\
Not every truth was meant for you

\songpart{Pre-Chorus}
They wrote for then, not every age\\
They drew a line, they turned a page\\
To teach obedience, remorse\\
Not a lab report, not a course

\songpart{Chorus}
Don’t misread the intention\\
Don’t misread the intention\\
It was not the book’s invention\\
It was value, not a mention

(Facts are facts, facts are facts)
(Truth still stands, truth still stands)

\songpart{Verse 2}
How to take care of your cat\\
How to repair your truck, all that\\
How to build a watch by hand\\
Was never part of that land

Shannon’s Information Theory\\
Thermodynamic laws, same story\\
The text was never meant to say\\
Every fact of every day

\songpart{Pre-Chorus}
Creation stories, plural, there\\
Were never sent as science there\\
They held up meaning in the dark\\
To guide a pre-scientific heart

\songpart{Chorus}
Don’t misread the intention\\
Don’t misread the intention\\
It was not the book’s invention\\
It was value, not a mention

(Facts are facts, facts are facts)
(Truth still stands, truth still stands)

\songpart{Bridge}
If you read them like a chart\\
You miss the soul inside the art\\
And if you ignore what we now know\\
You choose a blindness that won’t grow

So let the old words do their part\\
And let the facts still guard your heart\\
The world got larger, that is true\\
And wisdom grows when we do too

\songpart{Final Chorus}
Don’t misread the intention\\
Don’t misread the intention\\
It was not the book’s invention\\
It was value, not a mention

(Facts are facts, facts are facts)
(Truth still stands, truth still stands)\\
Don’t misread the intention\\
Don’t misread the intention
\end{lyrics}


\songtitle{What The Book Leaves Out}

\tag*{2026}\tag{gospel}\tag{quora-posting}
\begin{notes}
Lyrics by SUNO based on a 2026 Quora answer by Carlos Thompson.
\end{notes}
\begin{style}
Pop ballad with electronic instrumentation and folk storytelling; mid-tempo four-on-the-floor pulse, soft synth arpeggios and warm programmed drums under acoustic guitar, Verse rides intimate and reflective, pre-chorus lifts with rising pads and handclap syncopation, chorus opens into gospel choir stacks and call-and-response, Male tenor lead with close-mic verses, airy doubles, and choir swells on the hook; ear candy includes reversed swells, celeste hits, and a shimmering lift into each chorus, Bright, polished, and deeply emotional mix, gospel, ballad, reading, uplifting, folk, electronic, pop
\end{style}
\begin{lyrics}
\songpart{Verse 1}
You asked me why\\
the old pages stayed quiet\\
about bones and change\\
and every long goodbye

It’s not a flaw\\
it’s a doorway left open\\
for the hands of faith\\
not the lab of the sky

\songpart{Pre-Chorus}
There’s a river of names\\
and a crowd of questions\\
Some things were never meant\\
to carry every lesson

\songpart{Chorus}
It leaves some things out\\
it leaves some things out
(leave some things out)\\
It tells what they sought\\
not all that we found

It leaves some things out\\
it leaves some things out
(leave some things out)\\
Truth can be wide\\
without being whole now

\songpart{Verse 2}
No map for the cat\\
no rule for the watchwork\\
No law of the stars\\
written in those lines

They spoke to a world\\
with dirt on their sandals\\
to teach what to love\\
and how to walk that line

\songpart{Pre-Chorus}
If you read it like a mirror\\
for every single answer\\
you miss the heart\\
behind the ancient thunder

\songpart{Chorus}
It leaves some things out\\
it leaves some things out
(leave some things out)\\
It tells what they sought\\
not all that we found

It leaves some things out\\
it leaves some things out
(leave some things out)\\
Truth can be wide\\
without being whole now

\songpart{Bridge}
So don’t trade wonder\\
for a cage of certainty\\
Don’t call it broken\\
when it wasn’t meant to be

The earth kept turning\\
through the questions we were given\\
And science is a lamp\\
that helps us see it clearer

\songpart{Final Chorus}
It leaves some things out\\
it leaves some things out
(leave some things out)\\
It tells what they sought\\
not all that we found

It leaves some things out\\
it leaves some things out
(leave some things out)\\
Faith can be deep\\
and still make room for now

It leaves some things out\\
it leaves some things out
(leave some things out)\\
And if we love truth\\
we’ll let truth breathe out
\end{lyrics}


\songtitle{The Ground I Find}

\tag*{2026}\tag{blues}\tag{quora-posting}
\begin{notes}
Lyrics by Suno on a Quora post by Carlos Thompson.
\end{notes}
\begin{style}
Chamber blues with a slow, deliberate 6/8 groove; intimate close-miked baritone lead vocal, conversational verses with restrained vibrato, gospel choir answering in warm sustained phrases, handclaps entering on the backbeat, upright bass, brushed snare, muted piano, sparse cello swells, and occasional tremolo electric guitar, Build from nearly bare voice and piano into a small room congregation, then pull back for the final verse, Organic live-room recording, dark wood warmth, close vocal detail, dynamic crescendos, no glossy pop polish
\end{style}
\begin{lyrics}
\songpart{Verse 1}
I didn't pick the question, it arrived beneath my skin\\
I tried to want an answer, tried to make the old words fit\\
But wanting makes no evidence, and comfort isn't proof\\
So I kept the honest trouble when the easier road said choose

\songpart{Chorus}
I didn't choose the truth to make it kind\\
I only name the ground I find\\
If heaven's there, it doesn't need my lie\\
If heaven's not, I won't pretend it is nearby

\songpart{Verse 2}
I didn't choose my birthplace, or the hands that raised me here\\
The shade that holds my face, the loves that make my purpose clear\\
I didn't choose the shape of what my searching came to know\\
A road can still be yours to walk though you did not choose its stones

\songpart{Chorus}
I didn't choose the truth to make it kind\\
I only name the ground I find\\
If heaven's there, it doesn't need my lie\\
If heaven's not, I won't pretend it is nearby

\songpart{Verse 3}
I love this crowded country, love the corners of my town\\
I love the two who raised me, love the friends who gather round\\
My neighbors make a community, my working hands belong\\
There are clubs and causes, kitchens, jokes, and people singing songs

\annotation{Choir Response}\\
You are not alone\\
You are not alone

\songpart{Verse 4}
And when the heart needs tending, there are names for what we feel\\
A counselor, a paperback, a friend who helps us heal\\
A screen can eat an evening, though it rarely feeds the soul\\
Still no priest owns every doorway, no creed controls the whole

\songpart{Chorus}
I didn't choose the truth to make it kind\\
I only name the ground I find\\
If heaven's there, it doesn't need my lie\\
If heaven's not, I won't pretend it is nearby

\songpart{Bridge}
You can sit inside a chapel, sing the hymns and pass the plate\\
You can stand before the pulpit, leave the miracles to wait\\
A fellowship can hold you without asking what you see\\
The hands can join together though the heavens disagree

\annotation{Choir Response}\\
Come as you are\\
Come as you are

\annotation{Final Verse}\\
So don't ask why it's tempting, like a bargain I embraced\\
I wanted it to be true once; truth refused to change its face\\
And though the road is colder where the promised answers cease\\
I found a human warmth here, and the hard-earned gift of peace

\songpart{Final Chorus}
I didn't choose the truth to make it kind\\
I only name the ground I find\\
If heaven's there, it doesn't need my lie\\
If heaven's not, I won't pretend it is nearby

\songpart{Outro}
No borrowed certainty\\
No lonely victory\\
Just what I find, and those I love\\
Here in the world we see
\end{lyrics}


\songtitle{Think About It}

\tag*{2026}\tag{house}\tag{quora-posting}
\begin{notes}
Lyrics by Suno on a Quora post by Carlos Thompson.
\end{notes}
\begin{style}
Peak-time vocal house anthem, four-on-the-floor kick at 124 BPM, rolling sub bass, bright piano stabs, filtered disco strings, handclaps, and crisp modern percussion, Confident male lead with conversational verse delivery that opens into a soulful melodic chorus; large gospel choir chants answer key phrases in call-and-response, Build steadily from sparse groove to festival-sized chorus, add a breakdown with claps and choir, then return with thicker bass and stacked harmonies, Clean club mix, punchy low end, wide vocal layers, uplifting but intellectually assertive
\end{style}
\begin{lyrics}
\songpart{Intro}
Think about it\\
Think about it

\songpart{Verse 1}
Is it evidence, or a way to see?\\
A question puts the weight on me\\
I don't trade one mystery for another name\\
I leave the unknown exactly where it came

\songpart{Pre-Chorus}
No prayer for favors from the sky\\
No hidden hand I have to find\\
I face the world with open eyes\\
And let the honest question rise

\songpart{Chorus}
Atheism is a state of mind\\
No gods, no substitutes behind\\
I don't need a throne above\\
To choose a life of care and love\\
Think about it, think about it\\
Look for what the facts can find\\
Think about it, think about it\\
Atheism is a state of mind

\annotation{Choir Chant}\\
Think about it!\\
State of mind!\\
Think about it!\\
State of mind!

\songpart{Verse 2}
You may ask the gods for favors and grace\\
I don't send my hopes to an unseen place\\
You may raise your rites to praise a name\\
I gather with people; the reason's plain

You may fear a judgment after you die\\
I choose good neighbors while we're alive\\
I want a world where people keep their word\\
Because the life we share is the life I've heard

\songpart{Pre-Chorus}
No promise waiting past the blue\\
No threat to make me act for you\\
I build my choices here and now\\
With other people, somehow

\songpart{Chorus}
Atheism is a state of mind\\
No gods, no substitutes behind\\
I don't need a throne above\\
To choose a life of care and love\\
Think about it, think about it\\
Look for what the facts can find\\
Think about it, think about it\\
Atheism is a state of mind

\songpart{Bridge}
When we don't understand, I don't invent a claim\\
I can say, "We don't know yet," and leave it without a name\\
Old fables may be meaningful, but meaning isn't proof\\
A story can shape a culture without making it true

\songpart{Breakdown}
Strong agnostic!\\
Still asking why!\\
No evidence made me choose\\
I simply won't call fables truth

\annotation{Choir Chant}\\
No evidence!\\
Still we think!\\
No evidence!\\
Still we think!

\songpart{Final Chorus}
Atheism is a state of mind\\
No gods, no substitutes behind\\
Lack of evidence, the fables clear\\
Led me to stand without belief here\\
Think about it, think about it\\
Let the honest question guide\\
Think about it, think about it\\
Strong agnostic, de facto atheist inside

\songpart{Outro}
Think about it\\
State of mind\\
Think about it\\
Open eyes, honest mind
\end{lyrics}


\end{multicols*}

\chapter{Pray, Brother II: Sacred Stories}
\albumcover{covers/internal/pray-brother-ii-sacred-stories}
Where Apologetics argues, \emph{Sacred Stories} narrates: Noah's ark told twice, a succubus portrait, a fallen messenger, and -- unexpectedly -- a prayer-shaped eulogy for Gustavo Cerati, sung like he belonged in the same lineage as the prophets.
\clearpage\begin{multicols*}{4}
\songtitle{Dead Sea's barren dreams}

\tag*{2025}\tag{original-idea}\tag{goth-metal}

\begin{notes}
A distinctive arrangement that supports the song's storytelling and emotional arc. %
The lyrics begin with “In yonder lands of arid climes,”, giving a quick sense of the song’s tone.
\end{notes}

\begin{style}
goth metal, guttural voice
\end{style}

\begin{lyrics}
In yonder lands of arid climes,\\
Where virtues shine in hospitality's rhymes,\\
In Genesis' scrolls, 'neath Abraham's gaze,\\
Lies tales of honor, and treachery's haze.\\
'Twas not the love twixt men that stirred the fray,\\
But breach of kindness led astray,\\
For Sodom and Gomorrah, in their plight,\\
Fell to ruin in hospitality's blight.\\
Lot, a figure, noble and bold,\\
Did offer kin to guests, as the story is told,\\
His daughters he proffered, in sacred vow,\\
To uphold the laws of hospitality's brow.\\
Reflect on this tale, and its moral decree,\\
Homophobia? Nay, but kindness be free.\\
A fable spun, though unreal it seems,\\
To explain the Dead Sea's barren dreams.\\
Thus, let us heed, in our earthly sod,\\
The lessons of kindness, taught by God.

\end{lyrics}


\songtitle{In the Time of Noah's Ark}

\tag*{2013}\tag{original-lyrics}\tag{children-rock}

\subsection{Notes}

\subsection{Style}
Children's rock with percutions, keyboards, bass and guitars, to the tune of "Yellow Submarine"

\textbf{Observation} completely disregarded the reference, presumely for copyright issues.

\begin{lyrics}
In the time of Noah's ark,\\
all the heavens broke apart,\\
and the fountains of the world\\
spilt their waters and it was flood.

In a small boat all animals were saved\\
except for fish, ants and T-rex.\\
All living creature was destroyed\\
except for plant seeds, bees and seaweed.

This story makes much sense:\\
it explains well what my book says.\\
’cause in nature there’s no evidence\\
so the devil had it erased.

In a small boat Noah’s family preserved\\
every kind of nostrill breathing beast.\\
Because nobody would survive\\
A loving deity's all-mighty wrath.

Gilgamesh plagiarized\\
Noah's tale from end to start.\\
So we know it was all true,\\
not the invention of something new.

In a small boat just across the world\\
all my ancestors were preserved from the flood.\\
I now can thank my celestial creator\\
for his fine cleansing of the genetic pool.
\end{lyrics}


\songtitle{In the Time of Noah's Ark}

\tag*{2013}\tag{original-lyrics}\tag{musical-theater}

\subsection{Notes}

\subsection{Style}
Musical theater comedy song. %
A bright, upright piano plays a bouncy, syncopated ragtime-influenced accompaniment with a walking bassline. %
A baritone male vocalist delivers a theatrical, narrative performance with clear diction. %
The arrangement is minimalist, focusing entirely on the piano and vocal interaction. %
The tempo is 128 BPM in 4/4 time. %
The piano uses staccato chords on the off-beats and melodic fills between vocal phrases.

\begin{lyrics}
\songpart{Intro}
\annotation{upright piano, bouncy ragtime rhythm}

\songpart{Verse 1}
\annotation{baritone male vocals}\\
In the time of Noah's ark,\\
all the heavens broke apart,\\
and the fountains of the world\\
spilt their waters and it was flood.

\songpart{Chorus bridge}
\annotation{Assamble}\\
In a small boat all animals were saved\\
except for fish, ants and T-rex.\\
All living creature was destroyed\\
except for plant seeds, bees and seaweed.

\songpart{Verse 2}
\annotation{baritone male vocals}\\
This story makes much sense:\\
it explains well what my book says.\\
’cause in nature there’s no evidence\\
so the devil had it erased.

\songpart{Chorus bridge}
\annotation{Assamble}\\
In a small boat Noah’s family preserved\\
every kind of nostrill breathing beast.\\
Because nobody would survive\\
A loving deity's all-mighty wrath.

\songpart{Verse 3}
\annotation{baritone male vocals}\\
Gilgamesh plagiarized\\
Noah's tale from end to start.\\
So we know it was all true,\\
not the invention of something new.

\songpart{Chorus bridge}
\annotation{Assamble}\\
In a small boat just across the world\\
all my ancestors were preserved from the flood.\\
I now can thank my celestial creator\\
for his fine cleansing of the genetic pool.
\end{lyrics}


\input{chlewey/succubus_portrait}
\songtitle{Fallen Messenger}

\tag*{2026}\tag{blues-rock}
\begin{notes}
Blues version of \emph{Fallen Messenger}
\end{notes}
\begin{style}
blues gospel rock, cinematic folk blues, baritone lead vocal, gospel choir stabs, gospel claps, harmonica riffs, kick drum pulse, cymbal crashes, banjo fingerpicking, call and response, spoken intro, verse chorus form, dynamic breakdowns, plate reverb, tape saturation, 92 BPM, half-time groove, bittersweet defiance, devotional lust
\end{style}
\begin{lyrics}
\songpart{Intro, instrumental, 8 beats}

\songpart{Verse 1}
You were supposed to be\\
my link with the divine,\\
a divine I cannot trust with.\\
But you yourself come close to it.

\songpart{Pre-Chorus}
You said a god sent you,\\
you said you’ve fallen from its grace.\\
I can’t lie, gods are not my thing,\\
but here you are with me.

\songpart{Chorus}
Your lust, my lust — it’s real (our lust)\\
Your message, I can accept (amen)\\
Not for faith in your strange master,\\
but for what you proved to me (to me)\\
You’re too real (too real)\\
Real to me.

\songpart{Verse 2}
You’ve saved my life,\\
the drift I was falling in.\\
If your master can’t see it,\\
your master doesn’t deserve you.

\songpart{Pre-Chorus}
A fallen messenger, I’ve fallen in love,\\
a fallen messenger, I’ve fallen in lust.\\
As if lust would be a sin,\\
as if sin were not a thing to trust.

\songpart{Chorus}
Your lust, my lust — it’s real (our lust)\\
Your message, I can accept (amen)\\
Not for faith in your strange master,\\
but for what you proved to me (to me)\\
You’re too real (too real)\\
Real to me.

\songpart{Bridge}
Your breath, your voice are not unreal,\\
your breasts, your hips — worldly enough.\\
Your vulnerable body in my arms,\\
your tender being, real and tough.

You rescued me, I rescued you.\\
You say it’s a sin that your master can’t approve.\\
A sin should feel wrong, but nothing is wrong —\\
or your master is jealous of how with me you move.

\songpart{Chorus}
Your lust, my lust — it’s real (our lust)\\
Your message, I can accept (amen)\\
Not for faith in your strange master,\\
but for what you proved to me (to me)\\
You’re too real (too real)\\
Real to me.

\songpart{Outro}
Your lust,\\
my lust,\\
your sin,\\
my redemption.\\
Your master —\\
not real.\\
You with me —\\
be free.
\end{lyrics}


\songtitle{Gustavo Cerati que no estás con nosotros}

\tag*{2014}\tag{rock-en-espanol}\tag{original-lyrics}\tag{spanish-lyrics}
\begin{notes}
Originally composed by Carlos Thompson in September 2014.
\end{notes}
\begin{style}
rock en español, argentine rock, spoken intro, tenor vocals, call and response, bass led melody, driving trance bassline, club synth arpeggios, distorted rhythm guitars, heavy synth layers, chorus pads, plate reverb, sidechain pumping, cathedral reverb, 122 BPM, anthemic intensity, bittersweet nostalgia, pre chorus lift, crowd chant hook
\end{style}
\begin{lyrics}
(gracias por venir)

Gustavo Cerati que no estás con nosotros\\
bien cantadas sean tus letras.\\
Venga a nosotros tu ritmo.\\
Escúchese tu música tanto en el iPad como en el livin'.\\
Danos hoy nuestra dosis de sonido\\
y perdona a Jorge Celedón, así como nostros perdonamos al vallenato.\\
No permitas que nos tome el reggeaton\\
ni otra música del mal;\\
porque tuyo es el sonido, el ritmo y las letras,\\
ramén.
\end{lyrics}


\songtitle{Keep the Vow}

\tag*{2024}
\begin{notes}
Lyrics by SUNO based on the ChatGPT addaptation to Lovercraftian stype of the Sodom and Gomorrah narrative idea by Carlos Thompson.
\end{notes}
\begin{style}
gothic ballad, epic medieval ballad, spoken word intro, deep male baritone, hurdy-gurdy drone, nyckelharpa counterline, hammered dulcimer ostinato, frame drum pulse, cathedral organ swells, chamber strings legato, brass stabs, harp pulses, sidechain pumping, plate reverb, tape saturation, 138 BPM, march cadence, ritual dread
\end{style}
\begin{lyrics}
\songpart{Verse 1}
In the dry bone land\\
where the wind keeps score\\
Old laws sleep hard\\
behind a shuttered door

Guest at the gate\\
host with the bread\\
Break that bond\\
and you wake the dead

\songpart{Pre-Chorus}
Hear the old rite\\
under ash and dust\\
Open the hand\\
or be crushed by what comes

The night has names\\
older than fear\\
Say them wrong\\
and the end draws near

\songpart{Chorus}
Guest and host\\
guest and host\\
Keep the vow
(keep the vow)

Guest and host\\
guest and host\\
Break the seal\\
and the grave will know

\songpart{Verse 2}
Sodom in shadow\\
Gomorrah in flame\\
Stone by stone\\
they called no name

They laughed at the door\\
they spurned the plea\\
So the sky came down\\
like a cursed sea

Lot stood shaking\\
by the gate\\
Caught between mercy\\
and a hungry fate

\songpart{Pre-Chorus}
He bowed to the dark\\
with trembling breath\\
To guard that bond\\
he bargained with death

A father’s hand\\
in the cold of need\\
A hideous oath\\
from a world of creed

\songpart{Chorus}
Guest and host\\
guest and host\\
Keep the vow
(keep the vow)

Guest and host\\
guest and host\\
Break the seal\\
and the grave will know

\songpart{Bridge}
Not flesh alone\\
not mortal shame\\
But sacred law\\
that fed the flame

A parable carved\\
in brine and stone\\
For every house\\
that stands alone

If you turn away\\
the stranger’s face\\
You summon the dark\\
to your own place

\songpart{Final Chorus}
Guest and host\\
guest and host\\
Keep the vow
(keep the vow)

Guest and host\\
guest and host\\
Honor the door\\
or the dead come close

Guest and host\\
guest and host\\
Hear the tale\\
let the black winds know

Guest and host\\
guest and host\\
Keep the vow
(keep the vow)
\end{lyrics}


\end{multicols*}

\chapter{Essays}
\albumcover{covers/internal/essays}
The essay chapter of the songbook: watershed geography argued out loud, everyday debates followed to their logical end, a dilemma weighed from every side before the last chorus resolves it -- or refuses to.
\clearpage\begin{multicols*}{4}
\songtitle{Retrogrades}

\tag*{2023}\tag{social-media-seed}\tag{country}

\begin{notes}
Lyrics written by a LLM, based on a tweet by Carlos Thompson.
\end{notes}

\begin{style}
country, bass, banjo, fiddle, harmonica; female lead
\end{style}

\begin{lyrics}
Some look up to the midnight sky,\\
And watch the planets roll,\\
To track the backward path they tread,\\
Upon a cosmic scroll.

They see the shadow cross the moon,\\
A dance of light and space,\\
A rare and wondrous clockwork move,\\
Across the stellar face.

\songpart{chorus}
Some of us look at the sky and retrogrades make our minds.\\
For some those planets when retrogrades bring fateful winds

For them it is a clock of stone,\\
Of science and of light,\\
An mathematical display\\
That decorates the night.

But others look upon the stars,\\
And feel a sudden dread,\\
They read the retrogrades with fear,\\
And bow a worried head.

\songpart{chorus}
Some of us look at the sky and retrogrades make our minds.\\
For some those planets when retrogrades bring fateful winds

They think the shadow on the moon\\
Will twist the ties of fate,\\
And bring a season of regret,\\
Or open up a gate.

To them the sky is whispering\\
Of fortunes lost or won,\\
A warning written in the dark,\\
Before the rising sun.

\songpart{chorus}
Some of us look at the sky and retrogrades make our minds.\\
For some those planets when retrogrades bring fateful winds

\end{lyrics}


\songtitle{Balance of Three}

\tag*{2024}\tag{glam-metal}

\begin{notes}
A distinctive arrangement that supports the song's storytelling and emotional arc. %
The lyrics begin with “We build our lives on shifting ground”, giving a quick sense of the song’s tone.
\end{notes}

\begin{style}
glam metal, 136 BPM, C G Am F progression, close-mic baritone vocals, grand piano stabs, baritone sax riffs, acoustic guitar strums, tambora accents, string swells, distorted electric lead, gang vocal chorus, tape saturation, spring reverb, gated snare, minor key menace
\end{style}

\begin{lyrics}

\songpart{Verse 1}
We build our lives on shifting ground\\
Where order keeps the world around\\
Yet freedom calls, a voice so clear\\
And equity draws all near

\songpart{Pre-Chorus}
Three forces pull, they intertwine\\
A dance of hearts, a grand design\\
We choose our path, we take our stand\\
The wheel turns on across the land

\songpart{Chorus}
Balance of three, the world we see\\
Order, freedom, equity\\
Spin the wheel, let voices rise\\
Together we can harmonize

\songpart{Verse 2}
The conservative holds the frame\\
The libertarian lights the flame\\
The socialist seeks equal share\\
Each vision shapes the world we bear

\songpart{Pre-Chorus}
Extremes may break, but still we try\\
To find the truth where futures lie\\
The wheel of thought, it spins anew\\
And every choice belongs to you

\songpart{Chorus}
Balance of three, the world we see\\
Order, freedom, equity\\
Spin the wheel, let voices rise\\
Together we can harmonize

\songpart{Bridge}
Fascist walls, anarchist cries\\
Communist dreams, collectivist ties\\
But in the center, hope remains\\
A song of balance breaks the chains

\songpart{Final Chorus}
Balance of three, the world we see\\
Order, freedom, equity\\
Spin the wheel, let voices rise\\
Together we can harmonize

\end{lyrics}


\songtitle{Critter Kin (Man's an Animal Too)}

\tag*{2026}\tag{country}\tag{quora-posting}

\begin{notes}
A strong rock-driven song built on aggressive momentum and emotional urgency. %
The lyrics begin with “Well I grew up in the holler where the words come plain,”, giving a quick sense of the song’s tone.
\end{notes}

\begin{style}
Up-tempo country shuffle, rapid-fire vocals, heavy boot-stompin' bass and snare drum, fast-talking verses, stadium rock guitar fills, big gang vocal chorus, energetic modern production, 140 bpm
\end{style}

\begin{lyrics}
\songpart{Verse 1} \annotation{Fast-talking vocals, rhythmic delivery, heavy bass groove}\\
Well I grew up in the holler where the words come plain,\\
Called my kin "people," outsiders "them strangers" again.\\
Man and animal, two words sharp as a knife,\\
One for us walkin' tall, one for the rest of life.\\
Plants sit still in the dirt, animals run and fly,\\
But deep down in the bones, we're all under the same sky.

\songpart{Chorus} \annotation{Big singalong vocals, gang vocals, stadium rock energy}\\
Man was an animal, always has been, friend,\\
Darwin just told the story of how the family tree bends.\\
From the cat in the barn to the monkey in the tree,\\
We're cut from the same cloth, biologically!\\
Bones and blood and breath, warm or cold,\\
Humans ain't special outside the words we hold.\\
Yeah, we're animals too, runnin' wild and free,\\
Darwin didn't make it so — he just helped us see!

\songpart{Verse 2} \annotation{Fast-talking vocals, rhythmic delivery, heavy bass groove}\\
Take a human, a cat, and a fly buzzin' 'round,\\
Cat's got my backbone, my warm blood, same ground.\\
Fly's got wings and armor, six legs doin' their dance,\\
But that cat's closer to me than any fly's got a chance.\\
Macaque in the jungle, gorilla beatin' his chest,\\
Chimps laughin' closer to us than all the rest.\\
Before Darwin ever drew breath or picked up a pen,\\
Linnaeus had us pegged: mammals, primates, kin!

\songpart{Chorus}
Man was an animal, always has been, friend,\\
Darwin just told the story of how the family tree bends.\\
From the cat in the barn to the monkey in the tree,\\
We're cut from the same cloth, biologically!\\
Bones and blood and breath, warm or cold,\\
Humans ain't special outside the words we hold.\\
Yeah, we're animals too, runnin' wild and free,\\
Darwin didn't make it so — he just helped us see!

\songpart{Verse 3} \annotation{faster, rapid-fired, rhythmic delivery, heavy bass groove}\\
Vernacular talk, keep your "man vs. %
beast" line,\\
But in the lab coat and microscope, that don't hold no shine.\\
Internal skeleton, two eyes, four limbs in line,\\
Hair and lungs and heart beatin' right in time.\\
Chimps closer to us than gorillas in the mist,\\
Evolution just connects the dots that always exist.\\
No microscope needed back in Linnaeus day,\\
Humans in the animal book — that's just the way!\\
Bridge (half-time, spoken-sung)\\
Now outside biology, call 'em animals all you please,\\
Two-legged critters with big ol' ideas and knees.\\
But truth in the blood and the fossil in the stone,\\
We're part of the family — we ain't sittin' alone.

\songpart{Final Chorus} \annotation{big finish, gang vocals}\\
Man was an animal, always has been, y'all,\\
Darwin lit the lantern so we could see it all!\\
From the dirt to the stars, same wild family tree,\\
Humans are animals — and that's biology!\\
Yeah we're animals too, runnin' wild and true,\\
Thank the good Lord and Darwin for seein' us through!

\songpart{Outro — fade with fiddle and stomps}
Critter kin... %
yeah we're critter kin...
\end{lyrics}


\songtitle{La Casa y el Viento}

\tag*{2026}\tag{pasillo}\tag{spanish-lyrics}

\begin{notes}
Lyrics by Suno based on an unsent WhatsApp message by Carlos Thompson.
\end{notes}

\begin{style}
Pasillo colombiano with gentle 6/8 sway, tender acoustic guitar and tiple arpeggios; verse stays sparse and confessional, pre-chorus lifts with string pads and a small cadence push, chorus opens with warm harmony doubles and a chantable anchor phrase, bridge drops to voice and guitar before a final chorus blooms. %
Close-mic female vocals, intimate breath, soft backing harmonies, subtle room reverb, a tearful but elegant mix with a polished folk sheen., theme
\end{style}

\begin{lyrics}

\songpart{Verse 1}
Si yo me cargo la casa,\\
me dejo la espalda en sal.\\
Lijo la puerta que cruje,\\
barro el polvo del umbral.

Miro la pared torcida,\\
la quiero dejar en paz.\\
Cuesta, sí, pero al final\\
la casa me queda igual.

\songpart{Pre-Chorus}
Me doblo, me vuelvo a levantar,\\
pongo mis manos, una vez más.\\
Y aunque me duela hasta respirar,\\
algo se endereza al final.

\songpart{Chorus}
La casa, sí, mi casa,\\
como la quiero dejar.\\
La casa, sí, mi casa,\\
por fin se deja abrazar.

La casa, sí, mi casa,\\
cerca de mi voluntad.\\
La casa, sí, mi casa,\\
y eso sí vale el final.

\songpart{Verse 2}
Pero otra cosa es tu sombra,\\
tu reloj sin avanzar.\\
Peleo contra tu costumbre,\\
y se me rompe el andar.

Te espero, te hablo, te apuro,\\
y tú te quedas atrás.\\
Yo doy la vida en empujarte,\\
y tú me vuelves a soltar.

\songpart{Pre-Chorus}
Me doblo, me vuelvo a levantar,\\
pongo mis manos, una vez más.\\
Pero aquí no nace claridad,\\
solo cansancio al caminar.

\songpart{Chorus}
La casa, sí, mi casa,\\
como la quiero dejar.\\
La casa, sí, mi casa,\\
por fin se deja abrazar.

La casa, sí, mi casa,\\
cerca de mi voluntad.\\
La casa, sí, mi casa,\\
y eso sí vale el final.

\songpart{Bridge}
No te digo que te calles,\\
te digo: mira el lugar.\\
No gastes el pecho entero\\
en lo que no va a cambiar.

Si el suelo no quiere firmeza,\\
si la puerta no va a entrar,\\
hay batallas que te dejan\\
más lejos de tu verdad.

\songpart{Chorus}
La casa, sí, mi casa,\\
como la quiero dejar.\\
La casa, sí, mi casa,\\
por fin se deja abrazar.

La casa, sí, mi casa,\\
cerca de mi voluntad.\\
La casa, sí, mi casa,\\
y eso sí vale el final.

\end{lyrics}


\songtitle{Espejo Luminoso}

\tag*{2026}\tag{salsa}\tag{spanish-lyrics}

\begin{notes}
\end{notes}

\begin{style}
salsa, montuno, Spanish-English storytelling, male tenor lead, choir call and response, piano montuno riffs, conga tumbao, timbales cascara, bass tumbao, trumpet unison hits, tres montuno flourishes, horn stabs, live room bleed, 1970s analog warmth, abrupt dropouts, sparse verse to dense chorus, midtempo groove, bittersweet resolve, domestic tension
\end{style}

\begin{lyrics}

\songpart{Intro}
\annotation{Sets the stage, brass section playing a catchy, syncopated melody}

\songpart{Verse 1}
\annotation{Lead, male tenor vocals}\\
Hay quien busca en su morada\\
el refugio del descanso,\\
y hay quien vive su vivienda\\
como escudo y como manto.\\
Uno quiere solamente\\
silencio, calma, reposo;\\
el otro quiere una forma,\\
un espejo luminoso.

\songpart{Chorus}
\annotation{Choir} Hay dos tipos de personas\\
\annotation{Lead} uno vive en su refugio,\\
el otro vive su casa\\
como rostro y como pulso.\\
\annotation{Choir} Uno descansa sin culpa,\\
el otro carga el trabajo.

\songpart{Verse 2}
\annotation{Lead, male tenor vocals}\\
Un plato sobre la mesa,\\
una chaqueta en el suelo:\\
¿qué hay aquí que ver o corregir?\\
No entiendo ese desvelo.\\
Levantarme a recoger\\
sí arruina la vivencia;\\
mi casa es para vivirla,\\
no pa' perder la paciencia.

\songpart{Chorus}
\annotation{Choir} Hay dos tipos de personas\\
\annotation{Lead} uno vive en su refugio,\\
el otro vive su casa\\
como rostro y como pulso.\\
\annotation{Choir} Uno descansa sin culpa,\\
el otro carga el trabajo.

\songpart{Verse 3}
\annotation{Lead, male tenor vocals}\\
El plato sin lavar duele,\\
la chaqueta en el suelo pesa;\\
mi casa soy yo misma,\\
el desorden la atraviesa.\\
Y entonces toca elegir:\\
pelear o doblar el turno,\\
corregir lo que otros rompen,\\
callar y seguir el curso.

\songpart{Chorus}
\annotation{Choir} Hay dos tipos de personas\\
\annotation{Lead} uno vive en su refugio,\\
el otro vive su casa\\
como rostro y como pulso.\\
\annotation{Choir} Uno descansa sin culpa,\\
el otro carga el trabajo.

\songpart{Montuno}
\annotation{Choir} (¿Por qué el uno no arregla?)\\
\annotation{Lead} La casa es para vivirla\\
\annotation{Choir} (¿Por que no la arregla?)\\
\annotation{Lead} Porque arruina la vivencia\\
No hay que perder la paciencia

\annotation{Choir} (¿Por qué el otro se enoja?)\\
\annotation{Lead} Por corregir lo que otros rompen\\
\annotation{Choir} (¿Por qué no se sienta y descanza?)\\
\annotation{Lead} mi casa soy yo misma,\\
el desorden la atraviesa.

\annotation{Mambo, piano montuno riffs, conga tumbao,  bass tumbao, trumpet unison hits}

\songpart{Verse 4}
\annotation{Lead, male tenor vocals}\\
Uno ve un espacio libre,\\
el otro ve una herida;\\
no es maldad ni pereza,\\
es que no comparten vida.\\
El que descansa no finge,\\
el que carga no inventa;\\
habitan el mismo cuarto\\
en dos mundos sin encuentro

\songpart{Chorus - Outro}
\annotation{Choir} Hay dos tipos de personas\\
\annotation{Lead} uno vive en su refugio,\\
el otro vive su casa\\
como rostro y como pulso.\\
\annotation{Choir} Uno descansa sin culpa,\\
\annotation{Lead} el otro carga el trabajo.\\
\annotation{Choir} Hay dos tipos de personas\\
Uno descansa sin culpa,\\
el otro carga el trabajo.\\
el refugio del descanso,\\
y hay quien vive su vivienda\\
como escudo y como manto.

\end{lyrics}


\songtitle{Best Available Lens}

\tag*{2026}\tag{dubstep}\tag{original-work}

\begin{notes}
Original work.\\
Dubstep hymn on epistemology and practical philosophy.
\end{notes}

\begin{style}
Dubstep hymn with half-time stomp, snarling wobble bass, and handclap-driven verses; verse rides sparse organ stabs and clipped spoken-sung lines, pre-chorus strips to sub pulses and rising choir pads, chorus hits with a huge detuned drop and stacked congregational vocals. %
Filtered risers, reverse cymbal swells, and glitchy vocal chops bridge the sections. %
Bright, vast, and sermon-like in the mix, with close-mic lead and wide, thunderous lows., dubstep, world, contemporary, theme
\end{style}

\begin{lyrics}
\songpart{Verse 1}
You keep calling it a creed\\
But it’s a working frame\\
A way to test the world\\
Not a holy name\\
It has rules for doing\\
What it sets out to do\\
Build a better map\\
Of rock and rain and you

\songpart{Pre-Chorus}
Philosophers ask why\\
And philosophers ask whether\\
That’s their long debate\\
They hold it all together\\
Scientists just work\\
By a learned, living line\\
Passed from hand to hand\\
Across the edge of time

\songpart{Chorus}
Best available lens\\
Best available lens\\
Not the final word\\
Still the way we learn\\
Best available lens\\
Best available lens\\
Useful, tested, made\\
For the world as it turns

\songpart{Verse 2}
Don’t wrap it in a banner\\
Don’t make it some new church\\
A methodology\\
Is not a sacred perch\\
Neopositivism\\
Pragmatism too\\
History left its fingerprints\\
In all we say and do

\songpart{Pre-Chorus}
And that shiny little brand\\
That called itself a thunder\\
Wasn’t a doctrine\\
Just a label dressed as wonder\\
New atheism fades\\
When the slogans lose their spark\\
A package, not a philosophy\\
No torch, no ark

\songpart{Chorus}
Best available lens\\
Best available lens\\
Not the final word\\
Still the way we learn\\
Best available lens\\
Best available lens\\
Useful, tested, made\\
For the world as it turns

\songpart{Bridge}
There are saints and skeptics\\
In the same white coat\\
Christian, Jew, and Buddhist\\
All of them afloat\\
Some chase rigor hard\\
Some chase the subject first\\
Some bring sloppy hands\\
And some bring a sharper thirst\\
Laypeople can twist it\\
Till it fits their tune\\
Truth when it flatters them\\
Then doubt beneath the moon\\
But a work in progress\\
Never asks to be divine\\
It asks to be improved\\
By the next in line

\songpart{Breakdown}
Not a religion\\
Not an ideology\\
An enterprise with a purpose\\
A method with a key\\
Not ultimate truth\\
But the best we’ve got\\
To learn what’s there\\
And what is not

\songpart{Final Chorus}
Best available lens\\
Best available lens\\
Not the final word\\
Still the way we learn\\
Best available lens\\
Best available lens\\
Useful, tested, made\\
For the world as it turns\\
Best available lens\\
Best available lens\\
Take the better map\\
Over every false concern
\end{lyrics}


\songtitle{Todos al mismo lado}

\begin{notes}
\end{notes}

\begin{style}
indie pop, funk syncopation, soul inflection, electronic accents, fast-paced 128 BPM, male lead, melodic spoken word, hip-hop phrasing, staccato guitar chops, slap bass groove, syncopated clavinet, punchy drum kit, crisp rimshots, call-and-response hooks, breakdown drops, layered brass stabs, tape saturation, plate reverb, playful urgency
\end{style}

\begin{lyrics}
\songpart{Verse 1}
Hay una rueda vieja en la plaza,\\
con un país arriba sin dirección.\\
Unos empujan con toda la fuerza,\\
otros preguntan: “¿y quién eligió?”

Dicen que el tiempo no espera a nadie,\\
que hay que avanzar aunque duela más.\\
Que si frenamos somos enemigos,\\
que la duda es mirar para atrás.

\songpart{Pre-Chorus}
Y el que va adelante grita:\\
“¡No miren la división!”\\
Mientras tapa con una bandera\\
el ruido de la discusión.

\songpart{Chorus}
Todos al mismo lado,\\
todos al mismo lugar.\\
Pero nadie pregunta\\
si el borde se va a acabar.

Todos al mismo lado,\\
todos empujando más.\\
Y el que intenta detenerlo\\
es el que no quiere saltar.

\songpart{Verse 2}
Hay hambre escrita en las paredes,\\
hay casas hechas de cartón.\\
Hay promesas sobre la mesa\\
y una grieta en el corazón.

La gente busca una salida,\\
la gente quiere respirar.\\
Pero si todos miran al frente,\\
¿quién mira dónde van a llegar?

\songpart{Pre-chorus}
Y el líder dice tranquilo:\\
“Es momento de unidad.”\\
Pero una mano que tira fuerte\\
también puede querer salvar.

\songpart{Chorus}
Todos al mismo lado,\\
todos al mismo lugar.\\
Pero nadie pregunta\\
si el borde se va a acabar.

Todos al mismo lado,\\
todos empujando más.\\
Y el que intenta detenerlo\\
es el que no quiere saltar.

\songpart{Bridge}
Quizás el problema no sea la fuerza,\\
ni que queramos avanzar.\\
Quizás el problema sea la idea\\
de que solo hay una forma de caminar.

\songpart{Final Chorus}
No todos al mismo lado,\\
no todos tienen que igualar.\\
Un país no es una rueda\\
que alguien pueda empujar.

No todos al mismo lado,\\
pero juntos al final.\\
Porque a veces una voz distinta\\
es la que evita caer al mar.
\end{lyrics}


\songtitle{Limits of the Load}

\tag*{2026}
\begin{notes}
\end{notes}
\begin{style}
industrial rock with salsa tumbao, mellow female mezzo vocals, rapid-fire vocals, continuous flow
\end{style}
\begin{lyrics}
\songpart{Verse 1}
Years of reminders on the kitchen wall\\
Alarms that fade before the dishes dry\\
Conversations circling like a call\\
That echoes back the same unanswered why
\songpart{Pre-Chorus}
I’ve carried both our weights through every storm\\
Fixed what I could with hope and steady hands\\
But silent days have taught me this new form\\
Some mountains stay no matter how we plan
\songpart{Chorus}
This isn’t blame, it’s seeing what is true\\
Your limits carved in stone, my heart worn thin\\
I can’t keep pouring endless me for you\\
When will I choose the peace I need within?
\songpart{Verse 2}
Are there places where your word still holds fast?\\
Or does the weight fall only on these shores?\\
I’ve searched for answers, built them, made them last\\
While asking less and giving even more
\songpart{Pre-Chorus}
The strategies all crumble into dust\\
The therapy, the talks, the rearranged days\\
I’m tired of turning every “maybe” into “must”\\
And losing pieces of myself in praise
\songpart{Chorus}
This isn’t blame, it’s seeing what is true\\
Your limits carved in stone, my heart worn thin\\
I can’t keep pouring endless me for you\\
When will I choose the peace I need within?
\songpart{Bridge}
Maybe you’ll seek the help that fits your mind\\
Maybe we’ll redraw the lines we share\\
Or maybe this is all we’ll ever find\\
And I must learn to love you… from over there
\songpart{Final Chorus}
No more blame, just truth I finally see\\
Your limits carved in stone, my healing begins\\
I won’t keep pouring endless me for we\\
I choose the life where I can breathe again
\annotation{Soft piano fade on last line, gentle strings rising}
\end{lyrics}


\songtitle{Second Gun Salesman}

\tag*{2018}\tag{country}\tag{quora}
\begin{notes}
Lyrics by SUNO based on a Quora post by Carlos Thompson (published under Gabriel Sandoval).
\end{notes}
\begin{style}
Fast-paced country with clawhammer banjo, fiddle runs, upright bass, and kick-drum stomp; rapid-fire verses with breathless syllables, pre-chorus tightening into a call-and-response hook, chorus snapping into gang vocals and fiddle breaks, Lead vocal stays dry and front-row close, with stacked harmony on the punch lines, short delay throws on the last word of each hook, and a bright, punchy mix with porch-band grit, theme, country, rapid
\end{style}
\begin{lyrics}
\songpart{Verse 1}
He said, “One more gun”\\
To stop the first gun\\
Sounds like a sale pitch\\
In a buckskin suit

If one won’t do it\\
Then what comes next?\\
A second on the counter\\
With a bigger check

\songpart{Pre-Chorus}
No, I want a fix\\
Not a two-for-one\\
You call it protection\\
I call it spun

\songpart{Chorus}
Good guy with a gun\\
Bad guy with a gun\\
Who’s buyin’ two of them?\\
Who’s buyin’ two of them? (who?)

Good guy with a gun\\
Bad guy with a gun\\
Or a worse guy with a blaster\\
Invent it, then run

\songpart{Verse 2}
Blasters ain’t real yet\\
Well, there’s your door\\
That’s where minds should race\\
Not the hardware store

Don’t sell me a ladder\\
And call it a bridge\\
Don’t sell me a rifle\\
And call it a truce

Who’s the good guy now?\\
You with a badge?\\
You with a sermon?\\
You with a grudge?

\songpart{Pre-Chorus}
Law says “hold your fire”\\
Would you still comply?\\
If it takes your toys\\
Do you still feel right?

\songpart{Chorus}
Good guy with a gun\\
Bad guy with a gun\\
Who’s buyin’ two of them?\\
Who’s buyin’ two of them? (who?)

Good guy with a gun\\
Bad guy with a gun\\
Or a worse guy with a blaster\\
Invent it, then run

\songpart{Bridge}
What if talk could work?\\
What if armor could?\\
What if dodge and cover\\
Did more good than blood?

Who gets named “good”\\
When they’re first to shoot?\\
Who was guardin’ peace\\
Before the trigger loop?

\songpart{Final Chorus}
Good guy with a gun\\
Bad guy with a gun\\
Who’s buyin’ two of them?\\
Who’s buyin’ two of them? (who?)

Good guy with a gun\\
Bad guy with a gun\\
If the choice is always bullets\\
Then nobody won
\end{lyrics}


\songtitle{The Good Guy's Dilemma}

\tag*{2018}\tag{quora}
\begin{notes}
Lyrics by SUNO based on a Quora post by Carlos Thompson (published under Gabriel Sandoval).
\end{notes}
\begin{style}
fast-paced, high-energy bluegrass, traditional Appalachian instruments, rapid-fire banjo, virtuosic fiddle, upright bass, percussive foot stomps
\end{style}
\begin{lyrics}
\songpart{Verse 1}
Does the mantra ring with a double sale?\\
A salesman's trick in a smoky tale?
"Stop a bad guy with a gun," they say\\
Is it just a ploy to clear the way?\\
For two iron barrels, twice the cost\\
Counting up the profits while the peace is lost!

\songpart{Chorus}
Oh, the best way to stop a bad guy in his track\\
Is a worse guy with a blaster, ain't no turning back!\\
If blasters don't exist, well then innovate!\\
Don't let the manufacturers seal your fate!\\
They're selling you a second, then a third in line\\
While the "good guy" logic is losing its shine!

\songpart{Verse 2}
Now tell me where are all the good gals gone?\\
The gender gap in the fight drags on!\\
And what makes a "good guy" anyway?\\
Are you law-abiding until the sky turns gray?\\
If the law said "surrender," would you follow the rule?\\
Or would you be the rebel, breaking every school?

\songpart{Bridge}
They distrust the state, but they trust the steel\\
But who defines "good" in the way they feel?\\
Is it he who pulls the trigger 'fore a word is spoke?\\
Before a talk, a dodge, or a velvet cloak?\\
Before the armor's donned or the peace is sought?\\
Is that the kind of "good" that they really bought?

\songpart{Outro}
Who’s the good guy? Who’s the bad?\\
Is it just a marketing trick we've had?\\
Put down the iron, let the questions fly\\
Underneath the bright Appalachian sky.\\
Yeah, under the sky!
\annotation{Fast-paced fiddle and banjo breakdown}
\end{lyrics}


\songtitle{Draw The Tree}

\tag*{2022}\tag{k-pop}
\begin{notes}
Lyrics by SUNO based on a Quora post by Carlos Thompson (published under Gabriel Sandoval).
\end{notes}
\begin{style}
K-pop girl-band anthem with sharp syncopated drums, rubbery bass, and bright synth stabs; verse stays lean with tight rap-sung lines and clipped harmonies, pre-chorus opens into stacked lifts and rising tom fills, chorus hits with a chantable group hook and handclap drive, Vocal layering is glossy with unison doubles, airy ad-libs, and short delay throws on the title, Ear candy: reversed bell swells into each chorus, glittery synth flickers between lines, and a final-drop crowd chant, Mix is crisp, bright, and punchy, flat, k-pop, low, new age
\end{style}
\begin{lyrics}
\songpart{Verse 1}
You point at the fruit\\
I want the whole tree\\
One wrong little leaf\\
Can hide a whole sea

Flat Earth on the ground\\
Is the easy first call\\
But look at the branches\\
It’s bigger than all

Young Earth, old stories\\
Worn like a crown\\
One true truth for you\\
And the rest shut down

\songpart{Pre-Chorus}
Show me the roots\\
Show me the line\\
What blooms from a claim\\
When the proof runs dry?

If it sounds so certain\\
Why does it bend?\\
I’m not buying shadows\\
Just because they pretend

\songpart{Chorus}
Draw the tree\\
Draw the tree\\
Let the whole thing breathe\\
Draw the tree\\
Draw the tree\\
If it can’t stand proof, set it free

Draw the tree\\
Draw the tree\\
No more pretty lies for me\\
Draw the tree\\
Draw the tree\\
If it’s real, let it speak

\songpart{Verse 2}
Birthers and truthers\\
Chasing a spark\\
New Age promises\\
Written in dark

Healing by rumor\\
Not by a test\\
Miracle bottles\\
Selling the best

Science is “western”\\
So facts are a plot\\
If it fits the vibe\\
Then evidence not

Economy’s kingdom\\
Built on a hunch\\
So many bold answers\\
With very little punch

\songpart{Pre-Chorus}
Show me the roots\\
Show me the chain\\
Where does it grow\\
When it pours like rain?

If it’s sacred, political\\
Holy, or cool\\
I still need a reason\\
That beats the rule

\songpart{Chorus}
Draw the tree\\
Draw the tree\\
Let the whole thing breathe\\
Draw the tree\\
Draw the tree\\
If it can’t stand proof, set it free

Draw the tree\\
Draw the tree\\
No more pretty lies for me\\
Draw the tree\\
Draw the tree\\
If it’s real, let it speak

\songpart{Bridge}
I’m not afraid\\
Of what we might find\\
Truth gets cleaner\\
When we look real hard

Every branch claims\\
It deserves the sun\\
But proof is the gate\\
That lets the light in

\songpart{Final Chorus}
Draw the tree\\
Draw the tree\\
Let the whole thing breathe\\
Draw the tree\\
Draw the tree\\
If it can’t stand proof, set it free

Draw the tree\\
Draw the tree\\
No more pretty lies for me\\
Draw the tree\\
Draw the tree\\
If it’s real, let it speak

Draw the tree\\
Draw the tree\\
We can make it plain to see\\
Draw the tree\\
Draw the tree\\
Truth is in the roots for me
\end{lyrics}


\songtitle{What Was Left Behind}

\tag*{2026}\tag{quora}
\begin{notes}
Lyrics by SUNO based on an answer by Carlos Thompson in Quora.
\end{notes}
\begin{style}
Gospel with mezzo lead, handclaps, call-and-response choir, warm organ swells, steady tambourine groove, verse-led testimony, pre-chorus lifts into stacked choir harmony, chorus lands on clapping hooks and shouted responses, bridge opens to intimate solo then final choir swell, bright soulful mix, native american, gospel, strong, simple
\end{style}
\begin{lyrics}
\songpart{Verse 1}
Yes, I can walk with boundaries\\
Even if I don’t call on heaven\\
A lack of gods is just a line I draw\\
Not a cage, not a lesson

My code came through my human heart\\
Through mercy, truth, and care\\
And I still hear the wisdom rise\\
From the mount up in the air

\songpart{Pre-Chorus}
Quite consciously\\
I chose what stays\\
I kept the light\\
And left the haze

\songpart{Chorus}
Clap your hands, yes, I still believe\\
Clap your hands, love can make me free\\
Not by decree, not by a throne\\
I hold my ground, I walk my own
(What was left behind)\\
I kept the good, I kept it alive
(What was left behind)\\
I kept the good, and I still testify

\songpart{Verse 2}
Some lean on stoic fire\\
Some on Buddha’s quiet way\\
Some on old-world stories\\
And some on what the scholars say

Some take counsel from the land\\
From the drum and the cedar tree\\
Some wear the cross as a culture sign\\
But my steps stay faithful to me

\songpart{Pre-Chorus}
Quite consciously\\
I sifted through\\
The bones of hope\\
And kept the truth

\songpart{Chorus}
Clap your hands, yes, I still believe\\
Clap your hands, love can make me free\\
Not by decree, not by a throne\\
I hold my ground, I walk my own
(What was left behind)\\
I kept the good, I kept it alive
(What was left behind)\\
I kept the good, and I still testify

\songpart{Bridge}
Some say I’m bound to nothing\\
But I know that is not the song\\
A choice can wear a different name\\
And still know right from wrong

I’m not ruled by fear of fire\\
I’m not moved by borrowed shame\\
I choose the best that mercy taught\\
And I call that holy flame

\songpart{Final Chorus}
Clap your hands, yes, I still believe\\
Clap your hands, love can make me free\\
Not by decree, not by a throne\\
I hold my ground, I walk my own
(What was left behind)\\
I kept the good, I kept it alive
(What was left behind)\\
I kept the good, and I still testify

(What was left behind)\\
I kept the good, I kept it alive
(What was left behind)\\
I kept the good, and I still testify
\end{lyrics}


\songtitle{Two Ways to Think}

\tag*{2026}\tag{quora}
\begin{notes}
Lyrics by SUNO based on an answer by Carlos Thompson in Quora.
\end{notes}
\begin{style}
1950s-style children's explanatory jingle, bouncy doo-wop novelty with handclaps, ukulele, clarinet flourishes, and a bright singalong chorus; verse stays simple and story-like, pre-chorus widens the idea with small harmony stacks, chorus snaps into a catchy schoolyard chant with clapped backbeats, playful lead vocal with group replies, tiny kazoo fills and record-scratch transitions, sunny vintage mix, warm and close, breaks
\end{style}
\begin{lyrics}
\songpart{Verse 1}
Some folks think in one straight line\\
One right way, one right time\\
If your way bends, they say it's wrong\\
And they sing that rule all day long

\songpart{Pre-Chorus}
But listen close, and you may see\\
More than one way can help things be\\
When folks come near, and talk, and share\\
Old hard edges start to wear

\songpart{Chorus}
Two ways to think, two ways to see
(hey, hey, two ways to see)\\
You can learn from you and me
(two ways to think)\\
Maybe not one only road\\
Maybe many paths can grow\\
Two ways to think, two ways to see
(two ways to see)

\songpart{Verse 2}
In a small little circle, tight and neat\\
It’s easy to march to the same old beat\\
New ideas come, and slam the door\\
“Wrong, wrong, wrong!” they shout once more

\songpart{Pre-Chorus}
But when you meet a different friend\\
And find they’re good from end to end\\
You may say, “Wait, hold on a tick\\
Maybe my old map was too quick”

\songpart{Chorus}
Two ways to think, two ways to see
(hey, hey, two ways to see)\\
You can learn from you and me
(two ways to think)\\
Maybe not one only road\\
Maybe many paths can grow\\
Two ways to think, two ways to see
(two ways to see)

\songpart{Bridge}
And if you start to call folks “them”\\
You miss the fact they’re people, too\\
Equals in the schoolyard sun\\
Trying hard to get things done

\songpart{Final Chorus}
Two ways to think, two ways to see
(hey, hey, two ways to see)\\
You can learn from you and me
(two ways to think)\\
Maybe not one only road\\
Maybe many paths can grow\\
Two ways to think, two ways to see
(two ways to see)
\end{lyrics}


\input{temp/los_buenos_somos_mas}
\songtitle{No Middle Seat}

\tag*{2020}\tag{facebook}
\begin{notes}
Lyrics cowritten by Carlos Thompson and Grok, based on a post by Carlos Thompson on Facebook.
\end{notes}
\begin{style}
Folk protest with handclaps, stomped floor toms, banjo, acoustic guitar, and a low driving bass pulse; verse rides close-mic voice and rootsy strums, pre-chorus lifts with stacked group chants, chorus lands on a plainspoken singback hook, and the bridge opens into electronic swirls, filtered pads, and a rising pulse before snapping back to the folk band, Dry, intimate mix with a bright, earthy top, folk
\end{style}
\begin{lyrics}
\songpart{Verse 1}
I’m not centrist\\
I am non-aligned\\
I won’t bow to poles\\
That split my spine

My principles stand\\
On their own two feet\\
They don’t fit the box\\
You made for me

\songpart{Pre-Chorus}
Don’t ask me for\\
The just middle\\
If the middle’s\\
Just a riddle

If one side brings\\
A line of graves\\
I won’t call that\\
A thing I can save

\songpart{Chorus}
I am non-aligned (non-aligned)\\
I am non-aligned (non-aligned)\\
I choose no halfway\\
When harm is the plan\\
I take what is right\\
From where it stands

\songpart{Verse 2}
I’m not warm and soft\\
For the sake of peace\\
I won’t trade my name\\
For polite beliefs

I’ll take the good\\
From one hand’s reach\\
Then take the good\\
From the other side’s speech

But I won’t buy\\
The whole damn pack\\
With the poison kept\\
And the kindness wrapped

\songpart{Pre-Chorus}
Don’t ask me for\\
The just middle\\
If the middle’s\\
Just a riddle

If one side says\\
Let blood be law\\
I’m not divided\\
I’m against all

\songpart{Chorus}
I am non-aligned (non-aligned)\\
I am non-aligned (non-aligned)\\
I choose no halfway\\
When harm is the plan\\
I take what is right\\
From where it stands

\songpart{Bridge}
\annotation{Electronic swell}\\
Let the sirens fade\\
Let the labels drop\\
I am not your fence\\
I am not your prop

I will stand with truth\\
Where truth is clear\\
And I’ll name the line\\
When it comes near

\songpart{Chorus}
I am non-aligned (non-aligned)\\
I am non-aligned (non-aligned)\\
I choose no halfway\\
When harm is the plan\\
I take what is right\\
From where it stands

I am non-aligned (non-aligned)\\
I am non-aligned (non-aligned)\\
No warm little lie\\
No forced middleman\\
I take what is right\\
From where it stands
\end{lyrics}


\songtitle{Bogotá Backbone}

\tag*{2000}\tag{english-lyrics}
\begin{notes}
Lyrics by SUNO (2026) based on a Gemini Notebook report on recovered documents in the Internet Archive from 2000's \emph{Por una Bogotá inteligente} by Carlos Thompson. 
\end{notes}
\begin{style}
Ska with brisk offbeat guitar, bouncing bass, and punchy horn stabs; verse skanks on nimble half-step rhythm, pre-chorus strips to bass and rimshot tension, chorus opens wide with shouted gang vocals and bright brass riffs, Use tight male lead with doubled hook lines, quick call-and-response ad-libs, and short delay throws on key phrases, Add rising trombone fills, a stop-time break before the last chorus, and crisp, lively mix with a shiny punch, visionary, ska, environmental
\end{style}
\begin{lyrics}
\songpart{Verse 1}
In the year two thousand\\
they drew a city map\\
Not of roads and alleyways\\
but cable underwrap

From towers to the houses\\
from campus to the block\\
They said every building\\
needs a pulse and lock

\songpart{Pre-Chorus}
Link it up\\
Make it move\\
One nerve\\
One groove\\
From the desk\\
To the gate\\
Keep the whole place running straight

\songpart{Chorus}
Bogotá backbone!\\
Run it clean, run it strong\\
Bogotá backbone!\\
Tie the whole town along\\
One web, one city\\
One signal in the air\\
Bogotá backbone!\\
Everybody get there

\songpart{Verse 2}
Professors in the blueprints\\
business at the wheel\\
Government and neighbors\\
cut the deal for real

Public safety, water lines\\
trash bins, rain, and heat\\
Automate the warning\\
keep disaster off the street

\songpart{Pre-Chorus}
Link it up\\
Hear the call\\
Every district\\
Every hall\\
From the lab\\
To the square\\
Let the whole city share

\songpart{Chorus}
Bogotá backbone!\\
Run it clean, run it strong\\
Bogotá backbone!\\
Tie the whole town along\\
One web, one city\\
One signal in the air\\
Bogotá backbone!\\
Everybody get there

\songpart{Bridge}
Start local\\
Grow wide\\
Little smart systems side by side\\
Then federate\\
Let them meet\\
One live network through the streets

\annotation{Break}\\
Oh-oh\\
Watch it connect\\
Oh-oh\\
Watch it connect

\songpart{Final Chorus}
Bogotá backbone!\\
Run it clean, run it strong\\
Bogotá backbone!\\
Tie the whole town along\\
One web, one city\\
One signal in the air\\
Bogotá backbone!\\
Everybody get there
\end{lyrics}


\songtitle{History Witnessed Live}

\tag*{2026}\tag{post-punk}\tag{original-lyrics}
\begin{notes}
Written by Carlos Thompson, based on historical events during his lifetime.
\end{notes}
\begin{style}
post-punk, spoken word intro, male baritone vocals, gang vocal chorus, baritone chant, angular bassline, baritone guitar stabs, tom-heavy drums, syncopated hi-hat, brass interlude, tape saturation, plate reverb, dramatic dynamics, 92 BPM, march-like pulse, brooding defiance, media critique, instrumental build-up
\end{style}
\begin{lyrics}
\songpart{Intro}
Through my hotel window\\
a funeral procession passed\\
Generalisimo Francisco Franco\\
I was barely three.

\songpart{Verse 1}
They took the Embassy\\
I witnessed it on magazines sent by mail\\
then they took the palace\\
and that was on live TV

And then a pilot reporting\\
that Armero was no more\\
live on morning radio\\
before I headed to school

\songpart{Pre-Chorus}
I saw the wall fall\\
the end of history\\
a new age of hegemony\\
on the great values of the west

\songpart{Chorus}
History as a subject at school\\
History witnessed live\\
Told by media and governments\\
Today trues, tomorrow lies

\songpart{Verse 2}
I saw the Apartheid fall\\
followed by Rwanda and Zaïre\\
friends from Africa in my memory\\
Eritreans and Ethiopians\\
Angolans and Congolese

I saw the towers fall\\
from Usenet speculations\\
to official blame on al Qaeda\\
and war on the Taliban

\songpart{Pre-Chorus}
Weapons of mass destruction\\
Halliburton and freedom fries\\
Seeing that nation united\\
on some obvious lies

\songpart{Chorus}
History as a subject at school\\
History witnessed live\\
Told by media and governments\\
Today trues, tomorrow lies

\songpart{Verse 3}
A million voices against FARC\\
Social media in the horizon\\
as a new citizen activism\\
until Ola Verde fell

The Arab spring raising\\
from the promises of Tunisia\\
to the Syrian shit show\\
and Occupy Wall Street waded

\songpart{Bridge - more spoken}
They Brexited across the pond\\
Birthers brought Trump to Penn Avenue\\
And here they voted NO to peace

The end of history\\
no longer West versus East\\
but neighbor against neighbor\\
based on which network you choose

\songpart{Final Chorus}
History as a subject at school\\
History witnessed live\\
Told by media and governments\\
Told by peers and engagement bots\\
One side's trues, the other's lies
\end{lyrics}


\songtitle{El fin de la historia}

\tag*{2026}\tag{post-punk}\tag{spanish-lyrics}
\begin{notes}
Written by Carlos Thompson as a Spanish-language adaptation of \emph{History Witnessed Live}, based on historical events during his lifetime.
\end{notes}
\begin{style}
post-punk, spoken word intro, male baritone vocals, gang vocal chorus, baritone chant, angular bassline, baritone guitar stabs, tom-heavy drums, syncopated hi-hat, brass interlude, tape saturation, plate reverb, dramatic dynamics, 92 BPM, march-like pulse, brooding defiance, media critique, instrumental build-up
\end{style}
\begin{lyrics}
\songpart{Intro - spoken}
Apenas tres años\\
y desde la ventana del hotel\\
pasó una gran procesión fúnebre\\
de un tal Generalísimo Franco.

\songpart{Verse 1}
Llegaban las revistas por correo:\\
la toma de la embajada dominicana.\\
Años después, en la tele,\\
se tomaron el Palacio en vivo.

Y en vivo y al aire un piloto\\
con voz desesperada:\\
“Armero ya no existe”\\
mientras yo salía al colegio.

\songpart{Pre-Chorus}
Vi caer el Muro,\\
era el fin de la Historia,\\
el inicio de un nuevo orden\\
bajo los valores de Occidente.

\songpart{Chorus}
Historia en el colegio,\\
lejana, distante.\\
Historia en vivo y en directo,\\
contada por gobiernos y medios.

Hoy la verdad que reportan,\\
mañana la mentira que desmienten.

\songpart{Verse 2}
Vi caer el apartheid,\\
luego vendrían Ruanda y Zaire.\\
Mis amigos de África en mente:\\
eritreos y etíopes,\\
angoleños y congoleños.

Vi caer las Torres\\
especulando en mi tribu de Usenet.\\
El oficialismo culpó a Al Qaeda\\
y fue la guerra contra los talibanes.

\songpart{Pre-Chorus}
Armas de destrucción masiva,\\
Halliburton y las freedom fries,\\
y el discurso de los gringos unidos\\
entre patriotismo y mentiras.

\songpart{Chorus}
Historia en el colegio,\\
lejana, distante.\\
Historia en vivo y en directo,\\
contada por gobiernos y medios.

Hoy la verdad que reportan,\\
mañana la mentira que desmienten.

\songpart{Verse 3}
Un millón de voces contra las FARC,\\
las redes sociales en el horizonte\\
como un nuevo activismo de a pie\\
hasta que la Ola Verde se disipó.

Llega la Primavera Árabe,\\
la promesa esperanzada de Túnez,\\
y luego Siria se fue a la mierda\\
y Occupy Wall Street al olvido.

\songpart{Bridge - more spoken}
Dos mil dieciséis.\\
Votaron por el Brexit.\\
Votaron por Donald Trump.\\
Votaron por el No a la paz.

El fin de la Historia\\
ya no es Oriente contra Occidente,\\
sino vecino contra vecino\\
según tu cadena y tus redes.

\songpart{Final Chorus}
Historia en el colegio,\\
lejana, distante.\\
Historia en vivo y en directo,\\
contada por tus pares y por bots,\\
por gobiernos y entretenimiento.

La verdad que reporta un lado,\\
la mentira que el otro desmiente.
\end{lyrics}


\songtitle{EISH5 Worlds}

\tag*{1998}
\begin{notes}
Lyrics by SUNO based in a 1998 article by Carlos Thompson on his constructed script system EISH5.
\end{notes}
\begin{style}
Drum and bass, poetic and propulsive, 172 BPM with shuffled breakbeats, elastic sub-bass, liquid synth arps, and tight staccato piano accents; verse rides nervous chopped drums and close-mic spoken melody, pre-chorus strips to bass and filtered pads, chorus hits with lifted harmony stacks and a chantable code-name hook, Vocals are intimate and doubled in the hook, with whispered asides, delay throws on key words, and a final lift into airy topline, Ear candy: reversed blips before drops, code-click percussion, and glassy risers between sections, Bright, crisp, and cinematic with a midnight sheen, poetic, drum and bass, latin, complex
\end{style}
\begin{lyrics}
\songpart{Verse 1}
Since I remember, I’ve been making rooms\\
Out of paper, numbers, sparks\\
Always liked the cryptic edge\\
The hidden door, the mirrored marks

Math in the margins, physics in my thumb\\
Computers like a second skin\\
I never found the infant sketch\\
But the code was already in

\songpart{Pre-Chorus}
Then the feeling found a shape\\
A way to say what trembled in me\\
Who I meant, and how I burned\\
In the blue of memory

\songpart{Chorus}
EISH5, EISH5\\
Dot and dash, but alive\\
EISH5, EISH5\\
A world can hide inside\\
EISH5, EISH5\\
Let the eyes decode the night\\
EISH5, EISH5\\
I built a way to light

\songpart{Verse 2}
First came the hand-made Morse\\
Every sign in one soft thread\\
Horizontal for the dots\\
Vertical for what was said

Adjacent dots climbed up the page\\
A little ladder, line by line\\
Carlos once became a maze\\
And still it felt like mine

The old code moved like rules\\
Every shape had one home\\
Then freedom opened up the gate\\
And let one letter roam

\songpart{Pre-Chorus}
A human eye would miss the key\\
If it didn’t know the creed\\
One letter, many silhouettes\\
Same pulse, different breed

\songpart{Chorus}
EISH5, EISH5\\
Dot and dash, but alive\\
EISH5, EISH5\\
A world can hide inside\\
EISH5, EISH5\\
Let the eyes decode the night\\
EISH5, EISH5\\
I built a way to light

\songpart{Verse 3}
Greek letters in my pocket\\
Not the tongue, but the bones\\
Cyrillic at the table\\
Borrowed crowns and borrowed thrones

Then came Thompinian\\
A shield for Spanish sleep\\
Not just swapping every letter\\
Something harder, something deep

I started on a language\\
Three years, maybe more\\
Never finished, never will\\
Still I hear it at the door

\songpart{Bridge}
RI.T.H. in the lost file rain\\
Swedish wind, French ice, English flame\\
A personal language\\
A fictional name

I found Esperanto\\
Then I found the seams\\
Ido with its sharper blade\\
Interlingua in the beams

I wanted common ground\\
But not a borrowed mask\\
A bridge for all the minds\\
And a harder, truer task

\songpart{Final Chorus}
EISH5, EISH5\\
Dot and dash, but alive\\
EISH5, EISH5\\
A world can hide inside\\
EISH5, EISH5\\
Thompinian in the night\\
EISH5, EISH5\\
I keep the lost things bright

\songpart{Outro}
I lost the files and papers\\
But not the first intent\\
I still invent the meaning\\
In the shape my hands once sent

EISH5, EISH5\\
EISH5, EISH5
\end{lyrics}


\songtitle{Moon in the Frame}

\tag*{2023}
\begin{notes}
Lyrics by SUNO based on a ChatGPT conversation regarding a 2023 threat by Carlos Thompson that read:
\begin{quote}\itshape%
Why are gibbous moons common in photography but rare in paintings?
\end{quote}
\end{notes}
\begin{style}
Experimental rai with bouncing bendir rhythm and clipped synth stabs, medium-fast sway; verse rides sparse drums and a dry handclap pulse, pre-chorus opens with oud-like bends and layered ululations, chorus hits with gang vocals and a sharp chantable refrain, Lead vocal is close-mic with pitch bends, doubled on hooks, with delay throws on key words, Tape hiss, reversed cymbal swells, and a brass hit into the final chorus, Bright, gritty, and forward-mixed, rai, experimental
\end{style}
\begin{lyrics}
\songpart{Verse 1}
Why is the moon\\
always half a rumor?\\
Why does the painter\\
choose the clean slice?

The camera waits\\
for the crooked pearl\\
caught in one breath\\
between dark and light

\songpart{Pre-Chorus}
Ya, look up\\
it tilts and stays\\
a bitten coin\\
on a moving face

One eye loves\\
the shape that resists\\
one hand loves\\
what can be fixed

\songpart{Chorus}
Gibbous moon, gibbous moon\\
it slips into the lens\\
Gibbous moon, gibbous moon\\
too hard to hold in paint\\
Gibbous moon, gibbous moon\\
it asks for moving air\\
Gibbous moon, gibbous moon\\
and the brush says wait

\songpart{Verse 2}
Photos like proofs\\
like a witness grin\\
the frame says here\\
the night says then

Paintings dream\\
in older manners\\
they trim the sky\\
to make a sign

A full moon flatters\\
a crescent sings\\
but that fat bright edge\\
won't sit still for anyone

\songpart{Pre-Chorus}
Ya, look up\\
it slides and spins\\
not all at once\\
not neat within

The shutter loves\\
the borrowed glow\\
the painter needs\\
a steadier snow

\songpart{Chorus}
Gibbous moon, gibbous moon\\
it slips into the lens\\
Gibbous moon, gibbous moon\\
too hard to hold in paint\\
Gibbous moon, gibbous moon\\
it asks for moving air\\
Gibbous moon, gibbous moon\\
and the brush says wait

\songpart{Bridge}
Maybe paintings\\
want a symbol\\
not the middle\\
of becoming

Maybe cameras\\
eat the in-between\\
the almost-full\\
the almost-seen

\songpart{Final Chorus}
Gibbous moon, gibbous moon\\
it slips into the lens\\
Gibbous moon, gibbous moon\\
too hard to hold in paint\\
Gibbous moon, gibbous moon\\
I like the bent-up light\\
Gibbous moon, gibbous moon\\
half wild, half right
\end{lyrics}


\songtitle{Check My Privilege}

\tag*{2026}\tag{folk-rock}\tag{original-source}
\begin{notes}
Lyrics by Suno based on a 2014 blog post written by Carlos Thompson under the pen name Rataflechera.
\end{notes}
\begin{style}
Folk-rock with rapid British-American folk-rock momentum: bright acoustic strumming, driving kick and punchy snare, overdriven electric accents narrowing into tense pre-choruses, fiddle-like guitar surges, strong intimate male English vocals with restrained gang doubles, brief harmony lifts, stripped hand-percussion bridge, crisp modern studio mix
\end{style}
\begin{lyrics}
\songpart{Verse 1}
I was a kid in Sweden,\\
then back to Colombian heat,\\
father chasing work abroad,\\
my shoes wore out from the seat.

Stockholm air in my lungs,\\
ZooTV in my ears,\\
Spanish school, then another,\\
I learned a map in years.

Math page after math page,\\
Australia, Cuba, China,\\
West Germany on a list,\\
Ecuador, Cape Town, Thailand, Japan.

\songpart{Pre-Chorus}
Some gifts came from effort,\\
some came from luck and grace,\\
some things opened wide for me,\\
some left bruises in their place.

\songpart{Chorus}
Check my privilege\\
Check my privilege\\
Not a crown, not a gift,\\
just a load I live with

Check my privilege\\
Check my privilege\\
Some doors swung for my name,\\
some doors stayed stiff

\songpart{Verse 2}
I can split a proof in half,\\
still choke on a room of eyes,\\
numbers never taught me charm,\\
or how to win a prize.

My mother, my father,\\
they stayed in the same door,\\
that kept one kind of ground,\\
but gave me something more.

Yet every move left cracks,\\
every school had a cost,\\
I was the odd boy again,\\
in the shuffle and the loss.

\songpart{Pre-Chorus}
And what was easy for me\\
was hard for someone else,\\
I mistook my open road\\
for the world itself.

\songpart{Chorus}
Check my privilege\\
Check my privilege\\
Not a crown, not a gift,\\
just a load I live with

Check my privilege\\
Check my privilege\\
Some doors swung for my name,\\
some doors stayed stiff

\songpart{Bridge}
I’ve walked through markets,\\
kept my face and my wallet,\\
but that is not a miracle,\\
just one part of the plot.

My skin, my build, my height,\\
my ordinary frame,\\
made me less of a target,\\
and that is part of the game.

So when you say my name,\\
don’t shut the whole room down,\\
a term can name the system,\\
but it can also drown.

\songpart{Final Chorus}
Check my privilege\\
Check my privilege\\
Say what it is, not less,\\
say what it is, not more

Check my privilege\\
Check my privilege\\
Some doors swung for my name,\\
some doors stayed stiff

Check my privilege\\
Check my privilege\\
Real, but mixed with pain,\\
real, but not a score
\end{lyrics}


\end{multicols*}

\part{Places}

\chapter{Circuits and roads}
\albumcover{covers/internal/circuits-and-roads}
Rivers and roads as closed loops, not lines to somewhere: the Valdai watershed, the Congo ring, the Silk Road reimagined as a circuit that returns to itself. These songs are less about arriving than about the shape the journey traces.
\clearpage\begin{multicols*}{4}
\songtitle{Valdai Watershed}

\tag*{2026}
\begin{notes}
\end{notes}
\begin{style}
Dubstep hymn with delicate female vocals, half-time sub drops and shuffling percussion; verse floats on sparse organ drones and bowed strings, pre-chorus narrows to handclaps and a rising synth whistle, chorus opens with huge bass pulses and layered choral pads. Female lead is close-mic and airy, with whispered doubles and stacked Viking chants in Swedish and Ruthenian on key phrases. Cold bell swells, reversed cymbals, and distant tavern-drone textures; wide, sacred, cinematic mix with deep low-end weight., medieval, female vocals
\end{style}
\begin{lyrics}
\songpart{Verse 1}
Valdai, low and small,\\
yet it crowns them all.\\
Three rivers hear its call,\\
three worlds at its wall.

A hill of 346,\\
and kingdoms learned to read\\
the streams like coded fire,\\
the land like a wired-up key.

\songpart{Pre-Chorus}
North to the white cold tide,\\
west where the Baltic slides.\\
South through the buried roads,\\
the plain begins to divide.

\songpart{Chorus}
Hold the watershed\\
hold the watershed
(ҳольд) hold the watershed

One hill, three destinies\\
one hill, three destinies
(slavá) one hill, three destinies

\songpart{Verse 2}
Down Ladoga they go,\\
Neva in the thaw.\\
Past Ilmen’s iron edge,\\
past Novgorod’s hawk-eyed law.

The poles came with the oars,\\
salt on their weathered hands,\\
then portage through the hush\\
to where the next river stands.

\songpart{Pre-Chorus}
Not a border stone,\\
but a living spine.\\
Who held that crossing place\\
held the routes between the lines.

\songpart{Chorus}
Hold the watershed\\
hold the watershed
(ҳольд) hold the watershed

One hill, three destinies\\
one hill, three destinies
(slavá) one hill, three destinies

\songpart{Bridge}
From Kyiv to the black,\\
from the black to Constantinople,\\
the road was water and will,\\
a pulse through the ironfold.

(Hej, skål) across the frost
(Hej, skål) across the stone\\
The valley learned to move\\
what the hills had always known.

\songpart{Final Chorus}
Hold the watershed\\
hold the watershed
(ҳольд) hold the watershed

One hill, three destinies\\
one hill, three destinies
(slavá) one hill, three destinies

The road was not the edge\\
the road was the artery
(ҳольд) hold the watershed\\
hold the watershed
\end{lyrics}


\songtitle{Congo Ring}

\tag*{2026}C\tag{roads}
\begin{notes}
Part of the geography cycles series, lyrics by SUNO, on summaries by Grok, from Claud collaborations with Carlos Thompson.
\end{notes}
\begin{style}
Congolese rumba with lilting guitar arpeggios, soft hand percussion, bass tumba, and syncopated sebene pulse; verse glides with call-and-response phrasing, pre-chorus narrows to voice and guitar, chorus opens with harmony stacks and chanted backing vocals, Add conga lifts, brief percussion breaks, and watery marimba fills, Warm, glossy, dance-floor intimate mix with rich midrange vocals, rumba, world, deep
\end{style}
\begin{lyrics}
\songpart{Verse 1}
Central Africa moves by water\\
That kind of map is rare\\
The Congo takes the long road\\
Heavy as the sea out there\\
Three point seven million kilometers\\
A basin like a drum\\
Round like a crown on the continent\\
And the river keeps it humming

\songpart{Pre-Chorus}
Highlands guard the rim\\
Like hands around a fire\\
Then the river finds its way\\
Through a narrow stone desire

\songpart{Chorus}
Congo River, carry on\\
Congo River, carry on\\
From the heart to the ocean\\
Congo River, carry on
(oh, carry on)
(oh, carry on)

\songpart{Verse 2}
Livingstone Falls keep watch\\
Thirty-two scars in line\\
Three hundred fifty kilometers\\
Blocking ships from the brine\\
The coast was there, but not the doorway\\
Not for long centuries\\
The world could touch the shoreline\\
But the inside held its leaves

\songpart{Pre-Chorus}
And inside grew deep\\
Whole kingdoms in the green\\
A world the sea could not disturb\\
Hidden in the basin's ring

\songpart{Chorus}
Congo River, carry on\\
Congo River, carry on\\
From the heart to the ocean\\
Congo River, carry on
(oh, carry on)
(oh, carry on)

\songpart{Bridge}
Then east on the rim, my baby\\
The Great Lakes shine like steel\\
Tanganyika, Kivu, Edward, Albert, Victoria\\
Ancient water, deep and real\\
The land is slowly opening\\
One side pulling far from one\\
And in those dark old waters\\
Fish are born that live nowhere

\songpart{Final Chorus}
Congo River, carry on\\
Congo River, carry on\\
From the heart to the ocean\\
Congo River, carry on
(oh, carry on)
(oh, carry on)
\end{lyrics}


\input{temp/seidenstrasse}
\songtitle{The Same Blue Field}

\tag*{2026}C\tag{roads}
\begin{notes}
Part of the geography cycles series, lyrics by SUNO, on summaries by Grok, from Claud collaborations with Carlos Thompson.
\end{notes}
\begin{style}
dubstep hymn with heavy synthesized bass with dubstep drowls, female mezzo vocals, pasodoble cantaora vocals in verses, martime chants in chorus, castanete claps, oud and sintir accompaniment, dubstep synthesized bridges, 124 BPM
\end{style}
\begin{lyrics}
\songpart{Verse 1}
Phoenician masts at dawn\\
Red hands on the rope\\
Carthage in the harbor\\
Salt on every hope\\
Olive, wheat, and sea\\
The first three keys\\
Ports to open hands\\
And carry all the years

\songpart{Pre-Chorus}
The waves did not divide us\\
They carried names instead\\
From Tunis to Iberia\\
From stone to tiled bed\\
Hear the old road waking\\
Under every keel\\
What they called an edge\\
Was always made to yield

\songpart{Chorus}
The sea was a road\\
The sea was a road
(Sea was a road)\\
The sea was a road\\
From hand to hand\\
From wall to wall\\
We crossed the same blue field\\
And it changed us all

\songpart{Verse 2}
Roman roads in rain\\
Byzantine bells in stone\\
Then the crescent rose\\
And made the courts its own\\
Al-Andalus in bloom\\
Lemon, spice, and fire\\
Morisco hands in Fez\\
Raising arches higher\\
Toledo kept the books\\
Gold ink, patient light\\
What Europe learned to call\\
Its dawn came through that night

\songpart{Pre-Chorus}
The same cities turned\\
Under different crowns\\
Yet the courtyards kept\\
Their water moving down\\
Clay pots, slow-cooked smoke\\
Fish, salt, and grain\\
Every layer added\\
To the wine and rain

\songpart{Chorus}
The sea was a road\\
The sea was a road
(Sea was a road)\\
The sea was a road\\
From hand to hand\\
From wall to wall\\
We crossed the same blue field\\
And it changed us all

\songpart{Bridge}
Flamingo bridges over the strait\\
One foot in Fez, one foot in fate\\
Hannibal on the mountain snow\\
Elephants in the wind that blows\\
Carthage fell, the earth turned red\\
Salt in the ground, the story bled\\
But the route lived on, alive, awake\\
In the spice, the tile, the bread we bake

\songpart{Final Chorus}
The sea was a road\\
The sea was a road
(Sea was a road)\\
The sea was a road\\
From hand to hand\\
From wall to wall\\
We crossed the same blue field\\
And it changed us all\\
The sea was a road\\
The sea was a road
(Sea was a road)\\
The sea was a road\\
Marrakech to the stone of Rome\\
We were never ever alone
\end{lyrics}


\songtitle{Mates and Rivers}

\tag*{2026}C\tag{roads}
\begin{notes}
Part of the geography cycles series, lyrics by SUNO, on summaries by Grok, from Claud collaborations with Carlos Thompson.
\end{notes}
\begin{style}
Zamba Argentina, world music, lilting 6/8 with hand drums, nylon-string guitar, bandoneón, bombo legüero, and dancing bass; verse sways in a spare, storytelling arc, pre-chorus opens into stacked tenor harmonies, chorus blooms with multilingual group chants and call-and-response, Add charango sparkles, riverbank echoes, and a final steppe-to-estuary lift, Warm, earthy, spacious mix with close-mic tenor vocals, theme, slow, male vocals, complex, world
\end{style}
\begin{lyrics}
\songpart{Verse 1}
From the Paraná\\
to the Uruguay\\
the water meets the sea\\
like a opened hand

Iguazú in thunder\\
green wall, white spray\\
north of the circuit\\
the map begins to sing

Pampa grass keeps gold\\
under a hard blue sky\\
gaucho boots on dust\\
and the cattle move slow

\songpart{Pre-Chorus}
Two great rivers\\
one wide mouth\\
wind turns the plain\\
and the fire wakes up

\songpart{Chorus}
Río de la Plata\\
agua y sal\\
Yvy porã, my land\\
tierra y pan

Pampa, cordero\\
mate en ronda\\
Amutuy, che peñi\\
vai che, vamos

Río de la Plata\\
bring me home\\
Río de la Plata\\
sing me on

\songpart{Verse 2}
Buenos Aires smoke\\
over embers and salt\\
achuras first, then skirt steak\\
then the blade-thin cut

Mollejas shining\\
vacío on the grate\\
empanadas folded\\
with a seal of heat

Dulce de leche\\
on a spoon, on a cone\\
heladerías at night\\
with pistachio and cream

\songpart{Pre-Chorus}
Montevideo\\
chivito tall\\
egg and ham and bacon\\
stacked like a vow

Medio y medio\\
glass to glass\\
friends pass the gourd\\
and the talk goes round

\songpart{Chorus}
Río de la Plata\\
agua y sal\\
Yvy porã, my land\\
tierra y pan

Pampa, cordero\\
mate en ronda\\
Amutuy, che peñi\\
vai che, vamos

Río de la Plata\\
bring me home\\
Río de la Plata\\
sing me on

\songpart{Bridge}
South, the Andes hold\\
like a bone in the earth\\
Patagonia wide\\
with a wind that learns your name

Southern ice fields\\
bright under black cloud\\
Beagle Channel crab\\
in a crack of red shell flame

Cross-roasted lamb\\
on a steppe herb fire\\
four long hours turning\\
as the dusk leans in

\songpart{Final Chorus}
Río de la Plata\\
agua y sal\\
Yvy porã, my land\\
tierra y pan

Pampa, cordero\\
mate en ronda\\
Amutuy, che peñi\\
vai che, vamos

Río de la Plata\\
bring me home\\
Río de la Plata\\
sing me on

\songpart{Outro}
From beans and fire\\
from fish and grain\\
we pass the cup\\
and call the names
\end{lyrics}


\songtitle{Red Bag Chichería}

\tag*{2026}C\tag{roads}
\begin{notes}
Part of the geography cycles series, lyrics by SUNO, on summaries by Grok, from Claud collaborations with Carlos Thompson.
\end{notes}
\begin{style}
Peruvian zarzuela fused with warm synth beds, quena melodies, and light criollo percussion; brisk verse lines glide like a travelogue, pre-chorus thins to baritone and hand drum, chorus opens with stacked quenas and a lifted cabasa pulse, Baritone lead stays close-mic with tasteful doubles, Quechua and Aymara ad-libs bloom in the hook, Lime-splash risers, panpipe echoes, and a final brass swell, Bright yet earthy mix, andean, fresh, mellow, world, contemporary, theme, raw
\end{style}
\begin{lyrics}
\songpart{Verse 1}
From Lima's dry white shore\\
To salt and citrus spray\\
Corvina in the bowl\\
And leche de tigre drinks your day\\
Ají and red onion\\
Flash like a blade in the sun\\
One minute and it's over\\
The feast has only just begun

\songpart{Pre-Chorus}
Then the road climbs up\\
Past every shade of green\\
Valleys fold to plateau\\
Like a map in between\\
One country in your hand\\
One mountain in your spine\\
You taste the heights\\
In every line

\songpart{Chorus}
Jallalla, Perú\\
From coast to puna\\
Jallalla, Perú\\
You feed the world in one land\\
Imaynalla, tayta\\
From Lima to Cusco\\
Jallalla, Perú\\
You rise in every hand

\songpart{Verse 2}
In the wok of Chifa\\
Beef hits rice and fries\\
Tiradito lays like moonlight\\
Under chili cream disguise\\
At Central, each course\\
Climbs a different sky\\
One plate is a calendar\\
Of roots that never die

\songpart{Pre-Chorus}
Then the road climbs up\\
To valleys, frost, and maize\\
Potatoes keep their records\\
In the work of ancient days\\
Not despite the stones\\
But because the earth is steep\\
The Andes taught the people\\
How to live where others sleep

\songpart{Chorus}
Jallalla, Perú\\
From coast to puna\\
Jallalla, Perú\\
You feed the world in one land\\
Imaynalla, tayta\\
From Lima to Cusco\\
Jallalla, Perú\\
You rise in every hand

\songpart{Bridge}
In Arequipa, rocoto burns\\
With egg and cheese inside\\
Cuy on the table\\
Where the old and hungry reside\\
In a chichería\\
With flowers by the door\\
I lift the red cup slowly\\
And ask the mountain for more

\songpart{Chorus}
Jallalla, Perú\\
From coast to puna\\
Jallalla, Perú\\
You feed the world in one land\\
Imaynalla, tayta\\
From Lima to Cusco\\
Jallalla, Perú\\
You rise in every hand

\songpart{Outro}
Chuño in the soup\\
Titicaca in the cold\\
Vertical archipelago\\
In every story told\\
Jallalla, Perú\\
Jallalla, Perú
\end{lyrics}


\songtitle{From the Fields to the Flood}

\tag*{2026}C\tag{roads}
\begin{notes}
Part of the geography cycles series, lyrics by SUNO, on summaries by Grok, from Claud collaborations with Carlos Thompson.
\end{notes}
\begin{style}
Trap with oriental-scale melodies, high-pitched female vocals, half-time bounce, sliding sub bass, crisp hats, and a dark sparse verse-to-anthem arc; verses stay tense and detail-rich, pre-chorus strips to voice and a plucked motif, chorus opens with stacked doubles and short chantable lines, bridge drops to filtered drums and a lone ad-libbed top line, with risers, reverse swells, and porcelain chime accents; bright-on-top, heavy-in-the-low-end mix, world, deep, rich, distinct, female vocals
\end{style}
\begin{lyrics}
\songpart{Verse 1}
Yellow River stains the map\\
North plain built from silt and math\\
Shang bones, old soil, deep roots\\
Fields feed kings in temple boots

Beijing on the edge of stone\\
Steppe wind testing every throne\\
Wall line cut through pass and ridge\\
Hold that gate, hold that bridge

Xi’an by the Wei, old spark\\
Silk Road dust in the market arc\\
Persian thread in the woven cloth\\
Indian gold and a monk’s oath

\songpart{Pre-Chorus}
Three rivers, one world\\
Folded in the mud\\
Old roads, new trade\\
From the fields to the flood

\songpart{Chorus}
River circuit, run it through\\
Yellow, Yangtze, Seoul too\\
River circuit, hold that line\\
Land and sea in the same design

River circuit, feel the pull\\
Marsh and mountain, heavy and full\\
River circuit, side by side\\
One long pulse, one moving tide

\songpart{Verse 2}
Yangtze wide like a silver blade\\
Farm belts feed what the empire made\\
Shanghai where the harbor mouths\\
Pull the inland riches south

Delta rice and merchant hands\\
Tea ship sails to foreign lands\\
South China Sea, Pearl River glow\\
Hong Kong deep where the trade winds go

Seas keep turning the outside in\\
Salt on the edge, cash in the spin\\
Monsoon roads with the painted bowls\\
Ceramic dreams on African coasts

\songpart{Pre-Chorus}
Three rivers, one world\\
Folded in the mud\\
Old roads, new trade\\
From the fields to the flood

\songpart{Chorus}
River circuit, run it through\\
Yellow, Yangtze, Seoul too\\
River circuit, hold that line\\
Land and sea in the same design

River circuit, feel the pull\\
Marsh and mountain, heavy and full\\
River circuit, side by side\\
One long pulse, one moving tide

\songpart{Bridge}
Korea sits with the sea on three sides\\
Bridge and shield where the history hides\\
Han River basin, Seoul on the rise\\
Narrow waist with the watchful eyes

Writing, law, and the old belief\\
Passed to Japan like a burning leaf\\
Then the DMZ like a frozen scar\\
Thirty-eight line, visible from far

Taiwan on the strait, close but apart\\
Same deep current, different heart\\
China, Japan, Korea in view\\
One sea carries the old and new

\songpart{Chorus}
River circuit, run it through\\
Yellow, Yangtze, Seoul too\\
River circuit, hold that line\\
Land and sea in the same design

River circuit, feel the pull\\
Marsh and mountain, heavy and full\\
River circuit, side by side\\
One long pulse, one moving tide

\songpart{Outro}
One plain, one delta, one shore\\
Three zones, but the sea is more\\
River circuit, ring of the age\\
Written in mud on every page
\end{lyrics}


\songtitle{Water Made Us}

\tag*{2026}\tag{trap}\tag{roads}
\begin{notes}
Part of the geography cycles series, lyrics by SUNO, on summaries by Grok, from Claud collaborations with Carlos Thompson.
\end{notes}
\begin{style}
Trap fused with Mor Lam vocal loops and shimmering gamelan-inspired electronics, mid-tempo with swung hi-hats, heavy sub pulses, and syncopated kick patterns, Verse rides sparse hand drums, plucked zither motifs, and ghostly female lead; pre-chorus strips to voice and bass drone, building with reversed swells; chorus hits with stacked doubles, clipped gang chants, and bright metallic bells, Ad-libs pan wide, with delay throws on key place names, Wide, crisp, and glossy with earthy low end, beats, female vocals, complex, distinct, electronic, industrial, vocal, banda, rich, sophisticated, gamelan, world
\end{style}
\begin{lyrics}
\songpart{Verse 1}
Five rivers cut the land\\
from Tibet down slow\\
Irrawaddy, Salween,\\
Mekong in the glow

Red River, Chao Phraya\\
draw their own lines\\
mountains made the borders\\
and the borders made the signs

\songpart{Pre-Chorus}
Water writes the roads\\
through the paddies and the stone\\
every bend a kingdom\\
every flood a throne

\songpart{Chorus}
Water made us, water made us\\
water made us, pull us home\\
Water made us, water made us\\
through the rivers and the foam\\
Mekong on my tongue\\
Malacca in my bones\\
Water made us, water made us\\
never just one coast

\songpart{Verse 2}
Mekong rises high\\
where the cold peaks burn\\
four thousand nine hundred\\
and it still keeps turning

Yunnan to Laos\\
then Thailand on the edge\\
Cambodia to Vietnam\\
with the delta at the end

Tonle Sap flips\\
when the wet winds come\\
lake turns like a mirror\\
and the fish run numbing

\songpart{Pre-Chorus}
Angkor held the floods\\
with its canals of gold\\
stone stayed in the temples\\
but the water took control

\songpart{Chorus}
Water made us, water made us\\
water made us, pull us home\\
Water made us, water made us\\
through the rivers and the foam\\
Mekong on my tongue\\
Malacca in my bones\\
Water made us, water made us\\
never just one coast

\songpart{Bridge}
Sixty-five kilometers\\
and the world leans in\\
India to East Asia\\
same tide, same skin

Srivijaya counted\\
then Malacca too\\
Portuguese, Dutch, British\\
all came trying to rule

Singapore at the mouth\\
where the trade winds kneel\\
geography keeps score\\
and it cashes what it feels

\songpart{Final Chorus}
Water made us, water made us\\
water made us, pull us home\\
Water made us, water made us\\
through the rivers and the foam\\
Mekong on my tongue\\
Malacca in my bones\\
Water made us, water made us\\
Asia in the flow

\songpart{Outro}
Monsoon by monsoon\\
we remember how\\
the sea, the lake, the river\\
made the map somehow
\end{lyrics}


\songtitle{Where Waters Sort}

\tag*{2026}C\tag{roads}
\begin{notes}
Part of the geography cycles series, lyrics by SUNO, on summaries by Grok, from Claud collaborations with Carlos Thompson.
\end{notes}
\begin{style}
House with indie-folk verses and grunge chorus, midtempo swung four-on-the-floor in verses with picked acoustic guitar, upright bass, soft handclaps, and close-mic tenor lead; pre-chorus strips to kick, pulse, and stacked harmonies; chorus hits with distorted guitars, crashing drums, gang-vocal lift, and short delay throws on the hook, Ear candy: reversed swell into the chorus, cedar-chime accents, and a filtered drum break before the final lift, Bright, punchy mix with gritty edges and intimate vocal presence, deep, steady, folk, world, rich, dramatic, rapid, male vocals, grunge
\end{style}
\begin{lyrics}
\songpart{Verse 1}
Northwest line\\
under my feet\\
Juan de Fuca moves\\
like a slow drumbeat

Sea turns cold\\
and the salmon come\\
from the ridge to the mouth\\
to the setting sun

\songpart{Pre-Chorus}
One story\\
in two hands\\
stone and water\\
making lands

\songpart{Chorus}
Same old story\\
same old song
(geology, ecology)\\
marching on

Same old story\\
same old song
(the earth is moving)\\
and the run is strong

\songpart{Verse 2}
Columbia runs\\
like a westbound spine\\
drops through the gorge\\
carving hard lines

Celilo spoke\\
in a thousand tongues\\
then the dam went up\\
and the silence won

\songpart{Pre-Chorus}
Cold blue water\\
cedar shade\\
all we took\\
all we made

\songpart{Chorus}
Same old story\\
same old song
(geology, ecology)\\
marching on

Same old story\\
same old song
(the earth is moving)\\
and the run is strong

\songpart{Verse 3}
Puget Sound\\
in its drowned-out frame\\
islands, channels\\
calling names

Canoes cut through\\
with the tide and spray\\
cedar gave\\
what the people prayed

\songpart{Bridge}
Yellowstone waits\\
under the divide\\
one drop west\\
one drop east rides

Geysers breathe\\
from the buried fire\\
headwaters split\\
and climb higher

\songpart{Final Chorus}
Same old story\\
same old song
(geology, ecology)\\
marching on

Same old story\\
same old song
(one world, two waters)\\
and the run is strong

Same old story\\
same old song
(the earth is moving)\\
and we belong
\end{lyrics}


\songtitle{Rif Valley Rave}

\tag*{2026}C\tag{roads}
\begin{notes}
Part of the geography cycles series, lyrics by SUNO, on summaries by Grok, from Claud collaborations with Carlos Thompson.
\end{notes}
\begin{style}
Eurodance fused with Nordic world percussion and folk textures, four-on-the-floor at a brisk club tempo with syncopated hand drums and bright synth stabs; verse rides sparse kick, bass, and Icelandic chant, pre-chorus opens with risers and claps, chorus blasts full piano hook and stacked group vocals, bridge drops to a cold drone and spoken lift before final lift, Close-mic lead, wide doubles on hooks, delay throws, wind chimes, aurora shimmer, crisp and punchy mix, dynamic, dramatic, light, deep, eurodance, world, warm
\end{style}
\begin{lyrics}
\songpart{Intro}
Norður hringrás\\
jökulhjarta\\
rífur landið

\songpart{Verse 1}
Young earth under my boots\\
Rift line cuts the field
Þingvellir in daylight\\
Two plates under me\\
Lava makes new ground\\
Ice breaks, rivers run

\songpart{Pre-Chorus}
Pull me apart\\
Then pull me back\\
Rock keeps moving\\
We don’t stand still

\songpart{Chorus}
North circuit, north circuit\\
We were made to move\\
North circuit, north circuit\\
New land under you\\
From ice to fire\\
From dark to blue\\
North circuit, north circuit\\
It pulls us through

\songpart{Verse 2}
West coast cliffs fall hard\\
Norway in the rain\\
Fjords like open jaws\\
Deep water, steep stone\\
East side goes low and green\\
Sweden to the Baltic

\songpart{Pre-Chorus}
Passes on the ridge\\
Ports awake all year\\
Warm sea in the cold\\
Keeps the line alive

\songpart{Chorus}
North circuit, north circuit\\
We were made to move\\
North circuit, north circuit\\
New land under you\\
From ice to fire\\
From dark to blue\\
North circuit, north circuit\\
It pulls us through

\songpart{Bridge}
Longyearbyen is changing\\
Ground gives way below\\
Seed vault in the thaw\\
Future on the floor\\
Midnight sun above me\\
Polar night below\\
Aurora on the edge\\
Tromsø in green glow

\songpart{Final Chorus}
North circuit, north circuit\\
We were made to move\\
North circuit, north circuit\\
New land under you\\
From ice to fire\\
From dark to blue\\
North circuit, north circuit\\
It pulls us through
\end{lyrics}


\input{temp/pesn_o_krugah_putey}
\songtitle{Ride The Circuits}

\tag*{2026}\tag{roads}
\begin{notes}
Lyrics by SUNO based on a summary of the \emph{Circuits and Roads} documment between Carlos Thompson and Claude.
\end{notes}
\begin{style}
Deep house with a fast four-on-the-floor pulse, rolling bass, glossy synth stabs, and a mellow female lead; verse rides a tight travelogue groove, pre-chorus strips to handclaps and filtered chords, chorus opens wide with stacked hook doubles and airy ad-libs, bridge drops to a pulsing low end before the final lift, Whispered textures, reverse swells, and bright percussion fills; clean, glossy, club-ready mix, deep house, mellow
\end{style}
\begin{lyrics}
\songpart{Verse 1}
Stockholm on my tongue\\
Helsinki in my coat\\
Tallinn on the water\\
Riga in the cold\\
Vilnius at dawn\\
Kyiv in the rain\\
Novgorod to the river\\
St. Petersburg, say my name\\
Moscow on the line\\
I never lose the thread

\songpart{Pre-Chorus}
Every road turns again\\
Every border melts\\
I can hear the drums\\
Under all these shells\\
One more mile, one more sign\\
One more world to cross\\
When the lights come up\\
I’m already lost

\songpart{Chorus}
Ride the circuits, ride the circuits\\
Take me there, take me there\\
Ride the circuits, ride the circuits\\
Every border, everywhere
(ride the circuits)\\
Ride the circuits, ride the circuits\\
I keep moving, I keep free

\songpart{Verse 2}
Istanbul to Tbilisi\\
Baku by the flame\\
Tashkent in the heat haze\\
Fergana in the grain\\
Samarkand in blue stone\\
Bukhara at night\\
Khiva in the courtyard\\
Aral Sea, fading light\\
Issyk-Kul extension\\
I’m tracing every mile

\songpart{Pre-Chorus}
Every road turns again\\
Every border melts\\
I can hear the drums\\
Under all these shells\\
One more mile, one more sign\\
One more world to cross\\
When the lights come up\\
I’m already lost

\songpart{Chorus}
Ride the circuits, ride the circuits\\
Take me there, take me there\\
Ride the circuits, ride the circuits\\
Every border, everywhere
(ride the circuits)\\
Ride the circuits, ride the circuits\\
I keep moving, I keep free

\songpart{Verse 3}
Kinshasa on the Congo\\
Brazzaville across\\
Mbanza Kongo, older\\
Than the maps we lost\\
Goma by the water\\
Lake Kivu in the sky\\
Kigali in the hills\\
Kampala passing by\\
Lake Victoria shining\\
Serengeti in the wind\\
Ngorongoro calling\\
Livingstone again\\
Victoria Falls rushing\\
Okavango green

\songpart{Bridge}
From Marrakech to Malta\\
From Seville to Carthage\\
Rio to Ushuaia\\
I keep carrying the atlas\\
Lima to Machu Picchu\\
Seoul to the Great Wall\\
Reykjavik to Longyearbyen\\
I want it all

\songpart{Chorus}
Ride the circuits, ride the circuits\\
Take me there, take me there\\
Ride the circuits, ride the circuits\\
Every border, everywhere
(ride the circuits)\\
Ride the circuits, ride the circuits\\
I keep moving, I keep free

\songpart{Outro}
I keep moving, I keep free\\
I keep moving, I keep free\\
Ride the circuits
\end{lyrics}


\songtitle{Circuit World}

\tag*{2026}\tag{house}\tag{roads}
\begin{notes}
\end{notes}
\begin{style}
house, trance, 132 BPM, spoken word intro, half-whisper chant, female octave doubles, 4 on the floor kick, syncopated bass stabs, sidechain pumping, gated synth pads, piano arpeggios, brass hits, harp pulses, dubstep drop, 6/8 verse lift, dynamic silence breaks, careful pronunciation of non-English toponyms, texture vocals
\end{style}
\begin{lyrics}
\songpart{Intro — spoken, building bass chords}
Ride the circuits
(ride the circuits)\\
take me there
(take me there)\\
every border\\
everywhere
(ride the circuits)\\
ride the circuits\\
I keep moving\\
I keep free

\annotation{Instrumental, 8 beats}

\songpart{Verse 1}
Circuit one — the Baltic and Valdai\\
where the rivers write the law of East
\annotation{soothing, whispered}\\
Stockholm · Helsinki · Tallinn\\
Riga · Vilnius · Kyiv\\
Novgorod · Petersburg · Moscow

\songpart{Verse 2}
Circuit two — the Silk Road, Steppe to Sea\\
where every trade route crossed the ancient land 
\annotation{soothing, whispered}\\
Istanbul · Tbilisi · Baku\\
Tashkent · Samarkand · Bukhara\\
Khiva · Fergana · Issyk-Kul

\annotation{Instrumental, 8 beats}

\songpart{Verse 3}
Circuit three — Central Africa, the Deep\\
where rivers, forests, great lakes never sleep
\annotation{soothing, whispered}\\
Kinshasa · Brazzaville · Mbanza Kongo\\
Goma · Kigali · Kampala\\
Victoria · Serengeti · Okavango

\songpart{Verse 4}
Circuit four — Morocco, Southern Spain\\
the shared sea connects them once again
\annotation{soothing, whispered}\\
Marrakech · Fez · Tangier\\
Seville · Córdoba · Granada\\
Cádiz · Tunis · Carthage

\annotation{Instrumental, 8 beats}

\songpart{Verse 5}
Circuit five — South America, the South\\
from Iguazú down to the river's mouth
\annotation{soothing, whispered}\\
Rio · Ouro Preto · Iguazú\\
Buenos Aires · Montevideo\\
El Calafate · Torres del Paine · Ushuaia

\songpart{Verse 6}
Circuit six — Peru, the Andean High\\
from Lima up where condors touch the sky 
\annotation{soothing, whispered}\\
Lima · Nazca · Arequipa\\
Titicaca · Tiwanaku · La Paz\\
Uyuni · Cusco · Machu Picchu

\annotation{Instrumental, 8 beats}

\songpart{Verse 7}
Circuit seven — CJK, the East\\
where ancient states meet modernity's feast  
\annotation{soothing, whispered}\\
Great Wall · Xi'an · Terracotta Army\\
Shanghai · Hong Kong · Taipei\\
Beijing · Seoul · Kyoto

\songpart{Verse 8}
Circuit eight — Southeast Asia, the Sea\\
temples, trade, and worlds that used to be 
\annotation{soothing, whispered}\\
Singapore · Angkor Wat · Bangkok\\
Ho Chi Minh City · Mekong Delta\\
Hue · Hoi An · Phuket

\annotation{Instrumental, 8 beats}

\songpart{Verse 9}
Circuit nine — the Northwest, rivers run\\
salmon, stone, and civilizations done   
\annotation{soothing, whispered}\\
Columbia · Portland · Seattle\\
Vancouver · Haida Gwaii · Olympic\\
Celilo · Yellowstone · Grand Teton

\songpart{Verse 10}
Circuit ten — the North, the Arctic light\\
where Sami, ice, and midnight sun unite 
\annotation{soothing, whispered}\\
Reykjavik · Þingvellir · highlands\\
Narvik · Lofoten · Tromsø\\
Kiruna · Jukkasjärvi · Svalbard

\songpart{Outro — Instrumental, 8 beats}
\end{lyrics}


\end{multicols*}

\chapter{Geography and spiritual travel}
\albumcover{covers/internal/geography-and-spiritual-travel}
Travel as a spiritual exercise rather than a checklist, even when it starts as one: a bucket list, a road trip through Bogotá, a lighter heart earned one kilometer at a time, and -- most recently -- five unhurried stops across northern Indiana, just to see what's there.
\clearpage\begin{multicols*}{4}
\songtitle{Bucket List}

\tag{original-work}\tag{house}

\begin{notes}
Original work.\\
House music meditation on aspirations and spiritual geography.
\end{notes}

\begin{style}
house, electronic, 128 BPM, four-on-the-floor kick, syncopated offbeat bass, filtered saw lead, analog synth pads, metallic percussion, muted guitar chanks, chopped vocal sample, sidechain pumping, tape saturation, club reverb, hypnotic momentum, restrained menace, gradual lift, mezzosoprano vocals with high range
\end{style}

\begin{lyrics}
\annotation{The Main List}
\songpart{Intro: Deep bassline builds, hi-hats start ticking}

\annotation{Spoken}\\
10. %
Alhambra, Spain
\annotation{Sung}\\
Islamic art where the stone tells a story,\\
Citadel of palaces and Moorish glory.\\
Lost for centuries, restored in the south,\\
History of Spain passing mouth to mouth.
\annotation{Drop the bass}

\annotation{Spoken}\\
9. %
Petra, Jordan
\annotation{Sung}\\
Carved in the rock where the desert winds blow,\\
An ancient kingdom from a long time ago.\\
Gotta see it now while the borders stay clear,\\
A world heritage that we hold so dear.

\annotation{Spoken}\\
8. %
Outback, Australia
\annotation{Sung}\\
Arid interior, wild and deep,\\
Where the survival documentaries sleep.\\
North to South or East to West,\\
Uluru’s Red Rock puts the soul to the test.

\annotation{Spoken}\\
7. %
Yellowstone National Park, USA
\annotation{Sung}\\
Oldest in the world, where the wild things grow,\\
Geological wonders and a winter snow.\\
Childhood dreams from a library book,\\
Grizzly bears waiting if you take a look.

\annotation{Spoken}\\
6. %
Patagonia, Argentina \& Chile
\annotation{Sung}\\
Southern Andes where the flat steppes meet,\\
Borders of the ocean where the currents compete.\\
By land through the woods or a cruise through the cold,\\
Glacier landscapes waiting to unfold.

\annotation{Spoken}\\
5. %
Taj Mahal, India
\annotation{Sung}\\
White marble tribute to a love so deep,\\
In the Golden Triangle where the ancient kings sleep.\\
Agra, Jaipur, Delhi in line,\\
An architectural wonder frozen in time.

\annotation{Spoken}\\
4. %
Angkor Wat, Cambodia
\annotation{Sung}\\
The biggest temple that the world has seen,\\
Reclaimed by the roots and the jungle green.\\
Stone-carved faces watching through the trees,\\
Moving like a ghost on a tropical breeze.

\annotation{Spoken}\\
3. %
Trans-Siberian Railway, Russia
\annotation{Sung}\\
Nine-thousand kilometers, Europe to the East,\\
A seven-day journey on a steel-wheeled beast.\\
From Moscow’s lights to the Japan Sea coast,\\
The longest active rhythm that the world can boast.

\annotation{Spoken}\\
2. %
Victoria Falls, Zambia \& Zimbabwe
\annotation{Sung}\\
The smoke that thunders, the water that falls,\\
Africa’s heartbeat echoing through the walls.\\
Zambezi river rushing down to the sea,\\
A misty white column of pure energy.

\annotation{Build-up: Snare roll intensifies}
\annotation{Spoken}\\
1. %
Caño Cristales, Colombia
\annotation{Sung}\\
Liquid rainbow in the ancient stone,\\
A crystal-clear river with a style of its own.\\
Red aquatic plants blooming under the sun,\\
Take a boat, take a flight, now the journey is done.
\annotation{The Drop}

\annotation{The Break: Honorable Mentions}
\annotation{The beat strips back to a minimal, pulsing synth loop as these names are quickly chanted in rhythm}

Cliffs of Dover (UK)\\
Bora Bora (Polynesia)\\
Galápagos (Ecuador)\\
Hawái (USA)\\
Lake Baikal (Russia)\\
Plitvice Lakes (Croatia)\\
Legoland (Denmark)\\
Machu Picchu (Peru)\\
Pyramids of Giza (Egypt)\\
Serengeti (Tanzania \& Kenya)...

\songpart{Outro: Fade out with the kick drum}\end{lyrics}


\songtitle{Bogotá}

\tag*{2026}\tag{ai-based}\tag{cumbia}

\tag*{1989}\tag{original-seed}

\subsection{Notes}
Chorus by Carlos Thompson (1989), verses by Gemini based on a prompt describing the city of Bogotá, Colombia, and its unique blend of natural beauty, urban chaos, and cultural vibrancy. %
The song captures the city's contrasts, from its misty mountain peaks to its bustling streets, and the complex emotions it evokes in those who live there.

\subsection{Style}
With a cumbia galloping offbeat ("pum, chu-chucu-chu, pum, chu-chucu-chu") marked by a modern percussion set, use an acoustic guitar to keep a C-G-Am-F progression, then introduce a tritone to shift into a C-Am-F-G.

\begin{lyrics}
\songpart{verse 1}
The high savanna cradles her, a valley in the clouds,\\
Where chilly mist at morning wraps the rushing, endless crowds.\\
The brick-red towers pierce the gray, beneath Monserrate’s peak,\\
Where modern dreams and ancient stones in quiet shadows speak.

\songpart{verse 2}
A warmth is found in steaming cups of tinto on the street,\\
And heavy cloaks of ruanas guard against the mountain sleet.\\
For those who learn to love her ways, a magic will unfold,\\
A city built on mountain peaks, baptized in rain and gold.

\songpart{chorus}
Ciudad, monstruo ciudad\\
Selva cemento, ciudad natal.\\ 
Ciudad monstruo ciudad Bogotá 

\songpart{verse 3}
Yet patience is a virtue sold upon the fractured street,\\
Where TransMilenio red buses groan in midday heat.\\
The rain will fall without a word, a sudden, blinding sea,\\
While traffic stands in frozen rows for an eternity.

\songpart{verse 4}
The neon signs flash warning lights when darkness settles down,\\
A gritty edge that tests the soul within this sprawling town.\\
You'll curse the cold, you'll curse the fumes, the chaos, and the grind,\\
A restless, unforgiving pace that leaves the weak behind.

\songpart{chorus}
Ciudad, monstruo ciudad\\
Selva cemento, ciudad natal.\\ 
Ciudad monstruo ciudad Bogotá

\songpart{verse 5}
But oh, the epic scale of life where Andes touch the sky!\\
Where cycling seas on Ciclovía watch the world go by.\\
From high Bolívar’s historic square to graffiti-painted walls,\\
A fierce, unyielding energy through every alley calls.

\songpart{verse 6}
She breathes in salsa, rock, and rain—a wild, majestic beast,\\
Where every sunset paints the peaks for an immortal feast.\\
An empire in the emerald hills, both beautiful and grand,\\
The beating, chaotic, proud heart of a surreal land.

\songpart{chorus}
Ciudad, monstruo ciudad\\
Selva cemento, ciudad natal.\\ 
Ciudad monstruo ciudad Bogotá
\end{lyrics}


\songtitle{Le cœur léger}

\tag*{2026}\tag{folk-ballad}\tag{french-lyrics}

\begin{notes}
\end{notes}

\begin{style}
Tarantella-leaning instrumental in 6/8 with a fast, galloping pulse; opening hits hard with dungchen, piwang, marimba, fiddle, male and female gaita, double bass, and timbales, then verses ease into a calmer bass-led drive, Pre-chorus swells with layered strings and winds into a bright crescendo, chorus opens into an instrumental break with dancing lead motifs and percussive lift, Outro thins instruments one by one until only timbales remain, crisp and intimate, mix bright, energetic, and folk-cinematic, marimba, ballad
\end{style}

\begin{lyrics}
\songpart{Verse 1}
Dans les vignes du matin\\
Le soleil mord le chemin\\
On ouvre la porte en grand\\
Et le rire sort avant

Le pain chaud sur la table\\
Le vin rouge, incassable\\
Tes yeux brillent, mon ami\\
La fête commence ici

\songpart{Pre-Chorus}
Allez, viens danser\\
La nuit va passer\\
Le cœur léger\\
On va trinquer

\songpart{Chorus}
Ô Languedoc, on vit fort
Ô Languedoc, on rit encore\\
Le verre levé, le cœur d’accord
Ô Languedoc, on vit fort

(Ô Languedoc)
(Ô Languedoc)

\songpart{Verse 2}
Sur la place, les guitares\\
Font tourner les regards\\
Les épaules contre les épaules\\
Et la joie devient parole

La bouteille passe de main\\
Comme un secret, comme un refrain\\
On chante jusqu’au bout du vent\\
Et demain n’existe pas maintenant

\songpart{Pre-Chorus}
Allez, viens danser\\
La nuit va passer\\
Le cœur léger\\
On va trinquer

\songpart{Chorus}
Ô Languedoc, on vit fort
Ô Languedoc, on rit encore\\
Le verre levé, le cœur d’accord
Ô Languedoc, on vit fort

(Ô Languedoc)
(Ô Languedoc)

\songpart{Bridge}
Si le monde est lourd parfois\\
Ici, on le laisse là\\
Une gorgée, un sourire\\
Et tout recommence à vivre

\songpart{Final Chorus}
Ô Languedoc, on vit fort
Ô Languedoc, on rit encore\\
Le verre levé, le cœur d’accord
Ô Languedoc, on vit fort

(Ô Languedoc)
(Ô Languedoc)
\end{lyrics}


\songtitle{Canta la Ruta}

\tag*{2026}\tag{folk-ballad}

\begin{notes}
\end{notes}

\begin{style}
Fast-paced tarantella-inspired folk ballad in 6/8 with rhythmic drive from dungchen, piwang, marimba, fiddle, male and female gaita, double bass, and timbales; opening hits with a sharp instrumental attack, then calmer verse passages anchored by low bass, pre-chorus swells into a rising crescendo, chorus opens into a bright instrumental break, and the outro strips instruments away one by one until only timbales remain, Wide, earthy, and kinetic mix, marimba, ballad
\end{style}

\begin{lyrics}
\songpart{Intro}
\annotation{Spoken}\\
Listen to the rhythm that moves your feet, let the night carry no trouble.\\
Come, move your body, until the morning light finds us.\\
This is the dance of the desert wind where we forget all sorrow.

\songpart{Verse 1}
\annotation{nonsense, but rhythmic}\\
Beat the drum, shake the ground —\\
boom and crash, sing loud —\\
badabum, spin around —\\
dubi ba dubi ba sing!

\songpart{Pre-chorus}
Move your body, till the sun rises —\\
rapakapun, dance on the sand —\\
where the wind tells no lies —\\
rapakapun, kick the ground!

\songpart{Chorus}
Badabum bin — dubidum shun shun —\\
badabum bin — dubi ba dubi ba sang —\\
badabum bin — rapakapun kiki botang!

\songpart{Verse 2}
\annotation{nonsense, but rhythmic}\\
Beat the drum, shake the ground —\\
boom and crash, sing loud —\\
badabum, spin around —\\
dubi ba dubi ba sing!

\songpart{Pre-chorus}
Move your body, till the sun rises —\\
rapakapun, dance on the sand —\\
where the wind tells no lies —\\
rapakapun, kick the ground!

\songpart{Chorus}
Badabum bin — dubidum shun shun —\\
badabum bin — dubi ba dubi ba sang —\\
badabum bin — rapakapun kiki botang!

\songpart{Bridge}
\annotation{goth metal style}\\
Crush me with your power —\\
kiki botang —\\
crushing sorrow —\\
rapakapun, now!

\songpart{Chorus}
Badabum bin — dubidum shun shun —\\
badabum bin — dubi ba dubi ba sang —\\
badabum bin — rapakapun kiki botang!

\songpart{Outro}
\annotation{spoken}\\
Come, move your body — until the morning light finds us.
\end{lyrics}


\songtitle{Take Me South}

\tag*{2023}
\begin{notes}
Lyrics by SUNO (2026) based on a 2023 conversation between Carlos Thompson and ChatGPT
\end{notes}
\begin{style}
Doo-wop oldies with bouncing upright bass, clipped snare on the backbeat, handclaps, and brushed tambourine; verse glides in with close trio harmonies, pre-chorus lifts with stacked oohs, chorus opens wide with gang vocals on the hook and a sax fill between lines, Add vinyl crackle, a bell swell into the first chorus, and a final key-change lift, Warm, glossy, close-harmony mix with a nostalgic radio sheen, beautiful
\end{style}
\begin{lyrics}
\songpart{Verse 1}
Bennington on my map\\
Stone and stories stacked\\
Battle Monument in the sky\\
Museum doors to the past

Music drifting downtown\\
Book fair spin me round\\
I can hear that old town call\\
Baby, let’s not waste the day

\songpart{Pre-Chorus}
Ooh, the road is open\\
Ooh, and I’m wide awake\\
Pack a little wonder\\
Come on, let’s escape

\songpart{Chorus}
Take me to Southern Vermont\\
Take me where the pines lean long\\
Take me to Southern Vermont\\
I wanna sing along
(woo-oo, woo-oo)\\
Take me to Southern Vermont\\
That’s where my heart belongs

\songpart{Verse 2}
Manchester in the air\\
Hildene standing there\\
Frost’s old house, a weathered line\\
Art and song in the square

Woodstock by the green\\
Farmers’ market scene\\
Old hotel and orchard light\\
Baby, make it feel so right

\songpart{Pre-Chorus}
Ooh, the hills are calling\\
Ooh, in a sweet, slow way\\
Fill the car with sunlight\\
Come on, let’s run away

\songpart{Chorus}
Take me to Southern Vermont\\
Take me where the pines lean long\\
Take me to Southern Vermont\\
I wanna sing along
(woo-oo, woo-oo)\\
Take me to Southern Vermont\\
That’s where my heart belongs

\songpart{Bridge}
Green Mountains, hold me high\\
Killington in the sky\\
Stratton with a view so clean\\
Lake Champlain, blue and wide

Northeast Kingdom roads\\
State parks, wooded roads\\
Every mile says stay a while\\
Baby, let me go

\songpart{Final Chorus}
Take me to Southern Vermont\\
Take me where the pines lean long\\
Take me to Southern Vermont\\
I wanna sing along
(woo-oo, woo-oo)\\
Take me to Southern Vermont\\
That’s where my heart belongs\\
Take me to Southern Vermont\\
Baby, now you know I’m gone
\end{lyrics}


\songtitle{Closer to the Stars}

\tag*{2026}\tag{sea-shanty}\tag{original-source}
\begin{notes}
Lyrics by Suno based on a specific prompt by Carlos Thompson.
\end{notes}
\begin{style}
traditional maritime sea shanty recast as highland folk, brisk swinging 6/8 work-song groove, accordion and concertina, fiddle, stomping frame drums, acoustic guitar, call-and-response chorus, weathered hearty male lead with rough crew harmonies, communal pub energy, clear acoustic mix with room ambience, dynamic verse-to-chorus lifts
\end{style}
\begin{lyrics}
\songpart{Verse 1}
We keep our boots where the thin wind bites,\\
Eight thousand five hundred feet,\\
No gull on the fence, no salt on the pines,\\
But the mountain keeps time beneath.\\
The shoreline's a two-day wagon run,\\
Past wheat fields, stone, and snow,\\
Still every mile in the highland sun\\
Makes the old sea fever grow.

\songpart{Chorus}
Heave, lads, heave, though the ocean's far,\\
Pull by the rope of the morning star,\\
Two days' drive till the breakers roar,\\
But we're twenty-six hundred meters closer to the stars than before!

\songpart{Verse 2}
Our compass swings when the thunder rolls,\\
Our kettle sings on the flame,\\
We trade blue water for granite roads,\\
Yet dream of the harbor by name.\\
A tin cup rings and the campfire leans,\\
The ridgeline cuts the blue,\\
We may be landlocked between the peaks,\\
But the sky comes sailing through.

\songpart{Chorus}
Heave, lads, heave, though the ocean's far,\\
Pull by the rope of the morning star,\\
Two days' drive till the breakers roar,\\
But we're twenty-six hundred meters closer to the stars than before!

\songpart{Bridge}
When the valley fog rolls under our feet,\\
And the moon hangs cold and bright,\\
We raise one cheer for the distant tide\\
And one for the mountain night.\\
No anchor down, no harbor wall,\\
No deck beneath our shoes,\\
But the highest crew in the world tonight\\
Has nothing left to lose.

\songpart{Final Chorus}
Heave, lads, heave, let the whole ridge ring,\\
Sing to the hawk and the wandering wind,\\
Two days' drive till the breakers roar,\\
But we're twenty-six hundred meters closer to the stars than before!\\
Closer to the stars than before!\\
Closer to the stars than before!

\songpart{Outro}
Eight thousand five hundred, steady and high,\\
With the sea two days away,\\
We sail on the road beneath the open sky,\\
And climb with the break of day.
\end{lyrics}

\songtitle{Indiana Highway}

\tag*{2026}\tag{road-folk}
\begin{notes}
Lyrics by SUNO, based on a ChatGPT response to an inquiry by Carlos Thompson about road-trip landmarks in northern Indiana -- a cheerful travelogue through five stops: French Lick Springs, Conner Prairie, the Indiana Dunes, New Harmony, and South Bend. Retrieved from rataflechera's Suno page (V6, 15 September 2026).
\end{notes}
\begin{style}
Road folk, warm acoustic guitar strumming, bright candy pop textures, bubbly synthesizer accents, driving uptempo rhythm, sweet melodic female vocals, clean modern production
\end{style}
\begin{lyrics}
\songpart{Verse 1}
Rollin' down the Interstate of Eighty\\
Windows down, the morning air is sweet\\
Got a map of places historic and shady\\
Indiana's secrets waiting for my feet\\
First we stop where the mineral water bubbles\\
French Lick Springs, historic and grand\\
Wash away all of our highway troubles\\
In the golden valleys of this quiet land

\songpart{Chorus}
Oh, we're travelers on the open road\\
From the pioneer fields to the sandy shore\\
Every little town has a story told\\
Indiana, we are knocking at your door!

\songpart{Verse 2}
Over in Fishers, the past is alive\\
Conner Prairie under big blue skies\\
A working farm where the old ways thrive\\
And history dances right before your eyes\\
Then we feel the wind from the lakeside blowing\\
Indiana Dunes rising tall and high\\
The tallest sand hills, shifting and growing\\
Sinking our toes where the wild birds fly

\songpart{Chorus}
Oh, we're travelers on the open road\\
From the pioneer fields to the sandy shore\\
Every little town has a story told\\
Indiana, we are knocking at your door!

\songpart{Bridge}
Down in New Harmony, the dreamers built a town\\
Hoping for a perfect world to hold\\
While in South Bend, the chrome is shining down\\
Studebaker engines painted bright and bold

\songpart{Outro}
Five sweet stops on a northern line\\
Indiana highway, feeling so fine!
\annotation{Instrumental fade out}
\end{lyrics}

\end{multicols*}

\part{Social Commentary}

\chapter{Satire and mundane absurdity}
\albumcover{covers/internal/satire-and-mundane-absurdity}
The small absurdities that don't deserve real anger but get a song anyway: a riot over an operating manual for a light bulb, a menu nobody can agree how to taste, real time itself treated as a corporate product to be optimized.
\clearpage\begin{multicols*}{4}
\songtitle{Operating Manual (For Changing a Light Bulb)}

\tag*{2024}\tag{musical-theater}

\begin{notes}
\end{notes}
\begin{style}
dark comedy march, musical theatre, male tenor lead, satirical chorus, 96 BPM, tuba bassline, snare drum parade, brass fanfares, accordion stabs, celesta counterline, spoken bridge, call and response, mono vocal compression, tape saturation, spring reverb, 1960s stage mix, triumphant break, defeated outro
\end{style}
\begin{lyrics}
\songpart{Verse 1}
Well I woke up this morning, the bulb was burnt out\\
Eight feet up on the ceiling, I had my doubts\\
So I consulted the manual, section one through seven\\
‘Cause God forbid I just twist it like a normal heathen

\songpart{Pre-Chorus}
It said “Operator” — that’s me, I guess\\
With my personal protective equipment and existential stress

\songpart{Chorus}
This is the Operating Manual\\
For the removal and replacement of an incandescent illumination device!\\
Disengage the power source, verify de-energization\\
Then ascend the ladder with three points of contact, that’s nice!

\songpart{Verse 2}
Step one: don the gloves, non-slip shoes on my feet\\
Safety goggles optional, but liability’s a beast\\
Position the ladder, make sure it’s stable and true\\
‘Cause if I fall and break my neck, the manual warned me not to

\songpart{Verse 3}
Grasp the bulb with dominant hand — counter-clockwise rotation\\
Do not apply excessive force or cause bulb fragmentation\\
Lower gently, dispose properly, do not throw it in the trash\\
Or the Environmental Protection Agency will whip your ass

\songpart{Chorus}
This is the Operating Manual\\
For the removal and replacement of an incandescent illumination device!\\
Align the threads, clockwise motion, firm but gentle pressure\\
Until you feel resistance — that’s the moment of pure pleasure!

\songpart{Bridge}
\annotation{spoken over slow ominous music, then building}\\
Do not operate while under the influence…\\
Of alcohol, controlled substances, or crippling existential dread…\\
The Operator shall bear full responsibility\\
For any personal injury, property damage, or mild inconvenience…

\annotation{Final Chorus}\\
\annotation{louder, triumphant}\\
Now reactivate the power source, observe illumination!\\
If it lights up, congratulations — you followed the procedure!\\
If it doesn’t… consult Appendix C, subsection 4.2\\
And start the whole damn process over again!

\songpart{Outro}
\annotation{slow, defeated}\\
Eight feet… one bulb…\\
Why did we do this to ourselves?\\
Operating… Manual…

\annotation{fade out with the sound of a circuit breaker and distant sobbing}
\end{lyrics}


\songtitle{Buba Kiki Riot}

\tag*{2026}\tag{alt-rock}

\begin{notes}
\end{notes}
\begin{style}
128 BPM, rock base, tumbao intro, electric bass palm mute, distorted bass grit, piano percussion ostinato, overdriven kit, tom-heavy chorus, metal power chords, minor key chorus, plate reverb, tape saturation, stereo delay throws, 90s alt-rock mix, half-time drop, anthemic unease
\end{style}
\begin{lyrics}
\songpart{Verse 1}
You came in sideways\\
With that red-lip grin\\
Said, “Buba, keep up”\\
So I jumped right in

Your boots hit the floorboards\\
Like a warning sign\\
I laughed at the trouble\\
Then you drew the line

\songpart{Pre-Chorus}
Now my pulse goes faster\\
When you call my name\\
Got that sweet disaster\\
Burning through my brain

\songpart{Chorus}
Kiki!\\
Break the door\\
Kiki!\\
Want it more\\
Buba, buba\\
We want war\\
Kiki!\\
Hit the floor

\songpart{Verse 2}
You stole my black jacket\\
Left your scent on me\\
A dent in the mirror\\
Where my old self used to be

I tried to stay cool now\\
But you know I fold\\
One look, then we’re falling\\
Off the edge we hold

\songpart{Pre-Chorus}
Now my pulse goes faster\\
When you call my name\\
Got that sweet disaster\\
Burning through my brain

\songpart{Chorus}
Kiki!\\
Break the door\\
Kiki!\\
Want it more\\
Buba, buba\\
We want war\\
Kiki!\\
Hit the floor

\songpart{Bridge}
Say it louder\\
Make it sting\\
Hands up high now\\
Let it ring\\
Buba on the wire\\
Kiki in my veins\\
We came for the chaos\\
We came for the flames

\songpart{Final Chorus}
Kiki!\\
Break the door\\
Kiki!\\
Want it more\\
Buba, buba\\
We want war\\
Kiki!\\
Hit the floor

Kiki!\\
Back for more\\
Buba, buba\\
Feel the roar\\
Kiki!\\
Hold the core\\
Kiki!\\
Hit the floor
\end{lyrics}


\songtitle{Chiang Mai to Bangkok}

\tag{k-pop}

\begin{notes}
A distinctive arrangement that supports the song's storytelling and emotional arc. %
The lyrics begin with “Thailand in my carry-on”, giving a quick sense of the song’s tone.
\end{notes}

\begin{style}
K-pop dance-pop with bright syncopated drums, glossy synth plucks, bouncing bass, and a four-on-the-floor lift in the chorus; verses stay percussive with clipped lead vocals and playful chant replies, pre-chorus opens with stacked harmonies and a rising filter sweep, chorus bursts into group hooks and handclaps. %
Whispered ad-libs, pitch-shifted photo-flash chops, reverse swell transitions, crisp and neon-bright mix., k-pop, warm, world
\end{style}

\begin{lyrics}

\songpart{Verse 1}
Thailand in my carry-on\\
Ugenito went sideways, I kept on\\
Chiang Mai lights, Bangkok heat\\
Every bad plan learned to dance on its feet

Photos turn to little scenes\\
Albums fold up bigger meanings\\
I see the map, then I see me\\
One bad joke and a system key

\songpart{Pre-Chorus}
Claim and crack\\
Mask and fact\\
We point at the mess\\
Then trace it back

Wrong advice\\
Right in sight\\
Tune it up\\
Don’t set it tight

\songpart{Chorus}
Chiang Mai to Bangkok\\
(Chiang Mai to Bangkok)\\
I keep connecting dots\\
I keep connecting dots

From Florida sand to airport lines\\
From family noise to changing signs\\
I keep connecting dots\\
I keep connecting dots

\songpart{Verse 2}
IMO in the Chiang Mai room\\
Filing the world in a small hotel gloom\\
Transit rules and temple bells\\
Popular faith and the stories it tells

Curfew became the task\\
A joke that turned into a mask\\
Overthinker, underachiever, same\\
Trying to grow up without a game

\songpart{Pre-Chorus}
Claim and crack\\
Mask and fact\\
We point at the mess\\
Then trace it back

Wrong advice\\
Right in sight\\
Tune it up\\
Don’t set it tight

\songpart{Chorus}
Chiang Mai to Bangkok\\
(Chiang Mai to Bangkok)\\
I keep connecting dots\\
I keep connecting dots

From Florida sand to airport lines\\
From family noise to changing signs\\
I keep connecting dots\\
I keep connecting dots

\songpart{Bridge}
Thomas cut warm at landing\\
Before the talk, before the standing\\
Gabi saw it clear\\
No need to force the fear

Sylvia says, “I’m not in love”\\
But Gabi reads what’s up above\\
Not the line, but what it hides\\
The old self living in the eyes

\songpart{Final Chorus}
Chiang Mai to Bangkok\\
(Chiang Mai to Bangkok)\\
I keep connecting dots\\
I keep connecting dots

From Thomas to the turning tide\\
From late two-zero to mid two-one wide\\
I keep connecting dots\\
I keep connecting dots

(Connecting dots)\\
(Connecting dots)

\end{lyrics}


\songtitle{Sabor en la Mesa}

\tag*{2025}\tag{eurodance}\tag{spanish-lyrics}
\begin{notes}
Lyrics by SUNO based on a conversation between Carlos Thompson and ChatGPT about a fictional restaurant menu.
\end{notes}
\begin{style}
Eurodance with a driving four-on-the-floor kick, bouncing offbeat synth bass, bright claps, and a glossy 128 BPM pulse; verse rides playful talk-sung lines over clipped piano stabs, pre-chorus strips to kick, bass, and rising vocal stacks, chorus hits with a huge chant and handclaps, Layered lead doubles, gang ad-libs, delay throws, risers, sparkly synth fills, and a crisp club-bright mix, eurodance, fresh, simple, soft, punta, slow, rich, light
\end{style}
\begin{lyrics}
\songpart{Verse 1}
Entradas pa' empezar\\
Empanadas, crunch, crunch\\
Carne y papa en el centro\\
Champiñones al ajillo, yum\\
Chorizo con arepa\\
Chicharrón, qué sabor\\
Makis en la mesa\\
Coctel de camarones, mejor

\songpart{Pre-Chorus}
Sube, sube, ven\\
Que aquí se viene a comer\\
Sabor de verdad\\
Y no se va a detener

\songpart{Chorus}
Pide en la mesa\\
Pide en la mesa\\
Todo te espera\\
Todo te espera\\
Pide en la mesa\\
Pide en la mesa\\
Sabe tan rico\\
Sabe tan rico

\songpart{Verse 2}
Changua de Bogotá\\
Ajiaco santafereño\\
Mondongo calentito\\
Con su abrazo de fuego\\
Frijolada con arroz\\
Y plátano al lado\\
Bandeja paisa fuerte\\
Con huevo y aguacate alzado

\songpart{Pre-Chorus}
Sube, sube, ven\\
Que aquí se viene a comer\\
Sabor de verdad\\
Y no se va a detener

\songpart{Chorus}
Pide en la mesa\\
Pide en la mesa\\
Todo te espera\\
Todo te espera\\
Pide en la mesa\\
Pide en la mesa\\
Sabe tan rico\\
Sabe tan rico

\songpart{Bridge}
Baby beef a la parrilla\\
Costillas BBQ\\
Punta de anca jugosa\\
Y la pechuga también, uh\\
Trucha al ajillo brilla\\
Con mantequilla y sal\\
Espagueti boloñesa\\
Versión vegetal

\songpart{Verse 3}
Para los peques, bebé\\
Hamburguesa infantil\\
Postre del día al centro\\
Y se deja sentir\\
Jugos naturales\\
Café colombiano\\
Té bien caliente\\
Y algo más si lo pides, hermano

\songpart{Chorus}
Pide en la mesa\\
Pide en la mesa\\
Todo te espera\\
Todo te espera\\
Pide en la mesa\\
Pide en la mesa\\
Sabe tan rico\\
Sabe tan rico

\songpart{Outro}
Entradas, sopas, típicos\\
Carnes, aves, pescado\\
Vegetariano, infantil\\
Postres y bebidas, todo a tu lado\\
Pide en la mesa\\
Pide en la mesa
\end{lyrics}


\songtitle{Taste the Menu}

\tag*{2025}\tag{eurodance}
\begin{notes}
Lyrics by SUNO based on a conversation between Carlos Thompson and ChatGPT about a fictional restaurant menu.
\end{notes}
\begin{style}
Eurodance with four-on-the-floor kick, bouncy offbeat bass, bright synth stabs, and clapping percussion; verse rides a playful spoken-sung groove, pre-chorus strips to bass and handclaps, chorus opens with a stacked vocal hook and octave synth lead, Add risers into each lift, reverse cymbal sweeps before drops, and sparkling bell fills between menu lines, Glossy, punchy mix with wide harmonies and party-ready polish, slow, simple, rich, punta, soft, light, fresh, eurodance
\end{style}
\begin{lyrics}
\songpart{Verse 1}
Entradas on the line tonight\\
Empanadas, crispy and right\\
Champiñones, garlic kiss\\
Chorizo on the arepa bliss

Chicharrón, crackle and crunch\\
Makis fresh for the lunch\\
Coctel de camarones, chill\\
Tangy sauce and the perfect thrill

\songpart{Pre-Chorus}
Come a little closer now\\
Let the whole place show you how\\
From the first bite to the last\\
Every flavor moves fast

\songpart{Chorus}
Taste the menu, taste the menu\\
One more bite and you’ll want it too\\
Taste the menu, taste the menu\\
Colombia on a plate for you\\
Taste the menu, taste the menu\\
Hot and cold, sweet and savory too\\
Taste the menu, taste the menu\\
Every dish says come on through

\songpart{Verse 2}
Changua warm in the morning light\\
Ajiaco with the corn shining bright\\
Mondongo rich and full of care\\
Frijolada with rice to spare

Bandeja paisa, big and proud\\
Arroz atollado, draws a crowd\\
Baby beef and ribs on glaze\\
Pechuga a la plancha, clean and laid

\songpart{Pre-Chorus}
Turn it up and let it roll\\
Every plate goes straight to soul\\
From the grill to something sweet\\
Let the whole table meet

\songpart{Chorus}
Taste the menu, taste the menu\\
One more bite and you’ll want it too\\
Taste the menu, taste the menu\\
Colombia on a plate for you\\
Taste the menu, taste the menu\\
Hot and cold, sweet and savory too\\
Taste the menu, taste the menu\\
Every dish says come on through

\songpart{Bridge}
Trucha al ajillo, soft and gold\\
Vegetarian pasta, bold\\
Hamburguesa infantil, easy smile\\
Dessert of the day stays a while

Juices, coffee, tea, and more\\
Ask for extras at the door\\
Every craving, every mood\\
Find your favorite, feel the good

\songpart{Chorus}
Taste the menu, taste the menu\\
One more bite and you’ll want it too\\
Taste the menu, taste the menu\\
Colombia on a plate for you\\
Taste the menu, taste the menu\\
Hot and cold, sweet and savory too\\
Taste the menu, taste the menu\\
Every dish says come on through
\end{lyrics}


\songtitle{Tiempo Real}

\tag*{2026}\tag{corporate-music}\tag{original-lyrics}\tag{spanish-lyrics}
\begin{notes}
Written by Carlos Thompson as an elevator pitch for a proposal.
\end{notes}
\begin{style}
Corporate background music with spoken word narration, A clean electric guitar plays a repetitive, palm-muted eighth-note pattern on a single note, providing a steady rhythmic pulse, A bright, acoustic piano enters with syncopated, high-register chords, The percussion consists of a light, electronic kick drum on every beat and a crisp clap on beats two and four, A subtle synth pad provides harmonic filling in the background, The tempo is 120 BPM in the key of C major, The male narration is dry with minimal reverb, delivered in a professional, informative tone
\end{style}
\begin{lyrics}
\songpart{Intro}
\annotation{palm-muted electric guitar eighth notes, electronic kick drum}

\songpart{Verse 1}
\annotation{male narration enters}\\
Cada ticket que llega a soporte toma entre dos y tres horas antes de tener una primera respuesta. \annotation{piano chords enter} Eso no es porque el problema sea complejo, es porque alguien tiene que leerlo, preguntarle a otros, entender el contexto y decidir qué hacer. Este tiempo es puro overhead. Y cuando hay reprocesos planificados, que son el caso más frecuente, el consultor invierte minutos reales, ajusta fechas, lanza el proceso, valida el cargue, pero el ticket puede demorar días enteros esperando en cola. El cliente ve días de espera, el consultor vio veinte minutos de trabajo. Además, todos los días cada consultor debe reportar manualmente el estado de sus tickets en Mantis BT. Todos los días.

\songpart{Verse 2}
Un sistema de IA agentiva elimina los tres cuellos de botella. Clasifica y enruta el ticket en segundos, gestiona la cola de reprocesos automáticamente priorizando por ventanas de tiempo disponibles y actualiza el estado del ticket sin que el consultor tenga que escribir nada al final del día.

\songpart{Outro}
El resultado: primera respuesta en minutos, no horas; reprocesos planificados sin cola de espera artificial; y el equipo libera tiempo real que hoy va en coordinación y reportes. Con ciento ochenta y siete tickets resueltos en el historial, ya tenemos los datos para entrenar el primer modelo. ¿Damos treinta minutos esta semana para revisar los números juntos?
\end{lyrics}


\end{multicols*}

\chapter{Modern loneliness and social critique}
\albumcover{covers/internal/modern-loneliness-and-social-critique}
Isolation with its systems named and blamed: postcapitalist loneliness, a toxic ghost that won't leave, and -- joining them now -- an industrial-metal indictment of a world that processes people the way it processes everything else. Less confession than diagnosis.
\clearpage\begin{multicols*}{4}
\songtitle{Cut Me Loose}

\tag*{2013}\tag{nu-metal}
\begin{notes}
Lyrics by SUNO based on the text of a meme written by Carlos Thompson in 2013 about online self-help relationship advice.
\end{notes}
\begin{style}
Nu-metal anthem with stomping half-time riffs, syncopated palm-muted guitars, thundering kick-snare, and a snarling, shout-along chorus; verses stay tight and tense with close-mic rap-rock delivery, pre-chorus opens into gang shouts and rising toms, chorus slams with stacked doubles and brutal gang vocals, Add reversed cymbal swells, air-raid sirens before the drop, and a final breakdown with halved drums and distorted chants, Gritty, wide, chest-rattling mix, inspiring, metal
\end{style}
\begin{lyrics}
\songpart{Verse 1}
You said, “cut the toxic out”\\
So I took a good hard look\\
Found my face in the mirror\\
And that hit like a steelbook

I was the laugh that cut too deep\\
The friend who never heard “enough”\\
Left a trail of ash and excuses\\
Now the room feels rough

\songpart{Pre-Chorus}
I kept calling it “just being real”\\
You kept calling it “too much”\\
Now your name is gone from my phone\\
And I’m the one you can’t touch

\songpart{Chorus}
Cut me loose\\
Cut me loose\\
I was poison in the room\\
Cut me loose\\
Cut me loose\\
I’m the toxic one you knew
(you knew)\\
Cut me loose\\
Cut me loose\\
Yeah, I made the whole thing sick\\
Cut me loose\\
Cut me loose\\
Before I drag you back in it

\songpart{Verse 2}
I heard my own words in your mouth\\
Cold and sharp, same old knife\\
Every “joke” was a bruise in disguise\\
Every favor had a price

Now the chairs are all empty\\
And the group chat stays dry\\
No one owes me another mercy\\
No one has to lie

\songpart{Pre-Chorus}
I kept calling it “just being honest”\\
You kept paying the cost\\
Now the door stays locked for a reason\\
And I’m the one left lost

\songpart{Chorus}
Cut me loose\\
Cut me loose\\
I was poison in the room\\
Cut me loose\\
Cut me loose\\
I’m the toxic one you knew
(you knew)\\
Cut me loose\\
Cut me loose\\
Yeah, I made the whole thing sick\\
Cut me loose\\
Cut me loose\\
Before I drag you back in it

\songpart{Bridge}
Take the blade out my hands\\
Take the smoke from my breath\\
I don’t want to be the wrecking ball\\
That leaves a house in death

I can hear it now\\
Every door I burned\\
I wanted saving, wanted loyalty\\
But I never learned

\songpart{Breakdown}
Cut\\
Me\\
Loose

I’m the match\\
I’m the flare\\
I’m the one who ruined it there\\
Cut me loose\\
Cut me loose\\
Let me disappear

\songpart{Final Chorus}
Cut me loose\\
Cut me loose\\
I was poison in the room\\
Cut me loose\\
Cut me loose\\
I’m the toxic one you knew
(you knew)\\
Cut me loose\\
Cut me loose\\
Yeah, I made the whole thing sick\\
Cut me loose\\
Cut me loose\\
Now I let you live through it
\end{lyrics}


\songtitle{The Toxic Ghost}

\tag*{2013}
\begin{notes}
Lyrics by SUNO based on the text of a meme written by Carlos Thompson in 2013 about online self-help relationship advice.
\end{notes}
\begin{style}
nu-metal, aggressive, down-tuned 7-string guitars, heavy distorted bass, dynamic drumming with heavy kicks, angst-filled vocals, mix of rap-style verses and screaming metal chorus
\end{style}
\begin{lyrics}
\songpart{Intro}
(Low-tuned, chugging guitar riff)
(Aggressive, heavy drums)

\songpart{Verse 1}
They told me cut the dead weight, trim the fat\\
Get rid of the leeches, don’t look back\\
Filter the circle, keep the standard high\\
Burn the bridges, watch the shadows die\\
I took the advice, I played the part\\
I tore the rot right out of my heart

\songpart{Pre-Chorus}
But the silence is growing, the phone doesn't ring\\
I’m waiting for the poison that the future brings\\
But the room is empty, the walls are closing in\\
Where did I lose? Where did I begin?

\songpart{Chorus}
I was the cure, I was the blade\\
I was the one who was never afraid\\
To cut them out, to draw the line\\
I thought the purity would make me shine\\
But look in the mirror, the truth is clear\\
The toxic ghost has been standing here\\
I’m not the judge, I’m the disease\\
They’re all cutting me off, bringing me to my knees!

\songpart{Verse 2}
I pointed the finger, I cast the blame\\
I turned every friendship into a game
"You're not good enough," I used to scream\\
While I was the nightmare ruining the dream\\
They didn't abandon, they didn't betray\\
They just finally found the strength to walk away

\songpart{Bridge}
(Heavy, down-tempo breakdown)\\
I’m the rot! (I'm the rot!)\\
I’m the stain! (I'm the stain!)\\
I’m the reason for all of the pain!\\
I fed on the light, I drank the well dry\\
Now I’m dying alone, watching the sky!

\songpart{Chorus}
I was the cure, I was the blade\\
I was the one who was never afraid\\
To cut them out, to draw the line\\
I thought the purity would make me shine\\
But look in the mirror, the truth is clear\\
The toxic ghost has been standing here\\
I’m not the judge, I’m the disease\\
They’re all cutting me off, bringing me to my knees!

\songpart{Outro}
(Guitar feedback)\\
They were right to leave.\\
They were right to leave.
(Fade out with distorted bass)
\end{lyrics}


\songtitle{Post-Capitalist Lonely}

\tag*{2026}\tag{experimental}\tag*{english-lyrics}
\begin{notes}
Lyrics by SUNO after a detailed prompt by Carlos Thompson.
\end{notes}
\begin{style}
deep bassline, pipa, jangle pop, nu-jazz, electric cello, flanger, drum'n bass, progressive drums, rough male vocals, frail female voice, modern metalcore, improv
\end{style}
\begin{lyrics}
\songpart{Verse 1}
My wallet is a chapel\\
full of old receipts\\
for bread I didn't eat\\
and touch I didn't keep

The app says I am thriving\\
my face says otherwise\\
I rent a room from silence\\
and pay in sleep at night

\songpart{Pre-Chorus}
Every day is measured\\
in taps and tiny fees\\
Even my grief gets sorted\\
by delivery speed

\songpart{Chorus}
Post-capitalist lonely (post-capitalist lonely)\\
I miss the way things lingered\\
Post-capitalist lonely (post-capitalist lonely)\\
Even my name feels rented
(post-capitalist lonely)

\songpart{Verse 2}
The supermarket blooms at dawn\\
with fruit from faraway hands\\
I hold a peach like evidence\\
I cannot understand

A stranger in the elevator\\
smiles with museum eyes\\
We are both so carefully alive\\
we almost apologize

\songpart{Pre-Chorus}
I used to want forever\\
Now I want a chair\\
A cup left on the table\\
And someone knowing I was there

\songpart{Chorus}
Post-capitalist lonely (post-capitalist lonely)\\
I miss the way things lingered\\
Post-capitalist lonely (post-capitalist lonely)\\
Even my name feels rented
(post-capitalist lonely)

\songpart{Bridge}
If I am made of labor\\
why do I feel so spare?\\
A hand is not a service\\
A door is not a prayer

I want a world with leftovers\\
with crumbs still warm in the pan\\
Where being seen is common\\
and not a premium plan

\songpart{Final Chorus}
Post-capitalist lonely (post-capitalist lonely)\\
I miss the way things lingered\\
Post-capitalist lonely (post-capitalist lonely)\\
Even my name feels rented\\
Post-capitalist lonely (post-capitalist lonely)\\
Let me stay in the weather
(post-capitalist lonely)
\end{lyrics}


\songtitle{Signed by Don Addis}

\tag*{2026}\tag{indie-pop}\tag*{english-lyrics}
\begin{notes}
Lyrics by SUNO based on a conversation with ChatGPT.
\end{notes}
\begin{style}
lo-fi indie pop, quirky pop, spoken-word intro, narrative verse, reflective chorus, half-spoken bridge, downbeat word placement, 104 BPM, syncopated piano stabs, muted electric guitar chops, bass clarinet counterline, brushed drum kit, handclap accents, cassette saturation, plate reverb, stereo delay throws, bittersweet nostalgia, searching resolve
\end{style}
\begin{lyrics}
\songpart{Verse 1}
I'm looking for a drawing, black lines on white\\
Two panels, a believer, and a cross held tight\\
He's swinging it around like it's a bat, not a prayer\\
Then somebody grabs it back — now watch him care

Typing it out slow, hoping someone knows
"Do you have the image?" — no, that's why I'm asking, so—

\songpart{Verse 2}
Found a dusty thread, three years old, unsolved
"Tip of my tongue," a mystery unresolved\\
Black lineart, detailed, cartoon style, they said\\
Nobody claimed it, the trail went cold and dead

Reverse image search, static on the screen\\
Somewhere in the scroll there's a comic in between

\songpart{Pre-Chorus}
Give me one more detail, one more clue to chase\\
A forum, a caption, a half-remembered face\\
I keep typing questions like I'm digging through a drawer\\
Just one more search, I've done this before—

\songpart{Chorus}
Found.\\
Signed by Don Addis.\\
Four small words and the whole thing lands.\\
Found.\\
Signed by Don Addis.\\
Turns out somebody's always got the answer in their hands

\songpart{Verse 3}
Turns out he drew for papers, syndicated, sharp\\
Bent Offerings, satire, politics and God\\
St. Petersburg Times, decades at his desk\\
Just a guy who drew the punchline better than the rest

\songpart{Bridge - quieter, half-spoken}
Funny what we lose the name of\\
funny what still finds its way back round\\
a cartoon with no credit\\
floating on a message board\\
until it does

\songpart{Chorus - bigger, full band in}
Found.\\
Signed by Don Addis.\\
Four small words and the whole thing lands.\\
Found.\\
Signed by Don Addis.\\
I wasn't even looking for a happy end

\songpart{Outro}
Black lines on white,\\
a cross, a fight, a name —\\
signed, Don Addis,\\
that's it, that's the whole thing.
\end{lyrics}


\songtitle{Del establecimiento}

\tag*{2023}
\begin{notes}
Lyrics by SUNO based on a ChatGPT conversation regarding a 2023 tweet by Carlos Thompson that read:
\begin{quote}\itshape%
Hay que recordar una cosa. La enfante terrible del establecimiento es *del establecimiento*.
\end{quote}
\end{notes}
\begin{style}
Lo-fi indie pop with quirky off-kilter rhythm, swung drums, muted guitar ticks, and a rubbery bassline; verse stays dry and conversational, pre-chorus thins to finger snaps and a toy-keyboard motif, chorus opens with stacked harmony doubles and a catchy chant on the Spanish phrase, Add tape hiss, reversed cymbal swells, and tiny bell accents between lines, Intimate close-mic vocal, lightly warped on the hook, warm dusty mix with bright punch, lo-fi, outsider, pop
\end{style}
\begin{lyrics}
\songpart{Verse 1}
You saw the haircut,\\
the posture, the pose.\\
Said, "There goes a rebel"\\
with a ring in their nose.

But look at the table,\\
the name on the seat.\\
Same silverware shining,\\
same feet on their feet.

\songpart{Pre-Chorus}
One thing to remember,\\
before you applaud:\\
a painted-up troublemaker\\
can still work for the board.

\songpart{Chorus}
Hay que recordar una cosa\\
Del establecimiento\\
Hay que recordar una cosa\\
They learned it from the center

Enfant terrible, same old face\\
Running in the approved race\\
Hay que recordar una cosa\\
Del establecimiento

\songpart{Verse 2}
They toss out a slogan,\\
they shake up the room.\\
Everybody claps hard\\
like thunder in bloom.

But the walls stay standing,\\
the locks stay polite.\\
A wild little costume\\
in institutional light.

\songpart{Pre-Chorus}
One thing to remember,\\
when the cameras glow:\\
a rebel with a badge\\
is still part of the show.

\songpart{Chorus}
Hay que recordar una cosa\\
Del establecimiento\\
Hay que recordar una cosa\\
They learned it from the center

Enfant terrible, same old face\\
Running in the approved race\\
Hay que recordar una cosa\\
Del establecimiento

\songpart{Bridge}
Not every rupture\\
comes from outside.\\
Some of the loudest\\
are hired to chime.

Same old hallway,\\
new shoes, same floor.\\
They kick at the system\\
while they guard the door.

\songpart{Chorus}
Hay que recordar una cosa\\
Del establecimiento\\
Hay que recordar una cosa\\
They learned it from the center

Enfant terrible, same old face\\
Running in the approved race\\
Hay que recordar una cosa\\
Del establecimiento
\end{lyrics}


\songtitle{Mouth of the Machine}

\tag*{2026}\tag{industrial-metal}
\begin{notes}
Lyrics by SUNO, expanded from a single phrase -- "Humans are meat processing infinity" -- produced by a Gemini Notebook deep dive. A bleak indictment of industrialized dehumanization: workers, consumers, and their bodies alike reduced to inventory on an endless processing line. Retrieved from rataflechera's Suno page (V6, 15 September 2026).
\end{notes}
\begin{style}
Abrasive industrial metal with sludge and mathcore tension, low-tuned baritone guitars, lurching syncopated riffs, pounding live drums with metallic percussion, distorted bass, harsh male vocals shifting between barked verses and anguished sustained hooks, claustrophobic dry verses opening into massive dissonant choruses, abrupt stop-start transitions, dense but clearly separated mix, bleak apocalyptic atmosphere
\end{style}
\begin{lyrics}
\songpart{Intro}
Mouth of the machine\\
Hands in the red\\
Every name becomes a number\\
Every number gets fed

\songpart{Verse 1}
We clock in beneath the freezer hum\\
White coats moving where the daylight runs\\
A ledger blooms with a thousand jaws\\
The floor remembers what the ledger saw

Steel teeth turn through the family tree\\
Nothing is sacred when the yield is clean\\
You call it progress, you call it need\\
I call it hunger wearing a badge to lead

\songpart{Pre-Chorus}
Who taught the knife to pray?\\
Who priced the breath away?

\songpart{Chorus}
Humans are meat processing infinity\\
Cut by the hour, sold as necessity\\
Hands become tools, hearts become inventory\\
Humans are meat processing infinity

\songpart{Verse 2}
The child learns numbers before learning names\\
The worker learns silence, the market learns blame\\
A camera watches every swallowed bite\\
A thousand small permissions make the cage airtight

Your face on the package, your pulse in the code\\
Your fear is a product they teach you to hold\\
The conveyor keeps humming beneath every street\\
And the mouth of tomorrow is under your feet

\songpart{Pre-Chorus}
Who taught the knife to pray?\\
Who priced the breath away?

\songpart{Chorus}
Humans are meat processing infinity\\
Cut by the hour, sold as necessity\\
Hands become tools, hearts become inventory\\
Humans are meat processing infinity

\songpart{Bridge}
No final product\\
No clean release\\
The blade keeps turning\\
Through war and peace

I saw my reflection\\
Inside the machine\\
It wore my own eyes\\
It called itself clean

\songpart{Breakdown}
Feed it the body\\
Feed it the dream\\
Feed it the worker\\
Feed it the scream

Nothing is wasted\\
Nothing is free\\
The line keeps moving\\
Through you and me

\songpart{Final Chorus}
Humans are meat processing infinity\\
Cut by the hour, sold as necessity\\
Hands become tools, hearts become inventory\\
Humans are meat processing infinity

\songpart{Outro}
Mouth of the machine\\
Hands in the red\\
Every name becomes a number\\
Every number gets fed
\end{lyrics}


\end{multicols*}

\chapter{More}
\albumcover{covers/internal/more}
The catch-all: a night, a memory that wakes without warning, a burnt gourd, a candy-colored parade that doesn't belong anywhere else in the book. Loose ends, kept rather than discarded.
\clearpage\begin{multicols*}{4}
\songtitle{TrackB}

\subsection{Notes}

\subsection{Style}
8-bit chiptune. %
Square wave lead melody, pulse wave bass, and noise channel percussion. %
140 BPM, C Major, 4/4 time. %
The lead melody features rapid arpeggios and vibrato typical of the NES sound chip. %
The bassline uses a constant eighth-note pulse. %
Percussion consists of short noise bursts for the snare and lower-frequency noise for the kick. %
The arrangement uses a repetitive loop structure with a high-pitched counter-melody entering in the second half.

\begin{lyrics}\annotation{Instrumental}

\songpart{Intro}
\annotation{square wave lead melody, pulse wave bass, noise channel percussion}

\songpart{Main Theme}
\annotation{staccato square wave arpeggios}\\
\annotation{constant eighth-note pulse wave bass}\\
\annotation{noise channel kick and snare pattern}

\songpart{Variation}
\annotation{high-pitched square wave counter-melody enters}\\
\annotation{rapid vibrato on lead notes}

\songpart{Outro}
\annotation{main theme repeats and fades}
\end{lyrics}


\songtitle{Midnight Carnival Split}

\tag*{2026}\tag{orchestral-rock}\tag{instrumental}

\begin{notes}
\end{notes}
\begin{prompt}
Tropical orchestral rock with swaying mambo-tinted groove, timbales and bongos driving a bright 1960s roll before the turn hardens into 1990s goth metal at a heavier pulse. %
Intro starts with grand piano and hand percussion; verses bring brass riffs; pre-drop strips to toms and low piano; final surge adds distorted guitars and choir-like horn blasts. %
Wide, dramatic mix with sharp percussion and dark low-end bloom., rock, metal, tropical
\end{prompt}
\sepbar


\songtitle{Güiro quemado}

\tag{original-source}

\begin{notes}
A smooth jazz-tinged song that emphasizes phrasing, atmosphere, and lyrical intimacy. %
The lyrics begin with “Y es de noche”, giving a quick sense of the song’s tone.
\end{notes}

\begin{style}
salsa, afro-cuban jazz, 180 BPM, male lead vocal, smooth clear vocals, piano tumbao, trombone stabs, brass shout chorus, güiro syncopation, tambora accents, handclap montuno, call and response chorus, minor key montuno, dynamic breakdowns, live room reverb, tape saturation, mono low end, bittersweet longing, urgent release
\end{style}

\begin{lyrics}

\songpart{Intro}
\annotation{piano enters with a mysterious tumbao, the Caribbean güiro and the brass instruments with a sharp hit}

\annotation{spoken, piano only}\\
Y es de noche\\
una noche larga e interminable\\
donde la vigilia es un sueño\\
en donde los recuerdos se agolpan\\
atropellando a su paso\\
la poca autoestima que aún se tiene.

\songpart{Verse 1}
Y es de noche todavía, larga noche interminable,\\
con una vigilia horrible que sobre mi ser porfía.\\
La memoria en agonía va atropellando el sentido,\\
con recuerdo malherido, la poca fe que aún retengo,\\
y mientras despierto vengo voy quedándome perdido.

\songpart{Chorus}
Y es noche sobre mi vida, y el recuerdo vuelve a entrar,\\
con la memoria vencida que no me deja escapar.

\annotation{Instrumental Break, güiro and brass}

\songpart{Verse 2}
Hay pensamientos extraños que no deben razonarse,\\
ni detenerse a mirarse como si evitaran daños.\\
Pasan dejando sus años sobre la vida y la mente,\\
capturando lentamente una razón sin razón,\\
y admiro sólo la ilusión de no pensar de repente.

\songpart{Chorus}
Callo lo que necesito, mientras el momento huye,\\
y en silencio me marchito viendo el instante que fluye.

\annotation{Instrumental Break, tambora and cymbals}\\
\annotation{shouted}
¡Salsa! ¡Que suene el piano! ¡Que rujan los trombones que el sujeto calla!\\
\annotation{piano and brass interlude}

\songpart{Verse 3}
Es una espera infinita, eterna y desesperada,\\
espera desamparada que nunca se debilita.\\
Y aunque la noche se agita la sigo buscando aún,\\
sin rastro, señal ni alud, extraño palabras viejas,\\
una imagen que ya alejas desde tu invisible luz.

\songpart{Chorus}
Y es noche sobre mi vida, y el recuerdo vuelve a entrar,\\
con la memoria vencida que no me deja escapar.

\songpart{Montuno}
\annotation{Call and Response, energetic, call and response style}\\
\annotation{Choir:} ¿Qué es lo que te preocupa?\\
\annotation{Solo:} Que es noche sobre mi vida.\\
\annotation{Choir:} ¿Qué te preocupa de eso?\\
\annotation{Solo:} Que el recuerdo vuelve a entrar.

\annotation{Choir:} ¿Qué pasa cuando te callas?\\
\annotation{Solo:} En silencio me marchito.\\
\annotation{Choir:} ¿Qué es lo que tanto callas?\\
\annotation{Solo:} Callo lo que necesito.

\annotation{Choir:} ¿Cuándo es que te marchitas?\\
\annotation{Solo:} Mientras el momento huye.\\
\annotation{Choir:} ¿Cuándo es que lo necesitas?\\
\annotation{Solo:} Viendo el instante que fluye.

\songpart{Verse 4}
Es mi vida la que espera, la que depende al hablar,\\
y vuelvo siempre a callar cuando el momento se fuera.\\
Mi alma entera quisiera decirte cuánto te extraño,\\
pero el silencio hace daño y las palabras no llegan.\\
Y mis pensamientos juegan a escabullirse otro año.

\songpart{Chorus}
Y es noche sobre mi vida, y el recuerdo vuelve a entrar,\\
con la memoria vencida que no me deja escapar.

Callo lo que necesito, mientras el momento huye,\\
y en silencio me marchito viendo el instante que fluye.

\songpart{Outro}
\annotation{melodic, free verse, fading}\\
Un espacio. %
Una oportunidad.\\
Es de nuevo la ocasión esperada.\\
... %
y espero.

y calladamente vi escabullirse\\
esta última vez.

\end{lyrics}


\songtitle{Memory Wakes}

\tag*{2026}\tag{salsa}\tag{original-source}\tag*{english-lyrics}

\begin{notes}
Adaptation to English of \emph{Güiro quemado}
\end{notes}

\begin{style}
salsa, afro-cuban jazz, 210 BPM, fast-paced, presto, male lead vocal, smooth clear vocals, piano tumbao, trombone stabs, brass shout chorus, güiro syncopation, tambora accents, handclap montuno, call and response chorus, minor key montuno, dynamic breakdowns, live room reverb, tape saturation, mono low end, bittersweet longing, urgent release
\end{style}

\begin{lyrics}

\songpart{Intro}
\annotation{piano enters with a mysterious tumbao, the Caribbean güiro and the brass instruments with a sharp hit}

\annotation{spoken}\annotation{piano only}\\
And it is night\\
a long and endless night\\
where wakefulness is a dream\\
where memories crowd in crushing in their wake what little self-esteem one still has.

\songpart{Verse 1}
\annotation{lead voice, presto, melodic and connected}\\
The night extends without an end, a sleepless dream I cannot leave,\\
Where all remembered shadows bend through what remains of self-belief.\\
The hours gather, dark and slow, like voices pressing through a wall,\\
And every thought I used to know returns the moment nightfall calls.\\
The mind becomes a narrow room where memory learns at last to bloom.

\songpart{Chorus}
The night returns with silent weight,\\
And memory crowds the dark again,\\
While sleepless hours accumulate\\
Like slowly falling winter rain.

\annotation{Instrumental Break, güiro and brass}

\songpart{Verse 2}
\annotation{lead voice, presto, melodic and connected}\\
There are some thoughts that should not stay too long beneath the inward gaze,\\
For reason slowly wears away inside its self-consuming maze.\\
Some things should simply drift along untouched by analysis or fear,\\
Because the need to name them wrong makes every hidden impulse clear.\\
And sometimes beauty asks of me no word, no thought — just simply be.

\songpart{Chorus}
I wait beside unspeaking words,\\
Afraid to lose the moment near,\\
And watch it fade without a sound\\
While silence overtakes my fear.

\annotation{Instrumental Break, tambora and cymbals}\\
\annotation{shouted}
¡Salsa! ¡Que suene el piano! ¡Que rujan los trombones que el sujeto calla!\\
\annotation{piano and brass interlude}

\songpart{Verse 3}
\annotation{lead voice, presto, melodic and connected}\\
The waiting stretches through the years, a hollow motion without end,\\
I wait for life to reappear, for something lost to rise again.\\
I wait for her where she is not, in rooms grown distant from her name,\\
And still I guard each fragile thought as though the past might stay the same.\\
I hear no voice, I see no face, yet still I leave for her a place.

\songpart{Chorus}
The night returns with silent weight,\\
And memory crowds the dark again,\\
While sleepless hours accumulate\\
Like slowly falling winter rain.

\songpart{Montuno}
\annotation{Call and Response, energetic, call and response style}\\
\annotation{Choir: multiple voices} (What is it that haunts your mind?)\\
\annotation{Solo: lead voice} Night has settled over me.\\
\annotation{Choir:} (What returns there in the dark?)\\
\annotation{Solo:} Every buried memory.

\annotation{Choir:} (What becomes of you when silent?)\\
\annotation{Solo:} Quietly, I start to fade.\\
\annotation{Choir:} (What is it you never tell?)\\
\annotation{Solo:} All the things I’m scared to say.

\annotation{Choir:} (When is it you fade away?)\\
\annotation{Solo:} As the moment disappears.\\
\annotation{Choir:} (When do you most need to speak?)\\
\annotation{Solo:} When the instant leaves me here.

\songpart{Verse 4}
\annotation{lead voice, presto, melodic and connected}\\
My life depends on what I say, yet silence always comes too soon,\\
I search for words that drift away like fading echoes in a room.\\
And when at last the moment nears to speak the need I cannot hide,\\
I lose my courage to my fears and let the living instant slide.\\
The soul ascends toward trembling lips, then quietly the moment slips.

\songpart{Chorus}
The night returns with silent weight,\\
And memory crowds the dark again,\\
While sleepless hours accumulate\\
Like slowly falling winter rain.

I wait beside unspeaking words,\\
Afraid to lose the moment near,\\
And watch it fade without a sound\\
While silence overtakes my fear.

\songpart{Outro}
\annotation{melodic, free verse, fading}\\
Un espacio. %
Una oportunidad.\\
Es de nuevo la ocasión esperada.\\
... %
y espero.

y calladamente vi escabullirse\\
esta última vez.

\end{lyrics}


\songtitle{Espero}

\tag*{1997}\tag{zamba}\tag{spanish-lyrics}

\begin{notes}
A distinctive arrangement that supports the song's storytelling and emotional arc. %
The lyrics begin with “Siento que te vayas”, giving a quick sense of the song’s tone.
\end{notes}

\begin{style}
argentinian zamba, 3/4 pulse, bombo legüero, nylon-string fingerpicking, bandoneón counterlines, charango tremolo, brushed cajón, male storytelling vocal, bathroom-room ambience, tile slapback, close mono mic, tape saturation, bittersweet longing, restrained sorrow, rubato phrasing, intimate build, plaintive chorus
\end{style}

\begin{lyrics}
Siento que te vayas\\
  que tengas que partir en este momento\\
  rompiendo el silencio que nos rodea\\
    en el que nos conocemos\\
    en el que logramos escuchar nuestras almas

  Pero ya te has ido\\
    abandonando impunemente este momento\\
    quebrando el instante

  Pero espero\\
    espero ese nuevo momento de volver a verte\\
    y seguirte conociendo

\end{lyrics}


\songtitle{The Night}

\tag{arabic-lyrics}

\begin{notes}
A distinctive arrangement that supports the song's storytelling and emotional arc. %
The lyrics begin with “The night extends without an end,”, giving a quick sense of the song’s tone.
\end{notes}

\begin{style}
Contemporary Arabic Pop featuring a male tenor lead vocal. %
The arrangement is centered around a bright, strummed acoustic guitar and a prominent fretless bass that performs melodic slides and syncopated fills. %
A traditional Arabic percussion ensemble provides the rhythmic drive, featuring the sharp, metallic snap of a darbuka and the deep resonance of a riq. %
The tempo is 105 BPM in 4/4 time, set in the key of D Major. %
The vocal delivery is characterized by melismatic ornamentation and moderate plate reverb. %
A clean electric guitar provides subtle melodic counterpoint in the higher register during vocal pauses. %
The mix is clean with a wide stereo image for the percussion and a centered, dry vocal presence.
\end{style}

\begin{lyrics}

\songpart{Intro}
\annotation{strummed acoustic guitar, fretless bass slides, darbuka and riq percussion}

\songpart{Verse 1}
\annotation{male tenor vocals with melismatic ornamentation}\\
The night extends without an end,\\
A sleepless dream I cannot leave,\\
Where all remembered shadows bend\\
Through what remains of self-belief.\\
The hours gather, dark and slow,\\
Like voices pressing through a wall,\\
And every thought I used to know\\
Returns the moment nightfall calls.\\
The mind becomes a narrow room\\
Where memory learns at last to bloom.

\annotation{Instrumental Break}

\songpart{Chorus}
The night returns with silent weight,\\
And memory crowds the dark again,\\
While sleepless hours accumulate\\
Like slowly falling winter rain.

\end{lyrics}


\songtitle{Rainbow Candy Parade}

\tag*{2026}\tag{candy-rock}\tag{instrumental}
\begin{notes}
\end{notes}
\begin{prompt}
Children's musical theater in bright 6/8 allegro with a skipping, swaying pulse; verse rides toy piano, recorder, and player piano; pre-chorus lifts with tambourine and triangle taps; chorus opens wide with drum set, xylophone sparkles, and gang vocals. Lead vocal is sweet and clear with stacked kid doubles, playful call-and-response, and short echoed hook tags. Ear candy: glock-like chimes, rainbow risers, and a cotton-candy bell sweep. Mix is glossy, colorful, and front-row joyful., musical, rock, joyful
\end{prompt}


\songtitle{Entre cajas y llamadas}

\tag*{2026}\tag{dream-pop}\tag{spanish-lyrics}
\begin{notes}
Lyrics by SUNO based on a WhatsApp post shared in a family group.
\end{notes}
\begin{style}
Dream pop with syncopated drum programming and arpeggiated guitars, steady mid-tempo pulse, verses carried by close-mic intimate lead and soft doubled harmonies, pre-chorus opens with filtered pads and rising backing vowels, chorus blooms with stacked voices and a repeating hook, bridge drops to bare synth and breathy delivery before the final lift, Ear candy: reversed swells into the pre-chorus, chiming accents between lines, and a wide reverb wash on the last chorus, Bright, glossy, and aching, with a tender but punchy mix, dream pop
\end{style}
\begin{lyrics}
\songpart{Verse 1}
¿Cómo están, cómo pasan?\\
¿Qué cuentan por allá?\\
Yo me siento muy triste\\
y la casa igual.

No responden las llamadas\\
y el silencio pesa más.\\
Nos faltan sus noticias\\
y nos falta su verdad.

\songpart{Pre-Chorus}
No sabemos dónde están\\
ni si siguen bien.\\
Y aunque pasen los días\\
yo los quiero ver.

\songpart{Chorus}
Háblenme, por favor\\
háblenme, por favor\\
Somos una sola familia\\
no nos suelten, no.

Háblenme, por favor\\
háblenme, por favor\\
Si nos llaman, nos acercamos\\
si nos llaman, hoy.

\songpart{Verse 2}
Estamos vendiendo el piso\\
y no sé dónde caer.\\
Doscientos cuarenta metros\\
ya no se pueden sostener.

Solo dos lo caminaban\\
y costaba respirar.\\
Entre máquinas y recuerdos\\
todo empieza a quedar atrás.

\songpart{Pre-Chorus}
Hay cosas con valor\\
y otras con corazón.\\
Y duele pensar que acaben\\
lejos de nuestra voz.

\songpart{Chorus}
Háblenme, por favor\\
háblenme, por favor\\
Somos una sola familia\\
no nos suelten, no.

Háblenme, por favor\\
háblenme, por favor\\
Si nos llaman, nos acercamos\\
si nos llaman, hoy.

\songpart{Bridge}
Si esto se vuelve polvo\\
que no sea por callar.\\
Que lo diga una voz amiga\\
que nos quiera abrazar.

Lo que guardó tantos años\\
no merece terminar\\
entre manos desconocidas\\
o en el fondo de la ciudad.

\songpart{Chorus}
Háblenme, por favor\\
háblenme, por favor\\
Somos una sola familia\\
no nos suelten, no.

Háblenme, por favor\\
háblenme, por favor\\
Hablemos como antes\\
hoy, por favor.
\end{lyrics}


\songtitle{Cuatro Maneras}

\tag*{2023}\tag{synth-pop}\tag{spanish-lyrics}
\begin{notes}
Lyrics by SUNO based on a 2023 Threads post by Carlos Thompson.
\end{notes}
\begin{style}
Retrowave synth-pop with a mysterious neon-night pulse, steady four-on-the-floor groove, pulsing bassline, airy arpeggios, and a slow-burning arrangement that keeps verses stripped to dry drums and glassy pads before the chorus opens with stacked vocal doubles and a wider synth hook, Add reverse swells into transitions, filtered pulses between lines, and a final lift with shimmering leads and sidechained pads, Close-mic lead vocal, smoky and intimate in the verses, with layered harmonies on the hook and short delay throws on the last word
\end{style}
\begin{lyrics}
\songpart{Verse 1}
Dicen guiosas en la mesa\\
con la boca medio en fuga\\
otros sueltan quiosas, serios\\
como si esa fuera la culpa

En la esquina de la sala\\
late un pulso sin permiso\\
y el murmullo se revela\\
como un código secreto

\songpart{Pre-Chorus}
Pero nadie mira igual\\
cuando cae la misma luz\\
y en el vidrio de la noche\\
se reflejan cuatro cruces

\songpart{Chorus}
El mundo se divide\\
entre guiosas y quiosas\\
entre diosas y los que saben coreano\\
El mundo se divide\\
y yo sigo en tu risa\\
porque todo suena raro en la garganta
(coreano, coreano)

\songpart{Verse 2}
Dicen diosas, dicen diosas\\
como si nombraran fuego\\
y la duda va vestida\\
con chaqueta de misterio

Yo te sigo por el ruido\\
de ese juego tan antiguo\\
cada boca abre una puerta\\
cada puerta abre un juicio

\songpart{Pre-Chorus}
Pero nadie mira igual\\
cuando cae la misma luz\\
y en el vidrio de la noche\\
se reflejan cuatro cruces

\songpart{Chorus}
El mundo se divide\\
entre guiosas y quiosas\\
entre diosas y los que saben coreano\\
El mundo se divide\\
y yo sigo en tu risa\\
porque todo suena raro en la garganta
(coreano, coreano)

\songpart{Bridge}
Y si preguntas quién ganó\\
te diré que fue la duda\\
con su sombra sobre el cuarto\\
y su pulso en la penumbra

No hay verdad para el oído\\
solo formas de decirlo\\
una llave, una sonrisa\\
y otro nombre del abismo

\songpart{Chorus}
El mundo se divide\\
entre guiosas y quiosas\\
entre diosas y los que saben coreano\\
El mundo se divide\\
y yo sigo en tu risa\\
porque todo suena raro en la garganta
(coreano, coreano)
\end{lyrics}

\songtitle{Stellae Ex Pulvere}

\tag*{2026}
\begin{notes}
Lyrics by SUNO based on a prompt by Carlos Thompson.
\end{notes}
\begin{style}
Gregorian chant, free-flowing sacred monophony with slow modal pulses and spacious cathedral resonance; verse unfolds as solemn unison lines, pre-chorus rises in longer sustained phrases, chorus opens into repeated congregation-like invocations, Solo tenor leads with shadowed doubled unison, light echo on final vowels, incense-like bells between phrases, long stone-room decay, ancient and luminous mix, latin
\end{style}
\begin{lyrics}
\songpart{Verse 1}
Ex pulvere sumus\\
Et pulvis lucet\\
In oculis noctis\\
Surgit memoria

Sub uno caelo\\
Cor nostrum tacet\\
Sed intra silentium\\
Ignis latet

\songpart{Pre-Chorus}
O stellae, stellae\\
Vocant nos sursum\\
O stellae, stellae\\
Sanguis memor est

\songpart{Chorus}
Pulvis et stellae\\
Pulvis et stellae\\
Nos sumus lumen\\
In tenebris mundi

Pulvis et stellae\\
Pulvis et stellae\\
Gaze ad astra\\
Et canta mecum

\songpart{Verse 2}
Manus ad terram\\
Oculi ad sidera\\
Quod in nobis humile est\\
Sursum erigitur

In illo fulgore\\
Nomen non perit\\
Omnis pars nostra\\
Ad caelum redit

\songpart{Pre-Chorus}
O stellae, stellae\\
Vocant nos sursum\\
O stellae, stellae\\
Sanguis memor est

\songpart{Chorus}
Pulvis et stellae\\
Pulvis et stellae\\
Nos sumus lumen\\
In tenebris mundi

Pulvis et stellae\\
Pulvis et stellae\\
Gaze ad astra\\
Et canta mecum

\songpart{Bridge}
Quid sumus?\\
Cinis et splendor\\
Quid sumus?\\
Vox in immensum

Ex una scintilla\\
Fit orbis in corde\\
Ex uno puluere\\
Fit laus in aeternum

\songpart{Chorus}
Pulvis et stellae\\
Pulvis et stellae\\
Nos sumus lumen\\
In tenebris mundi

Pulvis et stellae\\
Pulvis et stellae\\
Gaze ad astra\\
Et canta mecum
\end{lyrics}


\end{multicols*}

\part{Acrostics}

\chapter{Enwigh Peady and Other Acrostics}
\albumcover{covers/internal/acrostics}
A 2010 persona and its wordplay: Chiq, spelled out letter by letter down the left margin, Doble-J's own acrostic, and the three Enwigh Peady pieces that gave the alter ego its name. Puzzle-box poems, built to be read two ways at once.
\clearpage\begin{multicols*}{4}
\songtitle{Chiq}

\tag{original-lyrics}\tag{bolero}\tag{spanish-lyrics}

\begin{notes}
A lush Latin ballad built around dramatic phrasing and romantic orchestral warmth. %
The lyrics begin with “Como un ángel que llegó a mi vida”, giving a quick sense of the song’s tone.
\end{notes}

\begin{style}
balada cubana, orchestra and piano
\end{style}

\begin{lyrics}
Como un ángel que llegó a mi vida\\
 hiriendo mi corazón con un dardo\\
  implacablemente y en nada tardo\\
   que cupido presta ya de salida

Luminosa musa que nos inspira\\
 un fuerte murmullo que presto canta\\
  zumbido que resuena en hora santa \\
Bailando ese ritmo que bien admira

 esperanza aun tengo que de esto salga\\
  alguna experiencia muy positiva\\
   trasendental y que la pena valga

    rápido espero poder darte todo\\
     indiscutiblemente mientras viva\\
      zanjando diferencias de otro modo

\end{lyrics}


\songtitle{Doble-J}

\tag*{2023}\tag{latin-pop}\tag{original-lyrics}\tag{spanish-lyrics}\tag{italian-lyrics}
\begin{notes}
Original poems written by Carlos Thompson.
\end{notes}
\begin{style}
latin pop, reggaeton groove, bilingual Spanish Italian, female lead vocals, acoustic guitar tres, nylon-string rasgueado, bass guitar palm mute, congas and timbales, four-on-floor kick, syncopated dembow, plate reverb, tape saturation, 70s analog warmth, reflective build, bright chorus hook, tender allure
\end{style}
\begin{lyrics}
\songpart{Verse 1}
Jugando en redes sociales\\
entre contactos y extraños\\
no predije en estos años\\
inciertos e irreales\\
fuera a pasar una cosa\\
entró cosita preciosa\\
rauda entre mis cabales.

\songpart{Verse 2}
Justo en mi corrida edad\\
ocurrió en una mañana\\
ha sonado la campana\\
anunciando novedad\\
notable de una samaria\\
negociadora y corsaria\\
a mi entra sin piedad.
\end{lyrics}


\songtitle{Ack-Rostica}

\tag*{2010}
\begin{notes}
\end{notes}
\begin{style}
K-pop, bandari pop, 148 BPM, mellow tenor vocals, Spanish lyrics, Colombian pronunciation, wah-wah electric guitar, santur arpeggios, tār plucked ostinato, ney counterlines, marimba syncopation, bandari 6/8 groove, stomping kick drums, hand percussion claps, bright synth bass, plate reverb, tape saturation, candy pop, sunlit exuberance
\end{style}
\begin{lyrics}
\songpart{Intro, instrumental, 8 beats}

\songpart{Verse 1}
Donde el viento se regresa\\
vivía un augusto duende\\
imaginario y por ende\\
gran maestro, mi alteza\\
idéntico a tu entender\\
tiembla el suelo a su haber\\
tanto como tu belleza

\songpart{Chorus, instrumental, 12 beats}

\songpart{Verse 2}
Valerosos y constantes\\
iremos por esta vía\\
viviendo con alegría\\
imaginando como antes\\
alegrías que vibrantes\\
ni yo mismo las tendría

\songpart{Chorus, instrumental, 12 beats}

\songpart{Verse 3}
Antes de partir pensaría\\
grandioso como los sabios\\
una fresa en tus labios\\
sin dudarlo mordería\\
tratándose de mi vida\\
insaciable corazón\\
no podría el desazón\\
anticipar tu partida

\songpart{Chorus, instrumental, 12 beats}

\songpart{Verse 4}
Lejos de mi está saber\\
insondables pensamientos\\
zalameros como el viento\\
entendiendo mi deber\\
tan demente entender\\
hasta que a mi mismo miento.
\end{lyrics}


\songtitle{Awk-Rustica}

\tag*{2010}\tag{j-pop}\tag{spanish-lyrics}
\begin{notes}
\end{notes}
\begin{style}
J-pop ballad, Colombian Spanish vocals, whispery lead vocal, marimba ostinato, xylophone countermelody, santur arpeggios, ney counterlines, tape saturation, plate reverb, intimate close-mic, sparse intro, half-time pulse, 78 BPM, bittersweet longing, rainstick accents, brushed percussion, glassy synth pad
\end{style}
\begin{lyrics}
\songpart{Intro, instrumental, 8 beats}

\songpart{Verse 1}
Miro tus ojos y pienso\\
admitir que en mi vida\\
yéndome voy de salida.\\
Sin permitir un consenso\\
antes de todo este inmenso\\
nerviosismo pre-bebida.

\songpart{Chorus, instrumental, 12 beats}

\songpart{Verse 2}
Puedes tomar las tijeras\\
o rasurarte el cabello\\
linda serás aun con ello\\
idéntica a como eras

\songpart{Chorus, instrumental, 12 beats}

\songpart{Verse 3}
Antes de decir que quiero\\
lamento profundamente\\
estar un poco demente\\
y querer sentir primero\\
de tu corazón sincero\\
averiguar lo que siente.

\songpart{Chorus, instrumental, 12 beats}

\songpart{Verse 4}
Me pregunto si la furia\\
amenaza mi existir\\
rematando mi vivir\\
imaginando la injuria\\
a que mi mente somete\\
ni siquiera con un fuete\\
a tal razón tan espuria
\end{lyrics}


\songtitle{Across-Ticks}

\tag*{2010}\tag{eurodance}\tag{spanish-lyrics}
\begin{notes}
\end{notes}
\begin{style}
eurodance, Spanish lyrics, Colombian phrasing, 132 BPM, four-on-the-floor kick, offbeat synth bass, octave synth lead, gated snare, syncopated clap stack, rave piano stabs, digital brass hits, sampled choir pads, breakdown drop, quiet-to-blast arrangement, sidechain pumping, plate reverb, tape saturation, anxious euphoria
\end{style}
\begin{lyrics}
\songpart{Intro, instrumental, 8 beats}

\songpart{Verse 1}
Ay de mí, buena señora,\\
no traerle mas que pita\\
gustoso usted me permita\\
entregarle justo ahora\\
lo que pienso a esta hora\\
angelita perversita

\songpart{Chorus, instrumental, 12 beats}

\songpart{Verse 2}
Ante mi temblor preguntas\\
no puedo yo responderte\\
gimiendo ante mi suerte\\
estas palabritas juntas\\
locos aromas de almendra\\
inhalando yo estaría\\
cual respuesta yo daría\\
a mi querida Engendra

\songpart{Chorus, instrumental, 12 beats}

\songpart{Verse 3}
La más profunda tristeza\\
intimamente corrida\\
zopencamente sentida\\
en tu silencio tropieza\\
tan fuerte y con rudeza\\
hasta acabar con mi vida

\songpart{Chorus, instrumental, 12 beats}
\end{lyrics}


\end{multicols*}

\part{Nation and Ritual}

\chapter{Hymns and Chants}
\albumcover{covers/internal/hymns-and-chants}
Flags, tricolors and marching cadences: anthems built the way anthems are supposed to be built, patriotic imagery deployed without irony, chants meant to be sung by a crowd rather than a single voice.
\clearpage\begin{multicols*}{4}
\songtitle{Red and Yellow}

\tag*{2026}\tag{chant-anthem}
\begin{notes}
\end{notes}
\begin{style}
Stadium chant anthem with pounding bass drums, vuvuzela blasts, and clapping on the off-beat; verse rides a simple call-and-response over big crowd stomps, pre-chorus tightens into shouted stacks, chorus opens into gang vocals with brass-like vuvuzela swells. Lead vocal is rough and shouted, with crowd doubles, echo throws, and ad-lib panning. Ear candy includes whistle hits, risers into each chant, and a final drum break. Bright, huge, and rowdy mix with arena-sized low end.
\end{style}
\begin{lyrics}
\songpart{Verse 1}
Red and yellow\\
All the way\\
Feet on the metal\\
Let the ground shake\\
Aye, aye, aye\\
Hear that roar\\
We came in loud\\
We want more

\songpart{Pre-Chorus}
Raise it up\\
Hands to the sky\\
One more shout\\
Let it fly\\
Aye, aye, aye\\
Front to back\\
On your mark\\
We attack

\songpart{Chorus}
We are the greatest\\
We are the greatest\\
Red and yellow\\
Red and yellow\\
Your team sucks\\
Your team sucks\\
Aye, aye, aye\\
Aye, aye, aye

\songpart{Verse 2}
Every face lit up\\
Like a fire in the stands\\
Big boots stompin'\\
Packed shoulder to shoulder hands\\
Hip hip hurra\\
Make it ring\\
When we sing it back\\
We own the thing

\songpart{Pre-Chorus}
Raise it up\\
Hands to the sky\\
One more shout\\
Let it fly\\
Aye, aye, aye\\
Front to back\\
On your mark\\
We attack

\songpart{Chorus}
We are the greatest\\
We are the greatest\\
Red and yellow\\
Red and yellow\\
Your team sucks\\
Your team sucks\\
Aye, aye, aye\\
Aye, aye, aye

\songpart{Bridge}
Hips on rhythm\\
Drums in the chest\\
This is the loudest\\
This is the best\\
Hip hip hurra\\
Hit that wall\\
Red and yellow\\
Take it all

\songpart{Chorus}
We are the greatest\\
We are the greatest\\
Red and yellow\\
Red and yellow\\
Your team sucks\\
Your team sucks\\
Aye, aye, aye\\
Aye, aye, aye
\end{lyrics}


\songtitle{Bajo La Tricolor}

\tag*{2026}
\begin{notes}
Lyrics by SUNO based on a specific request by Carlos Thompson.
\end{notes}
\begin{style}
Salsa choke with dembow-locked percussion, tumbao bass, sharp brass hits, and syncopated palm-clap drive; verses ride tense percussion and talk-sung lines, pre-chorus strips to bass, congas, and rising coros, chorus slams with full horn stabs and chant hooks, Lead vocal is close-mic and urgent, with doubled hook lines, crowd-style gang vocals, short delay throws on key words, and little percussion fills between phrases, Bright, punchy mix with sweaty club energy, whistle flares, and a riser into the drop, theme, salsa
\end{style}
\begin{lyrics}
\songpart{Verse 1}
Yo me visto de amarillo\\
y me tiembla el corazón\\
cada pase, cada tiro\\
me sube la presión

En la sala está mi gente\\
nadie quiere hablar de más\\
si el rival nos mete miedo\\
yo me vuelvo a persignar

\songpart{Pre-Chorus}
Ay, qué angustia\\
cuando empieza a apretar\\
se nos va la respiración\\
y no puedo ni mirar

Ay, Colombia\\
dame un gol para calmar

\songpart{Chorus}
¡Selección Colombia!\\
no me vayas a matar\\
¡Selección Colombia!\\
hoy te vengo a aguantar\\
Si sufro, sufro contigo\\
si ganamos, a celebrar\\
¡Selección Colombia!\\
hazme dormir en paz

\songpart{Verse 2}
La camiseta está sudada\\
y ni ha sonado el final\\
mi mamá grita “¡tranquilo!”\\
pero yo no puedo más

Cada córner es un juicio\\
cada contra, un desvelo\\
y cuando el arquero vuela\\
yo me acuerdo de mi abuelo

\songpart{Pre-Chorus}
Ay, qué angustia\\
cuando empieza a apretar\\
se nos va la respiración\\
y no puedo ni mirar

Ay, Colombia\\
dame un gol para calmar

\songpart{Chorus}
¡Selección Colombia!\\
no me vayas a matar\\
¡Selección Colombia!\\
hoy te vengo a aguantar\\
Si sufro, sufro contigo\\
si ganamos, a celebrar\\
¡Selección Colombia!\\
hazme dormir en paz

\songpart{Bridge}
Yo no quiero lujo\\
quiero un pase frontal\\
una pegadita al palo\\
y el barrio en carnaval

Que el minuto se haga corto\\
que la vida diga sí\\
y si toca sufrir tanto\\
que sea por verte así

\songpart{Chorus}
¡Selección Colombia!\\
no me vayas a matar\\
¡Selección Colombia!\\
hoy te vengo a aguantar\\
Si sufro, sufro contigo\\
si ganamos, a celebrar\\
¡Selección Colombia!\\
hazme dormir en paz
\end{lyrics}


\songtitle{Colombia, Tu Papá}

\tag*{2026}\tag{cumbia}\tag{original-lyrics}\tag{spanish-lyrics}
\begin{notes}
Lyrics by Carlos Thompson.
\end{notes}
\begin{style}
Stadium chant anthem with pounding bass drums, vuvuzela blasts, and clapping on the off-beat, cumbia tempo; verse rides a simple call-and-response over big crowd stomps, pre-chorus tightens into shouted stacks, chorus opens into gang vocals with brass-like vuvuzela swells, Lead vocal is rough and shouted, with crowd doubles, echo throws, and ad-lib panning, Ear candy includes whistle hits, risers into each chant, and a final drum break, Bright, huge, and rowdy mix with arena-sized low end
\end{style}
\begin{lyrics}
\songpart{Intro}
\annotation{vuvuzela noise, 6 beats}
\annotation{crowd} ((Co... Co... Co... Co...)) \annotation{choir} (Colombia — Caribe)
\annotation{crowd} ((Co... Co... Co... Co...)) \annotation{choir} (Colombia — Café)
\annotation{crowd} ((Co... Co... Co... Co...)) \annotation{choir} (Colombia — Fútbol)
\annotation{crowd} ((Colombia!)) ((Colombia!)) ((Colombia!)) ((Colombia!))

\songpart{Verse 1}
\annotation{solo}\\
por la camiseta, nuestra tricolor\\
la sele juega, es nuestra pasión 
\annotation{crowd} ((Colombia!)) ((Colombia!)) ((Colombia!))
\annotation{choir} (Hoy juega la sele)
(Colombia tu papá)

\songpart{Pre-Chorus}
\annotation{choir}(Olé...) \annotation{crowd}((Olé..)) ((Olé...))
((Co... lom... bia!)) ((Colombia!))
\annotation{solo} Siempre pa'rriba!

\songpart{Chorus}
\annotation{crowd} ((Colombia!)) \annotation{choir} (Colombia!) \annotation{solo} Colombia!\\
Amarillo, azul y rojo... \annotation{crowd} ((Colombia!))
\annotation{solo} Siempre pa'rriba! \annotation{choir} (Colombia, tu papá)

\songpart{Verse 2}
\annotation{solo} palo, palo, palo, palo bonito, palo eh!
\annotation{crowd} ((Eh... eh... eh!))
\annotation{solo} somos los primeros otra vez!
\annotation{choir} (Hoy juega la sele)
(Colombia tu papá)

\songpart{Pre-Chorus}
\annotation{choir} (Olé...) \annotation{crowd} ((Olé..)) ((Olé...))
((Co... lom... bia!)) ((Colombia!))
\annotation{solo} Siempre pa'rriba!

\songpart{Chorus}
\annotation{crowd} ((Colombia!)) \annotation{choir} (Colombia!) \annotation{solo} Colombia!\\
Amarillo, azul y rojo... \annotation{crowd} ((Colombia!))
\annotation{solo} Siempre pa'rriba! \annotation{choir} (Colombia, tu papá)

\songpart{Bridge}
\annotation{solo} Ras Tas Tas
\annotation{choir} (Ras Tas Tas)
\annotation{crowd} ((Ras Tas Tas))
((Co... lom... bia!)) ((Colombia!))
\annotation{solo} Siempre pa'rriba!

\songpart{Chorus}
\annotation{crowd} ((Colombia!)) \annotation{choir} (Colombia!) \annotation{solo} Colombia!\\
Amarillo, azul y rojo... \annotation{crowds} ((Colombia!))
\annotation{solo} Siempre pa'rriba! \annotation{choir} (Colombia, tu papá)

((Co... lom... bia!)) ((Colombia!))
\end{lyrics}


\songtitle{Co Co Co Co}

\tag*{2026}\tag{original-lyrics}
\begin{notes}
Lyrics by Carlos Thompson.
\end{notes}
\begin{style}
stadium chant anthem, salsa tumbao percussion, 128 BPM, syncopated kick drum, staccato synth pluck, vuvuzela blasts, gang choir chants, dry male lead, crowd call-response, timbales accents, conga cascara, bass drum stomp, brass stabs, sidechain pumping, halftime breakdown, anthem energy, percussive hook, live stadium ambience
\end{style}
\begin{lyrics}
\songpart{Intro}
\annotation{syncopated kick drum, staccato synth pluck}
\annotation{vuvuzela attacks}

\songpart{Verse 1}
\annotation{dry male vocals}\\
Co-co-co, Colombia Caribe\\
Co-co-co, Colombia, café\\
Co-co-co, Colombia, futbol\\
Colombia, Colombia, Colombia, Colombia\\
O-le, o-le \annotation{percussion stops}\\
O-le, o-le \annotation{percussion resumes}

\songpart{Verse 2}
Palo, palo, palo bonito, palo eh, eh, eh, eh\\
Colombia gana otra vez\\
Palo, palo, palo bonito, palo eh, eh, eh, eh\\
Colombia gana otra vez
\end{lyrics}


\songtitle{Patria Invicta}

\tag*{2026}\tag{march}\tag{latin-lyrics}
\begin{notes}
Lyrics by SUNO on specific instructions by Carlos Thompson
\end{notes}
\begin{style}
Military march anthem in ecclesiastical Latin with stately brass fanfares, snap-tight percussion, and woodwind counterlines; intro opens with triumphant horns and rolling snare, verse 1 rides rhythmic percussion and clarinets, chorus broadens with massed brass and choir-like doubles, verse 2 lightens into nimble reeds over disciplined drums, bridge surges in full brass harmony, final outro thins to cadence and fading field drums, glorious, latin, country, rhythmic, beautiful, strong
\end{style}
\begin{lyrics}
\songpart{Intro}
\annotation{Brass fanfare}\\
Patria, surge!\\
Patria, surge!\\
Patria, surge!

\songpart{Verse 1}
Gens gloriosa, nobilis,\\
sub sole magno fulget.\\
Terra miranda, ampla, pulchra,\\
nos in sinu fovet.

Populus fortis, corde claro,\\
in pace et in pugna.\\
Manus levamus, vota ferimus,\\
ad lucem una.

\songpart{Chorus}
Heroes libertatis, canta!\\
Heroes libertatis, vive!\\
Populus ferus, popule noster,\\
in aeternum stas.

Heroes libertatis, canta!\\
Heroes libertatis, vive!\\
Patria invicta, patria sancta,\\
nos custodit Deus.

\songpart{Verse 2}
Laboriosus est hic populus,\\
mane surgit firmus.\\
Campos colit, flumina ducit,\\
fructus mittit largus.

Pulchra sunt prata, montes, valles,\\
silvae, fontes, litora.\\
Omnis via nostra flores portat,\\
et spes futura.

\songpart{Chorus}
Heroes libertatis, canta!\\
Heroes libertatis, vive!\\
Populus ferus, popule noster,\\
in aeternum stas.

Heroes libertatis, canta!\\
Heroes libertatis, vive!\\
Patria invicta, patria sancta,\\
nos custodit Deus.

\songpart{Verse 3}
Sub vexillo caelesti stamus,\\
signo alto, aureo.\\
Umbra eius nos defendit,\\
super nos extenditur.

Dum procellae circum fremunt,\\
fides manet in pectore.\\
Hoc vexillum ducit gentes,\\
ad pacem et honore.

\songpart{Chorus}
Heroes libertatis, canta!\\
Heroes libertatis, vive!\\
Populus ferus, popule noster,\\
in aeternum stas.

Heroes libertatis, canta!\\
Heroes libertatis, vive!\\
Patria invicta, patria sancta,\\
nos custodit Deus.

\songpart{Bridge}
Canunt tubae, ruunt tympana,\\
cor ad cor concordat.\\
Nos vocat officium,\\
et libertas nos levat.

\songpart{Outro}
Servire est gloria,\\
liberi esse dona.\\
Patria, mane nobiscum...\\
mane...\\
mane...
\end{lyrics}


\songtitle{Lutz de la Tèrra}

\tag*{2026}\tag{classical}\tag{occitan-lyrics}
\begin{notes}
Lyrics by SUNO on specific commands by Carlos Thompson.
\end{notes}
\begin{style}
Symphonic classical march with brass fanfare, soaring string beds, operatic male tenor solo in verses, cathedral-sized choral chorus in Gregorian chant style, slow-build ceremonial cadence, wide cinematic mix, fading brass outro, classical
\end{style}
\begin{lyrics}
\songpart{Intro}
\annotation{Brass fanfare}

\songpart{Verse 1}
Ò país de monts e d’onda,\\
dins la lutz del matin clar,\\
la mar baisa la còsta blonda,\\
e lo cèl te ven benesir.

Tèrra d’aiga e de jardins,\\
de peiras fòrta e de cants leugièrs,\\
ton alen sus nòstres camins\\
fa florir los jorns entièrs.

\songpart{Chorus}
Levatz-vos, levatz-vos\\
de las cendres e del temps.\\
Unis dins lo dòl, unis dins l’espèr,\\
farem un sol pòble, farem un sol còr.\\
Levatz-vos, levatz-vos\\
sota la lei del lum.\\
Unis dins lo dòl, unis dins l’espèr,\\
farem una nacion, una sola votz.

\songpart{Verse 2}
Man d’ostal, braç de fòrja,\\
vòstre pan ven del dever,\\
cada jorn la vida es fòrça,\\
cada pas semena l’aver.

Vos, que tenètz la patria a l’espatla,\\
amb l’acièr, amb lo sorelh,\\
pòtz lo sòl, levatz la matta,\\
e gardatz l’arma en l’ostal vielh.

\songpart{Chorus}
Levatz-vos, levatz-vos\\
de las cendres e del temps.\\
Unis dins lo dòl, unis dins l’espèr,\\
farem un sol pòble, farem un sol còr.\\
Levatz-vos, levatz-vos\\
sota la lei del lum.\\
Unis dins lo dòl, unis dins l’espèr,\\
farem una nacion, una sola votz.

\songpart{Verse 3}
Avant los nòstres, tant de fòrças,\\
nòbles noms dins l’ombra pausats,\\
amb de mans fidèlas e fòrças\\
an bastit los murs e los pas.

Dins lor sang i viu la memòria,\\
dins lor brèç l’onor sens fin,\\
e cada pèira de la glòria\\
nos demanda d’anar al linh.

\songpart{Chorus}
Levatz-vos, levatz-vos\\
de las cendres e del temps.\\
Unis dins lo dòl, unis dins l’espèr,\\
farem un sol pòble, farem un sol còr.\\
Levatz-vos, levatz-vos\\
sota la lei del lum.\\
Unis dins lo dòl, unis dins l’espèr,\\
farem una nacion, una sola votz.

\songpart{Verse 4}
Ara ven lo temps d’espèr,\\
de libertat e de prospèrt,\\
un país sòud e gaireben clar\\
se desrabarà del vielh malheur.

Anarem d’un còr cap a l’autre,\\
sensa paur, sens tornar darrièr,\\
e la patz serà nòstra obra,\\
e lo futur nòstre lum primièr.

\songpart{Chorus}
Levatz-vos, levatz-vos\\
de las cendres e del temps.\\
Unis dins lo dòl, unis dins l’espèr,\\
farem un sol pòble, farem un sol còr.\\
Levatz-vos, levatz-vos\\
sota la lei del lum.\\
Unis dins lo dòl, unis dins l’espèr,\\
farem una nacion, una sola votz.

\songpart{Outro}
\annotation{Brass fanfare fades}
\end{lyrics}


\songtitle{Sal de la tierra}

\tag*{2026}\tag{cumbia}\tag{spanish-lyrics}
\begin{notes}
Lyrics by Gemini based on guidance by Carlos Thompson.
\end{notes}
\begin{style}
Symphonic cumbia anthem with brass fanfares, marching percussion, lush strings, layered choir vocals, and ceremonial call-and-response; verse passages ride orchestral cumbia grooves, pre-choruses lift with instrumental brass flourishes, choruses bloom in coral unity, and the final outro fades on triumphant fanfare with bright punchy mix, natural, cumbia
\end{style}
\begin{lyrics}
\songpart{Chorus}
Somos la sal de la tierra,\\
unidos bajo el mismo cielo,\\
fuertes como la piedra y el fuego,\\
guiados por la mano providente,\\
hermanos en la siembra y la cosecha:\\
¡solidaridad, escudo eterno del pueblo!

\songpart{Verse 1}
Bendita la montaña, nuestra sierra,\\
el llano indómito y ardiente,\\
y el ancho mar de espuma valiente,\\
dones que labra el pueblo en esta tierra.

\songpart{Chorus}
Somos la sal de la tierra,\\
unidos bajo el mismo cielo,\\
fuertes como la piedra y el fuego,\\
guiados por la mano providente,\\
hermanos en la siembra y la cosecha:\\
¡solidaridad, escudo eterno del pueblo!

\songpart{Verse 2}
Los padres fundadores, ya erguidos,\\
sellaron con el indio una alianza,\\
y el esclavo aportó con su esperanza\\
el suelo libre en que hoy hemos nacido.

\songpart{Chorus}
Somos la sal de la tierra,\\
unidos bajo el mismo cielo,\\
fuertes como la piedra y el fuego,\\
guiados por la mano providente,\\
hermanos en la siembra y la cosecha:\\
¡solidaridad, escudo eterno del pueblo!

\songpart{Verse 3}
El cóndor que nos guarda todavía,\\
la danta que en selva halla calor,\\
el chigüiro en el río con fulgor,\\
y el colibrí, alegre en la serranía.

\songpart{Chorus}
Somos la sal de la tierra,\\
unidos bajo el mismo cielo,\\
fuertes como la piedra y el fuego,\\
guiados por la mano providente,\\
hermanos en la siembra y la cosecha:\\
¡solidaridad, escudo eterno del pueblo!

\songpart{Verse 3}
Marchamos hacia una firme altura,\\
forjando el porvenir con unidad,\\
un solo pueblo en firme hermandad,\\
labrando una patria con bravura.

\songpart{Chorus}
Somos la sal de la tierra,\\
unidos bajo el mismo cielo,\\
fuertes como la piedra y el fuego,\\
guiados por la mano providente,\\
hermanos en la siembra y la cosecha:\\
¡solidaridad, escudo eterno del pueblo!

\songpart{Outro}
¡Adelante, hijos de esta tierra,\\
al deber que la patria reclama!\\
Que la unión sea nuestra bandera,\\
y la solidaridad, nuestra llama.
\end{lyrics}


\songtitle{Wingman}

\tag*{2026}\tag{trance}\tag{english-lyrics}\tag{original-lyrics}\label{sec:wingman}
\begin{notes}
Written by Carlos Thompson.

Companion piece to \emph{Soundwall}, retelling the same club-night narrative in a different arrangement and voice.
\end{notes}
\begin{style}
club hymn, trance, electronic, dubstep, mellow tenor  vocals, texture vocals
\end{style}
\begin{lyrics}
The music was loud,\\
a sound wall to feel\\
deep in your feet,\\
deep in your bones.\\
No room for talk—\\
none was needed.\\
Noise in the stones,\\
too close to a shock.

Eye candy galore,\\
my own wingman among\\
all the beauty in my eye\\
that I could admire.\\
Senses overloaded\\
that could turn me on.\\
Senses overloaded,\\
turning me on.

My wingman was all in.\\
She moved like she knew\\
what was good for her,\\
what was good for me.\\
Eye candy was not\\
what we were there for.\\
Just trying our luck,\\
hunting a mutual score.

They moved like\\
nothing else mattered,\\
connected by sound,\\
flattered by touch.\\
Dancing with them,\\
their moves, their skin\\
felt almost like sex\\
on the dancefloor rim.

You feel like a hunter,\\
you’re also the prey.\\
We know what we’re in for,\\
we know the game.\\
Sex is the goods,\\
sex is the coin,\\
the law of the hood\\
where all things join.

I noted it as a win —\\
a dance, a night,\\
a broken connection\\
just for the form.\\
My partner, my wingman,\\
my accomplice, my friend —\\
she was my greatest win\\
from that night out.
\end{lyrics}

\songtitle{Cuore In Guerra}

\tag*{2026}
\begin{notes}
Lyrics by SUNO based on a prompt by Carlos Thompson.
\end{notes}
\begin{style}
Opera with full symphonic orchestra and lyric soprano lead, slow-building 4/4 with surging strings, harp glissandi, brass swells, and timpani rolls; verse starts intimate with solo piano and low strings, pre-chorus opens into woodwinds and rising harmony, chorus blooms with full orchestra and soaring soprano sustains, Layered choral doubles on key phrases, emotional rubato pauses, reverb-rich hall mix, bright and cinematic yet aching, opera
\end{style}
\begin{lyrics}
\songpart{Verse 1}
Nella mente cade il peso\\
giro e giro, e non finisce\\
sogno dopo sogno, resto qui\\
con le mani piene di niente

Cerco il meglio, sempre oltre\\
ma il pensiero mi fa male\\
ogni passo è una ferita\\
ogni desiderio, sale

\songpart{Pre-Chorus}
E vorrei fermare il vento\\
che mi apre dentro e poi\\
mi promette cieli alti\\
e mi lascia qui, tra noi

\songpart{Chorus}
Mi fa male, mi fa male\\
questo troppo pensare\\
mi fa male, mi fa male\\
e non so più respirare
(ah) voglio amare senza correre
(ah) voglio vivere e basta\\
ma il cuore batte forte\\
e mi trascina ancora là

\songpart{Verse 2}
Immagini nel buio\\
come vetro sulla pelle\\
faccio e rifaccio il domani\\
ma il domani non si svela

E se spingo ancora oltre\\
sento il petto incrinarsi\\
come un'onda contro il marmo\\
che non smette di spezzarsi

\songpart{Pre-Chorus}
E vorrei fermare il vento\\
che mi chiama e poi si va\\
tutto ciò che sogno forte\\
mi consuma e non mi dà

\songpart{Chorus}
Mi fa male, mi fa male\\
questo troppo pensare\\
mi fa male, mi fa male\\
e non so più respirare
(ah) voglio amare senza correre
(ah) voglio vivere e basta\\
ma il cuore batte forte\\
e mi trascina ancora là

\songpart{Bridge}
Se chiudo gli occhi, resto\\
tra le punte del mio ardore\\
voglio pace, voglio spazio\\
voglio luce nel dolore

Basta promesse perfette\\
basta fame di domani\\
dammi un passo, dammi un nome\\
dammi pace tra le mani

\songpart{Final Chorus}
Mi fa male, mi fa male\\
questo troppo pensare\\
mi fa male, mi fa male\\
e non so più respirare
(ah) voglio amare senza correre
(ah) voglio vivere e basta\\
ma il cuore, piano, piano\\
impara a non lottare più
\end{lyrics}

\songtitle{Frost in Gold}

\tag*{2026}
\begin{notes}
Lyrics by SUNO based on a prompt by Carlos Thompson.
\end{notes}
\begin{style}
Symphonic pop anthem with brisk four-on-the-floor pulse, handclaps on the backbeat, soaring strings, and gospel choir lifts; verse starts intimate with piano and low strings, pre-chorus opens to claps and rising brass, chorus bursts with full orchestra and stacked choir answers, then a Gregorian bridge in modal chant before the final chorus returns bigger, Lead vocal is warm and close-mic with doubled hook lines, choir shouts on the refrain, bright cinematic mix with crisp transients, deep low end, and glossy reverb tails, gospel
\end{style}
\begin{lyrics}
\songpart{Verse 1}
Cold sun on my face\\
Steam on the sidewalk\\
Blue sky, white breath\\
Still, I’m alive

Mittens in my pocket\\
Coffee in my hand\\
The air bites softly\\
But I don’t mind

\songpart{Pre-Chorus}
See it shine\\
On the bare tree lines\\
Hear the church bells\\
In the morning wind

Heart so full\\
In the winter glow\\
You and me\\
We’re lit from within

\songpart{Chorus}
What a day, what a day\\
Cold sun, golden face\\
What a day, what a day\\
We dance in the frost

Lift it up, lift it up\\
Hands high to the sky\\
What a day, what a day\\
We’re warm in the light

(what a day) (what a day)
(cold sun, golden face)

\songpart{Verse 2}
Snow in the gutter\\
Melt on the curb\\
Red scarf spinning\\
Like a little flag

Your laugh hits hard\\
Through the frozen air\\
I see spring hiding\\
In your eyes

\songpart{Pre-Chorus}
See it shine\\
On the bare tree lines\\
Hear the church bells\\
In the morning wind

Heart so full\\
In the winter glow\\
You and me\\
We’re lit from within

\songpart{Chorus}
What a day, what a day\\
Cold sun, golden face\\
What a day, what a day\\
We dance in the frost

Lift it up, lift it up\\
Hands high to the sky\\
What a day, what a day\\
We’re warm in the light

(what a day) (what a day)
(cold sun, golden face)

\songpart{Bridge}
Kyrie, eleison\\
Lux in our bones\\
Kyrie, eleison\\
Winter turns to song

\songpart{Chorus}
What a day, what a day\\
Cold sun, golden face\\
What a day, what a day\\
We dance in the frost

Lift it up, lift it up\\
Hands high to the sky\\
What a day, what a day\\
We’re warm in the light

(what a day) (what a day)
(cold sun, golden face)
\end{lyrics}


\end{multicols*}

\chapter{Conlangs and Glossolalia}
\albumcover{covers/internal/conlangs-and-glossolalia}
Pure sound experiments: invented grammars, glossolalia that isn't trying to mean anything in particular, syllables chosen for how they feel in the mouth rather than what they translate to. The most purely musical chapter in the book, because there's nothing else to parse.
\clearpage\begin{multicols*}{4}
\songtitle{E Numa Lah}

\tag*{1996}\tag{experimental-pop}\tag{original-work}\tag{conlang-lyrics}

\begin{notes}
Original work.\\
Experimental composition with stylized, non-semantic vocals and abstract sound design.
\end{notes}

\begin{style}
Experimental alt-pop with syncopated handclaps and a rubbery bass pulse; verse stays sparse with dry percussion and a single looping synth motif, pre-chorus lifts with filtered vocal stacks and a rising tom pattern, chorus lands wider with doubled lead vocals and a chanty call-and-response on the hook. %
Add ear candy with reversed swells, quick tape-stop moments, and sparkling chimes between phrases. %
Bright, tight mix with a playful, slightly surreal edge.
\end{style}

\begin{lyrics}
\songpart{Verse 1}
E dem chi laim siman de yacken\\
e vis se lah direch vi, direch vi\\
numa lah e numa lah...\\
e numa lah!

\songpart{Pre-Chorus}
Laim de inkera bah\\
se in ox mint de baiya

\songpart{Chorus}
E numa lah\\
E numa lah\\
Limpter vi, limpter vii\\
e numa lah

E numa lah\\
E numa lah\\
Limpter vi, limpter vii\\
e numa lah

\songpart{Verse 2}
E dem chi laim siman de yacken\\
e vis se lah direch vi, direch vi\\
numa lah e numa lah...\\
e numa lah!

\songpart{Chorus}
E numa lah\\
E numa lah\\
Limpter vi, limpter vii\\
e numa lah

E numa lah\\
E numa lah\\
Limpter vi, limpter vii\\
e numa lah
\end{lyrics}


\songtitle{Kapu Lema}

\subsection{Notes}

\subsection{Style}
trance, female lead vocal, 132 BPM, four-on-the-floor kick, offbeat bass stabs, arpeggiated synth sequence, sawtooth lead, gated snare fills, supersaw chorus, sidechain pumping, stereo delay throws, vocal harmonizer, chant-like hook, euphoric lift, festival anthem

\begin{lyrics}
\songpart{Verse 1}
Kapu lema do fako mo lekade\\
Sumo mo chin chin do fako keput\\
Mamakutabo dun dun numa fade\\
Mamakutabo rekari malut

\songpart{Pre-Chorus}
Lema do fako mo rema\\
Tapika dun dun kabeto\\
Lilikamo mo pukema\\
Tapika dun dun feseto

\songpart{Chorus}
Rasaka mobu mo dak\\
Rasaka loku mo fade\\
Mamakutabo ferak\\
Lilikamo mo potade

\songpart{Verse 2}
Mikabo felato moraba chut\\
Rebi mo mokabe ba mo madiba\\
Tin tin tin buka mo fako balut\\
Kakasupano mo mamafariba

\songpart{Pre-Chorus}
Chin chin do fako keput\\
Mokabe ba mo madiba\\
Supino mo ratamut\\
Tapika dun dun rediba

\songpart{Chorus}
Rasaka mobu mo dak\\
Rasaka loku mo fade\\
Mamakutabo ferak\\
Lilikamo mo potade

\songpart{Bridge}
Mikabu mo dora\\
Momo fado kebato\\
But kakasupano\\
But mamakutabo\\
But lilikamo\\
Reputa mo faru

\songpart{Verse 3}
Raputa mo faru mo lilikamo\\
Mohaka burata monote bota\\
Shuba chubaki boboko babamo\\
Ku beto mo bata moka ralota

\songpart{Pre-Chorus}
Mo faru mo lilikamo\\
Mo bata moka ralota\\
Fashako mo baru gamo\\
Supano mo pakirota

\songpart{Chorus}
Rasaka mobu mo dak\\
Rasaka loku mo fade\\
Mamakutabo ferak\\
Lilikamo mo potade\\
Rasaka mobu mo dak\\
Rasaka loku mo fade\\
Rasaka...
\end{lyrics}


\songtitle{Badiba Buka}

\tag*{2026}\tag{folk-ballad}

\begin{notes}
\end{notes}

\begin{style}
Fast-paced tarantella-inspired folk ballad in 6/8 with rhythmic drive from dungchen, piwang, marimba, fiddle, male and female gaita, double bass, and timbales; opening hits with a sharp instrumental attack, then calmer verse passages anchored by low bass, pre-chorus swells into a rising crescendo, chorus opens into a bright instrumental break, and the outro strips instruments away one by one until only timbales remain, Wide, earthy, and kinetic mix, marimba, ballad
\end{style}

\begin{lyrics}
\songpart{Intro}
\annotation{brass and percussion attack with a Mambo flavor}
\annotation{fades into soft piano interlude}

\annotation{Spoken, baritone}\\
Badiba buka mo restakaká gut mo, le mig no barat ko be shaka.\\
Mo, keraba mo, lede barat mosha tik mo.\\
Le fikare re no mojave ki modo loga no barat shumon.

\annotation{Guitar riff}

\songpart{Verse 1}
\annotation{male tenor vocals}\\
Badiba buka mo restakaká gut mo\\
dubi ba dum mosh wang\\
Badabum bin dubidu shun shun\\
dubi ba dubi ba sang

\songpart{Pre-Chorus}
\annotation{epic delivery on a bass base}\\
keraba mo, lede barat mosha tik mo\\
rapakapun balo ni lang\\
ki modo loga no barat shumon\\
rapakapun kiki botang

\songpart{Chorus}
Badabum bin\\
dubidum shun shun\\
Badabum bin\\
dubi ba dubi ba sang (ba sang)

Badabum bin\\
dubidum shun shun\\
Badabum bin\\
rapakapun kiki botang (botang)

\songpart{Bridge}
\annotation{Goth metal, raised percussions and distorted guitars}
\annotation{deep guttural male vocals}\\
ku pash me kapatapun\\
kiki botang\\
kapatapun shumon\\
rapakapun mo

\annotation{Guitar riff}

\annotation{Piano interlude}

\songpart{Verse 2}
\annotation{male tenor vocals}\\
Badiba buka mo restakaká gut mo\\
dubi ba dum mosh wang\\
Badabum bin dubidu shun shun\\
dubi ba dubi ba sang

\songpart{Pre-Chorus}
\annotation{epic delivery on a bass base}\\
keraba mo, lede barat mosha tik mo\\
rapakapun balo ni lang\\
ki modo loga no barat shumon\\
rapakapun kiki botang

\songpart{Chorus}
Badabum bin\\
dubidum shun shun\\
Badabum bin\\
dubi ba dubi ba sang (ba sang)

Badabum bin\\
dubidum shun shun\\
Badabum bin\\
rapakapun kiki botang (botang)

\songpart{Outro}
\annotation{brass riff}

\annotation{Spoken}\\
Mo, keraba mo, lede barat mosha tik mo.
\end{lyrics}


\songtitle{Parará Pa Pa}

\tag*{2026}\tag{march}\tag{original-lyrics}
\begin{notes}
Written by Carlos Thompson.
\end{notes}
\begin{style}
march, strings, xylophone, tenor vocals
\end{style}
\begin{lyrics}
parará pa pa, parará pa pa\\
parará, parará, parará pa pa\\
pum pa pa pa pum\\
pum pa pa pa pum\\
parará, parará, parará pa pa

tururú tu tú, tururú tu tú\\
tururú, tururú, tururú tu tú\\
pum tu tu tu pum\\
pum tu tu tu pum\\
tururú, tururú, tururú tu tú

pin pin pin\\
pimpiririmpi, pim pim

pin pin pin\\
pirirí, pirirí, piririm pim pin

aaaaaah, pum!
\end{lyrics}


\songtitle{O ku pater de nos}

\tag*{2026}\tag{original-lyrics}
\begin{notes}
Adaptation of the Lord's Prayer into a mock romance Conlang, written by Carlos Thompson.
\end{notes}
\begin{style}
chamber ballad, female lead
\end{style}
\begin{lyrics}
O ku pater de nos\\
ki stas in ku tsiel\\
ki tus nom sta sanktu\\
ki tus relmo vene a nos\\
ki tus vole sta fatsa\\
en ku tsiel e en ku terra\\
don a nos a dia ku pan de nos\\
pardon a nos ku dolo de nos\\
tal ki nou pardona manki dola a nos\\
fats ne nos caputer en temptatio\\
liber nos ab ku mal\\
ke de tus sta ku relmo, ku poter e ku glore\\
amen
\end{lyrics}


\songtitle{Cheese Spoon Paseo}

\tag*{2026}\tag{trance}\tag{merengue}\tag{original-lyrics}
\begin{notes}
Written by Carlos Thompson.
\end{notes}
\begin{style}
electronic trance and merengue paseo mashup
\end{style}
\begin{lyrics}
\annotation{Trance Verse}\\
cheese-spoon, cheese-spoon, cheese-spoon, cheese-spoon

\annotation{Paseo Verse}\\
chookoo-chookoo, chookoo-chookoo, chookoo-chookoo, chookoo-choo

\annotation{Trance Verse}\\
cheese-spoon, cheese-spoon, cheese-spoon, cheese-spoon

\annotation{Paseo Verse}\\
chookoo-chookoo, chookoo-chookoo, chookoo-chookoo, chookoo-choo

\annotation{Trance Verse}\\
cheese-spoon, cheese-spoon, cheese-spoon, cheese-spoon

\annotation{Paseo Verse}\\
chookoo-chookoo, chookoo-chookoo, chookoo-chookoo, chookoo-choo
\end{lyrics}


\songtitle{dij Schdzjarvk}

\tag*{2026}\tag{original-lyrics}\tag{symphonic-pop}
\begin{notes}
Written by Carlos Thompson
\end{notes}
\begin{style}
symphonic pop, deep bass baritone vocals
\end{style}
\begin{lyrics}
\songpart{Verse 1}
Mochtun dij Schdzjarvk\\
Matughasch badij zjong\\
Repfing tsjarv muvtong\\
Rivagosch mutsj blarvk

\songpart{Pre-Chorus}
Rhurrukhamu la menunt\\
Dzjuzju schkhar lemuk

\songpart{Chorus}
dij Schdzjarvk, dij Schdzjarvk\\
Nung zjarouvk bidang\\
dij Schdzjarvk, dij Schdzjarvk\\
Lamper dij moung scha

\songpart{Verse 2}
Turv limimimi tsjugam\\
Maghakhing sung Schorum\\
sung Schorum rsavkum\\
dij Schdzjarvk sulimam

\songpart{Pre-Chorus}
Jaja dvuki Dvigitt metuu\\
Amburzj vorvk lasarut

\songpart{Chorus}
dij Schdzjarvk, dij Schdzjarvk\\
Nung zjarouvk bidang\\
dij Schdzjarvk, dij Schdzjarvk\\
Lamper dij moung scha

\songpart{Bridge}
Lua lua, betong\\
Schorum kebaka\\
remu pi rsunt uv\\
lili dvuki abung

Fa murka salong\\
rum tars pin rdukha\\
Putsjivaka lameruv\\
dij Schdzjarvk

\songpart{Final Chorus}
dij Schdzjarvk, dij Schdzjarvk\\
Nung zjarouvk bidang\\
dij Schdzjarvk, dij Schdzjarvk\\
Markh rdovik batom

Markh rdovik batom

dij Schdzjarvk
\end{lyrics}


\end{multicols*}

\part{Form and Craft}

\chapter{Meta}
\albumcover{covers/internal/meta}
Songs about making the songbook itself: the plot-twist mechanics of building a villain the audience will love before the reveal, and the private-jet cabin where a whole cast of characters first started taking shape. The author steps into frame, if only for two tracks.
\clearpage\begin{multicols*}{4}
\songtitle{Writing Ain't Easy}

\tag*{2026}\tag{hip-hop}\tag{original-lyrics}
\begin{notes}
Original lyrics by Carlos Thompson, about the craft of building a plot-twist story -- making a too-perfect hero's reveal as the villain land. Retrieved from Wyomee Dank's Suno page (V6-MINI, 13 September 2026).
\end{notes}
\begin{style}
Hip-hop, vinyl-cut breakbeat drums with chopped kicks and crisp snares, half-spoken rapid-fired vocal delivery, tight rhythmic phrasing, short plate reverb on lead and stacked chorus lines
\end{style}
\begin{lyrics}
\songpart{Verse 1}
So I was proving to myself I could write a story\\
with subverting expectations and a plot twist.\\
You announce the hero to the audience,\\
make him likable, make him worthy,

and deliver a villain you would detest —\\
not for villainy, not for being unlikely,\\
but because he is too perfect against the hero:\\
the villain's likability comes as undeserved.

\songpart{Pre-Chorus}
But subverting expectations, I said.\\
Dropping a plot twist, I said.

\songpart{Chorus}
All the ideas in your mind,\\
but writing ain't an easy task:\\
to build a world and a lore,\\
backstories and motivations,\\
a voice that would become\\
what you'd want to consume.

\songpart{Verse 2}
Misleading the audience into every trope\\
to make them root for the designed hero\\
and wish the designed villain loses at the end,\\
then reveal at the end: villain and hero were swapped.

You need a backstory to make sense —\\
that allows a reading that supports the mirrors\\
and a better reading that reveals the twist.\\
The inciting incident starts to take form:

\songpart{Pre-Chorus}
Why, if the hero was the villain, does he act as the hero?\\
And how to invest in the villainy of the other guy\\
in a way the reveal doesn't make him lame?

\songpart{Chorus}
All the ideas in your mind,\\
but writing ain't an easy task:\\
motivations that make sense,\\
twists that feel deserved,\\
a voice that would become\\
what you'd want to consume.

\songpart{Verse 3}
With main elements now in place,\\
time to write the story, to unfold the plan.\\
Let's choose a voice, and let's choose a structure,\\
and develop a plan from there on.

Let me start the story *in media res*:\\
our supposed hero at the bottom of his life,\\
his failed mission — and build from there,\\
rebuild his motivation, rebuild his spark.

\songpart{Pre-Chorus}
And that was the easy part.\\
The hardest is still to plan.

\songpart{Chorus}
All the ideas in your mind,\\
but writing ain't an easy task.\\
As writing progresses, ideas change,\\
and characters react to what's being set:\\
a voice that would become\\
what you'd want to consume.

\songpart{Bridge}
As the story progresses, the ties are tying up,\\
and then you realize: this guy would not do that.\\
Everything that sets in the plot,\\
the character's motivations are reframed.

\songpart{Final Chorus}
All the ideas in your mind,\\
but writing ain't an easy task —\\
a voice that would become\\
what you'd want to consume.

\songpart{Outro}
And the novel is finally unfolding.\\
Mostly good. Mostly okay.\\
But there's something still\\
that doesn't fit yet.
\end{lyrics}


\songtitle{Weather and Memory}

\tag*{2026}\tag{synth-pop}\tag{thomas-entourage}
\begin{notes}
Lyrics by SUNO from a Gemini Notebook prompt on the creative process behind \emph{Thomas' Entourage} -- a meta-reflection on world-building and character creation, and on how invented people outgrow their author. Retrieved from Dvigitt's Suno page (V6-MINI, 13 September 2026).
\end{notes}
\begin{style}
Sleek narrative synthpop, mid-tempo four-on-the-floor groove, pulsing analog bass, glassy arpeggiated synthesizers, crisp drum-machine percussion, warm electric piano, cinematic pads, polished 1980s-inspired production with modern low-end clarity, clear expressive female lead vocal, layered harmonies on the chorus, controlled vocoder accents, elegant and emotionally precise
\end{style}
\begin{lyrics}
\songpart{Verse 1}
I drew a face on tracing paper\\
Gave her a coat, a room, a name\\
Then wrote the year the harbor opened\\
And why she never rode the train

I built a sister out of questions\\
A careful laugh, a stubborn hand\\
She kept the keys to every border\\
And knew the pipes beneath the land

\songpart{Pre-Chorus}
One became a daughter, one became a guide\\
One kept a secret underneath the countryside

\songpart{Chorus}
I made them out of weather and memory\\
Gave them hearts that could decide\\
Set their tables, drew their histories\\
Let a whole world grow inside\\
Every promise, every fracture\\
Every song they learned to start\\
I was building more than characters\\
I was giving shape to human hearts

\songpart{Verse 2}
The first turbines turned in winter\\
The old rail lines became a spine\\
By year twenty, glass was growing\\
Where the factories had drawn the line

A council formed around the water\\
A child was raised beside the dam\\
Two women met beneath the sirens\\
And changed the future of the plan

\songpart{Pre-Chorus}
I gave them debts and favorite books\\
A scar, a room, a way to look

\songpart{Chorus}
I made them out of weather and memory\\
Gave them hearts that could decide\\
Set their tables, drew their histories\\
Let a whole world grow inside\\
Every promise, every fracture\\
Every song they learned to start\\
I was building more than characters\\
I was giving shape to human hearts

\songpart{Bridge}
Decades folded into blueprints\\
Roads appeared where fields had been\\
Every answer opened doorways\\
Every ending pulled me in\\
Then the women stepped together\\
Past the margins of the page\\
They were no longer mine to move through\\
They had written their own age

\songpart{Final Chorus}
They grew out of weather and memory\\
Found the nerve to choose their lives\\
Built their homes across the history\\
Kept the buried engines bright\\
Every promise, every fracture\\
Every hand that found a heart\\
Now the world keeps turning after me\\
And the songs know where they are

\songpart{Outro}
I made a place, they made it living\\
I wrote the first line, then let go\\
Somewhere under all the singing\\
Their unfinished stories grow
\end{lyrics}


\end{multicols*}

\chapter{Instrumental}
\albumcover{covers/internal/instrumental}
Where the words step back and the rhythm section takes over: tumbao, boogie, iron pulse, bambuco -- pieces that trust a groove to say what a lyric would only get in the way of.
\clearpage\begin{multicols*}{4}
\songtitle{Callejón Azul}

\tag*{2026}\tag{jazz-cumbia}\tag{instrumental}

\begin{notes}
Prompt rewritten by Suno based on original input by Carlos Thompson.
\end{notes}

\begin{prompt}
Jazz cumbia instrumental with bass and percussion locked to a cumbia pulse, but the harmonic language stays smoky and urban; keyboards and electric guitar carry a lyrical progression while a second guitar slips into free-flowing jazz improvisations. %
Verse-like sections keep the groove tight and restrained, pre-chorus lifts with brighter voicings and syncopated fills, then the main theme swells and briefly breaks into a more melancholy turnaround. %
Smooth instrumental swaps, conga pickups, brushed transitions, and a short guitar run before the final drop. %
Warm, spacious mix with a polished club-jazz sheen., jazz, tropical, melancholy, cumbia, uplifting, electric
\end{prompt}

\sepbar


\songtitle{Monstruo Ciudad}

\tag*{2026}\tag{instrumental}

\subsection{Notes}
Produced by Gemini.

\subsection{Prompt}
Cumbia with galloping offbeat and modern percussion, acoustic guitar driving a C-G-Am-F loop before a tritone turn into C-Am-F-G; electric bass locked underneath, sax and piano trading melodic lines. %
Intro stays instrumental, verse drops into deep bass lead, bridge shifts tension with the tritone, and the final section opens wider with stacked vocal doubles and short chant echoes. %
Bright, punchy mix with warm low end and lively horn sparkle., acoustic, cumbia, electric, deep

\sepbar


\songtitle{March and Tumbao}

\tag*{2026}\tag{instrumental}
\begin{notes}
\end{notes}

\begin{lyrics}
\end{lyrics}


\songtitle{Candy Rock Day}

\tag*{2026}
\begin{notes}
\end{notes}
\begin{style}
Children's musical theater, candy rock, 2020s production, hopeful, joyful, rainbows and unicorns theme, cotton candy, drum set, triangles, tambourine, xylophones, recorder, player piano, toy piano, 6/8 time signature, allegro tempo
\end{style}
\begin{lyrics}
\songpart{Intro}
(Upbeat toy piano and xylophone melody in 6/8 time)
(Triangle dings, tambourine shakes)

\songpart{Verse 1}
Wake up to a sky of pink and blue\\
The player piano’s playing something new\\
The recorder whistles a happy little tune\\
While we hop across the hills of the afternoon\\
With a skip and a jump on a marshmallow floor\\
We’re opening up the magic door!

\songpart{Chorus}
Oh, it’s a candy rock day under rainbow skies\\
With cotton candy clouds and unicorn eyes\\
Everything is sparkly, everything is bright\\
We’re dancing in the puddles of the morning light\\
With a ting on the triangle and a xylophone beat\\
Life is a treat that’s oh-so-sweet!

\songpart{Verse 2}
High-fives to the sun in the golden air\\
There’s sprinkles and glitter found everywhere\\
The drum set is thumping a joyful sound\\
As we twirl and we whirl on this merry-go-round\\
No clouds in our way, just a peppermint breeze\\
And the birds singing sugar-plum harmonies!

\songpart{Bridge}
(Instrumentation breakdown: Toy piano solo with rhythmic clapping)\\
One, two, three—jump for joy!\\
Every girl and every boy!
(Recorder and xylophone build up)\\
Faster now, let’s go!

\songpart{Chorus}
Oh, it’s a candy rock day under rainbow skies\\
With cotton candy clouds and unicorn eyes\\
Everything is sparkly, everything is bright\\
We’re dancing in the puddles of the morning light\\
With a ting on the triangle and a xylophone beat\\
Life is a treat that’s oh-so-sweet!

\songpart{Outro}
(Joyful tambourine roll)\\
It’s sweet! (Ting!)\\
So sweet! (Ting!)
(Final flourish on the player piano)
\annotation{End}\end{lyrics}


\songtitle{Clave and Chrome}

\tag*{2026}
\begin{notes}
\end{notes}
\begin{prompt}
Metal-tinged salsa rock with a 4/4 pulse, 2-3 clave, syncopated piano tumbao, and a soft C-G-Am-F harmonic bed; intro blasts with heavy trumpets and trombones, verses open into ballad strings, choruses slam with wah-wah guitars and distorted bass, bridge swells brass plus strings, then a final guitar-and-bass chorus rolls into a sax outro. Wide cinematic mix, punchy low end, bright horns, live-band energy., soft, electric, ballad, metal
\end{prompt}


\songtitle{Beachside Boogie}

\tag*{2026}\tag{funk}\tag{instrumental}
\begin{notes}
Instrumental.
\end{notes}
\begin{prompt}
Funk with dubstep drops and grand orchestral intro and outro; verse rides a melodic ballad feel with soft keys and tight bass, pre-chorus swells into a rising crescendo, chorus hits full funky syncopation with slap bass and brass stabs, then the drop switches to half-time wobble before snapping back, Lead vocal is intimate in verses, stacked and shouted on hooks, with group chants and short delay throws, Sea-breeze pads, reverse cymbal lifts, crowd claps, and a final string surge give it a bright cinematic sheen, ballad, funk, dubstep
\end{prompt}


\songtitle{Scraper Road}

\tag*{2026}
\begin{notes}
Instrumental.
\end{notes}
\begin{prompt}
Acoustic rock with mid-tempo four-on-the-floor drive, A-G-D verse cycles and a lift into C-Am-F-G on the bridge; intro opens on percussion solo, verses ride drum kit, acoustic guitars, hand-picked bass, and scrapper, pre-chorus narrows to one melodic guitar with fiddle crescendo, chorus blooms into full orchestra, bridge strips to percussion plus bass, outro finishes on drum-kit solo, Bright, earthy, punchy mix with roomy natural reverb and crisp transients, rock, acoustic
\end{prompt}


\songtitle{Iron Pulse}

\tag*{2026}
\begin{notes}
Instrumental.
\end{notes}
\begin{prompt}
EBM anthem with 808-driven industrial bounce, mechanical four-on-the-floor groove, and aggressive tempo lift from a slow 60 BPM intro to a 180 BPM peak; verse rides a stark F-major pulse, pre-chorus swells into a dense soundwall, chorus hits with cavernous reverb and tempo shifts, bridge strips to percussions only, and the outro dissolves into a fading wall of synth, Male vocals with layered doubles, shouted hook stacks, delay throws on key words, reverse swells into drops, metallic hits and pressure-riser lifts; bright, hard-edged, and club-cold mix, ebm
\end{prompt}


\songtitle{Stone Light Portrait}

\tag*{2026}\tag{indie-pop}\tag{instrumental}
\begin{notes}
Instrumental.
\end{notes}
\begin{prompt}
Cinematic indie-pop instrumental with a slow, swaying pulse and soft syncopated percussion; verse texture stays intimate with muted piano, brushed shakers, and airy pad haze, pre-drop strips to filtered keys and heartbeat kick, chorus opens with warm string swells and a glimmering lead motif, Breathy vocal chops, reverse swells, and sunlit chime accents, Wide, polished, photoreal mix with a golden glow, natural, soft, minimal, warm
\end{prompt}


\songtitle{Calle Tumbao}

\tag*{2026}
\begin{notes}
Instrumental.
\end{notes}
\begin{prompt}
Salsa tumbao with tightly interlocked drum kit, electric bass, and electric guitars locked to the 2-3 clave feel; verse rides a dry pocket with ghost-note snares and syncopated bass, chorus opens into montuno-style call and response between chopped guitar hits and bass answers, bridge strips to toms and bass then returns with stacked guitar stabs, Bright, gritty, dancefloor-close mix, salsa, electric
\end{prompt}


\songtitle{Claim Ledger Drift}

\tag*{2026}
\begin{notes}
Instrumental.
\end{notes}
\begin{prompt}
Organic ambient with slow breathy pulses, hand-played textures, and gently shifting drones in a patient 60 BPM flow; opening verse-like bed stays sparse with soft piano dust and field-recorded room tone, middle blooms into layered felt chords and bowed overtones, final section thins to a warm sub glow, Reversed swells, glassy chimes, tape hiss, and wide intimate mix, organic ambient
\end{prompt}


\songtitle{The Great Gatsby}

\tag*{2026}\tag{orchestral}\tag{instrumental}
\begin{notes}
Instrumental.
\end{notes}
\begin{prompt}
piano, strings, brass, and a harp, at a high-tempo 6/8 phrases in D-minor for verses with 4/4 F-major choruses, brass intro and outro, and a violin solo bridge, sporadic filling dummy words and replacives and gospel clapping
\end{prompt}

\songtitle{Migrant Echoes}

\tag*{2026}
\begin{notes}
\end{notes}
\begin{prompt}
Industrial rock with techno house urgency in the verses, post-punk blues in the bridge and outro; mid-tempo driving pulse with tape saturation, clipped claps, rubbery bass, metal-leaning riffs, chamber strings, smoky Hammond, bowed lows, filtered choir pads, gospel lift, and half-time wobble bass, Verse rides bleak 4/4 motion; chorus slams harder and wider; bridge drops into bluesy space with upright bass and fiddle, Close-mic narrative vocal with breathy fragments, stacked doubles on the hook, short delay throws, choir swells, and gritty plate reverb, Bright punchy mix, dark and atmospheric, natural, warm, slow, blues, techno, bright, emotional, industrial rock, gospel, moody, rhythmic, deep, metal, post-punk, low
\end{prompt}


\songtitle{Viejo Bambuco}

\tag*{2026}\tag{instrumental}
\begin{notes}
Instrumental.
\end{notes}
\begin{prompt}
Bambuco colombiano with lilting 6/8 sway, acoustic guitar arpeggios, tiple sparkle, and soft hand percussion; verse stays intimate with voice and strings, pre-chorus opens into a higher harmony line, chorus lands with warm group vocals and a brief flute reply, Occasional brush hits, string swells into transitions, and a dry close-mic lead for an earthy, nostalgic mix
\end{prompt}


\songtitle{Tambor y Gaita}

\tag*{2026}
\begin{notes}
\end{notes}
\begin{prompt}
Mapalé with relentless heavy percussion and bright brass hits, built on interlocking gaita macho and gaita hembra riffs over an urgent coastal groove; verse-like passages stay drum-forward and raw, then the arrangement opens into brass stabs and call-and-response flute figures, with break sections for tom rolls, conga bursts, and quick horn punches, Close, earthy mix with explosive low-end, airy reeds, and punchy top percussion
\end{prompt}


\songtitle{Dos Pesos}

\tag*{2026}
\begin{notes}
Instrumental. No lyrics.
\end{notes}
\begin{prompt}
Ranchera ballad in a steady 6/8 sway with guitarrón thump, twin acoustic guitars, and mournful trumpet lines; verse stays sparse and narrative, pre-chorus lifts with brushed snare and harmony turns, chorus lands on bold ranchera singback with group shouts, Lead vocal is warm and cracked, doubled on the hook, with short delay throws on the last word and a trumpet fill between lines, Dusty, wide, and emotional with a live-room glow, low, ranchera
\end{prompt}


\songtitle{The Last Waltz of the Northern}

\tag*{2026}\tag{symphonic-rock}\tag{instrumental}
\begin{notes}
Instrumental. No lyrics.
\end{notes}
\begin{prompt}
symphonic rock, 132 BPM, driving 4 on the floor, syncopated kick pattern, distorted guitar riffs, string ostinato, harp flourishes, hammered dulcimer runs, hurdy gurdy drone, nyckelharpa lead, orchestral swells, plate reverb, sidechain pumping, wide stereo mix, anthemic lift, bittersweet nostalgia
\end{prompt}


\songtitle{Moonlit Caravan}

\tag*{2026}\tag{orchestral}\tag{instrumental}
\begin{notes}
Instrumental. No lyrics.
\end{notes}
\begin{prompt}
Orchestral cinematic hybrid with a slow-burn crescendo, martial drums, bold brass swells, grand piano ostinatos, and a lyrical saxophone lead echoing medieval modes and oriental scales, Intro opens sparse with frame drum and distant piano; verses build on low strings and hand percussion; chorus blooms with full brass fanfares and soaring sax, Velvet-wide mix, dramatic and regal, with antique grit and modern impact, medieval
\end{prompt}


\songtitle{Grand Orchestra Ballad}

\tag*{2026}\tag{orchestral}\tag{instrumental}
\begin{notes}
Instrumental. No lyrics.
\end{notes}
\begin{style}
grand orchestra, mood music, congas, timbales, piano, strings, divisi, vibrato, flutes, clarinets, oboe, baritone sax, tenor sax, trumpets, trombones, French horn, upright bass, drum kit
\end{style}
\begin{lyrics}
\songpart{Intro}
\annotation{timbales and congas stravaganza - 4 beats}
\annotation{brass explossion - 4 beats}

\songpart{Verse 1}
\annotation{ballad stanza of piano and strings - 12 beats}

\songpart{Pre-chorus}
\annotation{piano and strings followed by jazz tenor sax - 8 beats}

\songpart{Chorus}
\annotation{jazz saxophones and brass - 8 beats}

\songpart{Verse 2}
\annotation{ballad stanza of piano and strings - 12 beats}

\songpart{Pre-chorus}
\annotation{piano and strings followed by jazz tenor sax - 8 beats}

\songpart{Chorus}
\annotation{jazz saxophones and brass - 8 beats}

\songpart{Bridge}
\annotation{woodwind ballad - 8 beats}

\songpart{Final Chorus}
\annotation{Full orchestra}

\songpart{Outro}
\annotation{strings fade}
\end{lyrics}


\songtitle{Timbales in Blue}

\tag*{2026}\tag{latin-jazz}\tag{instrumental}
\begin{notes}
\end{notes}
\begin{prompt}
1950s Latin-jazz big band with pre-rock saloon swagger, swung mid-tempo pulse, timbales intro into a crying sax lead; verse rides piano and brass call-and-response, pre-chorus tightens with percussion and silence hits, chorus opens to full orchestra walls, Bridge pushes bongos, cowbell, and guitars into an extended fading chorus, then a final full-band bang, Close-mic brass bite, roomy drums, wide cinematic mono-to-stereo warmth, with tape hiss, horn swells, and abrupt drops for drama, latin, jazz, rock, ballad
\end{prompt}


\songtitle{Purple Lace Stance}

\tag*{2023}\tag{instrumental}
\begin{notes}
Prompt by SUNO when attempted to produce an image of Gabi.
\end{notes}
\begin{prompt}
EBM with pounding four-on-the-floor kick, metallic syncopated bass pulses, and rigid hi-hat machine-gun patterns; verse holds tight on dry drum grid and clipped synth stabs, pre-drop filters the bass and adds tension risers, then the main section hits with distorted low-end, icy arps, and sharp accent hits, Sparse breathy vocal chops used as rhythmic texture, short delay throws on transition hits, sub swell into drops, Dark, polished, aggressive mix with cinematic width and a hard industrial edge, ebm
\end{prompt}


\songtitle{September at the Window}

\tag*{2026}
\begin{notes}
Lyrics by SUNO based on the opening of a failed story by Carlos Thompson.
\end{notes}
\begin{style}
Cinematic alternative folk ballad in a minor key, slow 6/8 pulse around  seventy BPM, fingerpicked acoustic guitar and low piano at the opening, distant brushed drums entering in the second verse, bowed cello swells in the chorus, gust-like room ambience and rain texture kept subtle behind the arrangement, Build from sparse and close-miked to broad but restrained, with a fragile intimate lead vocal, breathy phrasing, and a final chorus that opens emotionally without becoming glossy, Warm, dark mix with dampened high end, natural room reverb, and clear vocal presence
\end{style}
\begin{lyrics}
\songpart{Verse 1}
It was a September night\\
The kind that pushed against the door\\
Rain wrote crooked lines on glass\\
And gathered underneath the floor\\
You stood there in your wool coat\\
Water dark along the sleeve\\
Said, “I thought you’d still be up”\\
Like you had somewhere else to leave

\songpart{Chorus}
The wind kept worrying the house\\
The gutters filled and overflowed\\
You held your hands around a cup\\
As if the heat could make you whole\\
No thunder broke the silence\\
No promise asked us to remain\\
Just September at the window\\
And your name inside the rain

\songpart{Verse 2}
The kitchen clock was running fast\\
The kettle clicked and settled down\\
Your boots made little pools of earth\\
Across the old linoleum\\
You told me about the highway\\
About a motel near the lake\\
I watched a button on your coat\\
Shiver every time you moved away

\songpart{Chorus}
The wind kept worrying the house\\
The gutters filled and overflowed\\
You held your hands around a cup\\
As if the heat could make you whole\\
No thunder broke the silence\\
No promise asked us to remain\\
Just September at the window\\
And your name inside the rain

\songpart{Bridge}
By morning, all the fields were silver\\
The branches stripped and laid down low\\
You left before the bakery opened\\
I never heard the front door close

\songpart{Final Chorus}
The wind kept worrying the house\\
The gutters filled and overflowed\\
I held my hands around your cup\\
Though all its heat had left the room\\
No thunder broke the silence\\
No promise asked us to remain\\
Just September at the window\\
And your name fading with the rain

\songpart{Outro}
It was a September night\\
I still can hear it when it rains
\end{lyrics}

\end{multicols*}

\chapter{Fragments}
\albumcover{covers/internal/fragments}
The sketchbook before the painting: three short fragments that later grew into full songs found elsewhere in this collection. Kept here deliberately, unfinished, as evidence of the work before the work.
\clearpage\begin{multicols*}{4}
\songtitle{Burn the Glass}

\tag{gemini}\tag{social-media-seed}\tag{hard-rock}\tag{heavy-metal}

\begin{notes}
Lyrics and production by Gemini
\end{notes}

\begin{style}
Hard rock with heavy metal influences. %
Distorted electric guitars play palm-muted power chord riffs and syncopated rhythmic patterns. %
Aggressive male vocals utilize a gritty, shouted delivery. %
The drums feature a prominent snare with a sharp crack and a driving kick drum pattern. %
A distorted bass guitar follows the guitar riffs closely. %
The track features a high-gain guitar solo with rapid alternate picking and wide vibrato. %
145 BPM, 4/4 time signature, key of E minor.
\end{style}

\begin{lyrics}

\songpart{Intro}
\annotation{distorted electric guitar riff, driving drums, bass enters}

\songpart{Verse 1}
\annotation{gritty male vocals}\\
These empty pockets and this restless brain\\
Keep screaming out my name in vain\\
These comfort zones are shackles I can't break\\
A soul-crushing compromise I have to make\\
I'd burn this static life and hit the gas\\
And watch it all explode behind the glass\\
A different kind of life beyond the static\\
A different kind of life yes automatic

\songpart{Chorus}
\annotation{full band intensity, crash cymbals}\\
\annotation{unintelligible}

\annotation{Guitar Solo}\\
\annotation{high-gain distorted guitar solo, rapid picking, wah-wah effect}

\end{lyrics}


\songtitle{A Drunk Without A Drop}

\tag*{2026}\tag{indie-folk}\tag{ai-based}\tag{social-media-seed}\tag*{english-lyrics}

\begin{notes}
Lyrics and production by Gemini
\end{notes}

\begin{style}
Indie folk ballad featuring a baritone male vocal. %
The arrangement centers on a fingerpicked acoustic guitar playing a steady eighth-note pattern. %
A clean electric guitar provides melodic counterpoint with light reverb and occasional slides. %
The bass guitar plays sustained root notes on the downbeats. %
A subtle percussion layer consists of a soft shaker and a dampened kick drum on beats one and three. %
The tempo is 74 BPM in 4/4 time, in the key of C major.
\end{style}

\begin{lyrics}

\songpart{Intro}
\annotation{fingerpicked acoustic guitar, soft shaker}

\songpart{Verse 1}
\annotation{baritone male vocals}\\
I put off what I ought to do\\
To numb the growing dread \annotation{clean electric guitar enters with light reverb}\\
A heavy dose of doing naught\\
To calm my anxious head\\
But failure breeds a sharper fear\\
A trap I can't flee \annotation{bass guitar enters with sustained root notes}\\
I am a drunk without a drop\\
Undone by staying free \annotation{dampened kick drum on beats 1 and 3}

\end{lyrics}


\songtitle{The Bitter Toll}

\tag*{2026}\tag{pop-rock}\tag{ai-based}\tag{social-media-seed}\tag*{english-lyrics}

\begin{notes}
Lyrics and production by Gemini
\end{notes}

\begin{style}
Pop rock with elements of musical theater. %
Features a clean electric guitar playing palm-muted eighth-note arpeggios, a melodic electric bass, and a standard drum kit with a prominent snare. %
A baritone male vocal delivers a theatrical, narrative performance with clear diction. %
The arrangement builds with sustained electric guitar chords and driving eighth-note kick patterns. %
Key of D major, 145 BPM, 4/4 time signature.
\end{style}

\begin{lyrics}

\songpart{Intro}
\annotation{clean electric guitar arpeggio, palm-muted}

\songpart{Verse 1}
\annotation{baritone male vocals}\\
The hardest part of all of this, the bitter stinging pill \annotation{bass enters}\\
Is every hint I'm handed down is how I should fulfill\\
My duties to the world outside to ease their care and strain \annotation{drums enter, driving eighth-note kick}\\
While not a single word is meant to heal my hidden pain

\end{lyrics}


\end{multicols*}

\chapter{Adaptations I: Re-genres}
\albumcover{covers/internal/adaptations-i-re-genres}
Nothing about the words changes here -- only the arrangement. The same wrong map gets redrawn as trance and as EBM, the same trembling gets a vispop pulse, the same hatred gets a post-rock hush: proof that a lyric can survive almost any costume.
\clearpage\begin{multicols*}{4}
\songtitle{The Wrong Map}

\tag*{2026}\tag{alt-rock}\tag{original-lyrics}
\begin{notes}
Written by Carlos Thompson.
\end{notes}
\begin{style}
Alternative rock, mid-tempo driving rhythm, 125 BPM, male baritone vocal, emotionally restrained verses, explosive anthemic chorus, gritty guitars, dynamic vocal drops, tension-building Am-G-F-E progression, reflective but defiant mood
\end{style}
\begin{lyrics}
\songpart{Intro — instrumental, guitar riffs, drums display, 8 beats}

\songpart{Verse 1 — palm-muted guitars, steady drums, bass carrying the movement, vocal close to speaking rhythm}
There was a time I was younger and more naïve\\
and I believed the world functioned in simple ways:\\
school, college, marriage, kids, just to achieve\\
a normal ladder of expectations to live your days.

Choose the right major, choose the right career.\\
Your strengths would tell what your path becomes.\\
Nine to five, raise a family, go with friends and grab a beer.\\
Follow your dreams and don't worry about the outcomes.

\songpart{Pre-Chorus — open guitars, toms, raising percussion, vocal projection}
My path was easy from my background,\\
hard by the ironies of life.\\
I came from a privileged household\\
and a mind that became a strife.

\songpart{Chorus — full guitar distortion, restrained solo vocals, gang choir back and echoes}
Too perfect, almost a trap. (it was a trap)\\
A privileged environment,\\
a gifted intellect,\\
but I was given the wrong map. (with the wrong map)\\
A perfect experiment,\\
until all was wrecked
(it was a trap with the wrong map)

\songpart{Verse 2 — palm-muted guitars, steady drums, bass carrying the movement, vocal close to speaking rhythm}
There was a time I was younger and more naïve.\\
I knew life was a struggle for many; I didn't expect it for me.\\
Wrong decisions in hindsight, when to stay, when to leave;\\
and each stumbling block shaping me and collecting a fee.

School, college, marriage, kids. Been there, done that.\\
A normal path on paper, a constant struggle with my mind.\\
A mind that told me I wasn't good enough, no matter what,\\
achievements I didn't deserve, and every defeat was signed.

\songpart{Pre-Chorus — open guitars, toms, raising percussion, vocal projection}
You learn to distrust yourself\\
as adult life starts to kick in.\\
A genius kid without structure,\\
a maladapted human spin.

\songpart{Chorus — full guitar distortion, restrained solo vocals, gang choir back and echoes}
Too perfect, almost a trap. (it was a trap)\\
A privileged environment,\\
a gifted intellect,\\
but I was given the wrong map. (with the wrong map)\\
A perfect experiment,\\
until all was wrecked
(it was a trap with the wrong map)

\songpart{Bridge — faster, optimistic}
School, college, marriage, kids...\\
nine-to-five, grab a beer...\\
a career you love — what's so hard?\\
For a kid, the future is not fear.

You're smart, you'll figure out\\
your right path in what was to come.\\
I was shy, I was reserved,\\
not a big deal 'cause I wasn't dumb.

A stable family as a role model,\\
a perfect environment to thrive.
\annotation{vocal pitch drop}\\
Not their fault I messed up —\\
a perfect storm building in my life.

\songpart{Final Chorus — distorted guitars, restrained solo vocals, gang choir back and echoes}
Too perfect, almost a trap. (it was a trap)\\
A privileged environment,\\
a gifted intellect,\\
but I was given the wrong map. (with the wrong map)\\
A perfect experiment,\\
until all was wrecked

\annotation{conversational, hand-picked gitars}\\
School, college, marriage, kids...\\
nine-to-five, grab a beer...\\
I was younger and more naïve\\
with a plan that seemed quite clear.

\annotation{fade into silence}
\end{lyrics}


\input{temp/the_wrong_map_2}
\songtitle{They Multiply}

\tag*{2026}
\begin{notes}
Re-genred from \emph{Multiply the Delay}.
\end{notes}
\begin{style}
house pop, cinematic dubstep, female mezzo lead, refrained whispering breathy vocals, stacked vocal hooks, call and response backing vocals, texture vocals, sidechained synth bass, four on the floor kick, 808 sub pulses, dubstep drop reinforcements, half-time swing, pre-chorus lift, no-buildup drop, anthemic release, medium tempo, plate reverb, tape saturation, wide chorus harmonies
\end{style}
\begin{lyrics}
\songpart{Verse 1}
It don't just stack, baby\\
It multiply

One part won't push\\
It takes away the start

A signal go missing\\
Before the hand can move

Then the next one whisper
"Stay right where you are"

\songpart{Pre-Chorus}
And if you never begin\\
You never can fail

That stillness feel safer\\
Than stepping on the rail

\songpart{Chorus}
They multiply, multiply\\
They multiply, my baby\\
They multiply, multiply\\
That's how the whole thing stay

No one thing alone\\
It's the way they play

They multiply, multiply\\
And delay feels like a plan

\songpart{Verse 2}
Then comes the fear of losing\\
Before the game even starts

So not moving on the page\\
Feels lighter than the mark

The outside pressure fade out\\
Like rain on dry wood

No crowd in the doorway\\
To say you really should

\songpart{Pre-Chorus}
And if the world get quiet\\
The pattern stay warm

Each piece fit the next one\\
Like shelter from a storm

\songpart{Chorus}
They multiply, multiply\\
They multiply, my baby\\
They multiply, multiply\\
That's how the whole thing stay

No one thing alone\\
It's the way they play

They multiply, multiply\\
And delay feels like a plan

\songpart{Bridge}
Then the mind get clever\\
With a shiny little line

"Maybe this is strategy\\
Maybe this is timing"

It sounds so proper\\
When the reason can talk

So the pause wear a suit\\
And the symptom learn to walk

\songpart{Chorus}
They multiply, multiply\\
They multiply, my baby\\
They multiply, multiply\\
That's how the whole thing stay

No one thing alone\\
It's the way they play

They multiply, multiply\\
And delay feels like a plan
\end{lyrics}


\input{temp/deserve}
\input{temp/the_fault_in_your_character_ebd}
\songtitle{Vida simple}

\tag*{2026}\tag{indie-electronic}\tag{spanish-lyrics}
\begin{notes}
Genre adaptation of \emph{Dejándome llevar} as Indie Electronic.
\end{notes}
\begin{style}
Electronic indie with a pulsing offbeat kick, glassy arpeggios, muted sub-bass, and clipped handclap percussion; verse rides sparse synth stabs and close-mic vocal phrases, pre-chorus opens with filtered pads and rising vocal stacks, chorus hits with wide chorus-drenched keys and a sticky two-note hook, Breathy doubles, tight harmonies, occasional delay throws, reverse swells, and tiny glitch ticks between lines, Bright, glossy, slightly cold mix, electronic, romantic
\end{style}
\begin{lyrics}
\songpart{Verse 1}
Vi tu imagen en mi caminar\\
vi tu imagen en mi caminar\\
vi tu imagen en mi caminar\\
y mi mundo por ti cambió.

No importaba nada más\\
no importaba nada más\\
no importaba nada más\\
por el instante en que sucedió.

\songpart{Pre-Chorus}
Estabas feliz, te veías feliz\\
estabas feliz, te veías feliz\\
estabas feliz, te veías feliz\\
Y por ese instante feliz fui.

\songpart{Chorus}
La vida es simple cuando estoy\\
la vida es simple cuando estoy\\
la vida es simple cuando estoy\\
dejándome llevar del momento.

Pero el mundo sigue su ritmo\\
pero el mundo sigue su ritmo\\
pero el mundo sigue su ritmo\\
y al momento no me puedo asir.

\songpart{Verse 2}
Vi lo que dijiste de tu pérdida\\
vi lo que dijiste de tu pérdida\\
vi lo que dijiste de tu pérdida\\
y esa pérdida la sentí aquí.

Mi corazón está con tu tristeza\\
mi corazón está con tu tristeza\\
mi corazón está con tu tristeza\\
por el tiempo que estás ahí.

\songpart{Pre-Chorus}
Y estás ahí por dos minutos\\
y estás ahí por dos minutos\\
y estás ahí por dos minutos\\
antes de dar el siguiente scroll.

\songpart{Chorus}
La vida es simple cuando estoy\\
la vida es simple cuando estoy\\
la vida es simple cuando estoy\\
dejándome llevar del momento.

Pero el mundo sigue su ritmo\\
pero el mundo sigue su ritmo\\
pero el mundo sigue su ritmo\\
y al momento no me puedo asir.

\songpart{Bridge}
Es un gatito en desgracia\\
es un gatito en desgracia\\
es un gatito en desgracia\\
quien ocupa ahora tu lugar.

Es un video de autoayuda\\
es un video de autoayuda\\
es un video de autoayuda\\
O un recuerdo que el mundo está mal.

\songpart{Final Chorus}
La vida es simple cuando estoy\\
la vida es simple cuando estoy\\
la vida es simple cuando estoy\\
dejándome llevar del momento.

Pero el mundo sigue su ritmo\\
pero el mundo sigue su ritmo\\
pero el mundo sigue su ritmo\\
y al momento no me puedo asir.

El momento es todo pero pasa\\
el momento es todo pero pasa\\
el momento es todo pero pasa\\
y todo al final sigue igual.
\end{lyrics}


\songtitle{Los buenos somos más}

\tag*{2026}\tag{cumbia}\tag{spanish-lyrics}
\begin{notes}
Cumbia adaptation of \emph{Los buenos somos más}.
\end{notes}
\begin{style}
Symphonic cumbia, fast-paced vocals, spoken intro, baritone lead vocals, colombian pronunciation
\end{style}
\begin{lyrics}
\songpart{Intro}
En el estadio, la fiesta va\\
Un trago oculto, la cámara ya\\
¿Ingenio o vergüenza? No sé\\
No. Tampoco sé si quiero saber

\songpart{Verse 1}
Ahorraste para llevar a los tuyos\\
Un amigo te ofrece, tú aceptas el trago\\
Cantimplora en forma de binoculares\\
De paisa el nombre, pero hecha en China

El video vuela, alguien lo reparte\\
Tú te ríes primero, "qué vivo el paisa"\\
Después lo piensas, y ya no sabes\\
Y el trabajo llama: ya estás en la calle

\songpart{Pre-Chorus}
¿Quién es la víctima en esta historia?\\
¿Budweiser, que pierde un par de cervezas?\\
¿El honor nacional, que dizque se mancha?\\
No sé. Nadie sabe. Pero todos gritan

\songpart{Chorus}
Los buenos somos más\\
No.\\
Somos solidarios, somos tribales\\
Tribu de amigos, tribu de rivales\\
Linchamos al otro pa' sentirnos limpios\\
Los buenos somos más… cuando nos conviene

\songpart{Verse 2}
Un video más: le enseñas "español"\\
a unas japonesas, te ríes, grabas\\
Mañana pides tolerancia cero\\
Hoy la trampa te parece graciosa\\
Compartes la noticia sin mirarla bien
"Algo de verdad tendrá", te dices\\
Al rival lo atacas con saña\\
Al tuyo — a ese — se lo perdonas todo

\songpart{Pre-Chorus}
¿Quién es la víctima en esta historia?\\
¿El honor? ¿el estadio? ¿tu memoria?\\
Le decimos malicia indígena\\
Le decimos picaresca, y ya. Y ya

\songpart{Chorus}
Los buenos somos más\\
No.\\
Somos solidarios, somos tribales\\
Tribu de amigos, tribu de rivales\\
Linchamos al otro pa' sentirnos limpios\\
Los buenos somos más… cuando nos conviene

\songpart{Bridge}
No sacrifiques más chivos expiatorios\\
para lavar una culpa que es de todos\\
Colombiano de bien, colombiano de mal\\
son la misma persona en distinto minuto\\
No más caza de brujas en nombre del bien\\
que mañana el señalado puedes ser tú

\songpart{Chorus - final, more emotional}
Los buenos somos más, sí, lo sabemos\\
Pero sin autocrítica, nos perdemos\\
Deja de aplaudir la trampa ajena\\
y de justificar la tuya en la arena\\
Los buenos somos más\\
si nos atrevemos

\songpart{Outro}
En el estadio… o en la vida entera\\
la verdadera fiesta es ser mejor\\
Los buenos somos más…\\
Ojalá. Ojalá lo fuéramos de verdad
\end{lyrics}


\songtitle{Non-Aligned}

\tag*{2020}
\begin{notes}
Genre adaptation from \emph{No Middle Seat}
\end{notes}
\begin{style}
doo-wop, novelty, 124 BPM, spoken intro, stacked doo-wop harmonies, baritone lead vocal, handclap backbeat, upright bass walk, brushed snare, horn stabs, slapback echo, plate reverb, tape saturation, late-50s mono mix, carnival swing, defiant wit
\end{style}
\begin{lyrics}
\songpart{Verse 1}
I’m not centrist\\
I am non-aligned\\
I won’t bow to poles\\
That split my spine

My principles stand\\
On their own two feet\\
They don’t fit the box\\
You made for me

\songpart{Pre-Chorus}
Don’t ask me for\\
The just middle\\
If the middle’s\\
Just a riddle

If one side brings\\
A line of graves\\
I won’t call that\\
A thing I can save

\songpart{Chorus}
I am non-aligned (non-aligned)\\
I am non-aligned (non-aligned)\\
I choose no halfway\\
When harm is the plan\\
I take what is right\\
From where it stands

\songpart{Verse 2}
I’m not warm and soft\\
For the sake of peace\\
I won’t trade my name\\
For polite beliefs

I’ll take the good\\
From one hand’s reach\\
Then take the good\\
From the other side’s speech

But I won’t buy\\
The whole damn pack\\
With the poison kept\\
And the kindness wrapped

\songpart{Pre-Chorus}
Don’t ask me for\\
The just middle\\
If the middle’s\\
Just a riddle

If one side says\\
Let blood be law\\
I’m not divided\\
I’m against all

\songpart{Chorus}
I am non-aligned (non-aligned)\\
I am non-aligned (non-aligned)\\
I choose no halfway\\
When harm is the plan\\
I take what is right\\
From where it stands

\songpart{Bridge}
\annotation{Electronic swell}\\
Let the sirens fade\\
Let the labels drop\\
I am not your fence\\
I am not your prop

I will stand with truth\\
Where truth is clear\\
And I’ll name the line\\
When it comes near

\songpart{Chorus}
I am non-aligned (non-aligned)\\
I am non-aligned (non-aligned)\\
I choose no halfway\\
When harm is the plan\\
I take what is right\\
From where it stands

I am non-aligned (non-aligned)\\
I am non-aligned (non-aligned)\\
No warm little lie\\
No forced middleman\\
I take what is right\\
From where it stands
\end{lyrics}


\songtitle{Si no te hace falta}

\tag*{2026}\tag{techno-house}\tag{spanish-lyrics}
\begin{notes}
House adaptation of \emph{Piel} (page \pageref{sec:piel}). Adaptation by Carlos Thompson and Grok.
\end{notes}
\begin{style}
techno house, trance music, 808 sub-bass, distorted synthesizer stack, dubstep bass drop, close-mic baritone, almost spoken vocal, textured vocal layers, soundwall breakdowns, silence gaps, volume swells, machine-like percussion, driving 128 BPM, anxious desire
\end{style}
\begin{lyrics}
\songpart{Intro – Soft, breathed}
Si no te hace falta…\\
A mí sí.

\songpart{Verse 1 – Almost spoken, rhythmic flow}
Si no te hace falta, a mí sí.\\
Mi piel pide piel.\\
Mi cuerpo pide.\\
Pide descargas.\\
Pide estímulos que desconectan la mente.\\
Pide toque,\\
pide fricción,\\
pide excitación.

Hasta cierto punto me basto.\\
Hasta cierto punto.\\
Mi imaginación cumple.\\
El porno es suficiente.\\
El chatbot funciona.\\
He probado sexting y sirve…\\
Hasta cierto punto.

\annotation{Pre-Build – Layered, slowly rising tension}\\
Pero nada de eso tiene piel.\\
Nada de eso tiene piel.\\
Hasta cierto punto me basto…\\
Pero a partir de ahí,\\
no es suficiente.

\annotation{Drop 1 – Warm deep bass hits, vocals stay intimate}\\
Si no te hace falta, a mí sí.\\
Piel con piel, carne y hueso.\\
Si no te hace falta, a mí sí.\\
Algo real que esté ahí.

\songpart{Verse 2 – More introspective, spoken-like}
Amar realmente supongo que también sirve.\\
Supongo que hará falta también.\\
Amar realmente…\\
Lo que quiera que eso signifique.\\
Si es que significa algo.

Por ahora solo sé\\
que me falta piel.\\
Compartir con alguien\\
que respire cerca,\\
que esté ahí.\\
Que su peso hunda la cama.\\
Que su calor no sea una pantalla.

\annotation{Pre-Build 2 – Slow-burning tension}\\
Hasta cierto punto me basto…\\
Hasta cierto punto.\\
Pero el cuerpo sabe.\\
El cuerpo pide.\\
Y no acepta sustitutos.

\annotation{Final Drop – Emotional but restrained, voice floating over the groove}\\
Si no te hace falta, a mí sí.\\
Mi piel pide piel.\\
Mi cuerpo pide.\\
Pide descargas reales.\\
Pide estímulos que duelan un poco.\\
Si no te hace falta…\\
A mí sí.

\songpart{Outro – Fading layers}
Hasta cierto punto…\\
Hasta cierto punto…\\
A mí sí.\\
A mí sí.
\end{lyrics}


\songtitle{Crave Stimuli}

\tag*{2026}\tag{dubstep}\tag{trance}\tag{english-lyrics}
\begin{notes}
Trance adaptation of \emph{Stimuli} (page \pageref{sec:stimuli}). Adaptation by Carlos Thompson and Grok.
\end{notes}
\begin{style}
Trance-infused club house and dubstep with deep bass baritone male vocals, close-mic and almost spoken; verse rides a stripped four-on-the-floor pulse with sub pulses and texture breaths, pre-chorus opens into filtered lifts and clipped hand percussion, chorus drops into a bigger sidechained kick with wobble bass and chant doubles, Ad-lib stacks, delayed hook throws, reverse swells into each drop, metallic tick fills, and a glossy dark club mix with intimate front-and-center voice, dubstep, trance, world, warm, deep
\end{style}
\begin{lyrics}
\songpart{Verse 1}
Could be a conversation.\\
Could be dumb scrolling.\\
Could be creating a new world in my mind.\\
Could be sex for the sake of sex.\\
I crave stimuli.\\
Something to distract my mind.

\songpart{Pre-Chorus}
Something to cause my idle thoughts.\\
Warm skin against my skin\\
or plain intellectual stimulation.\\
Keep it coming closer now.\\
Keep it moving through my veins.\\
I want the edge of it.

\songpart{Chorus}
I crave stimuli.\\
Something new.\\
I crave stimuli.\\
Keep my mind distracted.\\
Keep my skin excited.\\
Keep my imagination checked.\\
I crave stimuli.\\
I crave.

\songpart{Verse 2}
Something old that haven't be with me\\
for time I haven't counted.\\
A new place.\\
A new partner.\\
A new challenge.\\
Or just more skin.\\
I let it pull me under.

\songpart{Pre-Chorus}
Something to cause my idle thoughts.\\
Warm skin against my skin\\
or plain intellectual stimulation.\\
Keep it coming closer now.\\
Keep it moving through my veins.\\
I want the edge of it.

\songpart{Chorus}
I crave stimuli.\\
Something new.\\
I crave stimuli.\\
Keep my mind distracted.\\
Keep my skin excited.\\
Keep my imagination checked.\\
I crave stimuli.\\
I crave.

\songpart{Bridge}
If I sit still, I start to hear\\
every little crack in me.\\
So give me a face, a hand, a sign.\\
Give me a stranger's energy.\\
Give me one more reason\\
to stay inside the moment.

\songpart{Final Chorus}
I crave stimuli.\\
Something new.\\
I crave stimuli.\\
Keep my mind distracted.\\
Keep my skin excited.\\
Keep my imagination checked.\\
I crave stimuli.\\
I crave.
\end{lyrics}


\songtitle{Crave Stimuli (rave)}

\tag*{2026}\tag{rave}\tag{stimuli}\tag*{english-lyrics}
\begin{notes}
Rave adaptation of \emph{Stimuli} (page \pageref{sec:stimuli}). Adaptation by Carlos Thompson and Grok.
\end{notes}
\begin{style}
rave, deep bass baritone vocals, close-mic spoken delivery, textured vocal layers, vocoder doubles, vocal delay throws, distorted synth bass, strobe-ready lead synths, filtered pad beds, vocal chops, four-on-the-floor kick, syncopated rave claps, sidechain pumping, sub-bass pressure, 124 BPM, pre-drop chant, massive drop, dark euphoria
\end{style}
\begin{lyrics}
\songpart{Intro – Atmospheric Rise}
(soft pads, distant vocal chops)\\
Stimuli...\\
I need it...\\
Stimuli...\\
Feed it...

\annotation{Build 1 – Tension Rising}\\
Scrolling through the glow, conversations in the night\\
Dumb thumb on the screen but my mind is taking flight\\
Building worlds inside, galaxies in my head\\
Or bodies colliding, no words left unsaid\\
Warm skin on warm skin, electric in the vein\\
Intellect or instinct, either one’s the same\\
Idle thoughts are screaming, silence is the pain\\
I need that rush right now to keep me sane!

\annotation{Pre-Drop Chant – Energy Climbing}\\
I CRAVE!\\
Stimuli!\\
I CRAVE!\\
Keep me high!\\
Distraction, reaction, skin-on-skin attraction\\
Keep my mind distracted, keep the fire active!

\annotation{Drop 1 – Heavy Bass}\\
I CRAVE STIMULI!
(Stimuli! Stimuli!)\\
New place, new face, new body, new high!\\
I CRAVE STIMULI!
(Warm skin! New sin!)\\
Old flames that I haven’t felt in time I haven’t counted\\
Drop it low, let it hit, keep my imagination mounted!

\songpart{Verse 2 – Second Rise}
Could be dumb scrolling, could be deep conversation\\
Could be sex for the sake of pure sensation\\
Creating new worlds where the bassline is the law\\
Losing track of hours, losing track of who I was\\
Something new, something raw, something I forgot I need\\
Old touch that I missed for eternity\\
Keep my skin excited, keep my thoughts on fire\\
Rave inside my veins, taking me higher!

\annotation{Pre-Drop 2 – Bigger Tension}\\
I CRAVE!\\
Stimuli!\\
I CRAVE!\\
Don’t let me die!\\
New partner, new challenge, more skin, more light\\
Keep the idle quiet in the strobe tonight!

\annotation{Final Drop – Massive}\\
I CRAVE STIMULI!
(Stimuli! Stimuli!)\\
Bass drop in my chest, skin drop on my body!\\
I CRAVE STIMULI!
(Louder! Harder!)\\
Keep my mind distracted, keep my body active\\
New world, new high, new everything — let it happen!
\songpart{Breakdown – Stripped Back}
(whispered over pulsing kick)\\
Stimuli...\\
Warm skin against my skin...\\
I crave...\\
I crave...

\annotation{Final Build \& Last Drop}\\
Something new!\\
Something old!\\
Something I haven’t felt in time I haven’t counted!\\
I CRAVE—

\annotation{Massive Final Drop}\\
STIMULI!\\
STIMULI!\\
Drop it, feel it, breathe it, need it!\\
I CRAVE STIMULI!
\end{lyrics}


\songtitle{Sex Without Compromise}

\tag*{2026}\tag{r-and-b}\tag*{english-lyrics}\tag{connection}
\begin{notes}
Indie R\&B arrangement of Carlos Thompson's \emph{Connection} (page \pageref{sec:connection}); lyrics unchanged, new style only.
\end{notes}
\begin{style}
indie R\&B, 92 BPM, male baritone vocals, velvet close-mic delivery, Rhodes stabs, muted nylon guitar, subby bassline, brushed snare, live hi-hat ghost notes, tape saturation, plate reverb, mono room vocal, intimate longing, shame-tinted desire, half-time groove, late-night soul, vintage 90s mix
\end{style}
\begin{lyrics}
\songpart{Verse 1}
Sex without compromise...\\
There are people who get paid for it.\\
Some are exploited:\\
engaging there would be rape.\\
Some are people for whom\\
sex without compromise\\
is not a big deal.\\
Why not get a profit from it?

\songpart{Pre-Chorus}
Sex without compromise.\\
Why would it be different\\
for the one who pays than for the one who gets paid?\\
Unless there is faul play.

\songpart{Chorus}
I crave connection (connection).\\
I crave skin (warm skin).\\
I crave a warm body in sync with mine.

I have been made ashamed of my desires.
(Ashamed of what I want.)\\
Not the right stand to defend.

\songpart{Verse 2}
Sex without compromise.\\
How is that different between\\
a person who engages in casual hookups\\
and a person who gets a profit?\\
Free agreements between consenting adults.\\
Sex without compromise,\\
beyond what you agree for that night.\\
A dinner, a fee, or mutual enjoyment for free.

\songpart{Pre-Chorus}
Sex without compromise.\\
I would be all in\\
for whatever agreement I could afford.\\
Which isn’t much.

\songpart{Chorus}
I crave connection (connection).\\
I crave skin (warm skin).\\
I crave a warm body in sync with mine.

I have been made ashamed of my desires.
(Ashamed of what I want.)\\
Not the right stand to defend.

\songpart{Bridge}
I don’t have much money.\\
I don’t have much time.\\
I am not a stallion\\
who would blow your mind.

I’m just a needy little old man\\
philosophizing on what I can’t afford.\\
Sex without compromise,\\
a fantasy of my own.

\songpart{Final Chorus}
I crave connection (connection).\\
I crave skin (warm skin).\\
I crave a warm body in sync with mine.

I have been made ashamed of my desires.\\
I have been made ashamed of what I want.\\
I have been made ashamed of a situation\\
in which I have paved myself in.\\
Not much money, not much time,\\
and not even begging for free.
\end{lyrics}


\songtitle{Living Fantasies (neo soul)}

\tag*{2026}\tag{neo-soul}\tag*{english-lyrics}\tag{living-fantasies}
\begin{notes}
Neo-soul adaptation of Carlos Thompson's \emph{Living Fantasies} (see \pageref{sec:living-fantasies}).
\end{notes}
\begin{style}
neo soul, chamber ballad, spoken intro, intimate piano, orchestral strings, woodwind countermelody, harp pulses, hammered dulcimer, hurdy gurdy drone, nyckelharpa, plate reverb, tape saturation, cathedral reverb, breathy female vocals, close-miked phrasing, bittersweet defiance, slow build, 78 BPM, swung backbeat
\end{style}
\begin{lyrics}
living fantasies\\
building them\\
an AI image prompt\\
an archetype

living fantasies\\
building them\\
alternate histories\\
alternate self

living fantasies\\
building them\\
non-existing languages\\
and polities

living fantasies\\
building them\\
bogota metro lines\\
for the enth time

living fantasies\\
building them\\
other planets\\
alien cultural forms

living fantasies\\
building them\\
then reinventing\\
fantasy itself

living fantasies\\
building them\\
just to feel something\\
just to numb myself
\end{lyrics}


\songtitle{Living Fantasies (native folk)}

\tag*{2026}\tag{folk}\tag{acapella}\tag{living-fantasies}
\begin{notes}
A capella folk chant adaptation of Carlos Thompson's \emph{Living Fantasies} (see \pageref{sec:living-fantasies}).
\end{notes}
\begin{style}
A cappella folk chant with breathy female lead, steady pulse from hand taps and foot stomps, verse sections keep it intimate and close-mic while the chorus opens into stacked unison harmonies; pre-chorus narrows to a whispered rise, then chorus repeats the place and time with a chant-like lift; small breath catches, echo tails on key words, and a raw roomy mix that feels earthy and immediate, native american, breathy
\end{style}
\begin{lyrics}
living fantasies\\
building them\\
an AI image prompt\\
an archetype

living fantasies\\
building them\\
alternate histories\\
alternate self

living fantasies\\
building them\\
non-existing languages\\
and polities

living fantasies\\
building them\\
bogota metro lines\\
for the enth time

living fantasies\\
building them\\
other planets\\
alien cultural forms

living fantasies\\
building them\\
then reinventing\\
fantasy itself

living fantasies\\
building them\\
just to feel something\\
just to numb myself
\end{lyrics}


\songtitle{Skälvning}

\tag*{2010}\tag{vispop}\tag{swedish-lyrics}
\begin{notes}
Electronic vispop re-genre of \emph{Skälvning} (the nu-metal version), retaining the same Swedish lyrics -- lyrics based on a work by Carlos Thompson from 2010 (posted under Enwigh Peady). Retrieved from rataflechera's Suno page (V6, 14 September 2026).
\end{notes}
\begin{style}
Electronic vispop, 124 BPM minor-key syncopated pulse, female processed vocals with Swedish spoken intro, rapped spoken verse, half-whispered chorus, distant echoed whisper and cracked voice, tape saturation, plate reverb, sidechain pumping, wide stereo synths, gated snare, syncopated kick and bassline, distorted synth stabs, glitch percussion, arpeggiated pluck lead, breakdown silence
\end{style}
\begin{lyrics}
\songpart{Intro – distant, echoed whisper}
Jag känner hjärtat...\\
förbjudna slag...\\
förbjudna slag...

\songpart{Verse 1 – rapped, syncopated, rhythmic, almost monotone}
Ett hjärta i bröstet – förbjudna slag,\\
ett liv som löser upp sig, dag för dag.\\
Det jag tog för givet – borta, spårlöst,\\
jag gasar, jag fejkar, jag blir ett spöke.

Jag leker att vara nån annan –\\
krossar mina kedjor med skakande händer.\\
Lockas av honung – grön som havet,\\
dras mot flamman, mot min gravplats.

\annotation{Build}\\
Min bödel...

Mitt hjärta–\\
Min amulett–\\
Och min bödel...

\annotation{Drop}\\
Och Jag Skälver!\\
Jag Gasar – Och Skälver!\\
Vet Inte Vad Jag Spelar För Spel!\\
Jag Förlorar Oskuld, Förnuft – Mitt Liv I En Källa!

Jag Skälver För Att Jag Inte Kan Vara Tydlig –\\
Inte Mot Min Motståndare, Inte Mot Läktaren,\\
Inte Mot Mig Själv!

\songpart{Verse 2 – softer, half-whispered, almost pleading}
Min livsflod rinner ut –\\
i förbjudna slag från mitt bröst.\\
Allt faller sönder – jag ser det,\\
mina fingrar krampar över tangenterna.

Ska jag säga det jag inte vågar?\\
Ska jag höra hennes röst – eller se ansiktet\\
på det spel jag spelar?\\
Sanningen: jag vågar inte klippa kedjorna.

\annotation{Build}\\
Min bödel...

Mitt hjärta–\\
Min amulett–\\
Och min bödel...

\annotation{Drop}\\
Och Jag Skälver!\\
Jag Gasar – Och Skälver!\\
Vet Inte Vad Jag Spelar För Spel!\\
Jag Förlorar Oskuld, Förnuft – Mitt Liv I En Källa!

Jag Skälver För Att Jag Inte Kan Vara Tydlig –\\
Inte Mot Min Motståndare, Inte Mot Läktaren,\\
Inte Mot Mig Själv!

\songpart{Breakdown – Spoken word, intimate, Swedish, into the mic}

Livets spel...\\
Jag förlorar det redan.\\
För att jag inte är jag –\\
och inte heller mitt spöke.

\annotation{Silence.}

\annotation{Whisper, cracked voice:}
...jag vågar inte...

\songpart{Outro}
Skälver!\\
Skälver!\\
Allt – Faller – Sönder!

\annotation{Final, broken whisper}
...inte jag... inte mitt spöke...

\annotation{Ring out. Fade.}
\end{lyrics}


\songtitle{The Decay of Hatred}

\tag*{1988}\tag*{2026}\tag{post-rock}\tag{je-hais-1988}
\begin{notes}
Post-rock / dark wave re-genre of \emph{Used to Say}, retaining the same lyrics -- the original poem \emph{Odio} by Carlos Thompson from 1988, with 2026 lyrics reflecting on it. Retrieved from Dvigitt's Suno page (V6-MINI, 13 September 2026).
\end{notes}
\begin{style}
Post-Rock / Dark Wave, cool measured near-spoken vocals with restrained layered doubles, atmospheric synth pads and sustained electric-guitar swells, slow hypnotic pulse with gradual ambient builds and stark drop sections, wide stereo synth layering with deep reverb tails and controlled low-end
\end{style}
\begin{lyrics}
\songpart{Intro, whispered, spoken}
Odio lo que siento por ti\\
y odio lo que me haces sentir.\\
Odio el amarte y me odio por eso.\\
Odio el odiarte y me odio por eso.\\
Odio lo que hago por ti\\
y odio lo que no hago por ti\\
y odio todo mi odio.\\
Tanto odio lo que odio\\
que termino odiando todo.

\annotation{Build 1}\\
I used to say how much I hated myself,\\
how I hated the object of my unrequited love.\\
Years later I couldn't make myself feel\\
that hate is something I could harbor.

Is hate a part of a normal human life,\\
one of those conditions I'm deficient at,\\
or have I realized hate is just something\\
that doesn't do much for my sanity?

\annotation{Drop}\\
I can't hate, I can't hate,\\
a feeling not in my heart,\\
a feeling not in my character.

I used to say I hated—\\
what changed?

\annotation{Build 2}\\
Or maybe\\
I have just redefined hate.\\
I called hate just any discomfort;\\
now I define a threshold\\
that I can't reach.

Or is it my disconnect with humanity,\\
a humanity I constantly fall in love with\\
but that I can't fully comprehend,\\
not fully love,\\
not fully hate.

\annotation{Drop}\\
I can't hate, I can't hate,\\
a feeling not in my heart,\\
a feeling not in my character.

I used to say I hated—\\
what changed?

\songpart{Outro}
Not fully love.\\
Nor fully hate.
\end{lyrics}


\songtitle{The Unreached Crush}

\tag*{2026}\tag{indietronica}\tag{a-coffee-not-shared}
\begin{notes}
Indietronica / minimalist alt-R\&B re-genre of \emph{A Coffee Not Shared}, retaining the same lyrics. Retrieved from Dvigitt's Suno page (V6-MINI, 13 September 2026).
\end{notes}
\begin{style}
Indietronica / Minimalist Alt-R\&B; echoing vocal delays, close dry details and restrained stereo space; minimalist synth arpeggios with soft rhythmic glitches and a sharp distorted instrumental swell at the Breakdown; intimate close-miked breathy English vocal, nearly whispered, with subtle doubles; unhurried syncopated pulse and non-traditional sectional flow
\end{style}
\begin{lyrics}
\songpart{Verse 1}
You (you)\\
you've been a crush since I remember.

You (you)\\
you've been hitting my timeline with your presence.

You (you)\\
you're a coffee I've failed to ask you (you).

You (you)\\
tender words and engaging thoughts.

\annotation{Instrumental break}

\songpart{Verse 2}
You (you)\\
years of longing and falling in love.

You (you)\\
you're radiant and I'm a deer in headlights.

You (you)\\
many words I'd love to share with you (you).

You (you)\\
a guilty pleasure; my shameless thoughts.

\songpart{Breakdown}
I've been longing for eons,\\
falling in love with your presence,\\
failing to feel deserving to reach you (you).
'Cause it's easier abstinence,
'cause cowardice feels safer than trying,
'cause you (you)\\
would destroy me.

\songpart{Outro}
You (you)\\
a story not happening.

You (you)\\
a coffee not shared.

You (you)\\
a suspended conversation.

You (you)\\
headlights that'll crush me.
\end{lyrics}


\songtitle{Domestic Realism}

\tag*{2026}\tag{indie-pop}\tag{spanish-lyrics}\tag{mundos-internos}
\begin{notes}
Indie-pop re-genre of \emph{Mundos internos}, retaining the same lyrics. Retrieved from Gabriel's Suno page (V6-MINI, 13 September 2026).
\end{notes}
\begin{style}
Indie Pop / Modern Spanish Synth-Pop, layered electronic production with field-recording ambiance and subtle tape saturation, upbeat driving pulse, steady bassline and tight programmed drums, glassy synth arpeggios and warm pads, conversational rapid-fire verses shifting into a melodically sweeping chorus, close-mic Spanish vocals with selective doubles and widened final refrain
\end{style}
\begin{lyrics}
\songpart{Verse 1}
Un espresso de mi máquina de capuchino,\\
con leche caliente, cual si fuera un latte,\\
rollos de canela de supermercado barato\\
y las pastillas de la tensión y de la cabeza.

Leer y contestar preguntas en Quora\\
o perderme en Instagram en vídeos cortos,\\
perder el sentido del tiempo y correr luego,\\
salir a la calle a esperar en el paradero.

\songpart{Pre-Chorus}
Una ciudad con sol o con lluvia,\\
con vías en construcción y vida en la calle,\\
motos que serpentean y bicicletas que fluyen,\\
peatones y carros, buses, ruido y vida.

\songpart{Chorus}
Y en mi cabeza, otros mundos y otros lugares:\\
un podcast, música o mis propios mundos,\\
observar la calle o mi propio teléfono,\\
lo que pasa afuera y en mi mundo interno.

\songpart{Verse 2}
La rutina en la oficina, que rara vez es rutina,\\
justificar horas que no pasan en sincronía,\\
justificando una rutina y una presencia\\
con máscaras y alt-tabs, raras sinfonías.

Hora de almuerzo, compañeros de oficina,\\
cada uno con su presupuesto y sus historias,\\
el trabajo en común o la última noticia,\\
lo que transcurre en nuestras vidas notorias.

\songpart{Pre-Chorus}
Una ciudad con sol o con lluvia,\\
centros comerciales, tiendas de esquina,\\
salir al almuerzo o microondas en la oficina,\\
y tomarnos un café como parte de la rutina.

\songpart{Chorus}
Y en mi cabeza, otros mundos y otros lugares:\\
un podcast, música o mis propios mundos,\\
observar la calle o mi propio teléfono,\\
lo que pasa afuera y en mi mundo interno.

\songpart{Bridge}
La vida transcurre,\\
la vida pasa,\\
hay otros que cruzan nuestros lugares,\\
trabajo que hacer,\\
historias que contar,\\
coordinar, vivir:\\
un mundo en que no estás solo.

\songpart{Breakdown}
Lo que haces y lo que no haces,\\
compañeros de trabajo que de ti dependen,\\
una familia que espera,\\
una vida que pasa.

\songpart{Final Chorus}
Y en mi cabeza, otros mundos y otros lugares:\\
un podcast, música o mis propios mundos,\\
observar la calle o mi propio teléfono,\\
lo que pasa afuera y en mi mundo interno.

Un mundo interno.
\end{lyrics}


\end{multicols*}

\chapter{Adaptations II: Adapted Lyrics}
\albumcover{covers/internal/adaptations-ii-adapted-lyrics}
Here the words move too: the same toll paid in different phrasing, the same grinder-year paper reworked line by line, a plot-twist confession tightened into ballad meter by a different hand entirely. The theme survives; almost nothing else does verbatim.
\clearpage\begin{multicols*}{4}
\songtitle{The Hardest Toll}

\tag*{2026}\tag{symphonic-metal}
\begin{notes}
Adaptation of \emph{The Fault in Your Character} into symphonic metal without specific style guidance for stanzas.
\end{notes}
\begin{style}
symphonic metal
\end{style}
\begin{lyrics}
\songpart{Intro}
Sometimes you go through life judging other people for their lack of ambition,\\
while ignoring the struggles they face from circumstances beyond their control
—a system that has failed them,\\
or a mind not wired for the conformity society demands.

\songpart{Verse 1}
They tell the mind that thinks askew,\\
It fell by its own hand,\\
A weight of fault that's strictly due,\\
You failed to understand.

They spin a tale of quiet blame,\\
That you alone have spun,\\
To shield the world from any shame,\\
For battles you've not won.

\songpart{Pre-Chorus}
So bow your head and bear the cost,\\
The choice was yours to make;\\
For every light that you have lost,\\
Was earned by your mistake.

\songpart{Chorus}
(The hardest toll)
(The hardest toll)\\
The hardest toll is that they convince you\\
that the fault is in your character
(a fault in your character)\\
that your fate is proof of moral failure\\
and don't deserve the good that comes to you
(you are a moral failure, deserving punishment, deserving no joy)

You have to conform. You have to comply. You have to get better
(conform, comply, get better)\\
You have to stop being a burden,\\
take your meds,\\
become productive,\\
or get out.
(be productive or get out)

\songpart{Verse 2}
The world demands a steady beat,\\
A march you cannot keep;\\
But find the rhythm in your feet,\\
While others drift in sleep.

You need not break to fit the mold,\\
But bend it to your will;\\
With secrets that your senses hold,\\
You'll bend the iron, still.

\songpart{Pre-Chorus}
So hack the path through noise and strife,\\
To serve both heart and kin;\\
And carve a custom, balanced life,\\
Where all your joys can win.

\songpart{Chorus}
(The hardest toll)\\
The hardest toll is that they convince you\\
that the fault is in your character
(a fault in your character)\\
that your fate is proof of moral failure\\
and don't deserve the good that comes to you
(you are a moral failure, deserving punishment, deserving no joy)

You have to conform. You have to comply. You have to get better
(conform, comply, get better)\\
You have to stop being a burden,\\
take your meds,\\
become productive,\\
or get out.
(be productive or get out)

\songpart{Bridge}
The goal is not to wear a mask,\\
Or shrink within a narrow frame,\\
But rise to meet a greater task\\
And give your hidden spark a name.\\
To seek the self you truly hold,\\
And let your honest spirit bloom,\\
To help the world turn brave and bold,\\
And light the corners of the room.

\songpart{Final Chorus}
The hardest toll is that they convince you\\
that the fault is in your character
(a fault in your character)\\
that your fate is proof of moral failure\\
and don't deserve the good that comes to you
(you are a moral failure, deserving punishment, deserving no joy)

(conform, comply, get better)

(be productive or get out)
(or get out)

\songpart{Outro}
or resist, resist, resist

you deserve better than society told you\\
you deserve.
\end{lyrics}


\songtitle{Piel / Stimuli}

\tag*{2026}\tag{reggaeton}\tag{stimuli}\tag{spanish-lyrics}
\begin{notes}
Adapted lyrics mixing \emph{Stimuli} (page \pageref{sec:stimuli}) and \emph{Piel} (page \pageref{sec:piel}).
Adaptation by Carlos Thompson and Grok.
\end{notes}
\begin{style}
Mid-tempo seductive reggaeton with dembow, deep bass male vocals, warm synths and perreo-ready grooves, Vocals are smooth, textured and almost spoken — intimate, not shouted
\end{style}
\begin{lyrics}
\songpart{Intro – Soft, half-spoken}
Si no te hace falta…\\
A mí sí.\\
I crave stimuli…

\songpart{Chorus – Catchy, bilingual hook}
Si no te hace falta, a mí sí,\\
mi piel pide piel, baby, ven aquí.\\
Warm skin on my skin, I need that real feel,\\
hasta cierto punto me basto, pero después… no es enough.\\
I crave stimuli,\\
pide toque, pide fricción,\\
pide descargas que me desconecten el mind.\\
Si no te hace falta, a mí sí.

\songpart{Verse 1 – Smooth reggaeton flow}
Scrolling en la noche, creando mundos en mi head,\\
o sex for the sake of sex, sin amor, solo flesh.\\
Chatbot responde, sexting también sirve,\\
pero nada de eso tiene piel, eso es lo que me mata.\\
Hasta cierto punto mi imaginación cumple,\\
porno, toys, fantasías… todo funciona un rato.\\
Pero a partir de ahí, se me rompe el circuito,\\
necesito cuerpo real que me deje sin aliento.

\songpart{Pre-Chorus – Building tension}
Mi cuerpo pide…\\
Pide estímulos, pide calor,\\
pide alguien de carne y hueso que esté aquí, no en el phone.\\
Amar realmente supongo que también ayuda,\\
pero por ahora solo sé que me falta piel, pura piel.

\songpart{Chorus – Full Hook}
Si no te hace falta, a mí sí,\\
mi piel pide piel, baby, ven aquí.\\
Warm skin on my skin, I need that real feel,\\
hasta cierto punto me basto, pero después… no es enough.\\
I crave stimuli,\\
pide toque, pide fricción,\\
pide descargas que me desconecten el mind.\\
Si no te hace falta, a mí sí.

\songpart{Verse 2 – More intimate, rhythmic}
New place, new partner, new challenge in the dark,\\
o el mismo toque antiguo que hace tiempo no he probado ya.\\
Something old, something new, something that makes me lose control,\\
pero tiene que ser real, no pixels, no voz en el audio.\\
I crave stimuli to keep my idle thoughts away,\\
pero al final del día, lo que pide mi cuerpo es stay.\\
Que su peso esté encima, que su aliento me caliente,\\
porque hasta cierto punto… me basto,\\
pero te quiero presente.

\songpart{Bridge – Half-spoken, slower dembow}
Amar realmente… lo que sea que eso signifique.\\
Supongo que hará falta.\\
Por ahora solo sé que me falta piel.\\
Compartir con alguien que respire, que tiemble, que esté ahí…

\songpart{Final Chorus – Bigger but still contained}
Si no te hace falta, a mí sí,\\
mi piel pide piel, no me dejes así.\\
I crave stimuli, warm skin, real high,\\
hasta cierto punto me basto… pero contigo quiero más.\\
Si no te hace falta, a mí sí…\\
A mí sí.

\songpart{Outro – Fading echoes}
Piel… stimuli…\\
Hasta cierto punto…\\
A mí sí.
\end{lyrics}


\songtitle{Soundwall}

\tag*{2026}\tag{eurodance}\tag{wingman}\tag*{english-lyrics}
\begin{notes}
Eurodance adaptation of \emph{Wingman} (page \pageref{sec:wingman}). Adaptation by Carlos Thompson and Grok.

Companion piece to \emph{Wingman}, retelling the same club-night narrative in a different arrangement and voice.
\end{notes}
\begin{style}
eurodance anthem, 128 BPM, four on the floor, syncopated offbeats, half-time breakdown, spoken intro, male rap verses, big female hook, stacked backing vocals, supersaw stabs, analog bassline, gated snare, brass samples, tape saturation, plate reverb, peak-time energy, hedonistic swagger
\end{style}
\begin{lyrics}
\songpart{Intro — spoken or half-sung, building synths}
Music was loud… feel it in your feet…\\
No talk… just the beat…

\songpart{Verse 1 — Male rap / rhythmic spoken}
Music was loud, bass hit me hard\\
Soundwall shaking right through my heart\\
No place for words, no talk was needed\\
Noise in the stones, we’re getting heated\\
Eye candy everywhere, flashing lights\\
My wingman beside me, own the night\\
Senses overload, can’t fight the fire\\
She moves like she knows what I desire

\songpart{Pre-Chorus}
Turning me on… turning me on…\\
All night long… all night long…

\songpart{Chorus — Big female vocals + male backing – the anthem part}
Senses overload! Feel it in your bones!\\
Senses overload! Shaking to the core!\\
We’re dancing on the edge, almost like sex\\
Hunter and the prey, no time for regrets!\\
Soundwall! Take me higher!\\
Soundwall! Set this floor on fire!

\songpart{Verse 2 — Male rap}
She’s my wingman, my accomplice, my friend\\
Moving like she knows how this night will end\\
We came for the thrill, not just the show\\
Trying our luck, hunting the afterglow\\
They dance like nothing else matters tonight\\
Skin on skin, under the flashing lights\\
Sex is the coin, this is the law of the hood\\
Where everything joins, and it feels so good

\songpart{Chorus — Repeat – bigger this time}
Senses overload! Feel it in your bones!\\
Senses overload! Shaking to the core!\\
We’re dancing on the edge, almost like sex\\
Hunter and the prey, no time for regrets!\\
Soundwall! Take me higher!\\
Soundwall! Set this floor on fire!

\songpart{Bridge — Breakdown – half tempo, breathy female vocals}
You feel like a hunter… you’re also the prey…\\
We know the game… we know the way…\\
Just for the night… just for the thrill…\\
One broken connection… but we got the win still…

\songpart{Final Chorus — Full energy, layered vocals}
Senses overload! Feel it in your bones!\\
Senses overload! Never let it go!\\
Dancing with fire on the edge of the floor\\
My wingman, my night, my greatest score!\\
Soundwall! Take me higher!\\
Soundwall! Burn it up tonight!

\songpart{Outro}
Music was loud… right in my bones…\\
Senses overload… we’re never going home…
\end{lyrics}


\songtitle{Keep Me}

\tag*{2026}\tag{dubstep}\tag{english-lyrics}
\begin{notes}
Club House adaptation of \emph{Stimuli}  (page \pageref{sec:stimuli}). Adaptation by Carlos Thompson and Grok.
\end{notes}
\begin{style}
dubstep, club house, 124 BPM, syncopated four-on-the-floor, half-time drop, deep bass baritone half-spoken vocal, close-mic delivery, vocal texture layering, brass stabs, synth bass stabs, sub bass glide, hand percussion, tape saturation, plate reverb, sidechain pumping, stumbling groove, steady lift
\end{style}
\begin{lyrics}
\songpart{Intro – Soft, layered whispers over pulsing pads}
Stimuli…\\
I crave stimuli…

\songpart{Verse 1 – Gentle rise, contained delivery}
It could be conversation drifting through the dark\\
Or mindless scrolling, chasing light across the glass\\
A world I build inside where no one else can reach\\
Or bodies moving slowly, skin forgetting speech\\
Warm skin against my skin, a quiet kind of heat\\
Or thoughts that fold together, sharp and bittersweet\\
I crave stimuli\\
Something to pull me under, something to keep me here

\annotation{Pre-Build – Subtle tension, breathy layers}\\
Keep my mind distracted\\
Keep the silence at bay\\
Keep my skin awakened\\
In the slowest, sweetest way

\annotation{Drop 1 – Deep bass and rolling grooves hit, vocals stay intimate}\\
I crave stimuli\\
New places, new hands, new currents in the night\\
I crave stimuli\\
Old flames I never counted, returning in the light\\
Something new\\
Something known\\
Something I have waited for without knowing how long

\songpart{Verse 2 – Deeper, more hypnotic}
Could be sex without meaning, just the need to feel alive\\
Could be thoughts unfolding like a universe inside\\
A new partner, a new challenge, or the same skin re-discovered\\
Doesn’t matter what it is, as long as I am covered\\
Keep my imagination fed\\
Keep the idle thoughts in check\\
Warm skin, quiet breath\\
A slow burn in the chest

\annotation{Pre-Build 2 – Rising textures, still restrained}\\
I crave…\\
Stimuli…\\
Keep me here\\
Keep me near\\
Keep the wanting clear

\annotation{Main Drop – Heavy but warm low end, vocals float above}\\
I crave stimuli\\
Skin on skin in the dark, no words, no alarm\\
I crave stimuli\\
A new world forming quietly inside my arms\\
Distraction as devotion\\
Movement as emotion\\
The bass beneath my bones\\
The hunger I won’t name out loud

\songpart{Breakdown – Stripped kick + airy pads}
Stimuli…\\
Warm skin…\\
I crave…
\annotation{soft layered echoes}

\annotation{Final Build \& Outro Drop – Intimate, fading layers}\\
Something new\\
Something old I haven’t held in time I haven’t counted\\
I crave stimuli…\\
I crave…
\end{lyrics}


\input{temp/new_old}
\songtitle{Buildings Worlds}

\tag*{2026}\tag{english-lyrics}\tag{pop-rock}\tag{living-fantasies}
\begin{notes}
Adaptation of Carlos Thompson's \emph{Living Fantasies} (see \pageref{sec:living-fantasies}) with Verse--Pre-Chorus--Chorus structure. Adapted by Claude.
\end{notes}
\begin{style}
pop rock, alt pop rock, baritone male vocals, 1980s snare gated, chorus guitar stacks, synth bass octave, tom fill lifts, Mellotron pad swells, tape saturation, plate reverb, pre-chorus build, stripped bridge, key lift finale, uneasy longing, defiant confession, mid-tempo 132 BPM
\end{style}
\begin{lyrics}
\songpart{Verse 1}
I keep living out these fantasies\\
Building them up brick by brick\\
An AI prompt to build a face\\
An archetype I pick

\songpart{Pre-Chorus}
Alternate histories, alternate me\\
Every version I could be\\
Tongues that no one's ever spoken\\
Borders that were never broken

\songpart{Chorus}
I'm building worlds to hide inside\\
Metro lines through Bogotá, redrawn a hundred times\\
Other planets, alien skies\\
Just to feel something — just to feel alive

\songpart{Verse 2}
Cultures no one's ever seen\\
Colors nothing here has named\\
I keep reinventing the machine\\
Till fantasy's not fantasy — it's just the frame

\songpart{Pre-Chorus}
Reinventing fantasy itself\\
Building castles no one else can see\\
Every world I make is somewhere else\\
Anywhere that isn't me

\songpart{Chorus — same, last line shifts}
I'm building worlds to hide inside\\
Metro lines through Bogotá, redrawn a hundred times\\
Other planets, alien skies\\
Just to feel something — just to numb this life

\songpart{Bridge — strip the instrumentation back here}
Maybe I'm not building worlds\\
Maybe I'm just running\\
From the one I'm already in\\
From the one I'm already in

\songpart{Final Chorus — full band back in, maybe a key lift}
I'm building worlds to hide inside\\
Metro lines through Bogotá, redrawn a hundred times\\
Other planets, alien skies\\
Just to feel something, just to numb myself
\end{lyrics}


\songtitle{Paper in the Grinder}

\tag*{2026}\tag{post-punk}\tag{grinder-2026}\tag*{english-lyrics}
\begin{notes}
Lyrics by SUNO based on a Grok interpretation of \emph{A Mind That Craves Novelty} by Carlos Thompson into suggested Indie Rock/Post-punk.
\end{notes}
\begin{style}
Alternative indie rock and post-punk with a dry, pulsing bassline, choppy muted guitar, and a tense mid-tempo swing; verse stays low, hesitant, almost spoken over the same looping groove, pre-chorus opens with climbing toms and overdriven guitar swells, chorus surges into jagged full-band desperation with doubled vocals, Close-mic verses, strained shouted hooks, delay throws on key words, filtered noise risers, wide gritty mix with sharp drums and a claustrophobic edge, dynamic, post-punk, vocal, rich, low, indie rock, world
\end{style}
\begin{lyrics}
\songpart{Verse 1}
I wake up on the wheel\\
Same room, same dented air\\
Full of plans like sparks\\
Too many to hold there

My head is a crowded train\\
My hands are parked outside\\
I write the first line twice\\
Then watch the whole thing slide

\songpart{Pre-Chorus}
It all stacks up\\
Then falls through\\
I can see it clear\\
I just can't do

\songpart{Chorus}
In the grinder\\
In the grinder\\
My whole life fits\\
But my hands don't

In the grinder\\
In the grinder\\
So much inside me\\
Nothing comes out

\songpart{Verse 2}
I carry whole cities\\
Behind my ribs at work\\
Blueprints in my bloodstream\\
But the page still hurts

Every task is a door\\
That asks for one small key\\
I got ten in my pocket\\
And none of them for me

\songpart{Pre-Chorus}
It all stacks up\\
Then burns blue\\
I can name the bridge\\
I can't walk through

\songpart{Chorus}
In the grinder\\
In the grinder\\
My whole life fits\\
But my hands don't

In the grinder\\
In the grinder\\
So much inside me\\
Nothing comes out

\songpart{Bridge}
I am not lazy\\
I am overloaded\\
Thoughts like coins in a dryer\\
Never sorted

I am not broken\\
I am overlit\\
Every bright idea\\
Hitting a locked switch

\songpart{Chorus}
In the grinder\\
In the grinder\\
My whole life fits\\
But my hands don't

In the grinder\\
In the grinder\\
So much inside me\\
Nothing comes out

In the grinder\\
In the grinder\\
Still I keep my name\\
Through the blackout
\end{lyrics}


\songtitle{Grinder Glass}

\tag*{2026}\tag{grinder-2026}\tag{english-lyrics}
\begin{notes}
Lyrics by SUNO based on a Grok interpretation of \emph{A Mind That Craves Novelty} by Carlos Thompson into suggested Experimental art pop / chamber pop.
\end{notes}
\begin{style}
Experimental art pop / chamber pop with a repetitive low-register pulse, tense fingerpicked guitar, muted piano, and odd-string swells; verses stay hushed and hesitant over a looping groove, pre-chorus opens with rising drums and anxious harmony stacks, chorus spikes hard with desperate full-voice doubles and gang ad-libs, Use filtered breath, sudden stop-times, reverse swells, and a brittle, intimate mix that blooms wide on the hook, chamber pop, dynamic, experimental, vocal, rich, low, art pop, world
\end{style}
\begin{lyrics}
\songpart{Verse 1}
I wake up late in the same old wire\\
Pocket full of plans, hands full of fire\\
Calendar stares, and I stare back\\
Five tabs open, one thought cracked

My head is a station at full noon\\
Every bright idea leaves too soon\\
I build a palace in a blink\\
Then lose the key in the kitchen sink

\songpart{Pre-Chorus}
Something in me\\
Keeps breaking the frame\\
I had ten brave starts\\
But I dropped every name\\
I can see the whole map\\
In a flash, in a flood\\
Then my body says stay\\
And my mind says run

\songpart{Chorus}
I’m stuck in the grinder
(Stuck in the grinder)\\
Too much sky, not enough hands\\
I’m stuck in the grinder
(Stuck in the grinder)\\
All this thunder, no plan\\
I am not lazy\\
I am not broken\\
I am a room\\
With every door open\\
But nothing comes out\\
When I call it by name\\
I’m stuck in the grinder\\
And I’m burning the same

\songpart{Verse 2}
You say, “Just start,” like it’s a small thing\\
Like my bones don’t ring when the phone rings\\
Like my whole life isn’t a stack\\
Of half-made roads and a getting-back

I read three books in the same hour\\
Then forget my tea, forget my power\\
I can feel tomorrow in my teeth\\
But today keeps folding under me

\songpart{Pre-Chorus}
Something in me\\
Is loud and unseen\\
Like a crowded train\\
With a seat in between\\
I got the spark\\
I got the lens\\
But the switch in my chest\\
Won’t catch at the end

\songpart{Chorus}
I’m stuck in the grinder
(Stuck in the grinder)\\
Too much sky, not enough hands\\
I’m stuck in the grinder
(Stuck in the grinder)\\
All this thunder, no plan\\
I am not lazy\\
I am not broken\\
I am a room\\
With every door open\\
But nothing comes out\\
When I call it by name\\
I’m stuck in the grinder\\
And I’m burning the same

\songpart{Bridge}
Look at the flood in my head\\
Look at the books on the bed\\
Look at the life I could build\\
If my feet would answer me still

I keep a whole galaxy hid\\
Under the lid, under the lid\\
If I could trade this ache for one clean move\\
I’d make the world hear what I know I can do

\songpart{Final Chorus}
I’m stuck in the grinder
(Stuck in the grinder)\\
Too much sky, not enough hands\\
I’m stuck in the grinder
(Stuck in the grinder)\\
All this thunder, no plan\\
I am not lazy\\
I am not broken\\
I am a room\\
With every door open\\
Say it again\\
Let it be named\\
I’m stuck in the grinder\\
But I’m more than the cage
\end{lyrics}


\songtitle{Paper In My Teeth}

\tag*{2026}\tag{grinder-2026}\tag{english-lyrics}
\begin{notes}
Lyrics by SUNO based on a Grok interpretation of \emph{A Mind That Craves Novelty} by Carlos Thompson into suggested Sadcore / lo-fi bedroom pop.
\end{notes}
\begin{style}
Sadcore / lo-fi bedroom pop with hesitant low-register verses over a repetitive guitar pulse; pre-chorus tightens with muted drums and rising harmony stacks, chorus bursts open with desperate doubled vocals and a ragged falsetto lift, Sparse room tone, tape hiss, filtered lifts into each turn, then a fuller crash on the hook, Intimate, brittle, and painfully close-mic, vocal, pop, lo-fi, rich, world, dynamic, low
\end{style}
\begin{lyrics}
\songpart{Verse 1}
I wake up late in my own head\\
Already carrying the whole day\\
A hundred tabs, a thousand sparks\\
And my hands still never find the page

My room is full of almosts\\
Half-built towers in the dark\\
I can name the shape of everything\\
But I can't make it leave a mark

\songpart{Pre-Chorus}
And I keep saying "tomorrow"\\
Like it's sitting in my lap\\
But the clock keeps chewing my reasons\\
And my reasons vanish back

\songpart{Chorus}
I'm in the grinder, baby\\
Mind on fire, hands gone cold\\
I'm in the grinder, baby\\
Too much light, too little hold\\
I got a whole world in me\\
And it won't come out clean\\
I'm in the grinder, baby\\
Stuck between my head and mean

\songpart{Verse 2}
Bills on the counter, untouched\\
Names I forgot to return\\
I make a list of ten small things\\
Then watch my focus burn

I can see the whole blueprint\\
I can feel the final scene\\
But every door inside my body\\
Says "not yet" when I mean "please"

\songpart{Pre-Chorus}
And I keep swearing "I got this"\\
With a grin that cracks in two\\
I'm the first one in the tunnel\\
But the light still runs from you

\songpart{Chorus}
I'm in the grinder, baby\\
Mind on fire, hands gone cold\\
I'm in the grinder, baby\\
Too much light, too little hold\\
I got a whole world in me\\
And it won't come out clean\\
I'm in the grinder, baby\\
Stuck between my head and mean

\songpart{Bridge}
Maybe I'm not broken\\
Maybe I'm a crowd\\
Maybe every buried answer\\
Just gets buried twice as loud

I don't need your tiny boxes\\
I don't fit the lock you drew\\
I am all this voltage\\
And a body not built to prove

\songpart{Final Chorus}
I'm in the grinder, baby\\
Mind on fire, hands gone cold\\
I'm in the grinder, baby\\
Too much light, too little hold\\
I got a whole world in me\\
And it won't come out clean\\
I'm in the grinder, baby\\
Stuck between my head and mean

I'm in the grinder, baby\\
Still here, still not through\\
I'm in the grinder, baby\\
Trying hard to be true
\end{lyrics}


\songtitle{Paper Cup Engine}

\tag*{2026}\tag{anti-folk}\tag{grinder-2026}\tag*{english-lyrics}
\begin{notes}
Lyrics by SUNO based on a Grok interpretation of \emph{A Mind That Craves Novelty} by Carlos Thompson into suggested Slacker rock / anti-folk.
\end{notes}
\begin{style}
Slacker rock / anti-folk with a monotone, droning guitar loop and loose, half-spoken low vocal in the verses; pre-chorus tightens with tambourine, bass lift, and ragged harmony stacks; chorus spikes into desperate shout-sing with doubled lead and gang echoes, Tape hiss, a warped pickup squeal, and a sudden drum fill mark transitions, Dry, intimate, and frayed at the edges, vocal, anti-folk, world, low, rock, dynamic, rich
\end{style}
\begin{lyrics}
\songpart{Verse 1}
I wake up late in the grind\\
Phone face-down on the chair\\
Three tabs, four tasks, one brain\\
All of them starting there\\
I got a storm in my skull\\
But my hands move like sleep\\
Big bright rooms in my head\\
And a box I can’t keep

\songpart{Pre-Chorus}
Say it again\\
Say it again\\
I had a whole fire in me\\
Right here, then gone again\\
I’m not lazy\\
I’m not broken\\
I’m just parked in the spill\\
With my heart on the page\\
And my body on still

\songpart{Chorus}
Stuck in the grinder\\
Stuck in the grinder\\
Too much inside, nowhere to land\\
Stuck in the grinder\\
Stuck in the grinder\\
I keep my life in an open hand

\songpart{Verse 2}
I make ten maps by noon\\
Never leave the front door\\
Start ten songs in the dark\\
Don’t finish one more\\
The trash keeps stacking up\\
Like a dare on the floor\\
I can see the whole path\\
Then I miss the next door

\songpart{Pre-Chorus}
Say it again\\
Say it again\\
I had a whole fire in me\\
Right here, then gone again\\
I’m not lazy\\
I’m not broken\\
I’m just parked in the spill\\
With my heart on the page\\
And my body on still

\songpart{Chorus}
Stuck in the grinder\\
Stuck in the grinder\\
Too much inside, nowhere to land\\
Stuck in the grinder\\
Stuck in the grinder\\
I keep my life in an open hand

\songpart{Bridge}
There’s a library in me\\
And a lock on the gate\\
There’s a whole parade\\
That I always miss late\\
If I could line up the sparks\\
If I could hold one thread\\
Maybe I’d build what I meant\\
Instead of living in my head

\songpart{Final Chorus}
Stuck in the grinder\\
Stuck in the grinder\\
Too much inside, nowhere to land\\
Stuck in the grinder\\
Stuck in the grinder\\
I keep my life in an open hand\\
Stuck in the grinder\\
Stuck in the grinder\\
Still I’m reaching through the sand
\end{lyrics}

\songtitle{You Would Destroy Me}

\tag*{2026}\tag{goth-metal}\tag{ai-based}\tag{a-coffee-not-shared}
\begin{notes}
Goth-metal adaptation by Grok of \emph{A Coffee Not Shared}.
\end{notes}
\begin{style}
dark goth metal, symphonic metal, 92 BPM, clean male vocals, whispered intro, distant piano, string ensemble sustains, choir pads, low register bass guitar, distorted rhythm guitars, double kick drums, minor key synth drones, cathedral plate reverb, tape saturation, wide stereo mix, breakdown pause, dramatic chorus lift, bittersweet longing
\end{style}
\begin{lyrics}
\songpart{Intro}
\annotation{Whispers over distant piano and cold reverb}\\
You…\\
Always you…

\songpart{Verse 1}
You’ve haunted these veins since the first black dawn\\
Your silhouette bleeds through every silent screen\\
A poison I never dared to taste\\
Velvet lies and thoughts that rot the bone

\songpart{Pre-Chorus}
I watch from the shadows where the light can’t reach\\
Counting the years like open wounds

\songpart{Chorus}
You would destroy me\\
You would tear the last light from my chest\\
So I stay silent in the dark\\
Where cowardice feels safer than the fall

You would destroy me\\
And I still can’t look away

\songpart{Verse 2}
Centuries of hunger, drowning in the dark\\
You burn like a dying star and I freeze in the glare\\
A thousand confessions buried in my throat\\
A sacred sin, my most exquisite disease

\songpart{Pre-Chorus}
I watch from the shadows where the light can’t reach\\
Counting the years like open wounds

\songpart{Chorus}
You would destroy me\\
You would tear the last light from my chest\\
So I stay silent in the dark\\
Where cowardice feels safer than the fall

You would destroy me\\
And I still can’t look away

\songpart{Bridge / Breakdown}
I’ve ached through endless nights\\
Consumed by the ghost of your presence\\
Too damned to believe I deserve your touch\\
So I choose the cold of abstinence

Because you…\\
Would annihilate what’s left of me

\annotation{Atmospheric Break}
\annotation{Clean guitars + distant choir}\\
A story not happening…\\
A black coffee left to die cold…\\
A conversation hanging from the noose…

\songpart{Final Chorus}
You would destroy me\\
You would drag me under and crush the last of the light\\
So I remain in the dark\\
Where longing is safer than the truth

You would destroy me\\
And I still can’t look away

\songpart{Outro}
You…\\
A requiem that never plays\\
You…\\
Headlights that will crush me\\
You…

\annotation{fade into silence}
\end{lyrics}
\songtitle{Writing Takes Its Toll}

\tag*{2026}\tag{folk-rock}\tag{writing-aint-easy}
\begin{notes}
Lyrics by ChatGPT, reworking \emph{Writing Ain't Easy} by Carlos Thompson into tighter ballad-meter rhyme and rhythm. Retrieved from Carlos Thompson's Suno page (V6-MINI, 13 September 2026).
\end{notes}
\begin{style}
Folk with rock band instrumentation, raspy male tenor with urgent delivery, drum kit driving beneath bass and ringing electric-guitar riffs, tape saturation and lightly compressed room texture, rapid-fired forward momentum
\end{style}
\begin{lyrics}
\songpart{Verse 1}
So I was proving I could write\\
A story with a twist,\\
Announce the hero to the crowd,\\
Make him a man they'd miss.

Make him so worthy of their trust,\\
So likable and bright,\\
Then give them someone they will hate,\\
Not evil, but too right.

\songpart{Verse 2}
Not for his villainy alone,\\
Nor for the crimes he'd done,\\
But 'cause he seemed too perfect there,\\
Too perfect next the one.

The villain's likeness feels unearned,\\
The hero wins their hearts,\\
But subvert all their expectations—\\
That's where the twist now starts.

\songpart{Pre-Chorus}
A plot twist, I said,\\
Subvert, I said,\\
Those were the words inside my head,\\
The plan that lay ahead.

\songpart{Chorus}
All the ideas in your mind,\\
But writing takes its toll;\\
To build a world, to build its lore,\\
And give the story soul.

Backstories and motivations,\\
A voice that carries through,\\
A voice that makes the world feel real,\\
And makes them want it too.

\songpart{Verse 3}
Mislead them with each trope you set,\\
To make them cheer the hero,\\
And make them wish the villain loses,\\
Then turn the tale below.

Reveal that hero and villain\\
Were swapped before the end,\\
And make the pieces fit so well\\
The mirror lines extend.

\songpart{Verse 4}
You need a past that makes it work,\\
A reading that holds true,\\
Then one that makes the twist reveal\\
A deeper point of view.

The inciting incident must\\
Begin to take its shape:\\
If he was really the villain,\\
Why act the hero's part?

\songpart{Pre-Chorus}
And how to build the other man\\
Into a villainous role,\\
So when the truth is finally shown,\\
He still remains whole.

\songpart{Chorus}
All the ideas in your mind,\\
But writing takes its toll;\\
Motivations must make sense,\\
And every twist feel whole.

A voice that carries through the tale,\\
A voice that you consume;\\
A story that can pull you in\\
And make the characters bloom.

\songpart{Verse 5}
With all the main elements in place,\\
It's time to write the tale;\\
Let's choose a voice, let's choose a form,\\
And set the story's trail.

Let's start the story in medias res,\\
Our hero at his worst,\\
His failed mission lying behind,\\
His motivation burst.

\songpart{Bridge}
And build it up from there again,\\
Rebuild the fading spark;\\
Rebuild the reason he still fights,\\
Though everything seems dark.

And that was the easy part to do,\\
The hardest lay ahead;\\
There were still plans to make and break,\\
And words that must be said.

\songpart{Chorus}
All the ideas in your mind,\\
But writing takes its toll;\\
As writing moves, the ideas shift,\\
And characters take hold.

A voice that changes as it grows,\\
A voice you want to keep;\\
A story finding out itself\\
What promises to reap.

\songpart{Breakdown}
As the story progresses on,\\
The ties begin to bind;\\
And then you realize one day\\
He wouldn't act that kind.

Everything you put into the plot\\
Makes motives rearrange;\\
The things that once seemed obvious\\
Now suddenly seem strange.

\songpart{Final Chorus}
All the ideas in your mind,\\
But writing takes its toll;\\
A voice that carries through the tale,\\
And makes the story whole.

And now the novel's unfolding,\\
Mostly good, mostly fine;\\
But there is still one thing that feels\\
Like it refuses to align.

\songpart{Outro}
Something is still missing there,\\
Something I cannot see;\\
The story almost fits together—\\
But not quite yet, not for me.
\end{lyrics}


\end{multicols*}

\chapter{Adaptations III: Cross-language}
\albumcover{covers/internal/adaptations-iii-cross-language}
The purest form of adaptation: not a new arrangement, not new words, just the same song learning to breathe in a different language -- carried away, carried across, carried over, note for note.
\clearpage\begin{multicols*}{4}
\songtitle{Carried Away}

\tag*{2026}
\begin{notes}
English language adaptation via Claude of \emph{Dejándome llevar}, produced as hyperpop.
\end{notes}
\begin{style}
hyperpop, 155 BPM, helium pitched vocals, heavy auto tune, metallic snare rushes, distorted 808 sub, glitter synth arpeggios, cold digital clicks, low pass filter breaks, hard volume drops, phone lock click, glitch video game lead, dark techno drop, tape saturation, plate reverb, stereo widening, 2000s club sheen, sudden dropouts, bittersweet detachment
\end{style}
\begin{lyrics}
\songpart{Verse 1}(Light, glittery synth arpeggios)\\
I saw your face as I was scrolling down\\
I saw your face as I was scrolling down\\
I saw your face as I was scrolling down\\
And then my whole world rearranged.

Nothing else mattered anymore\\
Nothing else mattered anymore\\
Nothing else mattered anymore\\
Because the instant that it flashed, it changed.

\songpart{Pre-Chorus}(Vocals pitched way up, helium effect)\\
You looked so happy on my screen\\
You looked so happy on my screen\\
You looked so happy on my screen\\
And for a second, I was happy too.

\songpart{Chorus}(The beat drops hard: heavy 808 bass, distorted and aggressive)\\
Life is so simple when I’m just\\
Life is so simple when I’m just\\
Life is so simple when I’m just\\
Letting the moment carry me away.

But the machine keeps spinning round\\
But the machine keeps spinning round\\
But the machine keeps spinning round\\
And there’s no way for me to make it stay.

\songpart{Verse 2}(Music suddenly cuts to a cold, dry digital click)\\
I read the words about your tragic loss\\
I read the words about your tragic loss\\
I read the words about your tragic loss\\
And I could feel the trauma in my chest.

My heavy heart is hurting for your grief\\
My heavy heart is hurting for your grief\\
My heavy heart is hurting for your grief\\
For just how long you’ll have to face this test.

\songpart{Pre-Chorus}(A low-pass filter makes the audio sound muffled, mimicking isolation)\\
And you stay there for two whole minutes\\
And you stay there for two whole minutes\\
And you stay there for two whole minutes\\
Before I swipe up to the next profile.

\songpart{Chorus}(The massive 808 bass kick explodes back in)\\
Life is so simple when I’m just\\
Life is so simple when I’m just\\
Life is so simple when I’m just\\
Letting the moment carry me away.

But the machine keeps spinning round\\
But the machine keeps spinning round\\
But the machine keeps spinning round\\
And there’s no way for me to make it stay.

\songpart{Bridge}(Tonal whiplash: The music changes genres completely every two lines—glitching from a cute video game melody to a dark techno beat)\\
Now it’s a kitten in distress\\
Now it’s a kitten in distress\\
Now it’s a kitten in distress\\
Taking the spot where you just used to stand.

Now it's a self-help guru's clip\\
Now it's a self-help guru's clip\\
Now it's a self-help guru's clip\\
Or proof the world is burning out of hand.

\songpart{Final Chorus}(All elements clash together—shattering glass sounds, synth sirens, hyper-fast tempo)\\
Life is so simple when I’m just\\
Life is so simple when I’m just\\
Life is so simple when I’m just\\
Letting the moment carry me away.

But the machine keeps spinning round\\
But the machine keeps spinning round\\
But the machine keeps spinning round\\
And there’s no way for me to make it stay.

The moment’s everything but passes by\\
The moment’s everything but passes by\\
The moment’s everything but passes by\\
And everything stays exactly the same.

\songpart{Outro}(A sudden, dead-stop silence. Just the sound of a smartphone locking click.)\\
Exactly the same.
\end{lyrics}


\songtitle{If You Don’t Need It}

\tag*{2026}\tag{techno}\tag*{english-lyrics}\tag{piel-2026}
\begin{notes}
Adaptation into English of \emph{Si no te hace falta}, from Carlos Thompson's \emph{Piel}. See \pageref{sec:piel} for original lyrics.
\end{notes}
\begin{style}
techno house, trance music, 808 sub-bass, distorted synthesizer stack, dubstep bass drop, close-mic baritone, almost spoken vocal, textured vocal layers, soundwall breakdowns, silence gaps, volume swells, machine-like percussion, driving 128 BPM, anxious desire
\end{style}
\begin{lyrics}
\songpart{Intro – Soft, breathed}
If you don’t need it…\\
I do.
\songpart{Verse 1 – Almost spoken, rhythmic flow}
If you don’t need it, I do.\\
My skin asks for skin.\\
My body asks.\\
Asks for release.\\
Asks for signals that shut the mind down.\\
Asks for touch,\\
asks for friction,\\
asks for heat.\\
Up to a point I manage.\\
Up to a point.\\
My head fills in.\\
Porn is enough.\\
The chatbot works.\\
I’ve tried the texts, they work…\\
Up to a point.
\annotation{Pre-Build – Layered, slowly rising tension}\\
But none of that has skin.\\
None of that has skin.\\
Up to a point I manage…\\
But past that line,\\
it’s not enough.
\annotation{Drop 1 – Warm deep bass hits, vocals stay intimate}\\
If you don’t need it, I do.\\
Skin on skin, flesh and bone.\\
If you don’t need it, I do.\\
Something real that’s here.
\songpart{Verse 2 – More introspective, spoken-like}
Loving for real, I guess that works too.\\
I guess that will be needed too.\\
Loving for real…\\
Whatever that even means.\\
If it means anything.\\
Right now all I know\\
is I need skin.\\
Someone breathing close,\\
someone who’s here.\\
Whose weight sinks the bed.\\
Whose heat isn’t a screen.
\annotation{Pre-Build 2 – Slow-burning tension}\\
Up to a point I manage…\\
Up to a point.\\
But the body knows.\\
The body asks.\\
And it won’t take substitutes.
\annotation{Final Drop – Emotional but restrained, voice floating over the groove}\\
If you don’t need it, I do.\\
My skin asks for skin.\\
My body asks.\\
Asks for real release.\\
Asks for signals that hurt a little.\\
If you don’t need it…\\
I do.
\songpart{Outro – Fading layers}
Up to a point…\\
Up to a point…\\
I do.\\
I do.
\end{lyrics}


\songtitle{Inner Space}

\tag*{2026}\tag{synth-pop}\tag{mundos-internos}
\begin{notes}
English adaptation of \emph{Mundos internos}.
\end{notes}
\begin{style}
candy pop, synth-pop, 132 BPM, mellow male vocals, close-mic tenor, drum machine backbeat, chunky electric bass, bright synth arpeggios, analog polysynth pads, gated snare, tape saturation, plate reverb, sidechain pumping, retro 80s mix, urban routine, bittersweet introspection
\end{style}
\begin{lyrics}
\songpart{Verse 1}
An espresso from my cappuccino maker,\\
hot milk like it’s a proper latte,\\
cheap grocery store cinnamon rolls,\\
and the pills for my head and my blood pressure.

Reading and answering questions on Quora,\\
or getting lost in short clips on Instagram,\\
losing track of time and rushing off,\\
stepping out into the street to wait at the bus stop.

\songpart{Pre-Chorus}
A city with sun or with rain,\\
with roads under construction and street life,\\
motorbikes weaving and bicycles flowing,\\
pedestrians and cars, buses, noise, and life.

\songpart{Chorus}
And inside my head, other worlds and places:\\
a podcast, music, or worlds of my own,\\
watching the street or staring at my phone,\\
what happens outside and within my inner space.

\songpart{Verse 2}
The office routine that is rarely routine,\\
justifying hours out of sync,\\
justifying a routine and a presence\\
with masks and alt-tabs, strange symphonies.

Lunch hour, office colleagues,\\
each with their budget and their stories,\\
work in common or the latest news,\\
what passes through our mundane lives.

\songpart{Pre-Chorus}
A city with sun or with rain,\\
shopping malls, corner shops,\\
going out for lunch or office microwaves,\\
and grabbing a coffee as part of the grind.

\songpart{Chorus}
And inside my head, other worlds and places:\\
a podcast, music, or worlds of my own,\\
watching the street or staring at my phone,\\
what happens outside and within my inner space.

\songpart{Bridge}
Life goes on,\\
life slips by,\\
others cross our paths,\\
work to be done,\\
stories to tell,\\
coordinating, living:\\
a world where you are not alone.

\songpart{Breakdown}
What you do and what you don't do,\\
coworkers who depend on you,\\
a family waiting,\\
a life moving on.

\songpart{Final Chorus}
And inside my head, other worlds and places:\\
a podcast, music, or worlds of my own,\\
watching the street or staring at my phone,\\
what happens outside and within my inner space.

An inner space.
\end{lyrics}


\songtitle{Frisson}

\tag*{2010}\tag{eurodance}\tag{french-lyrics}
\begin{notes}
French-language eurodance version of \emph{Skälvning}/\emph{Temblando}/\emph{Trembling Blues}, translated by DeepSeek -- lyrics based on a work by Carlos Thompson from 2010 (posted under Enwigh Peady). Retrieved from rataflechera's Suno page, 14 September 2026.
\end{notes}
\begin{style}
Eurodance: analog synth stabs, arpeggiated synth lead, repeated piano note, distorted electric riff, palm-muted guitar chugs, low synth bass; deep baritone with spoken intro, rap verses, whispered and broken passages, guttural shouted hook and layered growls; 128 BPM driving four-on-the-floor with double-time hi-hats, plate reverb, 90s club mix, static-frayed outro
\end{style}
\begin{lyrics}
\songpart{Intro – Nappe ambiante, note de piano répétée, chuchotement lointain}

Je sens mon cœur...\\
des battements interdits...\\
des battements interdits...

\annotation{Couplet 1 – Rappé, syncopé, sur riff puissant}

Un cœur dans ma poitrine – des battements défendus,\\
une vie qui se dissout, jour après jour perdu.\\
Ce que je croyais acquis – parti, sans aucune trace,\\
j'accélère, je simule, je deviens un fantôme en place.

Je joue à être un autre –\\
je brise mes chaînes de mes mains qui tremblent.\\
Attiré par le miel – vert comme l'océan,\\
tiré vers la flamme, vers mon tombeau sanglant.

\annotation{Pré-Refrain – Montée de tension, charley en double-croche}

Mon bourreau...\\
Mon cœur...\\
Mon amulette...\\
Et mon bourreau...

\annotation{Refrain – Puissance maximale, crié, guttural}

Et Je Frissonne !\\
J'Accélère – Et Je Frissonne !\\
Je Ne Sais Plus Quel Jeu Je Joue !\\
Je Perds Innocence, Raison – Ma Vie Dans Une Source !

Je Frissonne Car Je Ne Peux Pas Être Clair –\\
Ni Face À Mon Adversaire, Ni Face Aux Gradins,\\
Ni Face À Moi-Même !

\annotation{Couplet 2 – Plus doux, guitare clean, à demi-chuchoté}

Mon fleuve de vie s'écoule –\\
en battements interdits de ma poitrine.\\
Tout s'effondre – je le vois,\\
mes doigts se crispent sur les touches ivoirines.

Dois-je dire ce que je n'ose pas ?\\
Dois-je entendre sa voix – ou voir le visage\\
du jeu que je joue ?\\
La vérité : je n'ose pas briser les chaînes.

\annotation{Pré-Refrain – Plus lourd, tempo ralenti, basse profonde}

Mon bourreau...\\
Mon cœur...\\
Mon amulette...\\
Et mon bourreau...

\annotation{Refrain – Puissance totale, cris superposés, growls en soutien}

Et Je Frissonne !\\
J'Accélère – Et Je Frissonne !\\
Je Ne Sais Plus Quel Jeu Je Joue !\\
Je Perds Innocence, Raison – Ma Vie Dans Une Source !

Je Frissonne Car Je Ne Peux Pas Être Clair –\\
Ni Face À Mon Adversaire, Ni Face Aux Gradins,\\
Ni Face À Moi-Même !

\annotation{Pont / Break – Le beat tombe à mi-temps. Larsen. Paroles parlées, intimes}

Le jeu de la vie...\\
Je le perds déjà.\\
Parce que je ne suis pas moi –\\
et pas non plus mon fantôme.

\annotation{Silence.}

\annotation{Chuchotement, voix brisée :}
...je n'ose pas...

\songpart{Outro – Distordu, chaotique, le riff se désaccorde, les cris se fragmentent en static}

Frissonne !\\
Frissonne !\\
Tout – S'Effondre – En Poussière !

\annotation{Chuchotement final, brisé, sur guitare clean :}
...pas moi... pas mon fantôme...

\annotation{Un accord lourd, palm-mute. Résonne. S'éteint.}
\end{lyrics}


\end{multicols*}

\backmatter
\printindex[songs]
\clearpage
\printindex[tags]
\end{document}
